\documentclass[11pt,openany]{book}

\usepackage[T1]{fontenc}
\usepackage[utf8]{inputenc}
\usepackage{libertine}
\usepackage{microtype}
\usepackage{titlesec}
\usepackage{fancyhdr}
\usepackage{geometry}
\usepackage{textcomp}
\usepackage{multicol}
\usepackage{enumitem}
\usepackage[hidelinks]{hyperref}

\geometry{paperwidth=6.7in, paperheight=9.6in,
  inner=0.72in, outer=2.05in, top=0.82in, bottom=0.9in,
  marginparwidth=1.62in, marginparsep=0.2in}

\frenchspacing
\setlength{\parindent}{1.2em}
\setlength{\parskip}{0pt}
\setlength{\emergencystretch}{2em}
\setlength{\marginparpush}{2pt}
\linespread{1.05}

\hypersetup{pdftitle={Lingua Latina per se Illustrata},pdfauthor={Hans H. Ørberg},bookmarksnumbered=false}

\titleformat{\chapter}[display]{\normalfont\sffamily}{}{0pt}
  {\Large\bfseries\raggedright}[\vspace{8pt}{\titlerule[0.5pt]}]
\titlespacing*{\chapter}{0pt}{56pt}{30pt}
\titleformat{\section}[block]{\normalfont\sffamily\bfseries\large}{}{0pt}{}
\titlespacing*{\section}{0pt}{22pt}{10pt}
\titleformat{\tableofcontents}[block]{\normalfont\sffamily\bfseries\Large}{}{0pt}{}
\renewcommand{\contentsname}{Index Capitulōrum}

\newcommand{\bookchapter}[2]{%
  \clearpage\phantomsection
  \chapter*{{\normalfont\sffamily\small\mdseries #1}\\[6pt]#2}%
  \addcontentsline{toc}{chapter}{#1\quad #2}%
  \markboth{#1. #2}{#1. #2}\thispagestyle{plain}}

\newcommand{\booksection}[1]{%
  \phantomsection\section*{#1}%
  \addcontentsline{toc}{section}{#1}\markright{#1}}

\newcommand{\bookbackmatter}[1]{%
  \clearpage\phantomsection\chapter*{#1}%
  \addcontentsline{toc}{chapter}{#1}\markboth{#1}{#1}\thispagestyle{plain}}

\pagestyle{fancy}
\fancyhf{}
\renewcommand{\headrulewidth}{0.3pt}
\fancyhead[LE]{\sffamily\small\thepage}
\fancyhead[LO]{\sffamily\small\itshape\nouppercase{\rightmark}}
\fancyhead[RE]{\sffamily\small\itshape\nouppercase{\leftmark}}
\fancyhead[RO]{\sffamily\small\thepage}
\fancypagestyle{plain}{\fancyhf{}\renewcommand{\headrulewidth}{0pt}%
  \fancyfoot[C]{\sffamily\small\thepage}}

\newcommand{\navnote}[1]{\marginpar{\raggedright\scriptsize\linespread{1.0}%
  \selectfont #1}}
\newcommand{\contr}{$\leftrightarrow$}

\begin{document}

\begin{titlepage}
\centering
\vspace*{1.5in}
{\sffamily\bfseries\fontsize{30}{36}\selectfont Lingua Latīna\par}
\vspace{0.16in}
{\sffamily\large per sē illūstrāta\par}
\vspace{0.5in}
{\sffamily\bfseries\large Pars I\par}
\vspace{0.14in}
{\sffamily\Large Familia Rōmāna\par}
\vspace{0.6in}
{\sffamily\large Hans H. Ørberg\par}
\vfill
{\sffamily\small Domus Latīna\par}
\end{titlepage}

\pagenumbering{roman}
\pagestyle{fancy}
\phantomsection
\chapter*{Dē hāc editīone}
\addcontentsline{toc}{chapter}{Dē hāc editīone}
\markboth{Dē hāc editīone}{Dē hāc editīone}

The text is Hans H. \O{}rberg's \emph{Lingua Lat\={\i}na per s\=e ill\=ustr\=ata}, Pars I: \emph{Familia R\=om\=ana} (Domus Lat\=ina, 1991 and later printings), from the scan supplied with this collection. The book is in copyright: this edition is a private re-typesetting for the owner's own reading, not for distribution.

The scan carries the text layer made when it was digitised. That layer keeps the macrons of the printed book, which matter in a Latin reader, but it misreads badly wherever the book sets Latin in italic type --- in the grammar summaries, in the exercise paradigms and in the marginal glosses. The pages were therefore read a second time with Tesseract's Latin model and the two readings merged, word by word, checking each doubtful word against the book's own vocabulary. Reading errors remain, especially in the paradigms and in the outer margin, where the print is small and the OCR had least to work with; nothing has been invented or silently rewritten.

The layout follows the printed book: the running text with the Latin glosses beside it in the outer margin, the vocabulary of each chapter gathered at its end, then the grammar summary and the three pensa. The pictures that carry much of the book's teaching are not reproduced here --- a re-typeset edition cannot redraw them --- and neither are the printed line numbers, which refer to lines of the printed page. The declension tables (Tabula D\=ecl\=in\=ati\=onum), the Roman calendar, the word list and the grammatical index of the printed book, and the list of irregular forms (Formae m\=ut\=atae), are not reproduced: they are set in columns of very small type, and the scan's OCR cannot be trusted with them. The abbreviations the book uses are listed under \emph{Notae} at the end.

\tableofcontents
\clearpage
\pagenumbering{arabic}
\pagestyle{fancy}

\bookchapter{Cap. I}{Imperium Rōmānum}

\navnote{Italia... est sunt: Italia in Eurōpā est; Italia et Graecia in Eurōpā sunt}Rōma in Italiā est. Italia in Eurōpā est. Graecia in Eurōpā est. Italia et Graecia in Eurōpā sunt. Hispānia quoque in Eurōpā est. Hispānia et Italia et Graecia in Eurōpā sunt. S

Aegyptus in Eurōpā nōn est, Aegyptus in Āfricā est. Gallia nōn in Āfricā est, Gallia est in Eurōpā. Syria nōn est in Eurōpā, sed in Asiā. Arabia quoque in Asiā est. Syria et Arabia in Asiā sunt. Germānia nōn in Asiā, sed in Eurōpā est. Britannia quoque in Eurōpā est. Germānia et Britannia sunt in Eurōpā.

Estne Gallia in Eurōpā? Gallia in Eurōpā est. Estne Rōma in Galliā? Rōma in Galliā nōn est. Ubi est Rōma? Rōma est in Italiā. Ubi est Italia? Italia in Eurōpā est. Ubi sunt Gallia et Hispānia? Gallia et Hispānia in Eurōpā sunt.

Estne Nīlus in Eurōpā? Nīlus in Eurōpā nōn est. Ubi est Nīlus? Nīlus in Āfricā est. Rhēnus ubi est? Rhēnus est in Germāniā. Nīlus fluvius est. Rhēnus fluvius est. Nīlus et Rhēnus fluviī sunt. Dānuvius quoque fluvius: \navnote{Nīlus fluvius est; Nīlus et Rhēnus fluviī sunt}est. Rhēnus et Dānuvius sunt fluviī in Germāniā. Tiberis fluvius in Italiā est.

Nīlus fluvius magnus est. Tiberis nōn est fluvius \navnote{parvus --- \contr{} magnus fluvius magnus/parvus fluvit}magnus, Tiberis fluvius parvus est. Rhēnus nōn est fluvius parvus, sed fluvius magnus. Nīlus et Rhēnus nōn fluviī parvī, sed fluviī magnī sunt. Dānuvius quoque fluvius magnus est.

\navnote{Corsica Insula est; Corsica et Sardinia īnsulae sunt insula magna/parva īnsulae magnae/parvae}Corsica īnsula est. Corsica et Sardinia et Sicilia īnsulae sunt. Britannia quoque īnsula est. Italia īnsula nōn est. Sicilia īnsula magna est. Melita est īnsula parva. Britannia nōn Insula parva, sed insula magna est. Sicilia 30 et Sardinia nōn īnsulae parvae, sed īnsulae magnae sunt.

Brundisium oppidum est. Brundisium et Tusculum \navnote{Brundisium oppidum est; Brundisium et Tusculum oppida sunt oppidum magnum/ parvum oppida magnō/parva}oppida sunt. Sparta quoque oppidum est. Brundisium est oppidum magnum. Tusculum oppidum parvum est. Delphī quoque oppidum parvum est. Tusculum et Delphī nōn oppida magna, sed oppida parva sunt.

Ubi est Sparta? Sparta est in Graeciā. Sparta est oppidum Graecum. Sparta et Delphī oppida Graeca sunt. Tusculum nōn oppidum Graecum, sed oppidum Rōmānum est. Tusculum et Brundisium sunt oppida Rōmāna. Sardinia īnsula Rōmāna est. Crēta, Rhodus, Naxus, Samos, Chios, Lesbos, Lēmnos, Euboea sunt īnsulae Graecae. In Graeciā multae īnsulae sunt. In Italiā et in Graeciā sunt multa oppida. In Galliā et in Ger- 45 māniā multī sunt fluviī. Suntne multī fluviī et multa \navnote{paucī -ae -a --- multī -ae -a: multī/paucī fluviī, multae/paucae īnsulae, multa/pauca oppida num Creta...est? = est-ne}oppida in Arabiā? In Arabiā nōn multī, sed paucī fluviī sunt et pauca oppida.

Num Crēta oppidum est? Crēta oppidum nōn est! Quid est Crēta? Crēta īnsula est. Num Sparta insula est? \navnote{Rōma in Graeciā est? Rōma in Graeciā nōn est}Sparta nōn est insula! Quid est Sparta? Sparta oppidum est. Rhēnus quid est? Rhēnus est magnus fluvius. Num ōceanus Atlanticus parvus est? Nōn parvus, sed magnus est ōceanus. 55

Ubi est imperium Rōmānum? Imperium Rōmānum est in Eurōpā, in Asiā, in Āfricā. Hispānia et Syria et Aegyptus prōvinciae Rōmānae sunt. Germānia nōn est \navnote{imperium Rōmānum tn imperiō Rōmānō}prōvincia Rōmāna: Germānia in imperiō Rōmānō nōn est. Sed Gallia et Britannia sunt prōvinciae Rōmānae. In imperiō Rōmānō multae sunt prōvinciae. Magnum est imperium Rōmānum! NC

\booksection{Litterae et numerī}

I et II numeri sunt. III quoque numerus est. I, n, III numerī Rōmānī sunt. I et n sunt parvī numerl. CIO magnus numerus est.

A et B litterae sunt. c quoque littera est. a, b, c sunt trēs litterae. A est littera prima (I), B littera secunda (n), c littera tertia (in). r littera Graeca est. c est littera Latīna. c et d litterae Latīnae sunt. r et a sunt litterae \navnote{mānus}Graecae.

Fluvius et oppidum vocābula Latīna sunt. Ubi quoque \navnote{vocābulum: in vocābulō capitulum primum (cap. D): in capitulō prīmo}vocābulum Latīnum est. In vocābulō ubi sunt trēs litterae. In capitulō prīmō mille vocābula sunt.

In vocābulō īnsula sex litterae et trēs syllabae sunt: syllaba prima īn-, secunda -su-, tertia -la. In vocābulō 75 nōn sunt trēs litterae et ūna syllaba.

Quid est iii? III numerus Rōmānus est. r quid est? r littera Graeca est. Num c littera Graeca est? Nōn littera Graeca, sed littera Latīna est c. Estne B littera prima? b nōn littera prima, sed secunda est. Quid est nōn? Nōn 80 est vocābulum Latīnum. Nōn \contr{} sed, magnus, numerus vocābula Latīna sunt. Vocābulum quoque vocābulum Latīnum est!

\booksection{Grammatica Latīna}

\navnote{(numerus) singulāris: 1; plūrālis: II, III... fluvius magnus fluvit magn?}Singulāris et plūrālis [A] Nīlus fluvius magnus est.

Nīlus et Rhēnus fluviī? magnī sunt.

Fluvius singulāris

est. 'Fluvi? plūrālis est. Singulāris: -u\$. Plūrālis: -ī,

Exemplum: numerus, numeri.

\navnote{exemplum exempla insula magna īnsulae magnae}i parvus numerus est. i et II parvī numeri sunt. [B] Corsica Insula magna est

Corsica et Sardinia īnsulae magnae sunt.

Insule' singulāris est. ``Īnsulae'' plūrālis est. Singulāris: -a. Plūrālis: -ae.

Exempla: littera, litterae; prōvincia, prōvinciae.

a littera Latīna est. A et B litterae Latīnae sunt. Gallia est prōvincia Romāna. Gallia et Hispānia prōvinciae Rōmānae sunt. [C] Brundisium

oppidum magnum est.

\navnote{oppidum magnum oppida magna}Brundisium et Sparta oppida magna sunt.

Oppidum' singulāris est. 'Oppida' plūrālis est. Singulāris: -um. Plūrālis: -a.

Exempla: vocābulum, vocābula; exemplum, exempla.

Littera est vocābulum Latīnum, nōn Graecum. Littera et numerus nōn vocābula Graeca, sed Latina sunt.

\booksection{Pēnsum A}

\navnote{pensum pēnsa}Nīlus fluvi- est. Nīlus et Rhēnus fluvi

. Crēta īnsul Crēta et Rhodus Insul- sunt. Brundisium oppid

. Brundisium et Tusculum oppid

Rhēnus fluvi- magn- est. Tiberis est fluvi- parv-. Rhēnus et Dānuvius nōn fluvi- parv-, sed fluvi- magn- sunt. Sardinia Insul- magn- est. Melita īnsul- parv- est. Sardinia et Sicilia nōn Insul- parv-, sed īnsul- magn- sunt. Brundisium nōn oppid- parv-, sed oppid- magn- est. Tūsculum et Del- phi nōn oppid- magn-, sed oppid- parv- sunt.

Crēta Insul- Graec- est. Lesbos et Chios et Naxus sunt īnsul- Graec-. In Graeciā mult- īnsul- sunt. In Galliā sunt mult- fluvi-. In Italiā mult- oppid- sunt. In Arabiā sunt pauc- fluvi- et pauc- oppid-.

A et b litter- Latīn- sunt. c quoque litter- Latīn- est. Multī et paucī vocābul- Latīn- sunt. Ubi quoque vocābul- Latīnest. i et II numer- Rōmān- sunt. iii quoque numer- Rōmānest.

\booksection{Pēnsum B}

Sicilia--- --- est. Italia Insula --- est. Rhēnus --- est. Brundisium --- est. Sicilia et Sardinia --- magnae sunt. Melita īnsula --- est. Britannia nōn --- parva, sed --- ---

est. Brundisium nōn --- ---, sed --- magnum est. Est--- Brundisium in Graeciā? Brundisium --- est in Graeciā, --- in Italiā. --- est Sparta? Sparta est in Graeciā: Sparta oppidum --- est. Delphī--- --- oppidum Graecum est. Euboea, Naxus, Lesbos, Chios --- Graecae sunt. In Graeciā sunt --- īnsulae.

Quid est m? iii i11 --- est. --- est A? A littera est. A, B, c --- Latinae sunt. --- r littera Latina est? r --- littera ---, sed littera --- est. īnsula --- Latīnum est.

\booksection{Pēnsum C}

Ubi est Rōma? Estne Sparta in Italiā? Ubi est Italia? Ubi sunt Syria et Arabia? Estne Aegyptus in Asiā? Ubi sunt Sparta et Delphī? Ubi est Brundisium? Quid est Brundisium? Num Crēta oppidum est? Estne Britannia īnsula parva? Quid est Tiberis? Quid est D? Num a littera Latīna est? Estne II magnus numerus?

\booksection{Vocābula nova}
\noindent\small fluvius, īnsula, oppidum, ōceanus, imperium, prōvincia, numerus, littera, vocābulum, capitulum, syllaba, exemplum, pēnsum, magnus, parvus, Graecus, Rōmānus, Latīnus, multī, paucī, ūnus, duo, trēs, sex, mīlle, prīmus, secundus, tertius, est, sunt, in, et, sed, nōn, quoque, ubi?, num?, quid?, grammatica, singulāris, plūrālis.

\bookchapter{Cap. II}{Familia Rōmāna}

\navnote{DAVVS MEDVS SYRA DELIA}Iūlius vir Rōmānus est. Aemilia femina Rōmāna est. \navnote{ūnus (i) vir duo (n) virī ūnus puer duo pueri}Mārcus est puer Rōmānus. Quīntus quoque puer Rōmānus est. Iūlia est puella Rōmāna.

Mārcus et Quīntus nōn virī, sed pueri sunt. Viri sunt S [ūlius et Mēdus et Dāvus. Aemilia et Dēlia et Syra sunt feminae. Estne femina Iūlia? Nōn femina, sed parva puella est Iūlia.

Iūlius, Aemilia, Mārcus, Quīntus, Iūlia, Syra, Dāvus, Dēlia Mēdusque sunt familia Rōmāna. Iūlius pater est. Aemilia est māter. Iūlius pater Mārcī et Quīntī est. Iūlius pater Iūliae quoque est. Aemilia est māter Mārcī et Quīntī et Iūliae. Mārcus filius Iūliī est. Mārcus filius Aemiliae est. Quīntus quoque filius Iūlil et Aemiliae est. Iūlia est filia Iūliī et Aemiliae. 15

\navnote{quis? quae? quis est Mārcus? quae est Iūlia? quis est pater Mārcī? quae est māter Mārcī?}Quis est Mārcus? Mārcus puer Rōmānus est. Quis pater Mārcī est? Iūlius pater Mārcī est. Quae est māter Mārcī? Māter Mārcī est Aemilia. Quae est Iulia? Iūlia est puella Rōmāna. Quae māter Iūliae est? Aemilia māter Iūliae est. Pater Iūliae est Iūlius. Iūlia filia Iūliī est. \navnote{quī? quī sunt filii? Iūlia-que filiae-que}Quī sunt filiī Iūliī? Filii Iālir sunt Mārcus et Quīntus. Mārcus, Qulntus Iūliaque sunt trēs liberi. Liberi sunt filii filiaeque. Mārcus et Quīntus et Iūlia sunt liberi Iūliī et Aemiliae. In familiā Iūliī sunt trēs liberi: duo filii et ūna filia.

Estne Mēdus fīlius Iūliī? Mēdus fīlius Iūliī nōn est, Mēdus est servus Iūliī. Iālius dominus Mēdī est. Iūlius dominus servī est. Dāvus quoque servus est. Mēdus et Dāvus duo servī sunt. Iūlius est dominus Mēdī et Dāvī. Iūlius dominus servōrum est et pater līberōrum.

Estne Dēlia fīlia Aemiliae? Dēlia nōn est fīlia Aemi- 30 liae, Dēlia ancilla Aemiliae est. Aemilia domina Dēliae est. Aemilia domina ancillae est. Syra quoque ancilla \navnote{duo duae duo: duo servī duae ancillae duo oppida cuius? Iūli?, Aemiliae}est. Dēlia et Syra duae ancillae sunt. Aemilia domina ancillārum est.

Cuius servus est Dāvus? Dāvus servus Iūlit est. Cuius 35 ancilla est Syra? Syra est ancilla Aemiliae.

Quot līberi sunt in familiā? In familiā Iūliī sunt trēs \navnote{quot fīliī? quot fīliae? quot oppida?}līberī. Quot filiī et quot filiae? Duo filiī et ūna filia. Quot servī sunt in familiā? In familiā sunt centum servī. In familiā Iūliī sunt multī servī, paucī līberi. Iūlius est 40 dominus multōrum servōrum.

'Duo' et 'trēs' numeri sunt. 'Centum' quoque numerus est. Numerus servōrum est centum. Numerus līberōrum est trēs. Centum est magnus numerus. Trēs parvus numerus est. Numerus servōrum est magnus. Numerus līberōrum parvus est. In familiā Iūliī magnus \navnote{magnus numerus servōrum = multī servī parvus numerus llberōrum = paucīllberi}numerus servōrum, parvus numerus līberōrum est.

Mēdus servus Graecus est. Dēlia est ancilla Graeca. \navnote{multae-que}In familiā Iūliī sunt multī servī Graecī multaeque ancilso lae Graecae. Estne Aemilia femina Graeca? Aemilia nōn est femina Graeca, sed Rōmāna. Iūlius nōn vir Graecus, sed Rōmānus est.

\navnote{uēs tria: ues liberi tuēs litterae tria oppida magnus numerus oppidōrum = multa oppida magnus numerus fluviōrum = multī fluviī magnī-ne}Sparta oppidum Graecum est. Sparta, Delphī Tūscu- lumque tria oppida sunt: duo oppida Graeca et ūnum oppidum Rōmānum. In Graeciā et in Italiā magnus numerus oppidōrum est. In Galliā est magnus numerus fluviōrum. Fluviī Galliae magnī sunt. Magnine sunt fluvii Africae? In Āfricā ūnus fluvius magnus est: Nīlus; cēterī fluviī Āfricae parvī sunt. Suntne magnae īnsulae 60 Graecae? Crēta et Euboea duae īnsulae magnae sunt; \navnote{Cornēlius}cēterae īnsulae Graecae sunt parvae. --- HI!

Quis est Cornēlius? Cornēlius dominus Rōmānus est. Iūlius et Cornēlius duo dominī Rōmānī sunt. Mēdus nōn est servus Cornēliī. Mēdus servus Iūliī est. 65

Cornēlius: ``Cuius servus est Mēdus?''

Iūlius: ``Mēdus servus meus est.''

Cornēlius: ``Estne Dāvus servus tuus?''

Iūlius: ``Dāvus quoque servus meus est. Servī mel

Cornēlius: ``Estne Dēlia ancilla tua?''

Iūlius: ``Dēlia est ancilla mea, et Syra quoque ancilla mea est. Ancillae meae sunt Dēlia et Syra et cēterae multae. Familia mea magna est.''

Cornēlius: ``Quot servī sunt in familiā tuā?''

Iūlius: ``In familiā meā sunt centum servī.''

Cornēlius: ``Quid?''

Iūlius: ``Numerus servōrum meōrum est centum.''

Cornēlius: ``Centum servī! Magnus est numerus servōrum tuōrum!''

liber

\booksection{Lībertus Latīnus}

\navnote{ūnus liber duo libri}Ecce duo libri Latīnī: liber antīquus et liber novus. \navnote{antiquus -a -um * dre, y pāgina}lingva latina

est primus liber tuus Latīnus. Titulus libri tuī est 'lingva latina'.

Liber tuus nōn antīquus, sed novus est.

In lingva latina

sunt multae pāginae et multa capi- 85 tula: capitulum primum, secundum, tertium, cētera. 'imperivm romanvm' est titulus capitulī primi. Titulus capitulī secundi est 'familia romana'. In capitulō secundō sunt sex pāginae. In pāginā prīmā capitulī secundī multa vocābula nova sunt: vir, femina, puer, pu- 90 ella, familia, cētera. Numerus vocābulōrum Latīnōrum magnus est!

\booksection{Grammatica Latīna}

Masculīnum, fēmininum, neutrum 95 \navnote{masculīnum (m) \contr{} masculus = vir fēminīnum (f) \contr{} fēmina}[A] 'Servus' est vocābulum masculīnum. [B] 'Ancilla' est vocābulum fēminīnum. [C] 'Oppidum' est vocābulum neutrum.

Exempla: [A] Vocābula masculīna: fīlius, dominus, puer, vir; fluvius, 100 \navnote{-us}ōceanus, numerus, liber, titulus. Masculīnum: -us (-r). [B] Vocābula feminīna: fēmina, puella, fīlia, domina; īnsula, prōvincia, littera, familia, pāgina. Femininum: -a. \navnote{-um}[C] Vocābula neutra: oppidum, imperium, vocābulum, capitulum, exemplum, pēnsum. Neutrum: -um. Genetīznis Genetīvus \navnote{genetīvus (gen)}[A] Masculīnum:

Iūlius dominus serv? (Dāvī) est.

\navnote{servī servōrum}Iūlius dominus servōrum (Dāvt et Mēd?) est.

'Servī! genetīvus

est. 'Servōrum' quoque genetīvus est. 110 Serv?' genetīvus singulāris est. 'Servōrum' est genetīvus plūrālis. Genetīvus: singulāris -ī, plūrālis -ōrum. [B] Fēminīnum:

Aemilia domina ancillae (Syrae) est.

\navnote{ancillae ancillārum}Aemilia domina ancillārum (Syrae et Dēlia) est.

--- 'Ancillae' genetīvus singulāris est. fAncillarum' est genetīvus plūrālis. Genetīvus: singulāris -ae, plūrālis -ārum. [C] Neutrum:

d est prīma littera vocābulf dominus'.

\navnote{vocabuli vocābulōrum}Numerus vocāabulōrum magnus est.

120 --- 'Vocābul? genetīvus singulāris est. 'Vocābulōrum' est ge netivus plūrālis. Genetīvus: singulārisnetlvus plūrālis. Genetlvus: singulāris -f, plūrālis -ōrum.

\booksection{Pēnsum A}

Mārcus fīli- Iūliī est. Iūlia fīli- Iūliī est. Iūlius est vir Rōmān-. Aemilia femin- Rōmān- est. Iūlius domin-, Aemilia domin-

est. Mēdus serv- Graec- est, Dēlia est ancillGraec-. Sparta oppid- Graec- est.

Iūlius pater Mārc- est. Mārcus est filius Iūli- et Aemili-. Mēdus servus Iūli- est: Iūlius est dominus serv-. Iulius dominus Mēd- et Dāv- est: Iūlius dominus serv- est. Numerus serv- magnus est. Dēlia est ancilla Aemili-: Aemilia do-mina ancill- est. Aemilia domina Dēli- et Syr- est: Aemilia domina ancill- est. In familiā Iūli- est magnus numerus serv- et ancill-. Aemilia māter Mārc- et Quīnt- et Iūli- est. Mārcus, Quīntus Iūliaque sunt līberi Iūli- et Aemili-. Numerus līber- est trēs. Numerus serv- est centum.

In pāginā prīmā capitul- secund- multa vocābula nova sunt. Numerus capitul- nōn parvus est.

\booksection{Pēnsum B}

Mārcus --- Rōmānus est. Iūlius --- Rōmānus est. Aemilia est --- Rōmāna. Iūlius est --- Mārcī et Quīntī et Iūliae. In --- Iūliī sunt trēs ---: duo --- et ūna ---. --- līberōrum est Aemilia.

--- est Dāvus? Dāvus est --- Iūliī. Iūlius --- Dāvī est. --- est Syra? Syra --- Aemiliae est. Aemilia est --- Syrae.

Cornēlius: ``--- servī sunt in familiā tuā?'' Iūlius: aIn familiā --- sunt --- (c) servī.`` Cornēlius: Familia --- magna est!''

'lingva latina'

est titulus --- tuī Latīnī.

\booksection{Pēnsum C}

Quis est Qulntus? Qul sunt Mēdus et Dāvus? Mārcusne quoque servus Iūliī est? Cuius filia est Iūlia? Quot līberi sunt in familiā Iūliī? Quot servī in familiā sunt? Num Syra domina est? Quae est domina ancillārum? Estne Cornēlius vir Graecus? Num 'puella' vocābulum masculīnum est?

\booksection{Vocābula nova}
\noindent\small = et Mēdus Iūlius?, Aemilia, Iūlius, māter, pater, Mārcus Qulntus Iūlia, filius, filia, nova, vir, femina, puer, puella, familia, liberi, servus, dominus, ancilla, domina, liber, titulus, pāgina, antiquus, novus, cēteri, meus, tuus, centum, duae, tria, -que, quis?, quae?, quī?, cuius?, quot?, masculīnum, feminīnum, neutrum, genetīvus.

\bookchapter{Cap. III}{Puer improbus}

\booksection{Scaena prīma}

I \navnote{scaena}Persōnae: Iulia, Mārcus, Quīntus. Iūlia cantat: ``Lalla.'' Iūlia laeta est. Mārcus: ``St!'' Mārcus laetus nōn est. Iūlia cantat: ``Lalla, lalla.'' Mārcus: ``Ssst!'' Mārcus īrātus est. Iūlia cantat: ``Lalla, lalla, lalla.'' Mārcus Iūliam pulsat. \navnote{Mārcus Iūliam pulsat}Iam Iūlia nōn cantat, sed plōrat: ``Uhuhū!'' Mārcus ridet: ``Hahahae!''

\navnote{Iūlia plōrat Mārcus rīdet Quīntus Mārcum videt Mārcus Quīntum nōn videt}Quīntus Mārcum videt. Mārcus nōn videt Quīntum. Quīntus: ``Quid? Mārcus puellam pulsat --- et ridet!'' Quīntus īrātus est et Mārcum pulsat! Iam nōn ridet Mārcus. Mārcus īrātus pulsat Quīntum. Iūlia: ``Ubi est māter?'' Iūlia Aemiliam nōn videt. Iulia Aemiliam vocat: ``Māter! Mārcus Quīntum pulsat!'' Mārcus (īrātus): ``St!'' Mārcus Iūliam pulsat. \navnote{i e EN ES ``d d Quīntus Mārcum pulsat mamma = māter}Iūlia plōrat et Aemiliam vocat: ``Mamma! Mam-ma! Mārcus mē pulsat!''

\booksection{Scaena secunda}

\navnote{Mārcus Quīntum pulsat Iūlia Aemiliam vocat Aemilia venit}Persōnae: Aemilia, Iūlia, Mārcus, Quīntus. Aemilia interrogat: ``Quis mē vocat?'' \navnote{interrogat \contr{} \contr{} respondet}Quīntus respondet: ``Iūlia tē vocat.'' Aemilia Quīntum interrogat: ``Cūr Iūlia plōrat?'' \navnote{tam: Iūliam}Quīntus respondet: ``Iūlia plōrat, quia Mārcus eam pulsat.'' Aemilia: ``Quid? Puer parvam puellam pulsat? Fu! Cūr Mārcus Iūliam pulsat?'' Quīntus: Quia Iūlia cantat.``. Aemilia: ''Ō Iūlia, mea parva filia! Mārcus puer prō!bus nōn est; Mārcus est puer improbus!`` Quīntus: ''Iulia puella proba est.`` 35 --- Aemilia Quīntum interrogat: ''Ubi est Iūlius? Cūr \navnote{cūr Iūlius nōn venit? M PATER QN U d LL}nōn venit?`` Aemilia Iūlium nōn videt. Respondet Mārcus: ''Pater dormit.`` Quīntus: ''Māter nōn tē, sed mē interrogat!`` Aemilia: ''St, puerī! Ubi est pater?`` 40 --- Quīntus: ''Pater nōn hīc est, sed Mārcus hic est.`` Quīntus Iūlium vocat: ''Pater! Pa-ter!``

Iūlius dormit

\navnote{Quīntus Iūlium vocat ne-que = et nōn eum: Quīntum Iūlius dormit}Iulius Quīntum nōn audit neque venit. Cūr Iūlius Quīntum nōn audit? Iūlius eum nōn audit, quia dormit. Mārcus: ``Hahae! Pater dormit neque tē audit.'' Aemilia: ``Fū, puer!'' Aemilia īrāta est. Māter fīlium

\navnote{verberat = pulsat et pulsat (*tux-tax?)}Mārcus plōrat: ``Uhuhū!'' Iūlius eum audit. Iam nōn dormit pater. 007m e Su Pēs f cA Boon NI 224572) HL 28. 87 NH

\booksection{Scaena tertia}

Persōnae: Iūlius, Aemilia, Iūlia, Mārcus, Quīntus. Quīntus: ``Pater venit.'' Aemilia Quīntum nōn audit, quia Mārcus plōrat. \navnote{eum-que = et eum}Iūlius Quīntum videt eumque interrogat: ``Cūr Mārcus plōrat?'' Quīntus respondet: ``Mārcus plōrat, quia māter eum 55 \navnote{māter Mārcum}verberat.`` Iūlius: ''Sed cūr māter Mārcum verberat?`` Quīntus: ''Mārcum verberat, quia puer improbus est. \navnote{verberat, quia Mārcus puer improbus est}Mārcus parvam puellam pulsat!`` Iūlia: ''Mamma! Pater hic est.`` Aemilia Iūlium videt. Aemilia: ''Tuus Mārcus fīlius improbus est!`` Iūlius: ''Fū, puer! Puer probus nōn pulsat puellam. Puer quī parvam puellam pulsat improbus est!`` Iūlius īrātus puerum improbum verberat: tuxtax,

Mārcus plōrat. Quīntus laetus est et ridet. Iūlia laeta \navnote{lūlia nōn laeta est}nōn est neque rīdet. Cūr nōn laeta est Iulia? Nōn laeta est, quia Mārcus plōrat. Iūlia est puella proba! f s / Puer ridet. Puella plōrat. Quis est puer qul ridet? \navnote{quī quae puer quī ridet puella quae plōrat}Puer quī rīdet est Mārcus. Quae est puella quae plōrat? Puella quae plōrat est Iulia.

\navnote{quem puer quem Aemilia verberat puella quam Mārcus pulsat}Mārcus, quī puellam pulsat, puer improbus est. Pu-ella quam Mārcus pulsat est Iūlia. Iūlia Aemiliam vocat. Aemilia, quam Iūlia vocat, māter līberōrum est. Aemilia puerum verberat. Puer quem Aemilia verberat est Mārcus.

Quem vocat Quīntus? Quīntus Iūlium vocat. Iulius, \navnote{quem? lūlium Quīntum}quem Quīntus vocat, pater līberōrum est. Iūlius Quīntum nōn audit. Quem audit Iūlius? Iūlius Mārcum audit. Puer quem Iūlius audit est Mārcus.

Puella quae cantat laeta est. Puella quae plōrat nōn est laeta. Puer quī puellam pulsat improbus est!

\booksection{Grammatica Latīna}

Nōminātīvus et accūsātīvns \navnote{nōminātīvus (nōm) accūsātīvus (acc) Mārcus Mārcum Quīntum}[A] Masculīnum.

Mārcus ridet. Quīntus Mārcum pulsat.

Mārcus Quīntum pulsat. Quīntus plōrat.

Mārcus' nōminātīvus

est. 'Mārcum' accūsātīvus

est. \navnote{-ns -um}'Quīntum' est accūsātīvus, Quīntus! nōminātīvus. Nōminātivus: -us (-r). Accūsātīvus: -um.

Exempla: Iūlius, Tūlium; filius, filium; puer, puerum; eum. [B] Femininum.

Iūlia cantat. Mārcus Iūliam pulsat.

Iūlia Aemiliam vocat. Aemilia venit.

'Iūlia, Aemilia* nōminātīvus est. '``Iūliam, Aemiliam' accūsātīvus est. Nōminātīvus: -a. Accūsātīvus: -am.

\navnote{-am}Exempla: puella, puellam; parva, parvam; eam. Verbum

Iūlia cantat. Mārcus rīdet. Iūlius dormit.

'Cantat' verbum est. 'Cantat', 'ridet', 'dormit' tria verba sunt.

Exempla: cantat, pulsat, plōrat, vocat, interrogat, verberat (-at) ridet, videt, respondet (-et) dormit, venit, audit (-it). Nōminātīvus: Accūsātīvus: Mārcus Quīntus Iūlia Aemilia Quīntum ``Quis Quīntus ''Iūlia Mārcus Aemilia Mārcus Iūlius Iūlius ``Puer probus parvam puellam Iūlius Irātus

\booksection{Pēnsum A}

Cūr Mārc- Iūliam pulsat? Mārcus Iūli- pulsat, quia Iūlicantat. Iūlia plōr- quia Mārcus e- pulsat. Iūlia: ``Mamma! Mārcus --- pulsat.'' Aemilia puell- aud- et ven-. Māter Quīnt- videt et e- interrog-: ``Quis mē voc-?'' Quīnt- respond-: ``Iūlia --- vocat.''

Iūlius dorm-. Quīntus Iūli- voc-: ``Pater!'' Mārcus rīd-, quia Iūli- nōn venit. Aemilia Mārc- verber-. Iūlius ven-, quia Mārc- plōrat. Iūlius Aemili- et Mārc- et Quīnt- et Iūlividet. Iūlius: ``Puer quī parv- puell- pulsat improbus est.'' Iūlius puer- improb- verberat. Quem Iūli- verberat? Puer qu- Iūlius verberat est Mārcus. Mārcus plōr-. Puer qu- plōrat laet- nōn est. Puella qu- cantat laet- est.

\booksection{Pēnsum B}

Puella ---: ``Lalla.'' Puella --- cantat est Iūlia. Iūlia --- est.! Puer improbus puellam ---. Puella ---: ``Uhuhū!'' Puer ---:: ``Hahahae!'' Puer --- ridet est Mārcus. Iūlia Aemiliam ---: ---: ``Mamma!'' Aemilia ---, et Quīntum ---: ``Cūr Iūlia plōrat?'' Qulntus ---: ``Iūlia plōrat, --- Mārcus eam pulsat.'' Aemilia: ``Mārcus puer --- nōn est, puer --- est! Ubi est pater?'' Aemilia 1 Iūlium nōn ---. Quīntus: ``Pater nōn --- est.'' Qulntus Iūlium! ---: ``Pater!'' Iūlius Quīntum nōn ---. --- Iūlius Quīntum nōn ' audit? Iūlius eum nōn audit, quia ---. Mārcus plōrat, --- Ae-milia eum verberat. Iūlius Mārcum audit; --- Iūlius nōn dor!mit. --- Aemilia verberat? Aemilia Mārcum ---. Puer --- Ae!milia verberat improbus est. Iūlia laeta nōn est --- ridet. I

\booksection{Pēnsum C}

Quis Iūliam pulsat? Cūr Iūlia plōrat? Quīntusne quoque Iūliam pulsat? r Quem Quīntus pulsat? Cūr Aemilia venit? Quis Iūlium vocat? Cūr Iūlius Quīntum nōn audit?: Quem audit Iūlius? Cūr Mārcus plōrat?: Rīdetne Iūlia? Num 'Mārcus' accūsātīvus est? Num 'Iūliam' nōminātīvus est? Quid est 'dormit?

\booksection{Vocābula nova}
\noindent\small Aemilia, Iūliam, Aemiliam, cantat, ridet, dormit, -at, -et, -it, -um, -am, us, eum, eam, a, mē, scaena, persōna, mamma, laetus, īrātus, probus, improbus, pulsat, plōrat, rīdet, videt, vocat, venit, interrogat, respondet, audit, verberat, tē, neque, iam, cūr?, quia, hīc, quī, quae, quem, quam, nōminātīvus, accūsātīvus, verbum.

\bookchapter{Cap. IV}{Dominus et servī}

\booksection{Scaena prīma}

Persōnae: Iūlius, Aemilia, Mēdus. \navnote{pecūnia sacculus: in sacculō eius:: Iūliī}Sacculus Iūliī nōn parvus est. In sacculō eius est pecūnia. Iūlius pecūniam in sacculō habet.

Aemilia sacculum videt Iūliumque interrogat: ``Quot s nummī sunt in sacculō tuō?''

\navnote{nummus}Iūlius respondet: ``Centum.''

Aemilia: ``Num hic centum nummī sunt?''

Iūlius pecūniam numerat: ``Ūnus, duo, trēs, quattuor,

quīnque,

sex, septem,

octō, novem,

decem. 10 Quid? Decem tantum?``

Iūlius rūrsus pecūniam numerat: ``Ūnus, duo, trēs,

est centum, sed decem tantum.

Iūlius: ``Quid? In sacculō meō nōn centum, sed tan- 15 tum decem nummī sunt! Ubi sunt cēteri nummī? Servī meī ubi sunt?'' Mēdus: ``Servus tuus Mēdus hic est.'' Iūlius servum suum Mēdum videt, Dāvum nōn videt. Mēdus adest. Dāvus nōn adest, sed abest. Iūlius et \navnote{ad-est = hic est ab-est --- ad-est ad-sunt = hic sunt}Aemilia et Mēdus adsunt. Dāvus cēterique servī ab-sunt. Iulius: ``Quid? Ūnus servus tantum adest! Ubi est Dāvus? Dāvum vocā!'' Mēdus Dāvum vocat: ``Dāve!'' sed Dāvus Mēdum nōn audit neque venit. Mēdus rūrsus Dāvum vocat: ``Da-ā-ve! Venī!'' Dāvus venit. Iam duo servī adsunt.

\booksection{Scaena secunda}

Persōnae: Iūlius, Aemilia, Mēdus, Dāvus. Daāvus, quī dominum suum nōn videt, Mēdum interrogat: ``Quid est, Mēde?'' Mēdus: ``St! Dominus adest. Salūtā dominum!'' Servus dominum salūtat: ``Salvē, domine!'' \navnote{servus dominum salūtat NS}Dominus servum salūtat: ``Salvē, serve!'' Dāvus: ``Quid est, domine?'' Iūlius: St! Tacē, serve! Tacē et audi!`` Servus tacet. Iūlius: ''In sacculō meō sunt decem tantum nummī. Ubi sunt cēterī nummī meī?`` Dāvus tacet, neque respondet. Aemilia: ''Respondē, Dāve! Dominus tē interrogat.`` Dāvus respondet: ''Pecūnia tua hic nōn est. Interrogā Mēdum!`` Iūlius Mēdum interrogat: ''Ubi sunt nummī meī, \navnote{nūllum:: nōn ūnum verbum = vocābulum}Mēde?`` Mēdus nūllum verbum respondet. Iūlius rūrsus eum interrogat: ''Ubi est pecūnia mea? Respondē, serve!`` Mēdus Dāvum accūsat: ''Pecūnia tua in sacculō Dāvī est. Dāvus pecūniam tuam habet.`` Aemilia: ''Audī, Dāve! Mēdus tē accūsat.`` Dāvus: ''Quem Mēdus accūsat? mē?`` Iūlius: ''Tacē, Mēde! Servus quī servum accūsat improbus est!`` Mēdus tacet. Iūlius Dāvum nōn accūsat, sed interrogat eum: ''Estne pecūnia mea in sacculō tuō, Dāve?`` Dāvus: ''In sacculō meō nōn est pecūnia tua, do-mine.`` \navnote{Dāvus sacculum in mēnsā pōnit ecce sacculus = vidē: hīc est sacculus}Iūlius: ``Ubi est sacculus tuus? Dāvus: ''Hīc est. Ecce sacculus meus.`` Iūlius: ''Sacculum tuum in mēnsā pōne!`` \navnote{baculum eius:: Dāvī}Dāvus sacculum suum in mēnsā pōnit. Iam sacculus eius in mēnsā est. Iūlius baculum suum in mēnsā pōnit. Baculum dominī in mēnsā est. Dāvus: ``Vidē: in sacculō meō nūlla pecūnia est.'' Iūlius nūllam pecūniam videt in sacculō. In sacculō 65 \navnote{baculum et sacculus}Dāvī nūllī nummī sunt. Sacculus eius vacuus est. Dāvus pecūniam dominī nōn habet. Iūlius: ``Ō! Dāvus bonus servus est: pecūniam meam \navnote{bonus -a -um = probus}nōn habet. Ecce nummus tuus, Dāve!`` Iūlius ūnum nummum pōnit in sacculō Dāvī. Iam sacculus Dāvī nōn est vacuus: in sacculō dus est ūnus nummus. Dāvus laetus est. \navnote{sūmit \contr{} pōnit discēdit --- \contr{} venit quī quae quod: puer quī... puella quae... baculum quod.}Iūlius: ``Sūme sacculum tuum et discēde, bone serve!'' Dāvus sacculum suum sūmit et discēdit. Mēdus baculum, quod in mēnsā est, videt. Mēdus quoque discēdit! Cūr discēdit Mēdus? Mēdus discēdit, quia is pecūniam dominī in sacculō suō habet! Dāvus et Mēdus absunt.

\booksection{Scaena tertia}

Persōnae: lūlius, Aemilia. Iūlius: ``Dāvus bonus servus est. Is nōn habet pecūniam meam. --- Sed ubi est pecūnia mea, Mēde? Quis pecūniam meam habet?'' Mēdus nōn respondet. Iūlius: ``Ubi est Mēdus? Cūr nōn respondet?'' Aemilia: ``Mēdus nōn respondet, quia abest. Nūllus servus adest.'' Iulius Mēdum vocat: ``Mēde! Veni!'' sed Mēdus, quī abest, eum nōn audit neque venit. Iūlius rūrsus vocat: ``Mē-de! Veni, improbe serve!'' 90 Mēdus nōn venit. Iūlius: ``Cūr nōn venit Mēdus?'' Aemilia: ``Mēdus nōn venit, quia is habet pecūniam tuam! Eius sacculus nōn est vacuus!'' Aemilia ridet. Iūlius īrātus est --- is nōn rīdet! Iūlius: ``Ubi est baculum meum?'' Iūlius baculum, quod in mēnsā est, nōn videt. Aemilia: Ecce baculum in mēnsā.`` Iūlius baculum suum sūmit et discēdit.

\booksection{Grammatica Latīna}

\navnote{vocātīvus (voc) \contr{} vocat}Vocātīvus Mēdus Dāvum vocat: ``Dāve/'' 'Dāve' vocātīvus

est. Vocātīvus: -e (nōminātīvus: -n\$). Exempla: Mēde, domine, serve, improbe. Servus: ``Salvē, domiiu/'' Iūlius Mēdum vocat: ``Mēde/ Venī, improb strve!''

Imperāttvus et indicātīvus A Dominus: ``Vocā Dāvum!'' Servus Dāvum vocat. Dominus: ``Tac \contr{} et audi/'' Servus tacet et audit. Dominus: ``Discēd serve!'' Servus discēdit. Dominus imperat. Servus pāret. Vocā'imperātīvusest. 'Vocar'indicātīvusest. Tac, 'audf,

'discēd' est 'Vocat' indicātīvus est. Imperātīvus: -ā, -ē, -e, -ī. Indicātīvus: -at, -et, -ii, -iv Exempla: [1] salūtā, salūtat; [2] responde,

respondet; [3] 115 sūme, sūmit; [4] vent, venit. Mēdus: ``Salūta dominum!'' Dāvus dominum salūtar. Iūlius: Responde, serve!`` Servus respondet. Iūlius imperat: ''Sūm sacculum tuum et discēde/`` Dāvus sacculum suum sūmit et discēdit. Iūlius: ''Mēde! Veni/`` Mēdus nōn venit. Iūlius imperat. Dāvus pāret; Mēdus nōn pāret.

\booksection{Pēnsum A}

Mēdus ad-.

Dāvus ab-.

Iūlius imper-:

``Voc- Dāvum, Mēd-!u Mēdus Dāvum voc-: ''Dāv-! Ven-!`` Dāvus venneque Iūlium vid-. Mēdus: ''Salūt- dominum!`` Dāvus dominum salūt-: ''Salvē, domin-! Quid est?`` Dominus: ''Tac-, serv-! Nummī mel ubi sunt?`` Servus tac- neque respond-. sacculus Iūlius: ''Respond-!`` Dāvus: ''Interrog- Mēdum!`` Iūlius Mēpecūnia dum interrog-: ''Ubi est pecūnia mea, Mēd-?`` Mēdus: Da- tvus pecūniam tuam hab-.'' Iūlius: Tōn- sacculum tuum in 1 mēnsā, Dāv-!`` Dāvus pār-: sacculum suum in mēnsā pōn-. Dāvus: aVid-, domin-: sacculus meus vacuus est.'' Iūlius: Sūm- sacculum tuum et discēd-, bon- serv-!`` Dāvus saccu- \contr{} lum suum sūm- et discēd-. I

\booksection{Pēnsum B}

In sacculō Iūliī--- est. Iūlius pecūniam ---: ``Ūnus, duo, trēs, nummī sunt.

Dāvus dominum ---: ``Salvē, domine!'' Iūlius ---: ``Pōne: sacculum tuum in ---l!'' Dāvus sacculum --- in mēnsā ---.! Sacculus Dāvī--- --- est, in sacculō --- [: Dāvī]--- --- pecūnia est. \contr{} Dāvus sacculum suum --- et discēdit. 1

Iūlius: ``Mēde! Veni!'' Mēdus nōn venit, quia --- [: Mēdus]: pecūniam Iūliī ---. Iūlius baculum, --- in mēnsāā est, sūmit et ---. 1

Dominus imperat, bonus servus ---. t

\booksection{Pēnsum C}

Quot nummī sunt in sacculō Iūliī? Adestne Dāvus in scaenā primā? Quis Dāvum vocat? Suntne nummī Iūliī in sacculō Dāvī? Quid Iūlius pōnit in sacculō Dāvī? Quot nummī iam in sacculō Iūliī sunt? Estne vacuus sacculus Mēdī? Cūr Mēdus discēdit? Quem Iūlius vocat? Cūr Mēdus Iūlium nōn audit?

\booksection{Vocābula nova}
\noindent\small numerat \contr{} numerus I, ūnus, I, duo, II, III, quattuor, IV, quīnque, V, sex, VI, VII, septem vīll, octō, VIII, novem, X, decem, vocā, vocat, tacē, tacet, discēde, discēdit, audi, audit, imperat \contr{} \contr{} pāret, imperātīvus (imp), -at, -et, -it, mēnsa, baculum, vacuus, bonus, septem, habet, numerat, adest, abest, salūlat, accūsat, pōnit, sūmit, imperat, pāret, nūllus, eius, suus, is, quod, rūrsus, tantum, salvē, vocātīvus, imperātīvus, indicātīvus.

\bookchapter{Cap. V}{Vīlla et hortus}

hortus

Ecce villa et hortus Iūliī. Iūlius in magnā villā habitat. Pater et māter et trēs līberi in villa habitant. Iūlius et Aemilia trēs līberōs habent: duōs fīliōs et ūnam fīliam --- nōn duās filiās.

In villā multī servī habitant. Dominus eōrum est Iū- 5 \navnote{eōrum: servōrum}lius: is multōs servōs habet. Ancillae quoque multae in vīllā habitant. Domina eārum est Aemilia: ea multās \navnote{tārum: ancillārum}ancillās habet.

Iulius in villā sua habitat cum magnā familiā. Pater et \navnote{in hortō, in vīlla cum Mārcō, cum Iūliā}māter habitant cum Mārcō et Quīntō et Iūliā. Iūlius et 10 Aemilia in villā habitant cum līberis et servīs et ancillīs.

\navnote{in villis cum servīs, cum ancillīs AS CN d:}Vīlla Iūliī in magnō hortō est. In Italiā sunt multae vīllae cum magnīs hortīs. In hortīs sunt rosae et līlia. Iūlius multās rosās et multa līlia in hortō suō habet. \navnote{rosa līlium eō: hortō pulcher -chra -chrum: hortus pulcher (m) rosa pulchra (f) līlium pulchrum (/) eius: Syrae}Hortus Iūliī pulcher est, quia in eō sunt multae et pul- 15 chrae rosae llliaque.

Aemilia femina pulchra est. Syra nōn est femina pulchra, neque pulcher est nāsus eius, sed foedus est. Syra, quae bona ancilla est, nāsum magnum et foedum habet. Iūlius est vir Aemiliae, feminae pulchrae. Iūlius Aemiliam amat, quia ea pulchra et bona femina est. Aemilia Iūlium virum suum amat et cum eō habitat. Pater et māter līberōs suōs amant. Iūlius nōn sōlus, sed cum Aemiliā et cum magnā familiā in villa habitat.

In villā sunt duo ōstia: ōstium magnum et ōstium parvum. Villa duo ōstia et multās fenestrās habet. AQ EEESSEBNVEEBYEEN/A/ A, CBENANVER UR / A BEEBNNEBNUNEBVVVEREV/EEV/AEP, INL n ll

\navnote{ōstium et fenestra}\navnote{ātrium}In villa Iūlit magnum ātrium est cum impluviō. Quid est in impluviō? In eō est aqua. In ātriō nūllae fenestrae sunt.

\navnote{peristylum}NEN. - ---4 AN X IZ Bēs ---L ZZ -. T7. A. ez DOu------RVVVV. ``a z --- --- c ------7--- il T? 4 xm LU] EET xu ee]:u SN M] b Me - iem BH Etiam peristylum magnum et pulchrum in villa est. \navnote{etiam = quoque: etiam peristylum = p. quoque}*``Peristylum'

est vocābulum Graecum. In villis Graecis et Rōmānīs magna et pulchra peristyla sunt. Estne impluvium in peristylō? Id nōn in peristylō, sed in ātriō est. In peristylō parvus hortus est.

In vīllā sunt multa cubicula. Qulntus in cubiculō parvō dormit. Estne magnum cubiculum Mārcī? Id quoque parvum est. Iūlius et Aemilia in cubiculō magnō dormiunt. Ubi dormiunt servī? II quoque in cubiculīs dormiunt.

Suntne magna eōrum cubicula? Ea nōn magna sunt, et multī servī in ūnō cubiculō dormiunt. Etiam ancillae multae in ūnō cubiculō dormiunt, neque \navnote{iiae: ancillac}eae magna cubicula habent. ---

Aemilia in peristylō est. Estne sōla? Aemilia sōla nōn est: liberi cum eā in peristylō adsunt. Iūlius abest. Ae-milia sine virō suō Iūliō in villā est. Ubi est Iūlius? In oppidō Tūsculo est sine Aemiliā, sed cum servīs quattuor. II Aemilia cum Mārcō, Quīntō Iūliāque in peristylō est. \navnote{Iūlia cum Aemiliā est}Iūlia rosās pulchrās in hortō videt et ab Aemiliā discēdit. Iam ea cum Aemiliā nōn est. Aemilia eam nōn videt. Puella in hortō est. \navnote{Iūlia ab Aemiliā discēdit}Aemilia imperat: ``Mārce et Quīnte! Iūliam vocāte!'' Mārcus et Qulntus Iūliam vocant: ``Iūlia! Veni!'' sed Iūlia eōs nōn audit neque venit. \navnote{eōs: puerōs hortus Iūlia}Iūlia puerōs vocat: ``Mārce et Quīnte! Venīte! HIc multae rosae sunt.'' Puerī Iūliam audiunt, neque ii ab Aemiliā discēdunt. \navnote{Iūlia in hortō est Z}Quīntus: ``Carpe rosās, Iūlia!''

\navnote{Iūlia ex hortō venit z! d A d}\navnote{Iūlia rosās carpit}Iulia: ``Vidē, māter! Vidēte, pueri!

Vidēte rosās meās!`` Iūlia laeta est, rosae eam dēlectant. Aemilia: ''Ecce puella pulchra cum rosīs pulchrīs!`` Verba Aemiliae Iūliam dēlectant. Mārcus: ''Rosae pulchrae sunt; puella sine rosīs pulchra nōn est!`` Verba Mārcī Iūliam nōn dēlectant! Aemilia (īrāta): ''Tacē, puer improbe!

Iūlia puella pulchra est --- cum rosīs et sine rosīs.`` mm------ --- À--- Iūlia: Audīte, Mārce et Quīnte!'' Mārcus: ``Māter nōn videt nāsum tuum foedum!'' Mārcus et Qulntus rident: ``Hahahae!'' Iūlia: ``Audī, mamma: pueri etiam mē rident!'' Iūlia plōrat et cum ūnā rosā ab iīs discēdit. Aemilia: ``Tacēte, pueri improbī! Nāsus Iūliae foedus nōn est. Discēdite ex peristylō! Sūmite cēterās rosās eāsque in aquā pōnite!'' \navnote{eās: rosās its: īosīs}Pueri cēterās quattuor rosās sūmunt et cum iīs discē- 75 dunt. Aemilia, quae iam sōla est in peristylō, ancillās vocat:. ``Dēlia et Syra! Venīte!'' Dēlia et Syra ex ātriō veniunt. Aemilia eās interrogat: \navnote{tās: ancillās}``Suntne pueri in ātriō?'' Dēlia respondet: In ātriō sunt.`` Aemilia: ''Quid agunt Mārcus et Quīntus?``

Hic domina et ancillae puerōs audiunt ex ātriō: Quīn- 85 tus plōrat et Mārcus rīdet. Aemilia: ``Quid iam agunt pueri? Age, Dēlia! discēde et interrogā eōs!'' Dēlia ab Aemiliā et Syrā discēdit. Aemilia Syram interrogat: ``Ubi est Dāvus?'' Syra respondet: In oppidō est cum dominō.`` Dēlia ex ātriō venit et dominam vocat: ''Veni, o do-mina! Venī!`` Aemilia: ''Quid est, Dēlia?`` Dēlia: ''Quīntus est in impluviō!`` 95 Aemilia: ''In impluviō? Quid agit puer in impluviō?`` Dēlia: ''Aquam pulsat et tē vocat.`` Aemilia: ''Quid agit Mārcus?`` Dēlia: ''Is rīdet, quia Quīntus in aquā est!`` Aemilia: ''Ō, puer improbus est Mārcus! Agite! Iū 100lium vocāte, ancillae!`` Syra: ''Sed dominus in oppidō est.`` Aemilia: ''Ō, iam rūrsus abest Iūlius!`` Dēlia: ''Age! Veni, domina, et Mārcum verberā!`` Quid agit domina? Domina Irāta cum ancillīs ex peri 105stylō discēdit.

\booksection{Grammatica Latīna}

Accūsātīvus [A] Masculīnum. Iūlius nōn ūnum filium, sed duōs filios habet \navnote{filios}'FTln/ro' est accūsātīvus singulāris (vidē cap. III). 'Fīlis' accūsātīvus plūrālis est. Accūsātīvus: singulāris -um (nōminātīvus -4s), plūrālis -ōs (nōminātīvus -t). [B] Fēminīnum. Iūlius nōn duās filiās, sed ūnam filiam habet. 15 \navnote{filiam filiās}Filiam' est accūsātīvus singulāris (vidē cap. im). Filiās' accūsāūvus plūrālis est. Accūsātīvus: singulāris -am (nōminātīvus -a), plūrālis -as (noōminātivus -ae). [C] Neutrum. Villa nōn ūnum cubiculum, sed multa cubicula habet. \navnote{cubiculum cubicula}--- Cubiculum' accūsātīvus singulāris est. Cubicula'

accūsātīvus plūrālis est. Accūsātīvus: singulāris -um (= nōminātīvus), plūrālis -a (= nōminātīvus). Ablātīvus \navnote{ablātīvus (abl)}[A] Masculīnum. In hortō Iūliī. In hortīs Italiae.

\navnote{hortō hortīs}Hortō' ablātīvus

singulāris est. 'Hortfs' ablātlvus plūrālis est. Ablātīvus: singulāris -ō, plūrālis -i5. [B] Fēminīnum. In vīllā Iūliī. In villis Rōmānīs.

\navnote{villā villis}'Vīllā' ablātīvus singulāris est. 'Villis' ablātīvus plūrālis est. Ablātīvus: singulāris -ā, plūrālis -īs. [C] Neutrum.

\navnote{oppidō oppidīs}In oppidō Tusculō. In oppidīs Graecis.

Oppidoō' ablātīvus singulāris est. *'Oppidīs' est ablātīvus plūrālis. Ablātīvus: singulāris -ō, plūrālis -īs.

In, ex, ab, cum, sine cum ablātīvō: in ātriō, in cubiculīs, ex hortō, ex Italiā, ab Aemiliā, ab oppidō, cum servō, cum liberis, sine pecūnia, sine rosīs.

Imperātīvus et indicātīvus

Dāvum vocā, serve!`` Servus Dāvum vocat.

``Iūliam vocant.

'Vocā' est imperātīvus singulāris, Vocat' indicātīvus singulāris (vidē cap. iv). 'Vocāte' imperātīvus plūrālis est, vocant' indicātīvus plūrālis. Imperātīvus: singulāris -, plūrālis -te. Indicātīvus: singulāris -t, plūrālis -nt.

Exempla: \navnote{sing. -āte -at}Salūtā dominum, serve! Salūtāte dominum, servī! Servus dominum salūtat.

Servī dominum salūtant.

Tacēte, pueri! Pueri tacenr. Tac \contr{}, puer! Puer

Discēde, serve! Discēdite, servī! Discēd, serve! Servus discēdit.

Veni, Iūlia! Iūlia venit. Venīte, pueri! Pueri veniunt. Vem, Iūlia! Iūlia vemY.

\booksection{Pēnsum A}

Iūlius et Aemilia in vīll- habit- cum līber- et serv-. Dominus hortus mult- serv- et mult- ancill- habet.

in peristyl- est cum Mārc- et Quīnt- et Iūli-. Iūlia nāsus mult- ros- in hort- vid- et ab Aemili- discēd-. Iam Aemilia ōstium \navnote{fenestra}puell- nōn vid-, neque pueri eam vid-. Aemilia: ``Mārce et ātrium Quīnte! Voc- Iūliam!'' Pueri Iūli- voc-: ``Iūlia! Ven-!'' et aqua Iūlia puer- voc-: ``Mārce et Quīnte, ven-!'' Iūlia puerōs nōn \navnote{peristylum}cubiculum aud-, sed pueri Iūli- aud-. Iūlia: ``Cūr pueri nōn ven-?'' Iūlia pulcher ex hort- ven- cum v ros- pulchr-. Iūlia: ``Vid- ros- meās, foedus māter! Vid-, pueri!'' Mārcus: ``Rosae pulchrae sunt, puella habitat sine ros- pulchra nōn est!'' Iūlia cum ūn- ros- discēd-. Puerl amat rid-. Aemilia: ``Tac-, puerī! Sūm- ros- et discēd-!'' Puerī dēlectat ros- sūm- et discēd-; in ātri- aqu- sūm- ex impluvi- et rosagit in aquā pōn-.

\booksection{Pēnsum B}

Iūlius in magnā --- Aemiliam ---, quia --- [: Aemilia] bona et --- femina est.

Aemilia in peristylō est --- līberīs suīs, sed --- virō suō. iī Iūlia --- Aemiliā discēdit; iam puella in --- est. Iūlia rosās --- eae et --- hortō venit cum v ---. Puella laeta est: rosae eam ---. ea

Ubi est impluvium? --- [: impluvium] est in ---. In implueās viō --- est. In ātriō nūllae --- sunt.

\booksection{Pēnsum C}

Num Iūlius sōlus in villā habitat? Quot filios et quot filiās habent Iūlius et Aemilia? Ubi est impluvium? Ubi dormiunt servī? Adestne Iūlius in peristylō cum Aemiliā? Ubi est Iūlius? Estne Aemilia sōla in peristylō? Quid Iūlia agit in hortō? Cūr pueri Iūliam rīdent? vAriminum

\navnote{Placentia Padus Genua Via 4e a V. S ci u? e 3 Corsica e}coma V*

Mare Inferum o o eo, 9. 4 9 aut, Sicilia

\booksection{Vocābula nova}
\noindent\small in, ex, cum, ab, -is, sine, indicātīvus (ind), vocā, vocat, vocāte, vocant, -ete, -et, -ite, -it, -unt, -īte, -iunt.

\bookchapter{Cap. VI}{Via Latīna}

I In Italiā multae et magnae viae sunt: via Appia, via Latīna, via Flāminia, via Aurēlia, via Aemilia. Via Appia \navnote{disium (acc)}est inter Rōmam et Brundisium; via Latīna inter Rōmam et Capuam; via Flāminia inter Rōmam et Ariminum; via Aurēlia inter Rōmam et Genuam; via Aemilia inter Arīminum et Placentiam. Brundisium, Capua, Arīminum, Genua, Placentia, Ōstia magna oppida sunt. \navnote{prope Rōmam \contr{} procul ab Rōmā}Ubi est Ōstia? Ōstia est prope Rōmam. Tusculum quōque prope Rōmam est. Brundisium nōn est prope Rōmam, sed procul ab Rōmā: via Appia longa est. Via Latīna nōn tam longa est quam via Appia. Quam longa est via Flāminia? Neque ea tam longa est quam via Appia. Tiberis fluvius nōn tam longus est quam fluvius Padus.

\navnote{circum}Circum oppida mūri sunt. Circum Rōmam est mūrus antīquus. In mūrō Rōmānō duodecim portae sunt. Porta \navnote{porta = magnum ōstium}prima Rōmāna est porta Capēna. Circum oppidum Tusculum mūrus nōn tam longus est quam circum Rōmam.

\navnote{lectīca}\navnote{oppidum vīlla}Villa Iūliī est prope Tusculum. Ab oppidō Tūsculō \navnote{0---n ab oppidō ad vīllam}ad vīllam Iūliī nōn longa via est. Ecce Iūlius et quattuor servī in viā. Iūlius ab oppidō ad vīllam suam it. Domi- 20 \navnote{Iūlius it (sing) servī eunt (plūr)}nus et servī ab oppidō ad vīllam eunt. Dominus in lecticā est. Duo servī lectīcam cum dominō portant. Servī quī lectīcam portant sunt Ursus et Dāvus. Iūlius nōn in \navnote{ZR}viā ambulat, servī eum portant. Syrus et Lēander ambulant. Syrus saccum portat et Lēander quoque saccum 25 portat: Syrus et Lēander duōs saccōs in umeris portant. \navnote{V. umerus `` quī quōs: servī quī saccōs portant saccī quōs servī portant}Saccī quōs Syrus et Lēander portant magnī sunt, sed saccus quem Syrus portat nōn tam magnus est quam saccus Lēandrī. Quattuor servī dominum et duōs saccōs \navnote{Lēander -dri (gen) vehunt = portant (ab...}ab oppidō ad vīllam vehunt.

Iulius in lecticā est inter Ursum et Dāvum. Ursus est ante Iūlium, Dāvus post eum est. Syrus et Lēander nōn ante lectīcam, sed post lectīcam ambulant. Venitne Iūlius ā vīllā? Nōn ā villa venit. Unde venit Iūlius? Ab \navnote{āJab ante cēterās litterās}oppidō venit. Quō it Iūlius? Ad vīllam it. Post eum 35 Tusculum est, ante eum est villa.

Iūlius sōlus nōn est, nam quattuor servī apud eum \navnote{apud eum sunt = cum eō sunt}sunt. Mēdus nōn est apud dominum, nam is dominum īrātum timet. Mēdus est malus servus quī nummōs dominī in sacculō suō habet. Dominus servōs malōs baculō verberat; itaque servī malī dominum et baculum eius timent. Dāvus autem bonus servus est, neque is \navnote{autem = sed Dāvus}Mēdum amat. Dāvus amīcus Mēdī nōn est, nam servus \navnote{in-imīcus \contr{} amīcus}bonus et servus malus nōn amīcī, sed inimici sunt. Mēdus est inimīcus Dāvī. Ursus autem amīcus Dāvī est.

Mēdus abest ā dominō suō. Estne in oppidō Tusculō? Mēdus Tūsculi nōn est; neque Rōmae est Mēdus, sed in \navnote{in}viā Latīnā inter Rōmam et Tusculum. Unde venit Mēdus? Tusculō venit, neque is ad vīllam Iūlit it. Quō it so. Mēdus? Rōmam it. Tusculum post eum est, ante eum est Rōma. Mēdus viā Latinā Tūsculo Rōmam ambulat.

Etiam Cornēlius, amīcus Iūliī, in viā Latinā est inter \navnote{C ---À Op A Rōmae est B Tūsculi est \contr{} Roma Tūsculum it D Rōmam Tusculō it Tusculum --- ad oppidum Tūsculum}Rōmam et Tusculum. Unde venit Cornēlius? Is nōn Tusculō, sed Rōmā venit. Quō it? Cornēlius nōn Rōmam, sed Tusculum it. Rōma post eum, ante eum Tusculum est. Cornēlius in equō est. Equus quī Cornēlium vehit pulcher est. Iūlius et Cornēlius ad villas suās eunt. Villa ubi Iūlius habitat prope Tusculum est. Ubi habitat Cornēlius? Is Tusculī habitat.

Iam Iūlius prope vīllam suam est. Servī quī lectīcam portant fessī sunt. Dominus autem fessus nōn est, nam \navnote{Iūlius ab Ursō et Dāvō portātur = Ursus et Dāvus Iūlium portant}is nōn ambulat. Iūlius ab Ursō et Dāvō portātur, itaque is fessus nōn est. Fessī sunt Syrus et Lēander, nam iī duōs magnōs saccōs umeris portant, neque vacul sunt saccī! Saccī quī ā Syrō et Lēandrō

portantur

magnī 65 \navnote{saccī quī ā Syrō et Lēandrō portantur = saccī quōs Syrus et Lēander portant saccus quī ā Lēandrō portātur --- saccusquem Lēander portat is equō vehitur = equus}sunt, sed saccus quem Syrus portat nōn tam magnus est quam saccus quī ā Lēandrō portātur. Itaque Syrus nōn tam fessus est quam Lēander. Cornēlius nōn est fessus, nam is equō vehitur. Iūlius lecticā vehitur. Servī ambulant. Dominī vehuntur. Mēdus ambulat, nam is servus 70 est neque equum habet.

Iūlius prope vīllam suam est. Mēdus autem, quī dominum īrātum timet, procul ā villa Iūliī abest. Dominus ā servō malō timētur. Mēdus prope Rōmam est; iam mūri Rōmānī ab eō videntur et porta Capēna. (Is quī viā 75 Laūnā venit per portam Capēnam Rōmam intrat.) Cūr. \navnote{per portom n}Rōmam it Mēdus? Rōmam it, quia Lydia Rōmae habitat, nam Lydia amica eius est: Mēdus Lydiam amat et \navnote{Mēdus ab eā amātur = ea Mēdum amat}ab eā amātur. Mēdus Rōmam vocātur ab amīcā suā, quae femina est pulchra et proba. Itaque is fessus nōn 80 est et laetus cantat: \navnote{amica mea id ā Lydiā nōn audītur = Lydia id nōn audit}Nōn via longa est Rōmam, ubi amica habitat mea pulchra. Sed id quod Mēdus cantat ā Lydiā nōn audītur!

Iam Iūlius in villā est et ab Aemiliā līberīsque laetīs salūtātur. Cornēlius Tusculī est. Mēdus autem Rōmae 85 est ante ōstium Lydiae. Mēdus ōstium pulsat.

Lydia imperat: ``Intrā!''

Mēdus per ōstium intrat et amīcam suam salūtat:. ``Salvē, mea Lydia! Ecce amīcus tuus quī sōlus Rōmam ad tē venit.'' \navnote{Lydia verbīs Mēdī dēlectātur --- verba Mēdī Lydiam dēlectant}UU] Lydia verbīs Mēdī dēlectātur eumque salūtat: ``Ō amīce, salvē! Ubi est dominus tuus?'' Mēdus: ``Iūlius in villā est apud servōs suōs - neque is iam meus dominus est!'' 95 Verba Mēdī ā Lydiā laetā audiuntur.

\booksection{Grammatica Latīna}

\navnote{ūna praepositiō (prp) duae praepositiōnēs}Praepositiōnēs Iūlius ad vīllam it; ad oppidum; ad ancillās. Ursus ante Iūlium est; ante eum; ante vīllam. Dāvus post Iūlium est; post eum; post vīllam. Via inter Rōmam et Capuam; inter servōs. Ōstia est prope Rōmam; prope vīllam; prope eam. Circum oppidum mūrus est; circum mēnsam. Mēdus est apud amīcam suam, nōn apud dominum. 105 Mēdus per portam Capēnam Rōmam intrat; per ōstium. Ab, ad, ex, cēt. sunt praepositiōnēs.

Praepositiōnēs \navnote{ad, ante, apud, circum, inter, per, post, prope ablā, cum, ex, in, sine}cum accūsātīvō: ad, ante, post, inter, prope, circum, apud, per, cēt.; praepositiōnēs cum ablātīvō: ab/d, cum, ex, tn, sine, cēt. (vidē cap. v). Quō it Mēdus? Mēdus Rōmam it. Quō it Cornēlius? Is \navnote{quō? Rōmam/Tūsculum}Tusculum it. Unde venit Cornēlius? Cornēlius Rōma venit. Unde venit \navnote{unde? Rōma/TQsculō}Mēdus? Is Tūsculo venit. \navnote{ubi? Rōmai/Ttisculf oppidum}Ubi habitat Lydia? Lydia Rōmae habitat. Ubi habitat Cornēlius? Is Tūsculi habitat. Accūsātīvus:

Rōmam,

Tūsculum, Capuam,

Brundisium, Ōstiam = ad oppidum -am/-um. Ablātīvus: Rōmā, Tusculō, Capuā, Brundisiō, Ōstiā = ab oppidō -ā/-ō. Locātīvus \navnote{-ae}Brundisif, Ōstiae = in oppidō -ā/-o. Verbum āctīvum et passīvum, AN Servus saccum portat T = Saccus portātur ā servō

jS Servī saccōs portant iu = Saccī portantur d servīs

*``Portat portant'' verbum āctivum est. 'Portātur portantur \navnote{-nt -tur -ntur sing. plūr. -antur}est verbum passīvum.

Āctīvum: -t -nt. Passīvum: -tur -ntur.

Exempla: [1] -ātur -antur: Aemilia Iūlium amat et ab eō amātur. Pater et 130 māter līberōs amant et ab iīs amantur. Lydia verbīs Mēdī dēlectātur. \navnote{-entur}[2] -ētur -entur: Iūlia puerōs nōn videt neque ab iīs vidētur. Puerī Iūliam nōn vident neque ab eā videntur. \navnote{-untur}[3] -itur -untur: Dāvus sacculum in mēnsā pōnit = sacculus in 135 mēnsā pōnitur ā Dāvō. Pueri rosās in aquā pōnunt = rosae ā puerīs in aquā ponuntur. Cornēlius equō vehitur. [4] -ītur -iuntur: Iūlia puerōs audit neque ab iīs audītur = \navnote{-iuntur}pueri ā Iūliā audiuntur neque eam audiunt.

\booksection{Pēnsum A}

Iūlius ab oppid- Tusculō ad vlll- su- it. Vīlla eius prope Tuscul- est. Iūlius in lecticā est inter Urs- et Dāv-. Dominus ā servīs port-. Ursus et Dāvus nōn saccōs port-, sacc- ā Syrō et Lēandrō port-. Saccus quem Lēander port- nōn tam parvus est quam saccus qul ā Syr- port-.

Mēdus nōn est apud domin-, nam servus malus dominum um-. Dominus ā serv- mal- tim-. Dominus serv- mal- vocneque ab eō aud-. Serv- mal- ā domin- voc- neque eum aud-.

Quō it Mēdus? Rōm- it. Unde venit? Mēdus Tuscul- venit. Ante Mēd- est Rōma, Tusculum post e- est. Cornēlius nōn Tusculō Rōmam, sed Rōm- Tūscul- it, nam is Tusculvia habitat. Cornēlius nōn ambulat, sed equ- veh-.

Lydia, amica Mēdī, Rōm- habitat. Iam Mēdus Rōm- apud porta Lydi- est. Mēdus amīc- su- salūt- et ab amīc- su- salūt-, saccus \navnote{umerus}nam Lydia Mēd- am- et ab e- am-.

\booksection{Pēnsum B}

Ōstia nōn --- ā Rōmā, sed --- Rōmam est. --- venit Iūlius? \navnote{longus fessus}malus Tusculō venit et --- vīllam it. Duo servī eum ---. Syrus et! Lēander, qul --- lectīcam ambulant, duōs --- portant. Saccus Syri --- magnus est --- saccus quī--- --- Lēandrō portātur; --- portat Syrus nōn tam --- est quam Lēander.

Mēdus --- dominum nōn est, nam Mēdus dominum Irātum limet ---. Mēdus servus --- [= improbus] est; --- Mēdus et Iūlius intrai nōn ---, sed inimici sunt. Via Latīna, quae est --- Rōmam et quam Capuam, nōn tam --- est quam --- Appia. --- ambulat Mēdus? inter Is Rōmam ambulat, --- amica eius Rōmae habitat. Cornēlius circum --- [= sed C.] Rōmā Tusculum ---. Tusculum est --- eum, --- ad eum est Rōma. Cornēlius equō ---, is nōn ---. Iūlius et Cornēante lius ad vlllās suās ---.

\booksection{Pēnsum C}

\navnote{post}Ambulatne Iūlius? Quī Iūlium portant? ā Quid portant Syrus et Lēander? Unde venit Iūlius et quō it? Quō ii Mēdus? i Etiamne Cornēlius Tusculō Rōmam it? I Ubi habitat Cornēlius? Cūr Mēdus laetus est? Quae est Lydia? Quid habet Mēdus in sacculō suō? Suntne amlcī Iūlius et Mēdus? Num fortat verbum passīvum est?

\booksection{Vocābula nova}
\noindent\small oppidō Tūsculo Tūsculō --- ab oppidō, lusculō, Rōma0, Tūsculum, ante, -ōs, inter, -ās, prope, mē, apud, tē, per, -us, -am, -nt, -ae, tur, -um, -īs, -ntur, nova, nam, itaque, unde?, pracpositiō, locātlvus, ācūvum, passīvum.

\bookchapter{Cap. VII}{Puella et rosa}

\navnote{Iūlia rosam ante nāsum tenet illic: in peristylō oculus A lacrimat \contr{} ridet speculum f Az Jg}Ecce Mārcus et Quīntus ante ōstium vīllae. Puerī Iūlium exspectant. Māter nōn est apud filios, ea in peristylō est; illic virum suum exspectat. Aemilia laeta nōn est, quia Iūlius abest; nam Aemilia virum suum amat.

Ubi est Iūlia? In cubiculō suō est. Iūlia, quae sōla est 5 illic, rosam ante nāsum tenet. Puella lacrimat: in oculīs eius sunt lacrimae.

Iūlia speculum sūmit et ante oculōs tenet. Puella sē in speculō videt et sē interrogat: ``Estne foedus nāsus meus?'' Nāsus eius autem fōrmōsus nōn est. Iūlia rūrsus 10 \navnote{formōsus -a -um = pulcher}lacrimat.

Syra ōstium cubiculī pulsat.

Iūlia: ``Intrā!''

\navnote{tn cubiculum (acc) \contr{} ex cubiculō claudit --- \contr{} aperit hortus hortus}Syra ōstium aperit et in cubiculum intrat, neque ōstium post sē claudit. Iūlia Syram post sē in speculō videt. Syra nōn videt lacrimās Iūliae, nam puella sē nōn vertit.

Syra: ``Ō, hic est mea puella. Veni in hortum, Iūlia!'' Iūlia imperat: ``Claude ōstium!'' Ancilla pāret. Iūlia: ``Num nāsus meus foedus est, Syra?'' Syra: ``Foedus? Immō formōsus est nāsus tuus.'' \navnote{immō formosus: nōn foedus sed fōrmōsus est sunt}Iūlia sē vertit. Iam Syra lacrimās videt. Syra: ``Quid est, mea Iūlia? Tergē oculōs! Es laeta! Nāsus tuus tam fōrmōsus est quam meus.'' Iūlia: ``Sed nāsus tuus nōn fōrmōsus est!'' Syra: ``Quid? Nōnne formōsus est nāsus meus?'' Iūlia: ``Immō foedus est! Ecce speculum, Syra.'' Iūlia speculum tenet ante Syram, quae nāsum suum in speculō videt. Ancilla oculōs claudit et tacet. --- Ill Ecce Iūlius ad vīllam advenit. Servī lectīcam ante \navnote{ad-venit = ad... venit salvē! salvēte/ salvē, pater! (sing) salvēte, filii! (plūr) ōstiārius = servus quī ōstium aperit et claudit in vīllam \contr{} ex vīā}ōstium pōnunt. Pater filios salūtat: ``Salvēte, filii!'' et ā filiis salūtātur: ``Salvē, pater!'' Iūlius ambulat ad ōstium, quod ab ōstiāriō aperitur. Dominus per ōstium in vīllam intrat. Post eum veniunt Syrus et Lēander, quī duōs saccōs portant. Ōstiārius post eōs ōstium claudit. Ursus et Dāvus cum lectīcā vacuā discēdunt. Pueri saccōs plēnōs quī ā servīs portantur vident et \navnote{in-est= in...est}interrogant: ``Quid inest in saccīs?'' Iūlius respondet: ``In saccō quem Lēander portat māla Insunt. Hīc saccum pōne, Lēander!'' \navnote{eum:: saccum hic saccus = saccus quī hic (apud mē) est}Lēander saccum pōnit ante Iūlium, qul aperit eum. Iūlius: ``Vidēte, puerī: hic saccus plēnus mālōrum est.''

Iulius mālum ē saccō sūmit et sē vertit ad Mārcum: ``Ecce mālum tuum, Mārce.'' Iūlius Mārcō mālum dat. \navnote{ēlexante cēterās litterās}Pater filio suō magnum mālum dat. Iam Mārcus mālum habet, neque Quīntus mālum habet. Iūlius Quīntum ad sē vocat et ei mālum dat. Iulius Quīntō mālum dat. Iam et Mārcus et Quīntus māla habent.

Mārcus: ``Quid inest in saccō Syri?''

Iūlius: ``In saccō eius pira īnsunt. Aperi saccum, \navnote{pirum}Syre! Vidēte, pueri: hic saccus plēnus est pirōrum.`` Iūlius duo pira ē saccō sūmit: ''Ecce pirum tuum, Mārce, et tuum, Quīnte.`` Pater filiis suīs pira dat. Filit, quī iam nōn sōlum māla, sed etiam pira habent, laeti sunt. Servī \navnote{sōlum = tantum}autem neque māla neque pira habent.

\navnote{dat dant}Quīntus: ``Etiam servīs dā māla et pira, pater!''

Iūlius Syrum et Lēandrum servōs ad sē vocat et iīs māla et pira dat. Dominus servīs māla et pira dat.

Aemilia cum Dēliā ē peristylō in ātrium intrat, laeta ad Iūlium adit eumque salūtat. Aemilia virō suō ōsculum dat. Iūlius Aemiliae ōsculum dat.

Iūlius: Quid agit Iūlia?``

Aemilia: ``Rosās carpit in hortō.''

Iūlius: ``Currite in hortum, pueri, et vocāte eam!''

Quīntus currit. Mārcus nōn currit, sed ambulat.

Iulius imperat: ``Age! Curre, Mārce!''

Etiam Mārcus currit. Pueri per peristylum in hortum currunt. Illic autem puella nōn est. Pueri ex hortō per 70 peristylum in ātrium ambulant. Mārcus: ``Iūlia neque in hortō neque in peristylō est.'' Aemilia: ``Estne Syra in hortō?'' Quīntus: ``Nōn est. Nūlla ancilla illic est.'' 75 \navnote{it eunt}Aemilia: ``Nōn sōlum Iūlia, sed etiam Syra abest! Dēlia,1 1 ad cubiculum Iūliae!'' Dēlia ad cubiculum Iūliae it, ōstium pulsat, aperit, in cubiculum intrat. Illic nōn sōlum Iūlia, sed etiam Syra est. Oculī Iūliae plēnī sunt lacrimārum. Dēlia: ``Veni in ātrium, Iūlia! Illīc pater tuus tē exspectat.'' Iulia oculōs et nāsum terget, rosam sūmit, ē cubiculō exit. Syra et Dēlia post eam exeunt. Iūlia in ātrium ad \navnote{ex-it e \contr{} īntrat ex-it ex-eunt eī-que haec rosa = rosa quae hīc (apud mē) est}Iūlium currit eīque ōsculum dat. Iūlia: ``Ecce rosa. Nōnne pulchra est haec rosa?'' Iūlius: ``Nūlla rosa tam pulchra est quam filia mea!'' Iūlius fīliae suae ōsculum dat. Iamne lacrimat Iūlia? \navnote{īam-ne}Immō laeta est et rīdet. Iūlia: ``Num nāsus meus foedus est?'' Iūlius: ``Foedus? Immō tam fōrmōsus est quam---hoc \navnote{hocmālum- mālum quod hic (apud mē) est hic haec hoc: hic saccus (m) haec rosa (f) hoc mālum (n)}mālum! Ecce mālum tuum, Iūlia.`` Pater filiae mālum magnum et formōsum dat. Iūlia mālum terget et ante oculōs tenet. Iūlia: ''Ō, quam formōsum est hoc mālum!`` Puella laeta mālō suō ōsculum dat! \navnote{ef:}Iūlius: Hoc pirum etiam tuum est, Iūlia.`` Iūlius eī pirum dat. Iam puella et mālum et pirum habet. Aemilia: ''Etiam ancillīs meīs māla et pira dā!``

Iūlius ancillās ad sē vocat et iīs quoque māla et pira \navnote{iīs(= eīs):: ancillīs}dat. Ancillae laetae ex ātriō exeunt. ---

Cui Iūlius mālum dat? Puerō mālum dat. Puer cui \navnote{cui? puerō puellae}Iūlius mālum dat est filius eius.

Cui Iūlius ōsculum dat? Puellae ōsculum dat. Puella cui Iūlius ōsculum dat filia eius est.

\booksection{Grammatica Latīna}

Datīvus \navnote{datīvus (dat) \contr{} dat}[A] Masculīnum.

Iūlius servō (Syrō) mālum dat.

\navnote{servō servīs}Iūlius servīs (Syrō et Lēandrō) māla dat.

'Servā' datīvus singulāris est. 'Servīs' est datīvus plūrālis. Datīvus: singulāris -ō, plūrālis -?s. \navnote{ancillae ancillīs -ae -īs}[B] Fēminīnum.

Iūlius ancillae (Syrae) mālum dat.

\navnote{ancilkzi ancillfs}Iūlius ancillīs (Syrae et Dēliae) māla dat.

'Ancilla*' est datīvus singulāris. 'Ancillts' datīvus plūrālis \navnote{-ae}est. Datīvus: singulāris -ae, plūrālis -t\$. [C] Neutrum.

Fluvius oppidō aquam dat.

\navnote{oppidō oppidīs}Fluviī oppidīs aquam dant.

'Oppidō' datīvus singulāris est. Oppidizs' est datīvus plūrā- 120 lis. Datīvus: singulāris -ō, plūrālis -īs.

\booksection{Pēnsum A}

\navnote{-is}Iūlius Mārc-, fili- su-, mālum dat. Iūlius Mārc- et Quīnt-, fili- su-, māla dat. Iūlius etiam serv- su-, Syr- et Lēandr-, māla dat.

C- Aemilia ōsculum dat? Aemilia vir- su- Iūli- ōsculum dat. Iūlius Aemili- ōsculum dat. Iūlius Iūli-, fīli- su-, mālum dat, neque sōlum Iūliae, sed etiam Syr- et Deli-, ancillsu-. Iūlia māl- su- ōsculum dat!

\booksection{Pēnsum B}

Iūlius ad vīllam---. ---. Ōstiārius ōstium --- et post eum ---.

Saccī nōn vacuī, sed --- sunt. Iūlius: ``Vidēte, pueri: --- --- saccus plēnus --- est. Ecce mālum tuum, Mārce.'' Iūlius Mārcō mālum ---. Iūlius filiis nōn --- māla, sed etiam pira dat. i Iam puerī --- māla --- pira habent, sed servī--- --- māla --- pira habent. Dominus servōs ad --- vocat et --- quoque māla et --- 1 dat.

Aemilia ad Iūlium --- et eī--- --- dat. Iūlia abest. Puerī nōn! ambulant, sed --- in hortum. --- [: in hortō] Iūlia nōn est, ea in cubiculō suō est. Iūlia nōn ridet, sed ---: in --- eius sunt ---. Dēlia: ``In ātriō pater tuus tē ---, Iūlia.'' Iūlia --- [= ex] cubi culō ---, ad Iūlium currit et --- [: Iūliō] ōsculum dat. Iūlia rosam ante Iūlium ---. Iūlia: ``Nōnne --- rosa --- [= pulchra] \contr{} est?'' \contr{}

\booksection{Pēnsum C}

Quem puerī exspectant? I Venitne Iūlius Rōmā? Quis ōstium aperit et claudit? I Quid inest in saccīs?: Cui Iūlius mālum prīmum dat? s Cui Aemilia ōsculum dat? Estne Iūlia in hortō? \navnote{iīs cui}Quō it Dēlia? Estne Iūlia sōla in cubiculō suō? c Rīdetne Iūlia? Quō Iūlia currit? Quid Iūlius dat filiae suae?

\booksection{Vocābula nova}
\noindent\small Aemilia, ōsculum, dat, Mārcus, Quīntus, ambulat, currit, oculus, lacrima, speculum, ōstiārius, mālum, pirum, fōrmōsus, plēnus, exspectat, tenet, lacrimat, aperit, claudit, vertit, terget, advenit, inest, adit, exit, es, sē, hic, haec, hoc, immō, nōnne?, neque... neque, sōlum, illīc.

\bookchapter{Cap. VIII}{Taberna Rōmāna}

ALblNVS

Ecce taberna Rōmāna, in quā gemmae et margaritae / \navnote{margarita haec taberna, acc hanc tabernam}multae sunt. Cuius est haec taberna? Albīnī est. Albīnus hanc tabernam habet. Quī tabernam habet tabernārius est. Albīnus est tabernārius Rōmānus quī gemmās et margarītās vēndit. Aliī tabernāriī librōs vēndunt, aliī 5 māla et pira, aliī rosās et līlia.

\navnote{ānulus cum gemmā}Gemmae et margaritae sunt ōrnāmenta. Ānulus cum gemmā ōrnāmentum pulchrum est. Etiam līnea cum \navnote{līnea cum margaritis = līnea margaritāārum}margaritīs ōrnāmentum est. Līnea sine margarītīs nōn est ōrnāmentum! 10

\navnote{haec via, abl hāc viā}Multae feminae quae in hāc viā ambulant ante tabernam Albīnī cōnsistunt, nam feminae ōrnāmentīs dēlectantur. Eae quae magnam pecūniam habent multa ōrnā\navnote{quae = eae quae}menta emunt. Quae nūllam aut parvam pecūniam ha-bent ōrnāmenta aspiciunt tantum, nōn emunt. Etiam 15 viri multī ad hanc tabernam adeunt. Qul magnam pecūniam habent ōrnāmenta emunt et feminīs dant; cēteri rūrsus abeunt. Fēminae quārum viri magnam pecūniam \navnote{ab-it ab-eunt pecūniōsus -a -um: p. est = magnam pecūniam habet}habent multa ōrnāmenta ā viris suīs accipiunt. 20

Aemilia, cuius vir pecūniōsus est, multa ōrnāmenta ab eō accipit. Aemilia ānulum in digitō et margaritās in collō multaque alia ōrnāmenta habet. Ānulus digitum \navnote{collum digitus}Aemiliae ōrnat, margaritae collum eius ornant. Fēminae gemmīs et margaritis ānulisque ōrnantur.

In viā prope tabernam Albīnī vir et femina ambulant. Quī vir et quae femina? Est Mēdus, quī cum Lydiā, \navnote{quis? quis/quī vir? quid? quod ōrnāmentum?}amicā suā, ambulat. Mēdus est servus Iūliī, sed dominus eius Rōmae nōn est. Mēdus sine dominō suō cum feminā fōrmōsā in viīs Rōmae ambulat.

Lydia ōrnāmentum pulchrum in collō habet. Quod ōrnāmentum? Ōrnāmentum quod Lydia habet est līnea margarītārum. Collum Lydiae margaritis pulchrīs ōrnātur; Lydia autem nūllum aliud ōrnāmentum habet, quia pecūniōsa nōn est, neque pecūniōsus est amīcus eius. (Pecūniōsus est qul magnam pecūniam habet.) II

Albīnus clāmat: ``Ōrnāmenta!

Ōrnāmenta

feminārum! Ōr-nā-men-ta! Emite ōrnāmenta!``

Lydia cōnsistit oculōsque ad tabernam Albīnī vertit: Lydia tabernam aspicit. Mēdus nōn cōnsistit neque tabernam aspicit.

\navnote{illa taberna = taberna quae illic est}Lydia: ``Cōnsiste, Mēde! Aspice illam tabernam! Ō, quam pulchra sunt illa ōrnāmenta!'' Lydia tabernam Albīnī digitō mōnstrat. Mēdus sē vertit, tabernam videt, cum Lydiā ad tabernam adit. Mēdus et Lydia ante tabernam cōnsistunt. Albīnus eōs salūtat et margarītās 45 \navnote{Albīnus Lydiae (dar) margaritas ostendit hoc, abl hōc}in lineā ante oculōs Lydiae tenet: Albīnus Lydiae margaritās ostendit. Albīnus: ``In hōc ōrnāmentō vīgintī margaritae mag-nae sunt. Nōnne pulchrae sunt hae margaritae?'' \navnote{sing hic haec hoc plūr hi hae haec}Mēdus: ``Amīca mea multās margaritās habet.'' Lydia: ``In hāc tabernā multa alia ōrnāmenta sunt.'' \navnote{aspice/ aspicite/ aspicit. aspiciunt hi, acc hōs}Albīnus iīs trēs ānulōs sine gemmīs ostendit: ``Aspicite hōs ānulōs! Nōnne hi ānulī pulchri sunt?'' Mēdus et Lydia ānulōs aspiciunt. Lydia: ``In hīs ānulls gemmae nūllae sunt!'' \navnote{hi hae haec, abl hīs}Mēdus Albinum interrogat: ``Quot nummīs cōnstat ānulus in quō gemma est?'' \navnote{qul, abl quō ānulus gemmātus = ānulus cum gemmā}Albīnus Mēdō ānulum gemmātum ostendit. Albīnus: Hic ānulus centum nummīs cōnstat.`` Mēdus: Quid?'' Albīnus: ``Pretium huius ānulī est centum seēstertii.'' \navnote{hic haec hoc, gen huius sēstertius = nummus}Mēdus: ``Centum sēstertii? Id magnum pretium est!'' \navnote{hic ānulus, acc hunc ā.um, abl hōc ā.ō tantus -a -um = tam magnus}Albīnus: ``Immō parvum pretium est! Aspice hunc ānulum: in hōc ānulō magna gemma est. Tanta gemma sōla octōgintā sēstertiīs cōnstat.'' Mēdus: ``Num ānulus sine gemmāā vīgintī tantum sēs-

Mēdus, quī alium ānulum gemmātum post Albīnum \navnote{hic ānulus = ānulus quī hīc est ille ānulus = ānulus quī illic est}videt: Hic ānulus pulcher nōn est. Quot sēstertiīs cōnstat ille ānulus?``

Albīnus: ``Quī ānulus?''

\navnote{quantus -a -um = quam magnus ille, gen illīus,}Mēdus: ``Ille post tē. Quantum est pretium illīus ānu1?'' Mēdus ānulum post Albīnum digitō mōnstrat.

Alblnus: ``Ille quoque ānulus centum sēstertiīs cōnstat. Pretium illīus ānulī tantum est quantum huius; sed \navnote{acc illum tantus...quantus = tam magnus...quam}amica tua hunc ānulum amat, nōn illum.``

Lydia ānulum gemmātum ante oculōs tenet.

Lydia: ``Ō, quam pulchrum hoc ōrnāmentum est! Illud ōrnāmentum nōn tam pulchrum est quam hoc, \navnote{ille ānulus (m) illa}neque illa gemma tanta est quanta haec.``

Mēdus: ``Tanta gemma ad tam parvum ānulum nōn convenit.'' IH

Lydia, quae haec verba nōn audit, Mēdō digitōs suōs ostendit, in quibus nūllī ānulī sunt.

\navnote{illud ōrnāmentum (7) quī, abī plūr quibus}Lydia:

``Aspice,

Mēde! In digitīs meīs nūllī sunt ānulī. Aliae feminae digitōs ānulōrum plēnōs habent --- mel digitī vacuī sunt!``

Mēdus: ``Sacculus quoque meus vacuus est!''

Lydia ānulum in mēnsā pōnit. In oculīs eius lacrimae sunt. Mēdus, quī lacrimās videt, sacculum suum in mēnsā pōnit --- neque vacuus est sacculus, sed plēnus nummōrum! Quanta pecūnia est in sacculō Mēdī? In eō nōnāgintā sēstertii īnsunt.

Mēdus: ``Ecce sestertii nōnāgintā.''

\navnote{sēstertii nōnāgintā:}Alblnus: ``Sed nōnāgintā nōn satis est. Pretium ānulī est sēstertiī centum!''

\navnote{hic haec hoc, dat huic}Lydia: ``Dā huic tabernāriō centum sēstertiōs!''

Mēdus: ``Id nimis magnum pretium est! Aliītaber- nārii ānulum gemmaātum octōgintā sēstertiīs vēndunt.''

Albīnus: ``Quī sunt illī tabemāriī?''

Mēdus: ``Quī in aliīs viīs tabernās habent.''

\navnote{quī = ii quī sing ille illa illud plūr illīillae illa}Albīnus: ``Quae sunt illae viae in quibus illae tabernae sunt? Et quae sunt illa ōrnāmenta quae in illista- bernis parvō pretiō emuntur? Nōn sunt ōrnāmenta! Sed aspicite haec ōrnāmenta: hōs ānulōs, hās gemmās, hās 105 margaritās! Haec ōrnāmenta proba sunt! Neque pretium hōrum ōrnāmentōrum nimis magnum est!''

Mēdus: ``Accipe nummōs nōnāgintā --- aut nūllōs!''

Albīnus: ``Num hīc nōnāgintā sēstertiī sunt?''

Mēdus: ``Numerā eōs!''

Albīnus numerat sēstertiōs, quōrum numerus est nōnāgintā.

Albīnus: ``Sunt nōnāgintā.''

Mēdus: ``Satisne est?''

Albīnus nōn respondet, sed nummōs sūmit Mēdōque 115 ānulum dat. Albīnus pecūniam accipit et Mēdō ānulum vēndit sēstertiīs nōnāgintā.

Mēdus sē ad Lydiam vertit: ``Accipe hunc ānulum ab

quō digitō? In digitō mediō.

Mēdus: ``Hic ānulus ad digitum tuum nōn convenit. Ānulus nimis parvus est aut digitus nimis magnus!''

Lydia: ``Ō Mēde! Digitus medius nimis magnus est. Pōne ānulum in digitō quārtō!''

Mēdus ānulum in digitō Lydiae quārtō pōnit. Ānulus satis magnus est et ad digitum convenit, nam digitus quārtus nōn tantus est quantus digitus medius. Lydia laeta digitum suum aspicit et amīcō suō ōsculum dat.

Mēdus et Lydia ā tabernā abeunt. Lydia, quae Rōmae habitat, Mēdō viam mōnstrat.

Albīnus rūrsus clāmat: ``Ōrnāmenta! Ōr-nā-men-ta!'' et aliōs virōs pecuniōsōs, quōrum amicae nulla aut pauca ōrnāmenta habent, exspectat.

\booksection{Grammatica Latīna}

\navnote{duo prōnōmina Singulāris: nōminātīvus quislquīis ille is servus = ille servus accūsātīvus quem eum illum genetīvus cuius eius illīus datīvus cui ei ill ablātīvus quō eō illō Plūrālis: nōminātīvus illī accūsāativus quōs eōs illōs genetīvus quōrum eōrum illōrum}Prōnōmina 'quis 'quī', 'is', Hlley [A] Masculīnum.

Quis saccum portat? Servus saccum portat. Qul servus? Servus quī saccum portat est Syrus. Is/ille servus saccum portat.

lūlius servum vocat. Quem servum?

Servus quem Iūlius vocat est Syrus. Iūlius eum/illum servum vocat.

Iūlius dominus servī est. Cuius servī? Syrus est servus cuius dominus Iūlius est. Iūlius dominus e:?us/illtus servī est.

Iūlius servō mālum dat. Cui servō? Servus cui Iūlius mālum dat est Syrus. Iūlius ei/ill; servō mālum dat.

Saccus ā servō portātur. Ā quō servō? Servus ā quō saccus portātur est Syrus. Saccus ab eōlillō servō portātur.

Servī saccōs portant.

Quī servī? Servī quī saccōs portant sunt Syrus et Lēander. It/illī servī saccōs portant.

Iūlius servō5 vocat. Quōs servōs? Servī quōs Iūlius vocat sunt Syrus et Lēander. Iūlius eoōs/illōs servōs vocat.

Iūlius dominus servōrum est. Quōrum servōrum? Servī quōrum dominus est Iūlius sunt Syrus et Lēander. Iūlius dominus eōrum/illōrum servōrum est.

\navnote{datīvus quibus iīs illis ablātīvus quibus iīs illīs}Iūlius servīs māla dat. Quibus servīs? Servī quibus Iūlius māla dat sunt Syrus et Lēander. I. i1s/illis servīs māla dat.

Saccī ā servīs portantur.

Ā quibus servīs? Servī ā quibus saccī portantur

sunt Syrus et Lēander. Ab iīs/illis servīs saccī portantur. [B] Fēminīnum.

\navnote{Singulāris: nōminātīvus quae ea illa accūsātīvus quam eam illam genetīvus cuius eins illīus datīvus}Ancilla abest. Quae ancilla? Ancilla quae abest est Syra. Ealilla ancilla abest.

Iūlius ancillam vocat. Quam ancillam? Ancilla quam Iūlius vocat est Syra. Iūlius eam/illam ancillam vocat.

Iūlius dominus ancillae est. Cuius ancillae? Syra est ancilla cuius dominus Iūlius est. I. dominus eius/illtus ancillae est.

Iūlius ancillae mālum dat. Cui ancillae? Ancilla cui Iūlius mālum dat est Syra. Iūlius eī/illīancillae mālum dat.

\navnote{illī ablātīvus quā eā illā Plūrālis: nōminātīvus quae eae illae accūsātīvus quās eās illās genetīvus quārum eārum illārum}Iūlius ab ancillā salūtātur. Ā quā ancillā? Ancilla ā quā Iūlius salūtātur est Syra. Iūlius ab eā/illā ancillā salūtātur.

Ancillae absunt. Quae ancillae? Ancillae quae absunt sunt Syra et Dēlia. Eaelillae ancillae absunt.

Iūlius ancillās vocat. Quās ancillās? Ancillae quās Iūlius vocat sunt Syra et Dēlia. Iūlius eās/illās ancillās vocat.

Iūlius dominus ancillārum est. Quārum ancillārum? Ancil- 175 lae quārum dominus est Iūlius sunt Syra et Dēlia. Iulius do-minus eārumlillārum ancillārum est.

\navnote{datīvus quibus iīs illis ablātīvus quibus iīs illīs}Iūlius ancillīs māla dat. Quibus ancillīs? Ancillae quibus Iūlius māla dat sunt Syra et Dēlia. I. iīs/;lIis ancillīs māla dat.

Iūlius ab ancillīs salūtātur. Ā quibus ancillīs? Ancillae ā 180 quibus Iūlius salūtātur sunt Syra et Dēlia. Ab:is/illts ancillīs Iūlius salūtātur. [C] Neutrum.

\navnote{Singulāris: nōminātīvus quid/quod id illud}Quid est ānulus? Ānulus est ōrnāmentum. Quod ōrnāmentum? Ānulus est ōrnāmentum quod digitum ōrnat. IdJillud 185 ōrnāmentum pulchrum est.

\navnote{accūsātīvus quidlquod id illud}Quid Lydia in collō habet? Ōrnāmentum habet. Quod ōrnāmentum? Ōrnāmentum

quod Lydia in collō habet est līnea margarītārum. Lydia idlillud ōrnāmentum amat.

Pretium ōrnamentī est HS(= (7 sestertiī) c. Cuius ōrnamentī? \navnote{genetīvus cuius eius illīus}Ōrnāmentum cuius pretium est hs c est ānulus. Pretium eius Aillius ornāmenti est HS C.

\navnote{datīvus}Fluvius oppidō aquam dat. Cui oppidō? Oppidum cui fluvius aquam dat est Capua. Fluvius et/ill; oppidō aquam dat. 195

\navnote{ablātīvus quō eō iUo Plūrālis: nōminātīvus quae ea illa}Cornēlius in parvō oppida habitat. In quō oppidō? Oppidum in quō Cornēlius habitat est Tusculum. In eo/illō oppidō habitat Cornēlius.

Quae oppida prope Rōmam sunt? Ōstia et Tusculum sunt oppida quae prope Rōmam sunt. Ea/illa oppida prope Rōmam sunt.

\navnote{accūsātīvus quae ea illa Benetivus quōrum eōrum illōrum}Albīnus ōrnāmenta vēndit. Quae ōrnāmenta? Ōrnāmenta quae A. vēndit sunt ānulī. Ea/illa ōrnāmenta vēndit A.

Pretium ōrnāmentōrumest hs c. Quārum ōrnāmentōrum? Ōrnāmenta quōrumpretium est hs c Quōrum ōrnāmentōrum? Ōrnāmenta quōrum pretium est HS \contr{} ānulī sunt. Pretium eōrumJiūōrum ōrnāmentōrum est hs c.

Fluviī oppidīs aquam dant. Quibus oppidīs? Oppida quibus \navnote{datīvus quibus iīs illis}fluviī aquam dant sunt Capua et Brundisium.

Fluviī ilsliUxs oppidīs aquam dant.

\navnote{ablātīvus quibus iīs illis}Fēminae ōrnāmentts dēlectantur. Quibus ōrnāmentīs? Ōrnāmenta quibus feminae dēlectantur sunt margaritae et gem-mae. Iīs/illīsōrnāmentīs dēlectantur feminae. Prōnōmen hic Masculīnum Fēmininum ---

Sing. Nōm

hunc mūrum hanc viam

Gen.

huius verbī

Gen.

Dat.

huic mūrō

huic viae huic verbō

Abl

hōc mūrō hāc viā

.hl mūri hae viae Plūr. Nōm

hōs mūrōs haec verba

Gen.

Dat.

hīs mūris hīs viīs hīs verbīs

Abl. hīs mūrīs hīs viīs hīs verbīs

\booksection{Pēnsum A}

Qu- est Albīnus? Est tabernārius qu- ōrnāmenta vēndit. Quōrnāmenta? Ōrnāmenta qu- Albīnus vēndit sunt gemmae et margaritae. Qu- emit Mēdus? Ōrnāmentum emit. Qu- ōrnāmentum? Ōrnāmentum qu- Mēdus emit est ānulus c- pretium est Hs C. Digitus in qu- ānulus pōnitur est digitus quārtus.

H- servus Mēdus, ill- Dāvus est. Lydia h- servum amat, nōn ill-. Lydia amīca h- servī est, nōn ill-. Lydia h- servō ōsculum dat, nōn ill-. Lydia ab h- servō amātur, nōn ab ill---.

H- oppidum est Tusculum, ill- est Brundisium. Cornēlius in h- oppidō habitat, nōn in ill-. Viae h- oppidī parvae sunt.

\booksection{Pēnsum B}

Gemmae et margaritae --- pulchra sunt. Aemilia multa ōrnāmenta ā Iūliō ---. Aemilia --- in collō et ānulum in --- habet. Multae feminae ante tabernam Albīnī--- --- et ōrnāmenta eius aspiciunt. Viri ōrnāmenta --- et feminīs dant. --- gemmātus centum sēstertiīs --- --- ānulī est centum sēstertil, sed Mēdus --- (xc) tantum habet. Albīnus: ``Nōnāgintā nōn --- est!'' Mēdus: ``Accipe nōnāgintā sēstertiōs --- nūllōs!'' Ānulus ad digitum medium nōn ---: digitus medius --- magnus est. Sed ānulus convenit ad digitum --- (iv), quī nōn --- est quantus digitus ---. Lydia laeta digitum suum --- et cum Mēdō ā tabernā ---. Lydia Mēdō viam ---.

\booksection{Pēnsum C}

Quid Albīnus vēndit? Ā quō Aemilia ōrnāmenta accipit? Ambulatne Mēdus cum dominō suō? Ubi Mēdus et Lydia cōnsistunt? Cūr Mēdus margaritās nōn emit? Cūr Lydia nūllum ānulum habet? Estne vacuus sacculus Mēdī? Quot sēstertiīs cōnstat ānulus gemmātus? Ad quem digitum ānulus convenit?

\booksection{Vocābula nova}
\noindent\small accipe/, accipite!, accipit, accipiunt, digitus, medius, digitus N quārtus np A r, quārtus, sing, hic, haec, hoc, huius, huic, hāc, hōc, plūr, hi, hae, hōs, hās, hōrum, hārum, hīs, nova, taberna, gemma, margarita, tabernārius, ōrnāmentum, ānulus, līnea, collum, pretium, sēstertius, pecūniōsus, gemmātus, vlgintl, octōgintā, nōnāgintā, vēndit, cōnsistit, emit, aspicit, abit, ōrnat, clāmat, mōnstrat, ostendit, cōnstat, convenit, alius, ille, tantus, quantus, satis, nimis, aut, prōnōmen.

\bookchapter{Cap. IX}{Pāstor et ovēs}

I Hic vir, quī in campō ambulat, pāstor Iūliī est. Pāstor nōn sōlus est in campō, nam canis niger cum eō est et centum ovēs: ūna ovis nigra et ūndēcentum ovēs albae. Pāstor ūnam ovem nigram et multās ovēs albās habet. Is est dominus ovis nigrae et ovium albārum. Pāstor ovī nigrae et ovibus albīs aquam et cibum dat. Cum ūnā ove nigrā et ūndēcentum ovibus albīs pāstor in campō est.

Cibus ovium est herba, quae in campō est. In rivō est aqua. Ovēs in campō herbam edunt, et aquam bibunt ē \navnote{pānis}rivō, quī inter campum et silvam est. Canis herbam nōn ēst, neque pāstor herbam ēst. Cibus pāstōris est pānis, \navnote{ēst edunt vir pānem ēst viri pānem edunt}quī inest in saccō. Iūlius pāstōri suō pānem dat. Pāstor canī suō cibum dat: canis ā pāstōre cibum accipit. Ita-que canis pāstōrem amat.

In Italiā sunt multī pāstōrēs. Quī viā Appiā Rōmā 15 Brundisium it multōs pāstōrēs videt in campīs. Numerus pāstōrum magnus est. Dominī pāstōribus suīs ci-bum dant. Canēs ā pāstōribus cibum accipiunt.

\navnote{ūnus mōns ( \contr{} monts) duo montēs vallis --- \contr{} mōns ūna arbor duae arborēs ed TĒS lupus}Post campum montēs sunt. Inter montēs sunt vallēs. In campō est collis. (Collis est mōns parvus.) In colle 20 arbor est. In silvā multae arborēs sunt. Ovēs nōn in silvā neque in monte, sed in campō sunt. In silvā est lupus. Pāstor et ovēs lupum timent, nam lupus ovēs est. In silvīs et in montibus lupī sunt, in vallibus nūllī sunt lupī.

Sōl in caelō est suprā campum. In caelō nūlla nūbēs 25 \navnote{nūbēs}vidētur. Caelum est suprā terram. Montēs et vallēs, campī silvaeque sunt in terrā. In caelō sōl et nūbēs sunt, \navnote{ūna nūb*s duae nūbēs sub prp+abl \contr{} suprā}sed suprā hunc campum caelum sine nūbibus est. Ita-que sōl lūcet in campō.

Pāstor in sōle ambulat. Sub arbore autem umbra est. 30 Pāstor, quī nūllam nūbem videt in caelō, cum cane et ovibus ad arborem adit. Pāstor umbram petit. Etiam \navnote{umbram petit = ad um-}ovēs umbram petunt: post pāstōrem ad arborem adeunt. Pāstor ovēs suās ad arborem dūcit.

Ecce pāstor in umbrā arboris iacet cum cane et ovi- 35 bus. Arbor pāstōrī et canī et ovibus umbram dat; sed ovis nigra cum paucīs aliīs in sōle iacet. Pāstor, quī fessus est, oculōs claudit et dormit. Canis nōn dormit.

Dum pāstor in herbā dormit, ovis nigra ab ovibus albīs \navnote{pāstor in umbrā iacet ovēs albās relinquit = ab ovibus albīs abit}abit et ad rivum currit. Ovis nigra ovēs albās relinquit et 40 rivum petit; aquam bibit ē rīvō, et in silvam intrat!

Canis lātrat: ``Baubau!'' Pāstor oculōs aperit, ovēs aspicit, ovem nigram nōn videt. Pāstor ovēs numerat: ``Ūna, duae, trēs, quattuor, quīnque

ūndēcentum.`` Numerus ovium est undēcentum, nōn centum. Nūllae ovēs albae absunt, sed abest ovis nigra. Pāstor et canis ovēs albās relinquunt et silvam petunt. Pāstor saccum cum pāne in colle relinquit.

Dum cēterae ovēs ā pāstōre numerantur, ovis nigra in magnā silvā, ubi via nūlla est, errat. Ovis, quae iam procul ā pāstōre cēterisque ovibus abest, neque caelum neque sōlem suprā sē videt. Sub arboribus sōl nōn lūcēt. Ovis nigra in umbrā est.

In terrā inter arborēs sunt vestīgia lupī. Ubi est lupus \navnote{vestīgium lupus ipse, acc lupum ipsum}ipse? Nōn procul abest. Ovis vesūgia lupī in terrā videt, neque lupum ipsum videt. Itaque ovis lupum nōn ti-met. Parva ovis sine timōre inter arborēs errat.

\navnote{timor -ōris (gen) m \contr{} timet}Lupus autem prope ovem est. Pāstor et canis procul ab eā sunt. Lupus, qul cibum nōn habet, per silvam errat. Lupus in silvā cibum quaerit, dum pāstor et canis ovem nigram quaerunt.

Pāstor: ``Ubi est ovis nigra? Age! quaere ovem, canis, et reperi eam!'' Canis ovem quaerit et vestīgia eius in terrā reperit, neque ovem ipsam reperit. Canis lātrat.

\navnote{ovis īpsa, acc ovem īpsam dūcit dūcunt dūc! dūcite! (imp)}Pāstor: ``Ecce vestīgia ovis. Ubi est ovis ipsa? Dūc mē ad eam, canis!'' Canis dominum suum per silvam dūcit ad ovem, sed ovis procul abest.

Lupus ululat: ``Uhū!'' Et ovis et canis lupum audiunt. Canis currit. Ovis cōnsistit et exspectat dum lupus venit. Ovis bālat: ``Bābā!''

Ecce lupus quī ante ovem est! Iam ovis lupum ipsum ante sē videt. Oculī lupī in umbrā lucent ut gemmae et dentēs ut margaritae. Parva ovis oculōs claudit et dentēs

\navnote{ūnus dēns ( \contr{} dentsj duo dentēs}\navnote{ac-currit \contr{} ad-currit}Sed ecce canis accurrit! Lupus sē ab ove vertit ad 75 canem, quī ante lupum cōnsistit et dentēs ostendit. Lu-pus ululat. Canis lātrat. Ovis bālat.

\navnote{prope --- \contr{} procul clāmor -ōris (gen) m \contr{} clāmat p}Pāstor, quī iam prope est, clāmat: ``Pete lupum!'' Canis clāmōrem pāstōris audit, et sine timōre lupum petit. Lupus autem ovem relinquit et montēs petit.

Pāstor quoque accurrit et ovem suam, quae in terrā iacet, aspicit. In collō eius sunt vestīgia dentium lupī! Ovis oculōs aperit et ad pāstōrem suum bālat.

Pāstor laetus ovem in umerōs impōnit eamque portat ad cēterās ovēs, quae sine pāstōre in campō errant.

Procul in monte lupus ululat.

\booksection{Grammatica Latīna}

\navnote{pāstor ovem in umerōs im-pōnit ( \contr{} in-pōnit) dēclīnātiō -ōnis f (dēcl) \contr{} dēclīnat modus -Im: modō abl}Dēclīnātiō vocābulōrum [I] Dēclīnātiō prīma.

Vocābulum 'īnsula' dēclīnātur hōc modō:

Plūrahs Nōminānvus --- īnsulla īnsulae Accūsdtious --- īnsulam insulas Gencetivus īnsulae īnsulārum Datīvus īnsulae insulis 95 Ablātīvus

īnsula Accūsātīvus īnsulam īnsul ārum (ienetuus Dallvus Ablātīvus īnsulla \navnote{hōc modō: ut *īnsula'}femina, puella, fīlia, domina, ancilla, familia, littera, pecūnia, mēnsa, vīlla, aqua, via, silva, terra, cēt. [II] Dēclīnātiō secunda. Nōm.

[A] servus --- serv Acc. Gen.

Abl. [A] Ut 'servus' dēclīnantur vocābula masculīna: filius, dominus, fluvius, numerus, nummus, hortus, nāsus, mūrus, equus, saccus, umerus, amīcus, oculus, campus, rivus, lu-pus, cēt.; puer, vir, liber (nōm. sing. sine -ns). [B] Ut 'verbum' dēclīnantur vocābula neutra: oppidum, exemplum, baculum, ōstium, cubiculum, speculum, mālum, ōrnāmentum, collum, pretium, caelum, vestīgium, cēt. [III] Dēclīnātiō tertia. Nōm.

[A] pāstor

piāstorzs

[B] ovis ovēs Acc. Gen. pāstōr[is Dat.

[A] Ut 'pāstor' dēcllnantur haec vocābula: sōl sōlis, timor -ōris, clāmor -ōris (masc); arbor -oris (fem.); cēt. [B] Ut 4ovis' dēclīnantur haec vocābula: pānis, collis, vallis, canis (gen. plūr. canum); nūbēs -is (nōm. sing. -ēs); mōns montis, dēns dentis (nōm. sing. -s \contr{} -ts); et alia multa. Fēminlna sunt: ovis, vallis, nūbēs; masculīna: pānis, collis, mōns, dēns.

\booksection{Pēnsum A}

In Italiā sunt multī pāstōr-. Numerus pāstōr- magnus est. Pāstor Iūliī ūnum can- et multās ov- habet. Pāstor est dominus can- et ov-. Can- et ov- pāstōr- amant. Cibus ov- est herba, cibus pāstōr- est pān-. Pāstor pān- ēst.

In coll- ūna arbor est. Pastor cum can- et ov- ad arbor- it, Iam pāstor in umbrā arbor- iacet. Arbor pāstōr- et can- et ov- umbram dat, sed ov- nigra in sōl- iacet. Nūllae nūbante sōl- sunt. In silvā multae arbor- sunt, sub arbor- umbra est. Ov- nigra ā pāstōr- cēterisque ov- discēdit. Canis ovvidet.

\booksection{Pēnsum B}

Pāstor et centum --- in --- sunt. Pāstor ovibus aquam et --- dat. Cibus ovium est ---, cibus pāstōris est ---. Pāstor pānem ---. Ovēs herbam --- et aquam --- ē ---.

Sōl ---, nūlla --- in caelō --- hunc campum vidētur. In colle ūna --- est, in --- multae arborēs sunt. --- arboribus umbra est. Pāstor ovēs suās ad arborem ---. --- pāstor in --- arboris iacet, ovis nigra cēterās ovēs --- et silvam ---. In terrā sunt --- lupī; lupus --- nōn procul abest. Lupus in silvā cibum ---, dum pāstor et --- ovem quaerunt.

\booksection{Pēnsum C}

Num pāstor sōlus in campō est? Quot ovēs habet pāstor? Ā quō canis cibum accipit? Suntne montēs prope pāstōrem? Ubi sunt vallēs? Quid est collis? Quō it pāstor? Cūr pāstor umbram petit? Quō it ovis nigra? Quid ovis in terrā videt? Cūr lupus ovem nigram nōn ēst?

\booksection{Vocābula nova}
\noindent\small albus -a -um \contr{} niger, -gra -grum, plūr, sing, ovis, nōm, ovēs, acc, ovem, ovhm, ovts, gen --- ovis ovium, dat, ovi, ovibus, abl, ove, rivus, = parvus fluvius, nōm pāstor, pāstōrēs, acc pāstōrem pāstōrēs, gen pāstōris, pāstōrum, pāstōrf, pāstōribus, pāstōre, -ae, -um, -ās, -ārum, -ac, -īs, -ōrum, -ts, -is, -etn, -umJ-iim, -ibus, nova, campus, pāstor, canis, cibus, herba, pānis, mōns, vallis, collis, arbor, silva, lupus, sōl, caelum, terra, nūbēs, umbra, vestīgium, timor, dēns, clāmor, modus, niger, albus, ūndēcentum, ēst edunt, bibit, lūcet, petit, dūcit, iacet, relinquit, lātrat, errat, quaerit, reperit, ululat, bālat, accurrit, impōnit, ipse, procul, suprā, sub, dum, ut, dēclīnātiō, dēclīnat.

\bookchapter{Cap. X}{Bēstiae et hominēs}

\navnote{avēs in āere volant}\navnote{hominēs in terrā ambulant}\navnote{piscēs in aquā natant z RUN Viese \contr{} o Ita ``}Equus et asinus, leō et lupus, canis et ovis bēstiae sunt. \navnote{leō -ōnis m, pl leōnēs -ōnum}Leō et lupus sunt bēstiae ferae, quae aliās bēstiās capi\navnote{leō}unt et edunt. In Āfricā sunt multī leōnēs. Pāstōrēs Africae leōnēs timent, nam leōnēs nōn sōlum ovēs pāstōrum edunt, sed etiam pāstōrēs ipsōs! Nōn bēstia, sed homō \navnote{cauda homō -mis m, pl hominēs -inum fera = bēstia fera}est pāstor; leōnēs autem nōn sōlum aliās bēstiās, sed etiam hominēs edunt. Ferae et hominēs amīcī nōn sunt. Canis amīcus hominis est, ea bēstia fera nōn est.

Aliae bēstiae sunt avēs, aliae piscēs. Aquila est magna avis fera, quae parvās avēs capit et ēst. Avēs in āere \navnote{āēr āeris m pēs pedis m, pl pedēs -um Sed M y (ps āla d EE}volant. Piscēs in aquā natant. Hominēs in terrā ambulant. Avis duās ālās habet. Homō duōs pedēs habet. Piscis neque ālās neque pedēs habet. Avis quae volat ālās movet. Homō quī ambulat pedēs movet. Piscis quī natat caudam movet. Cum avis volat, ālae moventur. \navnote{cauda SEEN}Cum homō ambulat, pedēs moventur. Cum piscis na-tat, cauda movētur. Quī ambulat vestīgia in terrā facit. Qul volat aut natat vestīgia nōn facit.

In hortō et in silvā multae avēs sunt. Canis avēs, quae inter arborēs volant, aspicit. Canis ipse nōn volat, nam 20 canis ālās nōn habet. Canis volāre nōn potest. Neque \navnote{`` A --- quī ambulat vestīgia in terrā facit pot-est pos-sunt quod = quia petasus}pāstor volāre potest. Pāstor duōs pedēs habet, itaque pāstor ambulāre potest. Hominēs ambulāre possunt, quod pedēs habent, neque volāre possunt, quod ālās nōn habent.

\navnote{rcurius U Mercurius mercātor -ōris m imperium = id quod imperātur enim = nam īs nūntius = is quī verba portat (inter hominēs) i S}Mercurius autem ut avis volāre potest, nam in pedibus et in petasō eius ālae sunt. Mercurius nōn homō, sed deus est. Mercurius est deus mercātōrum. (Mercātor est homō quī emit et vēndit.) Mercurius imperia deōrum ad hominēs portat, is enim nūntius deōrum est. 30

Piscēs neque volāre neque ambulāre possunt. Piscēs in aquā natāre possunt. Etiamne hominēs natāre possunt? Aliī hominēs natāre possunt, aliī nōn possunt. Mārcus et Quīntus natāre possunt, Iūlia nōn potest.

Neptūnus natāre potest; is enim deus maris est. (Ōce-

IJ \navnote{Neptūnus mare -is n, pl maria -ium Rōmānī = hominēs Rōmānī Graecī = hominēs Graecī}anus est magnum mare.) Mercurius et Neptūnus deī Rōmānī sunt. Rōmānī et Graecī deōs multōs habent. Hominēs deōs neque vidēre neque audīre possunt. Deī ab hominibus neque vidēri neque audīri possunt.

\navnote{flūmen = fluvius flūmen -inis n, pl -ina -inum}Padus magnum flūmen est. In eō flūmine multī sunt 40 piscēs. Piscēs in aquā flūminis natant. In flūminibus et in maribus magnus numerus piscium est. Flūmina et maria plēna sunt piscium. Hominēs multōs piscēs capi\navnote{= nūllus homō}unt. Nēmō piscēs fluminum et marium numerāre po-test. Piscēs numerārī nōn possunt.

Piscēs in aquā vīvunt, neque in terrā vīvere possunt, nam piscēs in āere spīrāre nōn possunt. Homō sub aquā spīrāre nōn potest. Homō vīvit, dum spīrat. Quī spīrat \navnote{vīvus -a -um = qul vīvit}vīvus est, quī nōn spīrat est mortuus. Homō mortuus neque vidēre neque audīre, neque ambulāre neque cūrrere potest. Homō mortuus sē movēre nōn potest. Cum \navnote{pulmō}homō spīrat, anima in pulmōnēs intrat et rūrsus ex pulmōnibus exit. Anima est āēr quī in pulmōnēs dūcitur. Quī animam dūcit animal est. Nōn sōlum hominēs, sed \navnote{animam dūcere = spīrāre animal -ālis n, p! -ālia -ālium mare -is n, abl -i necesse est -}etiam bēstiae animālia sunt. Alia animālia in terrā vivunt, alia in marī.

Sine animā nēmō potest vīvere. Homō qul animam nōn dūcit vīvere nōn potest. Spīrāre necesse est hominī. \navnote{necesseest (+dat) ēst edunt: ēsse(= (--- edere) nēmō enim = nam nēmō}Ēsse quoque hominī necesse est, nēmō enim sine cibō vīvere potest. Necesse est cibum habēre. Pecūniam habēre necesse est, nam qul pecūniam nōn habet cibum emere nōn potest. Sine pecūniā cibus emī nōn potest. Necesse nōn est gemmās habēre, nēmō enim gemmās ēsse potest. Gemmae edī nōn possunt; sed quī gemmās suās vēndit pecūniam facere et cibum emere potest. Mercātor qul ōrnāmenta vēndit magnam pecūniam fa-cit. Fēminae quae pecūniam facere volunt ōrnāmenta sua vēndunt.

Avēs nīdōs faciunt in arboribus; nidi eārum inter rāmōs et folia arborum sunt. Aquilae in montibus nīdōs faciunt. In nidis sunt ōva. Avēs ōva pariunt, ex quibus \navnote{facere -it -iunt parere -it -iunt}parvī pullī exeunt. Canis nōn ōva, sed pullōs vīvōs pa-rit. Fēminae līberōs pariunt. ---

Mārcus et Qulntus et Iūlia in hortō sunt. Iūlia pilam tenet et cum pueris pilā lūdere vult, neque ii cum puellā 75 \navnote{vult volunt: Iūlia vult (--- volt) pueri volunt}lūdere volunt; pueri enim nīdōs quaerunt in arboribus. Itaque puella lūdit cum cane suā Margaritā.

\navnote{capere -it -iunt imp cape! capite!}Iūlia: ``Cape pilam, Margarīta!''

Canis pilam capit et caudam movet- Puella laeta ridet et canit. Pueri puellam canere audiunt.

\navnote{canere = cantāre vōx}Quīntus: ``Audī: Iūlia vōcem pulchram habet.''

Mārcus: ``Vōx eius nōn pulchra est!''

Canis avem suprā sē volāre videt eamque capere vult, neque potest. Canis īrātus lātrat: ``Baubau!'' Eius vōx pulchra nōn est, canis canere nōn potest! Avēs canere 85 possunt, piscēs nōn possunt: piscēs vōcēs nōn habent.

Cum aquila suprā hortum volat et cibum quaerit, \navnote{vūcis f. nōn audent: timent occultāre \contr{} \contr{} ostendere}parvae avēs canere nōn audent, et inter folia arborum sē occultant. Itaque aquila eās reperīre nōn potest. Neque avēs neque nidi avium ab aquilā reperīri possunt, quod 90 rāmīs et foliīs occultantur.

Mārcus autem nīdum reperit et Quīntum vocat: ``Veni, Quīnte! In hāc arbore nīdus est.'' Accurrit Quīntus.

Mārcus: ``Age! Ascende in arborem5 Quīnte!''

Quīntus in arborem ascendit; iam is suprā Mārcum 95 in arbore est. Mārcus ipse in arborem ascendere nōn audet! \navnote{ramus crassus}Mārcus interrogat: ``Quot sunt ōva in nīdō?'' Quīntus: ``Nūlla ōva, sed quattuor pullī/ \navnote{ramus tenuis}Nīdus est in parvō rāmō. Rāmus quī nīdum sustinet nōn crassus, sed tenuis est. Rāmus tenuis Quīntum sustinēre nōn potest, is enim puer crassus est. Ecce rāmus cum puerō et nīdō et pullīs ad terram cadit! Mārcus Quīntum ad terram cadere videt. Rīdetne Mārcus? Nōn ridet. Mārcus enim perterritus est. Iam Quīntus et pullī quattuor sub arbore iacent. Neque puer neque pullī sē movent. Pullī mortuī sunt. Quīntusne mortuus est? Nōn est. Quīntus enim spīrat. Quī spīrat mortuus esse nōn potest. Sed Mārcus eum spīrāre nōn \navnote{est sunt: esse neque enim = nōn enim facere = agere /--- A}videt, neque enim anima vidēri potest. Quid facit Mārcus? Mārcus perterritus ad vīllam cūrrit et magnā vōce clāmat: ``Age! Venī, pater!'' Iūlius puerum vocāre audit et exit in hortum. Pater filium perterritum ad sē accurrere videt eumque interrogat: ``Quid est, Mārce?''

\navnote{Mārcus perterritus est s}Iūlius: ``Quid? mortuus? Ō deī bonī!'' Pater, ipse perterritus, cum Mārcō ad Quīntum cūrrit. Iūlia quoque accurrit cum cane suā. Quīntus oculōs aperit. Iūlius eum oculōs aperīre videt. Iūlius: ``Ecce oculōs aperit: ergō vīvus est.'' Mārcus et Iūlia Quīntum vīvum esse vident. Puer autem ambulāre nōn potest, neque enim pedēs eum sustinēre possunt;

ergō necesse est eum

portāre. Quīntus ā Iūliō in vīllam portātur et in lectō pōnitur.

\navnote{PP}Aemilia filium suum ā Iūliō portāri videt, et interrogat: ``Cūr puer ipse ambulāre nōn potest?''

Iūlius: ``Quīntus ambulāre nōn potest, quod nōn est \navnote{Quīntus ā Iūliō portātur Aemilia Quīntum ā Iūliō portār? videt Quīntus in lectō pōnitur Aemilia Quīntum in lectō pōnī aspicit}avis neque ālās habet! Quī volāre vult neque potest, ad terram cadit!``

Aemilia Quīntum ā Iūliō in lectō pōnī aspicit.

\booksection{Grammatica Latīna}

īnfinītīvus [A] Āctīvum.

Iūlius Mārcum nōn vidēre, sed audire potest. Pater fīlium 135 vocāre audit, et accurrere videt.

Vocāre', vidēre', accurrere', audire' infinitivus est. īnfīnītīvus: -re. [1] -āre: cantāre, pulsāre, plōrāre, interrogāre, verberāre, numerāre, salūtāre, imperāre, habitāre, amāre, dēlectāre, por- 140 tāre, ambulāre, exspectāre, intrāre, ōrnāre, clāmāre, mōnstrāre, errāre, volāre, natāre, spīrāre; dare; cēt. [2] -ēre: ridēre, vidēre, respondēre, habēre, tacēre, pārēre, timēre, tenēre, iacēre, movēre, cēt. [3] -ere: pōnere, sūmere, discēdere, carpere, agere, vehere, claudere, vertere, currere, vēndere, emere, cōnsistere, ostendere, bibere, petere, dūcere, relinquere, quaerere, vīvere, lūdere, canere, ascendere, cadere; capere, facere, aspicere, accipere, parere; cēt. [4] -?re: venīre, audīre, dormīre3 aperīre, reperire; īre (ad-īre, 150 ab-Īre, ex-īre); cēt.

Mēdus servus esse nōn vult. Nēmō gemmās ēsse potest. 'Esse' quoque et *ēsse' infinitivus est. [A] Passīvum.

Mārcus ā Iūliō nōn vidēri, sed audiri potest. Mārcus Quīn- 155 tum ā Iūliō portān et in lectō pōnī videt.

Portār?,

vidēr?, pōn?, 'audīr? infinitivus passīvī est. Infinītīvus passīvī: -rīl-ī. [1] -ān: portārī, numerāri, vocāri. [2] -ēri: vidēri, tenērī. [3] -f; pōnī, emī, edī, claudī. [4] -m: audīrī, reperīri, aperīri.

\navnote{= Syra Iūliam rosam tenēre videt = parva puella magnum ōstium neque aperīre neque claudere potest}Piscēs numerān nōn possunt. Syra rosam ā Iūliā teneri videt. Gemmae edī nōn possunt. Magnum ōstium ā parvā puellā neque aperin neque claudī potest. Dēclīnātiō tertia

'Avis' (f) et 'piscis' (m) dēclīnantur ut 'ovis'. Ut 'pāstor' dēclīnātur 'mercātor' (m).

'Leō' (m) et 'homō* (m) dēclīnantur hōc modō:

Sing. Plār. Nōm. ]leō leonas homō hominēs Acc. --- leōnem leōnēs hominem --- hominēs Gen. leons leōnum hominis hominum Dat. leoni leōn[ibus Nām. \navnote{leō}leŌ Acc.

leōn*m leōn5 Gen.

leonu Dat.

leōnf

eōnibus AbL

Abī. leonle --- leōnibus homine hominibus

'Vōx' (f) et pēs'

(m) dēclīnantur hōc modō: Nōm.

vōx \navnote{pēs}Acc Gen. Dat.

vōcī Abl.

vōce

\booksection{Pēnsum A}

Av- in āēr- volant. Pisc- in aquā natant. Iūlia neque volneque nat- potest. Homō duōs ped- habet, itaque homō ambul- potest. Homō mortuus sē mov- nōn potest. Spīr- neces se est homini, nam sine animā nēmō vīv- potest. Cum homō spīrat, anima in pulmōn- intrat et ex pulmōn- exit. Homō quī spīrat mortuus es- nōn potest. Homin- deōs vid- nōn possunt. Deī ab homin- vid- nōn possunt. Nēmō piscēs numer- potest. Piscēs numer- nōn possunt. Sine pecūniā cibus em- nōn potest.

Puerī Iūliam can- audiunt. Mārcus Quīntum ad terram cad- videt. Iūlius Mārcum clām- audit. Aemilia saccum ā Iūliō in mēnsā pōn- et aper- videt.

\booksection{Pēnsum B}

Leō et aquila --- sunt. Viri et feminae --- sunt. Mercurius nōn homō, sed --- est. In aquā sunt ---. In āere sunt ---. Quid agunt piscēs et avēs? Piscēs in aquā ---, avēs in āere ---. Avis duās --- habet, itaque avis volāre ---. Avis quae volat ālās---. Cum homō ambulat, --- moventur. Cum homō ---, anima in --- intrat. Homō quī spīrat --- est, quī nōn spīrat--- --- est. Nam sine animā --- vīvere potest. Spīrāre homini --- est.

Iūlia --- [= cantat]. Iūlia --- pulchram habet. Pueri --- quaerunt. Nidi sunt inter --- et folia arborum. In nidis avium sunt --- aut pullī. Avēs nōn pullōs vīvōs, sed ōva ---. Quīntus in arborem --- et iv --- videt in nīdō. Rāmus quī nīdum --- tenuis est. Rāmus tenuis puerum --- sustinēre nōn potest: Qulntus ad terram ---. Mārcus eum cadere videt et --- est.

\booksection{Pēnsum C}

Num Neptūnus homō est? Quis est Mercurius? Quid agunt mercātōrēs? Num necesse est margaritās habēre? Quid est ōceanus Atlanticus? Cūr aquila ā parvīs avibus timētur? Ubi sunt nidi avium? Quid est in nidis? Quae bēstiae ōva pariunt? Quid agunt pueri in hortō? Cūr rāmus Quīntum sustinēre nōn potest?

\booksection{Vocābula nova}
\noindent\small infinitivus - 17 (īnf), vocāre, vidēre, accurrere, audīre, -re, esse, ēsse, portāri, videri, pōnī, audiri, \contr{} peds Vocābula nova, asinus, leō, bēstia, homō, fera, avis, piscis, aquila, āēr, āla, pēs, cauda, petasus, deus, mercātor, nūntius, mare, flūmen, anima, pulmō, animal, nīdus, rāmus, folium, ōvum, pullus, pila, vōx, lectus, ferus, vīvus, mortuus, crassus, tenuis, perterritus, capere, volāre, natāre, movēre, facere, vīvere, spīrāre, parere, lūdere, canere, audēre, occultāre, ascendere, sustinēre, cadere, potest possunt, vult volunt, necesse est, nēmō, cum, quod, enim, ergō, infinitivus.

\bookchapter{Cap. XI}{Corpus hūmānum}

caput

\navnote{capillus frōns oculus auns gena j auris oōsVz labrum'' collum hūmānus -a -um \contr{} homō corpus -oris n, pl -ora -ōrum crūs crūns n, pl crūra -rum ūna manus, duae manūs}Corpus hūmānum quattuor membra habet: duo bracchia et duo crūra. Bracchium membrum est et crūs membrum est. In bracchiō est manus, in crūre pēs. Duae manūs et duo pedēs in corpore hūmānō sunt.

In corpore ūnum caput est, nōn duo capita. In capite \navnote{caput -itis n, pl capita -itum super prp-- acc \contr{} sub frōns -onus f īnfrā prp+acc \contr{} sub}sunt oculī et aurēs, nāsus et ōs. Super caput capillus est. Capillus virōrum nōn tam longus est quam feminārum. Suprā oculōs frōns est. īnfrā oculōs genae sunt. Post \navnote{frōns-ontis/ īnfrāprp+acc \contr{} suprā}frontem est cerebrum. Qul cerebrum parvum habet stultus est. Ōs est inter duo labra. In ōre lingua et dentēs īnsunt. Dentēs sunt albī ut margarītae. Lingua et labra rubra sunt ut rosae.

\navnote{ruber -bra -brum}Hominēs oculīs vident et auribus audiunt. Homō quī oculōs bonōs habet bene videt, quī oculōs malōs habet male videt. Quī aurēs bonās habet bene audit, quī aurēs malās habet male audit. Syra male audit, ea enim aurēs malās habet.

\navnote{pectus -oris n cor cordis n sanguis -inis m fluere color -ōris m iecur -oris n venter -tris m}Caput est super collum. Sub collō est pectus. In pec-tore cor et pulmōnēs sunt. In corde est sanguis, quī per vēnās ad cor fluit. Color sanguinis est ruber. īnfrā pul- 20 mōnēs est iecur et venter. In ventre cibus est. Cor, pulmōnēs, iecur, venter sunt viscera hūmāna..

Homō quī ventrem mālum habet cibum sūmere nōn \navnote{cibum sūmere = ēsse}\navnote{-gra -grum}potest, neque is sānus, sed aeger est. Homō sānus ventrem bonum, pulmōnēs bonōs, cor bonum habet.

Medicus ad hominem aegrum venit eumque sānum facit. Medicus est vir quī hominēs aegrōs sānat, sed \navnote{sānāre = sānum facere}multī aegrī ā medicō sānārī nōn possunt. / ``d / y

\navnote{pulmō}\navnote{cor jecur NM MPa er /venter viscera hūmāna}Estne sānus Quīntus? Nōn est: pēs eius aeger est. Puer super lectum iacet. Aemilia et Syra apud puerum aegrum sunt. Māter apud eum sedet manumque eius tenet. Syra nōn sedet, sed apud lectum stat. Quīntus iacet. Aemilia sedet. Syra stat.

Aemilia: ``Ecce mālum, Quīnte.''

Aemilia puerō aegrō mālum rubrum dat, neque is 35 \navnote{pōculum aquae = pōculum cum aquā modo = tantum, sōlum}mālum ēsse potest. Māter ei pōculum aquae dat.

Aemilia: ``Bibe aquam modo!'' Māter caput Quintu sustinet, dum puer aquam bibit. Aemilia: Iam dormī, Quīnte! Dormī bene!`` \navnote{tangere}Māter manum pōnit in fronte filii: frontem eius tan-git. Quīntus oculōs claudit atque dormit. I! lūlius, quī cum Syrō servō in ātriō est, imperat: *Í ad oppidum, Syre, atque medicum arcesse!`` Medicus Tusculī habitat. Iūlius servum suum Tusculum īre iubet atque medicum arcessere. Syrus equum \navnote{iubēre = imperāre equum ascendere = in equum ascendere re-venīre = rūrsus venīre aegrōtāre = aeger esse}ascendit, ad oppidum it, medicum arcessit. Servus cum medicō ad vīllam revenit. Medicus interrogat: ``Quis aegrōtat?'' Iūlius: ``Meus Quīntus filius aegrōtat; ambulāre nōn potest.'' Medicus: ``Cūr ambulāre nōn potest?'' Iūlius: ``Quia pēs eius aeger est. Puer stultus est, me-dice: nīdum in arbore reperit, arborem ascendit, dē arbore cadit! Itaque pedem aegrum habet nec ambulāre \navnote{nec = neque pede aeger est = pedem aegrum habet}potest. Nec modo pede, sed etiam capite aeger est.`` Iūlius medicum ad cubiculum Quīntī dūcit. Medicus \navnote{cubiculum intrāre = in cubiculum intrāre quiētus -a -um = quī sē nōn movet}cubiculum intrat, ad lectum adit atque puerum aspicit. Quīntus quiētus super lectum iacet nec oculōs aperit. Medicus puerum dormīre videt. Medicus dīcit: ``Puer dormit.'' \navnote{dīcere id quod medicus dīcit = verba medicī}Syra, quae male audit, id quod medicus dīcit audīre nōn potest; itaque interrogat: ``Quid dīcit medicus?'' Aemilia (in aurem Syrae): ``Medicus 'puerum dormīre' dīcit.'' Quīntus oculōs aperit atque medicum adesse videt. \navnote{ad-esse verbum facere = verbum dīcere `` NN A as EN ODEQJAM}Puer, quī medicum timet, nūllum verbum facere audet. Medicus: ``Ōs aperī, puer! Linguam ostende!'' Syra: ``Quid dīcit medicus?'' Aemilia: ``Medicus Quīntum ōs aperire atque lin-guam ostendere iubet.'' Qulntus ōs aperit atque medicō linguam ostendit. Medicus linguam eius rubram esse videt. Medicus: ``Lingua eius rubra est.'' Syra: ``Quid dīcit?'' Aemilia: ``Dīcit 'linguam eius rubram esse'.'' Medicus etiam dentēs Quīntī spectat et inter dentēs \navnote{spectāre = aspicere}albōs ūnum nigrum videt. Nōn sānus est dēns qul colōrem nigrum habet. Medicus: ``Puer dentem aegrum habet.'' Qulntus: ``Sed dēns nōn dolet; ergō dēns aeger nōn 80 est. Pēs dolet --- et caput.'' Syra: ``Quid dīcunt?'' Aemilia: ``Medicus dīcit 'Quīntum dentem aegrum habēre', et Quīntus dīcit 'pedem et caput dolēre, nōn dentem'.'' Iūlius: Nōn dentem, sed pedem modo sānā, me-dice!`` Medicus pedem Quīnt spectat atque digitum ad pe-dem appōnit: medicus pedem eius tangit. Puer digitum \navnote{ap-pōnere \contr{} ad-pōnere sentīre}medicī in pede suō sentit. Quīntus: ``Ei, ei! Pēs dolet!'' Il

Medicus (ad Iūlium): ``Tenē bracchium puerī!'' (ad Aemiliam:) ``Tenē pōculum sub bracchiō!'' (ad Quīntum:) ``Claude oculōs, puer!'' Medicus Quīntum oculōs claudere iubet, quod puer cultrum medicī timet.

Ecce medicus cultrum ad bracchium pueri appōnit. Perterritus Qulntus cultrum medicī sentit in bracchiō, nec oculōs aperire audet. Capilli horrent. Cor palpitat. \navnote{en TN v ] DN NS, 9 P `` SN --- horrēre: capilli horrent = capilli stant}Medicus vēnam aperit. Ruber sanguis dē bracchiō in pōculum fluit. Quīntus sanguinem dē bracchiō fluere sentit atque horret. Frōns et genae albae sunt ut līlia

\navnote{horrēre = perterritus esse}Medicus puerum oculōs aperire iubet: Aperi oculōs, puer!`` neque Quīntus oculōs aperit. Puer quiētus super lectum iacet ut mortuus.

Syra: ``Cūr Quīntus oculōs nōn aperit? --- Ō deī bonī! Puer mortuus est!''

Quīntus autem spīrat, ergō mortuus nōn est. Sed Syra eum mortuum esse putat, quod eum splrāre nōn audit. Iūlius et Aemilia fīlium suum quiētum spectant --- atque tacent.

Medicus manum super pectus pueri impōnit eumque spīrāre et cor eius palpitāre sentit.

Medicus: ``Puer spīrat et cor eius palpitat.''

\navnote{gaudēre = laetus esse}Aemilia gaudet quod filius vīvit.

Syra: ``Quid dīcit medicus?''

Aemilia: ``Medicus dīcit 'Quīntum spīrāre et cor eius palpitāre.' Ergō Quīntus vīvit.''

\navnote{Syra Quīntum vīvere gaudet = Syra gaudet quod Quīntus vīvit}Syra Quīntum vīvere gaudet. Aemilia sanguinem dē bracchiō filii dēterget. Iam \navnote{dē-tergēre}puer oculōs aperit. Quīntus: ``Ei! Dolet bracchium!'' Māter fīlium vīvum esse videt et audit. Aemilia imperat: ``Aquam arcesse, Syra!'' neque ancilla verba dominae audit. Iūlius: ``Domina tē aquam arcessere iubet, Syrā!'' Syra abit, atque revenit cum aliō pōculō aquae plēnō. Aemilia pōculum tenet, dum Quīntus bibit. Medicus: ``Iam necesse est puerum dormīre.'' \navnote{horrēre = timēre ab-esse noster -tra -trum = meus}Exit medicus. Quīntus, quī medicum horret, eum abesse gaudet. Iūlius: ``Iam filius noster nōn modo pede, sed etiam bracchiō aeger est.'' Aemilia: ``Ille medicus crassus filium nostrum sānāre nōn potest.'' Aemilia nōn putat medicum puerum ae-grum sānāre posse. Syra: ``Stultus est medicus! Neque cor neque cerebrum habet!'' Syra medicum stultum esse' dīcit. Iūlius et Aemilia eum stultum esse putant, nōn dīcunt.

\booksection{Grammatica Latīna}

Dēclīnātiō tertia [A/B] Masculīnum et feminīnum. Vocābula masculīna: pāstor, mercātor, clāmor, timor, color, sōl, āēr, venter, leō, pulmō, homō, pēs, sanguis; pānis, collis, piscis, mōns, dēns, cēt.

Vocābula feminīna: ovis, vallis, avis, auris, nūbēs, frōns; arbor, vōx, praepositiō, dēclīnātiō, cēt. [C] Neutrum.

Vocābula neutra: flūmen, ōs, crūs, corpus,

pectus, iecur, caput, cor; viscera (p/); mare, animal, cēt.

Plūrālis: -a. Accūsātīvus = nōminātlvus.

'Corpus' et 'flūmen' dēclīnantur hōc modō:

PlŪK Nōm.

corpus

corpora Acc.

corpus

corpora

flūmen flūmina Gen. corporis corporum flūminis flūminum Dat.

flūminum Gen.

corports

corporum

flūminis

ūūminibus Dat.

corporī

corport\&M5

flūminf

ūūm'mibus Abl.

Ut 'corpus' dēclīnantur: pectus -oris, ōs ōris, crūs crūris, iecur -oris (cor cordis, caput -itis); plūrālis: viscera -um.

Mare' et animal hōc modō dēclīnantur: Sing. ---

Plūr. Nōm.

marle

maria Acc. Gen.

maris

marium Dat.

Accūsātīvus cum īnfinītīvō

\navnote{Iūlia dormit Iūliam dormīre}Iūlia dormit. Syra Iūliam dormīre videt.

dlcit.

Syra: ``Iūlia dormit.'' Syra 'Iūliam dovmīre9

'Iūliam dormīre' est accūsātīvus cum īnfinītīvō. Accūsātīvus cum īnfinītīvō pōnitur apud multa verba: [1] vidēre, audire, sentīre: Puer medicum adesse videt. Pueri Iūliam canere audiunt. Medicus puerum spīrāre sentit. [2] iubēre: Dominus servum discēdere iubet. dīcere: Quīntus 'pedem dolēre' dīcit. [4] putāre: Syra Quīntum mortuum esse putat. [5] gaudēre: Aemilia filium vīvere gaudet. [6] necesse esse: Puerum dormīre necesse est.

\booksection{Pēnsum A}

Membra corpor- hūmaāni sunt duo bracchia et duo crūr-. In corpor- hūmānō ūn- caput est, nōn duo capit-. In capitsunt duae aur- et ūn- ōs. In ōr- sunt dent-. In pector- ūncor et duo pulmōn- sunt.

Medicus Quīnt- super lectum iac- videt; medicus puerdorm- videt. Medicus: Quīnt- dorm-.`` Medicus 'Quīntdorm-' dīcit. Medicus puer- linguam ostend- iubet, et 'lingu- eius rubr- es-' dīcit. Puer dīcit 'ped- et caput dol-.' Medicus Aemili- pōculum ten- iubet. Syra Quīnt- spīr- nōn audit, itaque Syra e- mortu- es- putat. Sed Quīnt- vīv-. Māter fili- vīv- gaudet. Necesse est puer- aegr- dorm-.

\booksection{Pēnsum B}

--- hūmānum habet quattuor ---: duo --- et duo ---. In bracchils duae --- sunt, in crūribus duo ---. Super collum est ---. In capite sunt duo oculī, duae ---, ūnus nāsus, ūnum ---. In ōre sunt dentēs et ---. Sub collō est ---. In pectore sunt pulmōnēs et ---. In corde et in vēnīs--- --- est. Sanguis per --- ad cor ---.

Aemilia apud lectum Quīnt ---, Syra apud lectum ---. Quīntus nōn ---, sed aeger est. Syrus medicum ex oppidō ---. Medicus digitum ad pedem pueri ---: medicus pedem eius ---. Quīntus, quī digitum medicī in pede ---: ``Ei! Pēs ---!''

\booksection{Pēnsum C}

Quae sunt membra corporis hūmāni? Ubi est cerebrum? Quid est in pectore? Ubi est venter? Cūr Quīntus cibum sūmere nōn potest? Estne Quīntus sōlus in cubiculō suō? Unde medicus arcessitur? Quid videt medicus in ōre Quīntī? Quid Quīntus in bracchiō sentit? Cūr Syra Quīntum mortuum esse putat?

\booksection{Vocābula nova}
\noindent\small potest, sunt, possunt, posse (īnf), -is, -um, -ibus, -ia, -ium, nova, corpus, membrum, bracchium, crūs, manus, caput, auris, ōs, capillus, frōns, gena, cerebrum, labrum, lingua, pectus, cor, sanguis, vēna, color, iecur, venter, viscera, medicus, pōculum, culter, hūmānus, stultus, ruber, sānus, aeger, noster, bene, male, fluere, sānāre, sedēre, stāre, tangere, arcessere, iubēre, revenire, aegrōtāre, dīcere, spectāre, dolēre, appōnere, sentīre, horrēre, palpitāre, putāre, gaudēre, detergere, posse, modo, super, īnfrā, dē, atque, nec.

\bookchapter{Cap. XII}{Mīles Rōmānus}

\navnote{pīlum: scūtum:}\navnote{hasta}\navnote{arma Rōmāna}I Quīntus est frāter Mārcī. Iūlia soror eius est. Mārcus et \navnote{frāter -tris m soror -ōris f pater -tris mn māter -tris f. nōmen -inis n}Quīntus frātrēs Iūliae sunt. Mārcus patrem et mātrem, frātrem et sorōrem habet. Nōmen patris est 'Iulius', mātris 'Aemilia'; 'Quīntus' est nōmen frātris, 'Iūlia' so-S rōris.

Mārcō una soror est. Iūliae duo frātrēs sunt. Nōmina \navnote{Mārcō (dat) ūna soror est = Mārcus ūnam sorōrem habet}frātrum sunt 'Mārcus' et 'Quīntus'. Patri et mātrī ūna fīlia et duo filiī sunt.

Matri Aemilia! nōmen est. Quod nōmen est patri? Eī nōmen est 'Lūcius Iūlius Balbus'. Virō Rōmānō tria \navnote{prae-nōmen -inis n cognōmen -inis n}nōmina sunt. 'Lūcius' est praenōmen, id est nōmen prīmum; 'Balbus' cognōmen est. Filiis nomina sunt 'Mārcus Iūlius Balbus' et 'Quīntus Iūlius Balbus'. 'Mārcus' \navnote{praenōmina Latīna:}et 'Quīntus' praenōmina sunt filiōrum. Alia praenōmina Latīna sunt 'Aulus', 'Decimus', 'Gāius', 'Pūblius', 'Sextus', Titus'. \navnote{R Quīntus Sex.}Aemiliae est ūnus frāter, cui 'Aemilius' nōmen est \navnote{avunculus -I m = frāter mātris trlstis \contr{} +laetus}(praenōmen 'Pūblius', cognōmen Paulus''). Frāter Aemiliae est avunculus līberōrum. Aemilius autem procul ā sorōre suā abest. Itaque trīstis est Aemilia, quae frā- 20 trem suum amat. Mārcus et Iūlia mātrem suam trīstem in hortum exīre vident et patrem interrogant: ``Cūr māter nostra tristis est?'' \navnote{vester -tra -trum (= tuus}Iūlius: ``Māter vestra tristis est, quod Aemilius procul 25 ab eā abest. Aemilius avunculus vester est, id est frāter mātris. Māter tristis est, quod frātrem suum vidēre nōn potest.'' Mārcus: ``Ubi est avunculus noster?'' Iūlius: ``Avunculus vester est in Germāniā. Aemilius 30 mīles est. In Germāniā multī sunt mīlitēs Rōmānī.'' \navnote{mīles -itis m}Iūlia: ``Quid est mīles?'' Iūlius: Mīles est vir quī scūtum et gladium et pīlum \navnote{fert = portat fert ferunt arma fert}fert. Scutum et gladius et pīlum sunt arma mīlitis Rōmānī. Mīles est vir armātus.`` Iūlia: ''Quid agunt mīlitēs Rōmānī in Germāniā?`` Iūlius: Mīlitēs nostri in Germāniā pugnant.'' Iulia: Meī quoque frātrēs pugnant.`` Iūlius: ''Puerī pugnīs, nōn armīs pugnant.

Mīlitēs pugnant gladiīs, pīlīs, hasūs.``

Mārcus: ``Num Aemilius et hastam et pīlum fert?''

Iūlius: ``Aemilius pīlum tantum fert, is enim pedes \navnote{pedes -itis m eques -ius m}est, nōn eques. Eques est mīles quī ex equō pugnat; quī pedibus pugnat pedes est. Equitēs hastās, peditēs pila ferunt. Pīlum nōn tam longum est quam hasta, neque gladius peditis tam longus est quam gladius equitis. Pīlum Aemiliī sex pedēs longum est.``

\navnote{D}Mārcus: ``Quam longus est gladius eius?''

Iūlius: ``Duōs pedēs longus est,'' 50

\navnote{brevis --- \contr{} longus}Mārcus: ``Duōs pedēs tantum? Cūr tam brevis est gladius?''

Iūlius: ``Quod gladius brevis nōn tam gravis est quam gladius longus. Gladius equitis longior et gravior est \navnote{gladius brevis nōn tam gravis est quam gladius longus ferre (īnf)= --- portāre levis \contr{} \contr{} gravis fertur feruntur Germānī -ōrum m pl}quam peditis. Pedes, quī pedibus it multaque alia arma fert, gladium longum et gravem ferre nōn potest; itaque gladius eius brevis et levis est --- brevior et levior quam is quī ab equite fertur. Etiam gladiī quī ā Germānīs feruntur longiōrēs et graviōrēs sunt quam Rōmānōrum ac pila eōrum longiōra et graviōra quam nostra sunt.``

\navnote{atque ante a,6,1,0,u,À acdatque ante cēterās litterās barbarus -a -um = nec Rōmānus nec Graecus incolere: terram i. = in terrā habitāre}Iūlia: ``Quī sunt Germānī?''

Iūlius: ``Germānī sunt hominēs barbari quī Germāniam incolunt. Germānia est magna terra nōn procul ā Galliā; Gallia autem prōvincia Rōmāna est, ut Hispānia, Syria, Aegyptus. Prōvincia est pars imperiī Rōmānī, ut membrum pars corporis est; Rōma enim caput imperiī est, prōvinciae membra sunt. Germānia nōn est prōvincia Rōmāna. Flūmen Rhēnus Germāniam ā Galliā prō\navnote{dividere}vinciā dīvidit. Rhēnus ac Dānuvius flūmina, quae Germāniam ab imperiō Rōmānō dīvidunt, fīnēs imperiī nostri sunt. Germānia est patria Germānōrum, ut Rōma 70 \navnote{patria -ae/: f: p. nostra trum nostrōrum contrāā prp- acc}nostra patria est.``

Iūlia: ``Cūr mīlitēs Rōmānī contrā Germānōs pug-nant? Suntne Germānī hominēs improbī?''

Iūlius: ``Mīlitēs nostri contrā Germānōs pugnant, quod Germānī amīcī Rōmānōrum nōn sunt nec Rōmā- 75 \navnote{pārēre/imperāre + dat hostis -is m \contr{} amīcus}nis pāārent. Germānī hostēs Rōmānōrum sunt, ac bellum est inter Germānōs et Rōmānōs. Germānī exercitum nostrum oppugnant.``

\navnote{op-pugnāre = pugnāre contrā}Iūlia: ``Quid est 'exercitus'?''

Iūlius: ``Exercitus est magnus numerus mīlitum quī 80 contrā hostēs dūcitur. Quī exercitum dūcit dux exercitūs est. Dux exercitui imperat, exercitus dūcī suō pāret, nam dux ab exercitū metuitur.

``In Germāniā et in Britanniā sunt magnī exercitūs Rōmānī quī contrā exercitus hostium pugnant. Mīlitēs 85 et ducēs exercituum Rōmānōrum ab hostibus metuuntur. In Hispāniā et in Galliā nōn multī sunt mīlitēs \navnote{Hispānī -ōrum m pl Gallī -ōrum m pl}Rōmānī, nam Hispānī et Gallī, quī eās prōvinciās inco\navnote{mīlitāre = mīles esse sagitta -ae f}lunt, iam exercitibus nostrīs pārent. In exercitibus Rōmānīs etiam Hispānī et Gallī multī mīlitant, quī et alia 90 arma et arcūs sagittāsque ferunt.

Iūlia: \contr{} Ubi habitat Aemilius?``

\navnote{arcus -ūsm}Iūlius: ``Aemilius in castris habitat mille passūs ā fine \navnote{castra -ōrum npl (vocābulum plūrāle tantum)}imperiī. Castra sunt mīlitum oppidum.``

\navnote{-īn}Mārcus: ``Quam longus est passus?''

\navnote{passus -ūs m: i passus}Iūlius: ``Ūnus passus est quīnque pedēs, ergō mille passūs sunt quīnque milia pedum. In castris Aemiliī sex \navnote{mille (M) hominēs; duo milia (MM/il) hominum pueri = pueri et puellae}mīlia mīlitum habitant. Nūllae feminae aut pueri illīc habitant, nec enim feminae puerique mīlitāre possunt. Circum castra fossa et vāllum longum et altum est.``

Mārcus: ``Quam altum est vāllum castrōrum.''

Iūlius: ``Prope decem pedēs altum est, et duo mīlia passuum longum. Quattuor portae per vāllum in castra dūcunt. Inter duās portās est via lāta, quae castra in duās partēs dīvidit; ea via centum pedēs lāta est.

\navnote{in bellō = cum bellum est dēfendere ++ oppugnāre iacere -it -iunt}``In bellō portae castrōrum clauduntur. Cum exercitus Germānōrum castra oppugnat, Rōmānī castra dēfendunt: vāllum ascendunt ac pila in Germānōs iaciunt. Illī autem nec pīla in castra iacere possunt, quod fossa no nimis lāta et vāllum nimis altum est, nec vāllum ascendere, quod Rōmānī pīlīs et gladiīs vāllum dēfendunt. \navnote{ex-pugnāre equitatus -ūs m = equitum numerus (pars exercitūs) impetus -ūsm: impetum facere in hostem = hostem armīs petere, oppugnāre fugere -it -iunt}Hostēs castra nostra expugnāre nōn possunt.

``Ecce portae aperiuntur atque equitātus noster in hostēs

impetum

facit.

Barbarī perterritī,

quī impetum equitātūs sustinēre nōn possunt, arma ad terram iaciunt atque in magnās silvās fugiunt.

``Mīles Rōmānus, quī hostem armātum accurrere videt, nōn ab eō fugit, sed armīs sē dēfendit. Mīlitēs Romānī fortēs sunt. Mīles fortis hostem nōn metuit, sed sine metū impetum in hostem facit. Aemilius, avuncu- 120 lus vester, mīles fortis est.

``Circum imperium Rōmānum multī sunt hostēs. Castra et oppida nostra ab hostibus oppugnantur neque expugnantur, nam mīlitēs nostri prōvinciās ac patriam nostram ā Germānīs et ab aliīs hostibus dēfendunt.'' ---

Mārcus: ``Etiam Germānī suam patriam dēfendunt.''

īūlius: ``Sed patria nostra pulchrior est quam illōrum! Atque Germānī hominēs barbarī sunt.''

Mārcus: ``Nōnne fortēs sunt Germānī?''

\navnote{impetus hostium}Iūlius: ``Fortēs sunt illī, sed Rōmānī fortiōrēs sunt, nec arma Germānōrum tam bona sunt quam nostra. Scutum eōrum nimis parvum est, pīlum nimis longum et grave; nec enim pīlum tam grave procul iaci potest. Itaque pīlum nostrum breve et leve est --- brevius et levius quam pīlum Germānōrum. Mīlitēs Rōmānī bene 135 pugnant, quod pila eōrum brevia et levia sunt, nōn lon-ga et gravia ut Germānōrum. Patria nostra bonīs armīs dēfenditur. Nūllus hostis Rōmam expugnāre potest.''

Mārcus: ``Cūr Rōmānī Germāniam nōn expugnant?''

Iūlius: ``Germānia nōn sōlum armīs dēfenditur, sed 140 etiam altīs montibus, magnīs silvīs atque lātīs et altīs flūminibus.''

\booksection{Grammatica Latīna}

Dēclīnātiō quārta 145

\navnote{versus -ūs m: cap. xil 142versūs habet}'Exercitus' (m) dēclīnātur hoc modō (vidē exempla in versibus 80-89 huius capitull):

Nōminātīvus

exercitus Accūsātīvus Genetīvus Datīvus Ablātīvus

Ut 'exercitus' dēclīnantur masculīna: arcus, passus, equitātus, impetus, metus, versus, cēt.; femininum: manus. 155 Adiectīvum [III] \navnote{ad-iectīvum -īn (adi)}Dēclīnātiō prīma et secunda.

Mōns altus. Arbor alta. Vāllum altum.

``Altus -a -um' est adiectīvum

primae et secundae dēclīnatoōnis. Masculīnum altus'! dēclīnātur ut 'servus', feminīnum 'alta' ut 'femina', neutrum 'altum' ut 'oppidum'. Singulāris

Singulāris

Masc. Fēm. Neutr. Nōm.alus ---

alta Nōm.

Acc.

altum Acc.

altarum altlorum Dat. Abl.

altjō

Hōc modō dēclīnantur haec adiectīva: albus, bonus, fessus, foedus, īrātus, laetus, lātus, longus, magnus, malus, meus, novus, plēnus, primus, sānus, tuus, cēt.; (-er -ra -rum:) pulcher, aeger, niger, ruber, noster, vester, cēt. [III] Dēclīnātiō tertia.

Gladius brevis. Via brevis. Pīlum breve.

\navnote{brevis breve}Brevis -e' est adiectīvum tertiae dēclīnātiōnis. Masculīnum et feminīnum 'brevis' dēclīnātur ut ovis' (sed abl. sing. -f), neutrum 'brev' ut mare'. Nōm.

brevia Acc.. brevem brevle breves brevia Gen.

brevis Abl.

Hoc modō dēclīnantur haec adiectīva: brevis, fortis, gra-vis, levis, tenuis, trīstis, cēt.

Exempla:

Pedes fors Rōmānus gladium brevem et levem, scūtum magnum et grave fert; eques hastam fert longam et gravem. Peditl fortī Rōmānō est gladius brevis et levis, scūtum magnum et grave; equitī hasta longa et gravis est. Pedes gladiō 190 brevī et levi, scūtō magnō et gravī armātus est; eques hastā longa et gravi pugnat.

Peditēs forts Rōmānī gladiōs breves et levēs, scūta magna et gravia ferunt; equitēs hastās ferunt longās et gravēs. Peditibus fortibus Rōmānīs sunt gladiī breves et levēs, scūta magna et gravia; equitibus hastae longae et gravēs sunt. Peditēs gladiīs brevibus et levibus, scūtīs magnīs et gravibus armātī sunt; equitēs hastīs longīs et gravibus pugnant.

Comparātīvus

\navnote{comparātīvus (comp)}Hic mūrus altior est quam ille.

Hoc vāllum altius est quam illud.

\navnote{altior altius}Aluor -ius

comparāātivus est. Comparātivus est adiectivum dēclīnātiōnis tertiae:

Nōm.

altior Acc. aluōrem altius Gen. ahiōris Dat. Abl. altuōre

Exempla: brevior, fortior, gravior, levior, longior, pulchrior.

Gladius equitis longior et gravior est quam peditis. Eques gladium longiorem et graviōrem fert quam pedes. Eques gladiō longiōre et graviōre pugnat.

Gladiī equitum longiōrēs et graviōrēs sunt quam peditum. Equitēs gladiīs longiōribus et graviōribus pugnant.

Pīlum brevius et levius est quam hasta. Pila breviōra et leviōra sunt quam hastae.

Dēlia ancilla pulchrior

est quam Syra. Lēander, quī ancillae pulchriōris amīcus est, ancillae pulchr:ōr? rosam dat. Quid pulchrius est quam rosa?

Mīlitēs Rōmānī fortiōrēs sunt quam hostēs. Dux Rōmānus mīlitibus fortiōribus imperat quam dux hostium. Ille dux militum fortiōrum est.

\booksection{Pēnsum A}

Pīlum ē man- iacitur, sagitta ex arc-. In exercit- Rōmānīs multī Gallī mīlitant, quī arc- ferunt. Equitāt- sine metimpet- in hostēs facit, neque hostēs impet- equitāt- sustinēre possunt. Mille pass- sunt quīnque milia pedum. Via Latīna CL[150: [150: centum quīnquāgintā] milia pass- long- est.

Rāmus tenu- puerum crass- sustinēre nōn potest, nam puer crass- grav- est. Ūnum mālum grav- nōn est, nec duo māla grav- sunt, sed saccus plēn- mālōrum grav- est. Lēander saccum magn- et grav- portat In saccō grav- sunt māla. Servī saccōs grav- portant. In saccīs grav- sunt māla et pira.

Saccus Lēandri grav- est quam Syri, nam māla grav- sunt quam pira. Via Appia long- est quam via Latina. Via Appia et via Aurēlia long- sunt quam via Latīna. Quīntus crass- est quam Mārcus. Lēander saccum grav- quam Syrus portat. In saccō grav- māla sunt.

Hoc pīlum long- et grav- est quam illud. Haec pila long-et grav- sunt quam illa.

\booksection{Pēnsum B}

Mārcus --- Quīntī est. Quīntus ūnum frātrem et ūnam --- habet. 'Iūlia' --- sorōris est. Frāter Aemiliae est --- līberōrum. Quīntus: ``Ubi est avunculus noster?'' Iūlius: ``Avunculus --- est in Germāniā.''

Prōvincia est --- imperiī Rōmānī. Rhēnus Germāniam ab imperiō Rōmānō ---; Rhēnus --- imperiī est. In Germāniā mulū --- Rōmānī sunt, quī contrā Germānōs ---; Germānī enim --- Rōmānōrum sunt. Mīlitēs Rōmānī scūta et --- et pīla ---. Mīlitēs in --- habitant. Circum castra est --- et ---, quod x pedēs --- est. Mīlitēs patriam ab hostibus ---.

\booksection{Pēnsum C}

Num 'Mārcus' cognōmen est? Quot frātrēs habet Aemilia? Quid agit Aemilius in Germāniā? Quae arma pedes Rōmānus fert? Quam longum est pīlum Aemiliī? Ubi habitant mīlitēs Rōmānī? Quī sunt Germānī et Gallī? Estne Germānia prōvincia Rōmāna? Quod flūmen Germāniam ā Galliā dīvidit? Cūr hostēs castra expugnāre nōn possunt? Num mīles fortis ab hoste fugit? Cūr hasta procul iaci nōn potest?

\booksection{Vocābula nova}
\noindent\small sing nōm exercitus, acc, exercitum, gen, exercitūs, dat, exercitui, abl, exercitā, dux ducis m, metuere = timēre, plūr, nōm exercitās, exercituum, exercitibus, -ns, -um, -uum, -ūs, -ibus, -uī, altus, alta, altum, -am, -ae, -ās, -ōs, a, -ōrum, -ārum, -is, -īs, ĪS, -er, -Ta, -rum, -ia, -es, is, -em, -ium, -l0r -tus -lōrem, -ius, -iōnis, -toris, -Mri, -iōre, -tore, -iōrēs, -iōra, -iōrēs -iōra -ioōrum, -iōrum, -ribus --- -ionibus -ribus, -iōribus, -idribus, nova, frāter, soror, nōmen, praenōmen, cognōmen, avunculus, mīles, scūtum, gladius, pīlum, arma, arcus, sagitta, pugnus, hasta, pedes, eques, pars, finis, patria, hostis, bellum, exercitus, dux, castra, passus, fossa, vāllum, equitātus, impetus, metus, versus, mīlia, trīstis, armātus, brevis, gravis, levis, barbarus, lātus, fortis, vester, ferre, pugnāre, incolere, dividere, oppugnāre, metuere, mīlitāre, dēfendere, iacere, expugnāre, fugcre, contrā, ac, adiectīvum, comparatīvus.

\bookchapter{Cap. XIII}{Annus et mēnsēs}

\navnote{Iān. Capricornus Febr Aquārius Mārt. Piscēs Apr. Ariēs Māi. Taurus Iūn. Geminī Cancer Aug. Leō Sept. Virgō Oct. Lībra Nov. Scorpiō Dec. Sagittārius}MENSIS MEIMSIS DIĒS DIĒS XXXI. DIĒS XXx DIESXXXI

MENSIS MENSIS MENSIS MENSIS IVLIVS AVGVSTVSSEPTEMBER OCTŌ XXL DIĒS XXxt DIĒS XXX

I Annus in duodecim mēnsēs dīviditur, quibus haec sunt \navnote{v vi sextus -a-um VH septimus -a -um vii septimus-a-um viii octāvus -a -um IX x xi XIl duodecimus -a -um \contr{} primus}nōmina: Iānuārius, mēnsis primus; Februārius, secundus; Mārtius, tertius; Aprilis, quārtus; Māius, qulntus; Iūnius, sextus; Iūlius, septimus; Augustus, octāvus; September, nōnus; Octōber, decimus; November, ūndecimus; December, mēnsis duodecimus ac postrēmus.

Ūnus annus duodecim mēnsēs vel trecentōs sexāgintā quīnque diēs habet. Saeculum est centum annī. Centum annī vel saeculum est longum tempus. Duo saecula sunt \navnote{tempus -oris n}ducentī annī. Homō sānus nōnāgintā vel etiam centum annōs vīvere potest; ducentōs annōs vīvere nēmō po-test.

Mēnsī prīmō et mēnsī tertiō ā deīs nōmina sunt: Iānuāriō ā deō Iānō, Mārtiō ā deō Mārte. Iānus et Mārs sunt deī Rōmānī. Iānus est deus cui duae faciēs sunt. \navnote{Mārs}Mārs deus bellī est.

\navnote{nōmināre \contr{} nōmen; nōminātur- nōmen habet}Mēnsis September nōminātur ā numerō septem, Octōber, November, December ab octō, novem, decem. Nam tempore antīquō Mārtius mēnsis primus erat. Tunc \navnote{antiquus \contr{} ante tunc = illō tempore}September mēnsis septimus erat, Octōber, November, 20 December

mēnsēs octāvus, nōnus,

decimus erant. Nunc autem mēnsis prīmus est Iānuārius, September \navnote{nunc = hōc tempore ...igitur = ergō}igitur mēnsis nōnus est, Octōber decimus, November ūndecimus, December duodecimus.

Quam longus est mēnsis November? November tri- 25 gintā diēs longus est. December ūnum et trigintā diēs \navnote{ūnus et trīgintā}habet. Iānuārius tam longus est quam December, sed \navnote{duo-dē-trigintā ūn-dē-trīgintā}Februārius brevior est: duodētrīgintā aut ūndētrīgintā diēs tantum habet. Februārius brevior est quam cēterī undecim mēnsēs: is mēnsis annī brevissimus est. Mār- 30 \navnote{Mārtius)}tius ūnum et trigintā diēs longus est (et item Māius, Iūlius, Augustus, Octōber). Aprīlis trigintā diēs habet (item Iūnius et September). Sex mēnsēs sunt dīmidia \navnote{dīmidia pars = Vi quārta pars = lA}pars annī, trēs mēnsēs quārta pars annī.

Diēs est dum sōl in caelō est. Prima pars diēī est 35 māne, pars postrēma vesper. Diēs est tempus ā māne ad vesperum. Nox est tempus ā vesperō ad māne. Vesper \navnote{nox noctis/ \contr{} f diēs}est finis diēī atque initium noctis. Māne finis noctis est \navnote{finis = pars postrēma initium = pars prima}atque initium diēt. CONES ATNSS

Diēs in duodecim hōrās dīviditur Ab hōrā prīmā 40 diēs initium facit. Hōra sexta, quae hōra media est inter hōram primam et duodecimam, 'meridiēs' nōminātur. \navnote{: medius diēs}Hōra sexta vel meridiēs diem dīvidit in duās aequās partēs: 'ante merīdiem' et 'post merīdiem Meridiē sōl altissimus in caelō est. Sex hōrae sunt dīmidia pars diēī. Il

Nocte sōl nōn lūcet, sed lūna et stēllae lucent. Lūna \navnote{stēlla-aef lūxlūcis/}ipsa suam lūcem nōn habet, lūx lūnae ā sōle venit; ita-que lūna nōn tam clāra est quam sōl. Sōl est stēlla clārissima, quae lūce suā et terram et lūnam illūstrat. Neque \navnote{\contr{} pars lūnae obscūrus -a -um \contr{} eclārus exiguus -a -um = parvus}tōta lūna sōle illūstrātur, sed tantum ea pars quae vertitur ad sōlem; cētera pars obscūra est. Cum exigua pars lūnae tantum vidētur, lūna 'nova' esse dīcitur. Diē septimō vel octāvō post lūnam novam lūna dīmidia vidētur, quae fōrmam habet litterae D. Diē quīntō decimō post \navnote{lūna nova dīmidia plēna quīntus decimus}lūnam novam lūna plēna est et formam habet litterae O. Cum lūna nōn lūcet, nox obscūra est.

\navnote{dīcitur = nōminātur}Diēs mēnsis prīmus 'kalendae' nōminātur. Diēs primus mēnsis Ianuāārii dīcitur 'kalendae Iānuāriae'; is diēs annī prīmus est atque initium annī novī. Diēī prīmō mēnsis Iūliī kalendae Iūliae' nōmen est; is est diēs annī medius.

\navnote{tertius decimus īdūs}Diēs tertius decimus post kalendās 'īdūs' nōminātur. 'īdūs Iānuāriae' diēs tertius decimus est post kalendās Iānuāriās. Item 'īdūs Februāriae' dīcitur diēs tertius decimus mēnsis Februāriī. Sed 'īdūs Mārtiae' diēs est qulntus decimus mēnsis Mārtiī, nam mēnse Mārtiō \navnote{-uum f pl (mensis) -ber -bris m, abl -bri}(item Māiō, Iūliō, Octobri) īdūs nōn diēs tertius decimus, sed quīntus decimus post kalendās est. Diēs nōnus ante īdūs dīcitur 'nōnae' (nōnae Iānuāriae: diēs quīntus Iānuāriī; nōnae Februāriae: diēs quīntus Fe- 70 bruāriī; nōnae Mārtiae: diēs septimus Mārtiī; cēt.)

Diēs octāvus ante kalendās Iānuāriās, quī dīcitur 'ante diem octāvum kalendās Iānuāriās', est diēs annī \navnote{ante diem octāvum kalendās Iūliās (a. d. vim kal. Iūl.) = diēs octāvus ante kalendās Iūliās}brevissimus. Ante diem octāvum kalendās Iūliās diēs annī longissimus est. Ante diem octāvum kalendās Apri- 75 lēs (id est diē octāvō ante kalendās Aprīlēs) nox atque diēs aequi sunt; is diēs 'aequinoctium' dīcitur. Item ante diem octāvum kalendās Octōbrēs aequinoctium dīcitur, nam eō quoque diē nox aequa est atque diēs.

\navnote{aestās -ātis/*- hiems}Tempora annī sunt quattuor: aestās et hiems, vēr et 80 autumnus. Aestās est tempus ā mēnse Iūniō ad Augustum, hiems ā mēnse Decembrī ad Februārium. Mēnsis Iūnius initium aestātis, December initium hiemis est. Tempus ā Mārtiō ad Māium vēr dīcitur, ā mēnse Septembrī ad Novembrem autumnus. Vēr ā mēnse Mārtiō 85 initium facit. Autumnus ā Septembrī incipit.

\navnote{incipere -it -iunt - initium facere}Aestāte diēs longi sunt, sōl lūcet, āēr calidus est. Aestās est tempus calidum. Hiems tempus frīgidum est. \navnote{imber -bris m nix}Hieme nōn sōlum imber, sed etiam nix dē caelō cadit. Imber est aqua quae dē nūbibus cadit. Nix frīgidior est 90 quam imber. Hieme montēs et campī nive operiuntur. \navnote{nivis}Vēre campī novā herbā operiuntur, arborēs novīs foliīs ōrnantur, avēs, quae hieme tacent, rūrsus canere incipi\navnote{f operīre pu `` TN hieme montēs nive operiuntur}unt. Autumnō folia dē arboribus cadunt et terram operiunt. In Germāniā hiemēs frigidiōrēs sunt quam in Ita- 95 liā: altā nive operītur tōta terra et lacūs glaciē operiuntur; hieme Germānī pueri super glaciem lacuum lūdere possunt. Iānuārius mēnsis annī frigidissimus est. Mēn sēs calidissimī sunt Iūlius et Augustus. Eō tempore 100 multī Rōmānī urbem relinquunt et villas suās petunt, \navnote{magnum (: Rōma)}neque enim tōtam aestātem in urbe vīvere volunt. --- II Hōra nōna est. Aemilia apud fīlium aegrum sedet. Cubiculum lūce sōlis illustrātur. Aemilia ad fenestram it caelumque spectat. Aemilia: ``Hōra nōna est. Sōl altus in caelō est nec \navnote{z c 4T EN X3 ub ``! - Ü pueri super glaciem lacūs lūdunt}nūbibus operitur.`` Quīntus: ''Quandō sōl altissimus est?`` \navnote{quandō? = quō tempore?}Aemilia: ``Hōrā sextā vel merīdiē. Sed post merīdiem āēr calidior est quam ante merīdiem.'' 110 Aemilia fenestram claudit. Iam satis obscūrum est cubiculum. Quīntus: ``Ō, quam longae sunt hōrae, cum necesse est tōtum diem in lectō iacēre!'' Aemilia: ``Vēr est. Diēs atque hōrae longiōrēs sunt vēre quam hieme.'' Quīntus: ``Quandō longissimī sunt diēs?'' Aemilia: ``Mēnse Iūniō.'' Quīntus: ``Iamne mēnsis Iūnius est?'' \navnote{īam-ne}Aemilia: ``Māius est: hic diēs est mēnsis Māiīpostrēmus. Hōc annī tempore diēs nōn tam calidī sunt quam aestāte et noctēs frigidiōrēs sunt.'' Quīntus: ``Quī mēnsis annī calidissimus est?''

Aemilia: ``Iulius.''

Quīntus: ``Cūr ille mēnsis patris nōmen habet?''

Aemilia rīdet et respondet: ``Mēnsis Iūlius nōn ā pa- 125 tre tuō nōmen habet --- nōn tantus vir est ille! --- sed ā \navnote{AA SSSQU ZH E `` A ips VV ''SEM U/ \$7 X Iūlius Caesar}Iūliō Caesare. Ante tempora Caesaris ei mēnsī erat nōmen 'Quīntīlis' ā 'quīntō' numerō --- nōn ā nōmine \navnote{(mēnsis) Quīntīlis -is m}tuō!``

Quīntus: ``Estne Iūlius mēnsis quīntus?'' Puer digitīs 130 numerāre incipit: ``Mēnsis prīmus Iānuārius, secundus

Aemilia: Antiquis temporibus Mārtius nōn tertius, sed primus mēnsis erat; tunc igitur Quintlis, quī nunc Iūlius nōminātur, mēnsis quīntus erat, nōn septimus, \navnote{Caesar Augustus (mēnsis) Sextilis -is m}ut nunc est. Item mēnsī Augustō nōmen erat 'Sextīlis', quia sextus erat, nunc Augustus nōminātur ā Caesare Augustō. --- Sed hoc satis est. Iam necesse est tē dormīre.`` Māter enim faciem filii aspicit eumque oculōs claudere videt. Aemilia puerum dormīre velle putat. ---

\navnote{vult volunt, īnf velle}Quīntus: ``Sed diēs est. Sōl lūcet. Nocte iubē mē dormīre, cum sōl in caelō nōn est! Ubi sōl est nocte, cum hīc nōn lūcet?''

Aemilia: ``Cum hic nox est \contr{} sōl lūcet in aliīs terris procul ab Italiā et ab urbe Rōmā. Cum nox est illīc, hīc 145 in Italiā diēs est.''

Quīntus: ``Ergō nunc in aliīs terris et urbibus nox est atque hominēs dormiunt. Hīc diēs est, nec tempus est dormīre!''

\booksection{Grammatica Latīna}

Dēclīnātiō quinta Diēs brevissimus est ante diem viii kalendās Iānuāriās. Māne initium diēt est. Diē? primō Iānuāriī nōmen est kalendae Iānuāriae; ab eō diē incipit novus annus. Mēnse Iuniō diēs longi sunt. Iūnius xxx diēs habet: numerus diērum est xxx. Diēbus primīs mēnsium nōmina sunt: kalendae Iānuāriae, Februāriae, Mārtiae, Aprīlēs, cēt. Ab iīs diēbus mēnsēs incipiunt.

Nōminātīvus Accūsātīvus Genetīvus Datīvus Ablātīvus Ut 'diēs' m dēclīnātur 'meridiēs' m, item 'faciēs' f, 'glaciēs' f, et pauca alia vocābula feminīna. --- 'Māne' est vocābulum neutrum indeēclinābile.

\navnote{dēcl) = quī dēclinārī nōn potest superlātīvus (sup)}Superlātīvus Aetna mōns altus est, altior quam cēteri montēs Siciliae: is mōns Siciliae altissimus est. 170 Via Appia longa est, longior quam cēterae viae Italiae: ea via Italiae longissima est. 'īre' verbum breve est, brevius quam cētera verba Latīna: id verbum Latīnum brevissimum est. Altior, longior, brevius* comparātīvus est. 'Altissimus, longissima, brevissimum' superlātīvus

est. Superlātīvus est adiectīvum dēclīnātiōnis i et II: -issimus -a -um.

\booksection{Pēnsum A}

Hōrae di- sunt xn. Hōra sexta di- dīvidit in duās partēs: ante meridi- et post meridi-. Sex hōrae sunt dīmidia pars di-. Mēnsis Iūnius xxx di- habet: numerus di- est xxx. Novus annus incipit ab eō di- qul dīcitur kalendae Iānuāriae. Iūlius mēnsis annī calid- est; Iānuārius est mēnsis annī frigid-. Mēnsis annī brev- est Februārius. Padus est flūmen Italiae long- et lāt-. Sōl stēlla clār- est.

\booksection{Pēnsum B}

Iānuārius --- primus est. December est mēnsis --- ac ---. Tempore antīquō September mēnsis septimus ---, nam --- mēnsis primus erat Mārtius. Nunc September mēnsis--- --- est. Diēs est tempus ā --- ad ---. Māne est --- diēī, vesper --- diēī est et initium ---. Diēs in xil --- dīviditur. Nocte --- et --- lūcent, neque eae tam --- sunt quam sōl. Sōl lūnam --- suā ---. Ea lūnae pars quae sōle nōn illūstrātur --- est.

Vēr et ---, aestās et --- sunt quattuor --- annī. --- ā mēnse Iūniō ---. Hieme nōn sōlum ---, sed etiam --- dē nūbibus cadit: montēs et campī nive ---. Aqua nōn tam --- est quam nix.

\booksection{Pēnsum C}

Quot sunt mēnsēs annī? Ā quō Iānuārius nōmen habet? Quam longus est mēnsis Aprilis? Ā quō mēnsis Iūlius nōmen habet? Cūr mēnsis decimus Octōber nōminātur? Cūr lūna nōn tam clāra est quam sōl? Qul diēs annī brevissimus est? Quī diēs aequinoctia dīcuntur? Quod tempus annī calidissimum est? Quandō nix dē nūbibus cadit? Quid est imber? Ā quō diē incipit annus novus?

\booksection{Vocābula nova}
\noindent\small diēs, diem, diēī, diē, diērum, diēbus, -em, -ēs, -ei, -ēnnn, -ēbus, Vocāābula nova, annus, mēnsis, saeculum, tempus, faciēs, māne, vesper, nox, initium, hōra, meridiēs, lūna, stēlla, lūx, forma, kalendae, īdūs, nōnae, aequinoctium, aestās, hiems, vēr, autumnus, imber, nix, lacus, glaciēs, urbs, postrēmus, dīmidius, aequus, clārus, tōtus, obscūrus, exiguus, calidus, frigidus, indēclinābilis, ūndecim, trīgintā, sexāgintā, ducentī, trecentī, quīntus, sextus, septimus, octāvus, nōnus, decimus, ūndecimus, duodecimus, nōmināre, illūstrāre, incipere, operīre, velle, erat erant, vel, tunc, nunc, igitur, item, quandō?, superlātīvus.

\bookchapter{Cap. XIV}{Novus diēs}

\navnote{gallus canēns}I1 Nox est. Familia dormit. Villa Iūliī obscūra et quiēta est. Mārcus quiētus in lectō suō cubat; is bene dormit. \navnote{cubāre = iacēre (in lectō) dolēre + dat: pēs eī /Quīntō dolet}Quīntus dormīre nōn potest, quod et caput et pēs eī dolet. Bracchium quoque dolet Quīntō; itaque is nōn dormit, sed vigilat.

\navnote{vigilāre dormīre}Cubiculum in quō Quīntus cubat nōn magnum est, \navnote{uter- utra- utrum-que: uterque puer:: et M. et Q.; utrumque cubiculum: et cubiculum M.ī neuter -tra -trum: neuter}nec magnum est cubiculum Mārcī. Utrumque cubiculum parvum est. Uterque puer cubat in cubiculō parvō, neuter in cubiculō magnō. Neutrum cubiculum mag\navnote{alter -era -erum duo, abl duōbus valēre \contr{} \contr{} aegrōtāre uter -tra -trum: uter puer?}num est. Uterque puer quiētus est, neuter puer sē mo-vet. Alter puer dormit, alter vigilat. Alter ē duōbus pueris valet, alter aegrōtat. Uter puer aegrōtat, Mārcusne an Quīntus? Quīntus aegrōtat, Mārcus valet.

Āēr frigidus cubiculum Mārcī intrat, fenestra enim aperta est. Mārcus fenestrā aperta dormit. Fenestra Quīntī aperta nōn est. Altera ē duābus fenestrīs est aperta, altera clausa. Utra fenestra clausa est? Fenestra \navnote{M.ne an Q.? duae, abl duābus clausus -a -um \contr{} eapertus -ā -um}Quīnt. Is fenestrā clausā dormit, quia aeger est.

\navnote{canēns -entis = quī canit}Ecce gallus canit: ``Cucurrū! Cucucurrū!'' Gallus canēns novum diem salūtat.

\navnote{gallia canēns gallum canentem}Mārcus oculōs nōn aperit neque sē movet. Quīntus, quī oculīs apertīs iacet, super lectum sē vertit. Puer

\navnote{dormiēns -entis = qul dormit vigilāns -antis = quī vigilat quō-modo? \contr{} quō modō? ho-diē = hōc diē}lum audit, et Dāvum vocat.

Dāvus cubiculum intrāns interrogat: ``Quōmodo sē 25 habet pēs tuus hodiē?''

Quīntus respondet: ``Pēs male sē habet, nec pēs tan-tum, sed etiam caput et bracchium dolet. Ō, quam longa nox est! Sed iam māne est, nam gallus canit. Dā mihi aquam, Dāve!'' 30

Dāvus Quīntō aquam in pōculō dat. Puer aquam bi-bit. Servus puerum bibentem aspicit.

Dāvus Quīntum in lectō iacentem relinquit et cubicu\navnote{ad-hūc = ad hoc tempus, etiam nunc}lum Mārcī intrat. Mārcus adhūc dormit. Dāvus ad puerum dormientem adit eumque excitat. Quōmodo servus 35 puerum excitat? In aurem pueri dormientis magnā vōce clāmat: ``Mārce! Māne est!'' Eō modō excitātur Mārcus, et oculōs aperiēns servum apud lectum stantem videt. Iam neuter ē duōbus pueris dormit.

``Hōra prima est'' inquit Dāvus, ``Surge ē lectō!''

\navnote{inquit = dīcit surgere}Mārcus ē lectō surgit. Iam nōn cubat inlectō, sed \navnote{af-ferre \contr{} ad-ferre}ante lectum stat. Mārcus Dāvum aquam afferre iubet: ``Affer mihi aquam ad manus!''

Servus Mārcō aquam affert et ``Ecce aqua'' inquit, \navnote{lavāre}``Lavā faciem et manūs! Manūs tuae sordidae sunt.'' Mārcus prīmum manūs lavat, deinde faciem. Dāvus: ``Aurēs quoque lavā!'' ``Sed aurēs'' inquit Mārcus ``in faciē nōn sunt!'' Dāvus: ``Tacē, puer! Nōn modo faciem, sed tōtum \navnote{``Y -AVV T eR a DN 4 Mārcus faciem lavat}caput lavā! Merge caput in aquam!`` Mārcus caput tōtum in aquam mergit atque etiam aurēs et capillum lavat. Iam tōtum caput eius purum est. Aqua nōn est pūra, sed sordida. Aqua quā Mārcus lavātur frigida est; itaque puer manūs et caput sōlum lavat, nōn tōtum corpus. Māne Rōmānī faciem et manūs aquā frigidā lavant; post meridiem tōtum corpus lavant aquā calidā. Mārcus caput et manūs tergēns Dāvum interrogat: ''Cūr frāter meus tam quiētus est?`` Dāvus respondet: ''Qulntus adhūc in lectō est.`` Mārcus: ''In lectō? Quīntus, quī ante mē surgere solet, adhūc dormit! Excitā eum!`` ''Nōn dormit'' inquit Dāvus, ``Frāter tuus vigilat, nec surgere potest, quod pēs et caput el dolet.'' 65 \navnote{vir togātus}Mārcus: ``Mihi quoque caput dolet!'' Daāvus: ``Tibi nec caput nec pēs dolet! Caput valēns \navnote{tunica}nōn dolet nec membra valentia.`` Per fenestram apertam intrat āēr frigidus. Mārcus fri-\navnote{frigēre = frigidus esse}get, quod corpus eius nūdum est (id est sine vestīmentīs); itaque Mārcus vestīmenta sua ā servō poscit: ``Dā \navnote{poscere = dari iubēre vestīre}mihi tunicam et togam! Vesti mē!`` Dāvus puerō frigentī tunicam et togam dat, neque eum vestit: necesse est puerum ipsum sē vestīre. Mārcus prīmum tunicam induit, deinde togam. Puer iam nudus nōn est.

\navnote{induere}(Toga est vestīmentum album, quod viri et pueri Rō\navnote{gerere = (in corpore) ferre, habēre togātus -a -um = togam gerēns praeter prp--acc}mānī gerunt. Graecī et barbari togam nōn gerunt. Multīs barbarīs magna corporis pars nūda est. Virō togātō nūlla pars corporis est nūda praeter bracchium alterum. Utrum bracchium virō togātō nūdum est, dextrumne an \navnote{dexter -tra -trum \contr{} sinister -tra -trum c}sinistrum? Bracchium dextrum est nūdum, bracchium 80 sinistrum

togā operītur. Mīlitēs togam nōn gerunt, nēmō enim togātus gladiō et scūtō pugnāre potest. Utrā manū mīles gladium gerit? Manū dextrā gladium gerit, scūtum gerit manū sinistrā.)

Mārcus, quī pedibus nūdīs ante lectum stat, calceōs 85 poscit: ``Dā mihi calceōs! Pedēs frīgent mihi.'' Dāvus eī calceōs dat, et eum sēcum venīre iubet: ``Veni mēcum! Dominus et domina tē exspectant.''

Mārcus cum servō ātrium intrat, ubi parentēs sedent \navnote{= pater et māter}līberōs exspectantēs. Ante eōs in parvā mēnsā pānis et 90 māla sunt. Parentēs ā filio intrante salūtantur: ``Salvēte, pater et māter!'' et ipsī fīlium intrantem

salūtant: ``Salvē, Mārce!''

Māter alterum filium nōn vidēns Dāvum interrogat: Quīntus quōmodo sē habet hodiē?``

Dāvus: ``Quīntus dīcit 'nōn modo pedem, sed etiam caput dolēre'.''

Aemilia surgit et ad fīlium aegrōtantem abit. Māter fīliō suō aegrōtantī pānem et mālum dat, sed ille, quī multum ēsse solet, hodiē nec pānem nec mālum ēst. Puer aegrōtāns nihil ēsse potest.

Mārcus autem magnum mālum ā patre poscit: ``Dā mihi illud mālum, pater! Venter vacuus est mihi.''

\navnote{esse, imp ēs! ēste!}Iūlius Mārcō pānem dāns ``Primum ēs pānem'' in-quit, ``deinde mālum!''

\navnote{utra-que manus utram-que manum}Mārcus pānem suum ēst. Deinde pater ei in utramque manum mālum dat, et ``Alterum mālum nunc ēs'' inquit, ``alterum tēcum fer!''

\navnote{tē-cum = cum tē ferre, imp fer! ferte!}Dāvus Mārcō librum et tabulam et stilum et rēgulam afFert.

Iūlius: ``Ecce Dāvus tibi librum et cēterās rēs tuās affert. Sūme rēs tuās atque abī!''

``Sed cūr nōn venit Mēdus?'' inquit Mārcus, quī Mēdum adhūc in familiā esse putat. Is servus cum puerīs īre solet omnēs rēs eōrum portāns. Mārcus ipse nūllam rem portāre solet praeter mālum.

Iūlius: ``Mēdus tēcum īre nōn potest. Hodiē necesse est tē sōlum ambulāre.''

\navnote{tabula}Cūr ille servus mēcum venīre nōn potest ut solet? Etiamne Mēdō caput dolet?``

Iūlius: ``Immō bene valet Mēdus, sed hodiē aliās rēs agit.''

Mārcus: ``Quae sunt illae rēs?''

Iūlius nihil ad hoc respondet et ``Iam'' inquit ``tempus est discēdere, Mārce.''

Mārcus mālum, librum, tabulam, stilum rēgulamque \navnote{salvē! valē! ``salvē!'' dīcit quī advenit valē!`` dīcit quī discēdit}sēcum ferēns ē villā abit. Fīlius ā patre discēdēns ``Valē, pater!'' inquit.

``Valē, Mārce!'' respondet pater, ``Bene ambulā!''

Quō it Mārcus cum rēbus suīs? Vidē capitulum quīntum decimum!

\booksection{Grammatica Latīna}

Participium

\navnote{participium -īn (part) /Puella vigilāns... masc./fem. sing. plūr. nōm. -ns acc. gen. dat. abl. neutr. -ntia nōm. -ns -ntia acc. -ns gen. -ntis dat. -ntī -ntibus abl.}Puer vigilāns in lectō iacens gallum canentem audit. Gallus 135 canēns nōn audītur ā puerō dormiente. Puer dormiēns servum clamantem audit, nec ā gallō canente, sed ā servō clāmante excitātur. Nōn vōx gallī canentis, sed vōx servī clāmantis puerum dormientem excitat.

Pueri vigHantēs gallōs canentēs audiunt.

Gallī canentēs ā 140 pueris dormientibus nōn audiuntur. Pueri dormientēs servōs clāmantes audiunt, nec ā gallīs canentibus, sed ā servīs clāmantibus excitantur. Nōn vōcēs gallōrum canentium, sed vōcēs servōrum clāmantium puerōs dormientēs excitant.

Caput valēns nōn dolet nec membra valentia. Canis animal volāns nōn est; animālia volantia sunt avēs. Mīlitēs inter pīla volantia pugnant.

'Vigilāns', 'iacens', canēns', dormizns' participia

sunt. \navnote{-ēns -entis}Participium est adiectīvum dēclīnātiōnis III: 111:gen. sing. -antis, -entis (abl. sing. -e vel -ī).

\booksection{Pēnsum A}

Puer dorm- nihil audit. Dāvus puerum dorm- excitat: in aurem puerī dorm- clāmat: Mārce!`` Mārcus oculōs aperservum apud lectum st- videt. Servus puerō frīg- vestīmenta dat. Parentēs fīlium intr- salūtant et ā filiō intr- salūtantur. Fīlius discēd- ''Valē!`` inquit.

Corpus val- nōn dolet. Medicus caput dol- sānāre nōn: potest. Piscēs sunt animālia nat-.:

\booksection{Pēnsum B}

Pueri in lectīs--- --- [= iacent]. --- puer dormit, alter ---; alter

Quīntus aegrōtat.

Servus puerum dormientem --- et ei aquam ---. Mārcus ē lectō --- et primum manūs ---, --- faciem. Puer vestīmenta ā ) servō ---, et primum tunicam ---, deinde ---. Iam puer --- nōn est. Virō --- [= togam gerentī] bracchium --- nūdum est. t Dāvus Mārcum sēcum venīre iubet: Veni ---!`` ---!''

Mārcus Mēdum, qul cum eō īre ---, nōn videt. Mēdus librōs et cēterās --- Mārcī portāre solet, Mārcus ipse --- por tāre solet --- mālum. ``Hodiē'' --- Iūlius ``Mēdus --- īre nōn potest.'' Mārcus sōlus abit librum et --- et stilum --- ferēns.;!

\booksection{Pēnsum C}

Qulntusne bene dormit? Uter puer aegrōtat? I Estne clausa fenestra Mārcī? Uter ē duōbus pueris gallum canentem audit? t Quōmodo servus Mārcum excitat? ) Totumne corpus lavat Mārcus? c - I Cūr Mārcus friget?: Quid Mārcus ā servō poscit? Utrum bracchium togā operītur? I Utrā manū mīles scūtum gerit? ā Quās rēs Mārcus sēcum fert? I

\booksection{Vocābula nova}
\noindent\small rēs reī f, sing plūr, nōm, rēs, acc, rem, gen, reī, rērum, dat, rēbus, abl, rē, -iZns -ientis Vocābula nova, gallus, vestimentum, tunica, toga, calceus?, parentēs, tabula, stilus, rēgula, apertus, clausus, sordidus, pūrus, nūdus, togātus, dexter, sinister, omnis, cubāre, vigilāre, valēre, excitāre, surgere, afFerre, lavāre, mergere, solēre, frigēre, poscere, vestlre, induere, gerere, inquit, uterque, neuter, alter, uter?, mihi, tibi, mēcum, tēcum, sēcum, nihil, quōmodo, hodiē, adhūc, prīmum, deinde, praeter, an, valē, participium.

\bookchapter{Cap. XV}{Magister et discipulī}

SEXTVS

\navnote{lūdus}MAGISTER ET DISCIPVLI

Mārcus librum et tabulam et cēterās rēs ferēns Tuscu- / lum ambulat. Illīc est lūdus puerōrum.

Multī puerī māne in lūdum eunt.

\navnote{magister -tri m}Magister ludī est vir Graecus, cui nōmen est Diodōrus. Mārcus magistrum metuit, Diodōrus enim magis- 5 ter sevērus est, quī discipulōs suōs virgā verberat; eō modō magister sevērus discipulōs improbōs pūnīre so-let. Discipulī sunt puerī quī in lūdum eunt. Mārcus et Quīntus sunt duo discipulī. Aliī discipulī sunt Titus et Sextus.

Sextus, quī ante Mārcum et Titum ad lūdum advenit, \navnote{nōn-dum = adhūc nōn sella -ae f IL A}prīmus lūdum intrat. Sextus sōlus est, nam ceteri discipulī nōndum adsunt.

Magister intrāns Sextum sōlum in sellā sedentem videt. Sextus dē sellā surgēns magistrum salūtat: ``Salvē, 15 magister!''

\navnote{cōn-sīdere = sedēre in-}Magister: ``Salvē, Sexte! Cōnsīde!''

\navnote{tacitus -a -um = tacēns}Sextus in sellā cōnsīdit. Discipulus tacitus ante magistrum sedet.

Magister interrogat: ``Cūr tū sōlus es, Sexte?''

``Ego sōlus sum, quod cēteri discipulī omnēs absunt'' respondet Sextus.

\navnote{ex-clāmāre}Sextus: ``Num ego discipulus improbus sum?''

Magister: ``Immō tū probus es discipulus, Sexte, at Mārcus et Quīntus et Titus improbī sunt!''

\navnote{iānua -ae f --- ōstium}Hic Titus ad lūdum advenit et iānuam pulsat antequam intrat: discipulus nōn statim intrat, sed prīmum iānuam pulsat; tum lūdum intrat et magistrum salūtat: \navnote{tum = deinde}``Salvē, magister!''

Titus: ``At Mārcus et Quīntus nōndum adsunt!''

Magister: ``Tacē, puer! Claude iānuam et cōnsīde! Aperite librōs, puerī!''

Sextus statim librum suum aperit, sed Titus, qul librum nōn habet, ``Ego'' inquit ``librum nōn habeō.''

Magister: ``Quid? Sextus librum suum habet, tū librum tuum nōn habēs? Cūr librum nōn habēs?''

Titus: ``Librum nōn habeō, quod Mārcus meum librum habet.'' Il

Post Titum Mārcus ad lūdum advenit, neque is iānuam pulsat antequam intrat. Mārcus statim intrat, nec magistrum salūtat.

``Ō Mārce!'' inquit Diodōrus, ``Cūr tū iānuam nōn pulsās cum ad lūdum venis, nec mē salūtās cum mē 45 vidēs?''

At Mārcus ``Ego'' inquit ``iānuam nōn pulsō cum ad lūdum veniō, nec tē salūtō cum tē videō, quia nec Sextus nec Titus id facit.''

Sextus et Titus: ``Quid?''

\navnote{īn}Mārcus (ad Sextum et Titum): ``Vōs iānuam nōn pulsātis cum ad lūdum venītis, nec magistrum salūtātis \navnote{audltis-ne}cum eum vidētis. Audītisne id quod dīcō?``

Tum Sextus et Titus ``Id quod dīcis'' inquiunt ``vērum nōn est: nōs iānuam pulsāmus cum ad lūdum venī- 55 mus, et magistrum salūtāmus cum eum vidēmus. Nōnne vērum dīcimus, magister?''

\navnote{vērum -In = id quod vērum est quod = id quod}Magister: ``Vōs vērum dīcitis. Quod Mārcus dīcit nōn est vērum. Discipulus improbus es, Mārce! Necesse est tē pūnīre. Statim ad mē venī!''

Diodōrus, magister sevērus, tergum pueri virgā verberat. (Tergum est posterior pars corporis.) Tergum \navnote{posterior -ius comp lacrimāre puerō (dat) nōn convenit = lacrimae puerō nōn conveniunt dēsinere --- \contr{} incipere}dolet Mārcō, neque ille lacrimat, nam lacrimāre puerō Rōmānō nōn convenit.

Mārcus clāmat: ``Ei! Iam satis est! Dēsine, magister!''

Magister puerum verberāre dēsinit et ``Ad sellam \navnote{red-īre = rūrsus Īre}tuam redī'' inquit, ``atque cōnsīde!''

Mārcus ad sellam suam redit, neque cōnsīdit, sed tacitus ante sellam stat.

\navnote{audīs-ne}``Audīsne, Mārce? Ego tē cōnsīdere iubeō. Cūr nōn cōnsīdis?'' interrogat magister.

Mārcus: ``Nōn cōnsīdō, quod sedēre nōn possum.''

Diodōrus: ``Cūr sedēre nōn potes?''

pars tergī īnferior in quā sedēre soleō.``

\navnote{īnferior -ius comp \contr{} īnfrā}Haec verba audientēs Titus et Sextus rident.

Diodōrus: ``Quid rīdētis, Tite et Sexte?''

Diodōrus: ``Tacēte! Eam corporis partem nōmināre nōn convenit! --- Sed ubi est frāter tuus, Mārce?''

Mārcus: ``Is domī est apud mātrem suam. Quīntus \navnote{domi \contr{} in lūdō Q.: Ego aeger sum'' Q. dīcit sē aegrum esse}dīcit 'sē aegrum esse''

Diodōrus: ``Si aeger est, in lādum īre nōn potest. At vōs bene valētis. Iam aperite librōs!''

Titus: ``Mārcus meum librum habet.''

Diodōrus: ``Quid tū librum Titī habēs, Mārce?''

Mārcus: ``Ego eius librum habeō, quod is meum mālum habet. Redde mihi mālum meum, Tite!''

\navnote{red-dere = rūrsus dare}Titus ridēns ``Mālum'' inquit ``tibi reddere nōn possum, id enim iam in ventre meō est!''

Mārcus Irātus Titum pulsāre incipit, sed magister ``Dēsine, Mārce!'' inquit, ``Titus tibi mālum dare nōn \navnote{mālum = --- māla rēs; m. dare tibi:: tē pūnīre nisi = si nōn}potest, at ego tibi mālum dare possum, nisi hic et nunc Titō librum reddis!``

Mārcus Titō librum reddit. IH

Discipulī librōs aperiunt. Item librum suum aperit magister ac recitāre incipit. Puerī autem priōrem libri \navnote{prior -ius \contr{} \contr{} posterior; prior pars:: initium}partem tantum audiunt, nam antequam magister par\navnote{pars posterior:: finis}tem libri posteriōrem recitāre incipit, omnēs pueri dormiunt!

\navnote{magister recitat discipulī dormiunt}Magister recitāre dēsinit et exclāmat: ``Ō improbī discipulī! Dormītis! Quod recitō nōn audītis!'' \navnote{quod z id quod}Mārcus magistrō īrātō ``Ego'' inquit ``nōn dormiō. Immō vigilō et tē audiō, magister'' Item Sextus et Titus ``Neque nōs dormīmus'' inquiunt, ``Vigilāmus et omnia verba tua audīmus. Bene recitās, Diodōre! Bonus es magister!'' Magister laetus ``Vērum est quod dīcitis'' inquit, ``ego bene recitō ac bonus sum magister --- at vōs male recitātis ac malī discipulī estis!'' Discipulī: ``Immō bonī discipulī sumus! Bene recitāmus!'' Diodōrus: Tacēte! Ubi estis, pueri?`` Discipulī: In lūdō sumus.'' Diodōrus: ``Vērum dīcitis: in lūdō estis --- nōn domī 115 in lectulis! In lectulō dormīre licet, hic in lūdō nōn licet \navnote{lecta/us = parvus lectus}dormīre!`` Magister Irātus discipulōs virgā verberat.

Discipulī: ``Ei, ei! Quid nōs verberās* magister?'' 120

Diodōrus: ``Vōs verberō, quod aliō modō vōs excitāre nōn possum. Ad sellās vestrās redīte atque cōnsīdite!''

Discipulī ad sellās suās redeunt, neque cōnsīdunt.

\navnote{red-īre -it -eunt}Magister: ``Quid nōn cōnsīditis?''

Discipulī: ``Nōn cōnsīdimus, quod sedēre nōn possu-

Diodōrus: ``Quid? Sedēre nōn potestis? Ergō vōs hōram tōtam taciti stāte, dum ego sedēns partem libri posteriōrem recitō! Nec enim stantēs dormīre potestis!''

\booksection{Grammatica Latīna}

Persōnae verbī

Puerī rident. Puella nōn rider.

Puerī: ``Cūr tū nōn ridēs, cum nōs rldtmus?''

Puella: ``Ego nōn rideo, quia vōs mē rīdētis/''

Rideō rīdēmus' dīcitur persōna

prima, 'ridēs ridētis' persōna secunda, 'ridet rident! persōna tertia.

\navnote{sing. plūr. -tis II -nt}Persōna prīma

rideō Persōna secunda

ridēs Persōna tertia ------

ridet Exempla: [1] clāmo clamāmus; [2] rideō rīdēmus; dīcō dīcimus; faciō facimus; [4] audiō audīmus.

Malus discipulus in lūdō clāmat et rīdet. Quī id nōn facit, et omnia verba magistri audit, bonus est discipulus.

Bonus discipulus malō dīcit: Ego bonus discipulus sum: in lūdō nec clāmō nec rideo, et omnia verba magistri audi?. At tū malus discipulus es: in lūdō clāma5 et rides nec verba magistrī audts/ Vērum dīcō/ Nōnne hoc facis?`` Malus discipulus: ''Quod dīcis vērum nōn est: ego hoc nōn faciō/``

Malī discipulī in lūdō clāmant et rīdent. Quī id nōn faciunt, et omnia verba magistri audiunt, bonī sunt discipulī.

Bonī discipulī malīs dīcunt: ``Nōs bonī discipulī sumus: in lūdō nec clamaāmus nec rīdēmus, et omnia verba magistri audi-mus. At vōs malī discipulī estis: in lūdō clamatis et ridetis nec verba magistri audītis/ Vērum dīcimus! Nōnne hoc facitis?'' Malī discipulī: ``Quod dīcitis vērum nōn est: nōs hoc nōn 155 facimus/'' [1] Ut 'clām clamāmus' dēclīnantur verba quōrum īnfīnītīvus dēsinit in -āre. [2] Ut 'rideō rīdēmus' dēclīnantur verba quōrum īnfinītīvus dēsinit in -ere. [3] Ut 'dīcō dīcimus! dēclīnantur verba quōrum

infinitivus dēsinit in -ere, praeter ea quae -i- habent ante -ō et -unt, ut faciō faciunt': accipere, aspicere, capere, facere, fugere, ia-cere, incipere, parere, cēt. [4] Ut 'audiō audtmus' dēclīnantur verba quōrum infinitivus 165 dēsinit in -īre, praeter verbum īre (pers. I sing. eō, pers. III plūr eunt), item ab-īre, ad-īre, ex-īre, red-Īre, cēt.

Ut 'sum sumus* dēclīnantur verba quōrum infinitivus deēsinit in -esse, ut ab-esse, ad-esse, in-esse, cēt., et posse (pers. i pos-sum pos-sumus, n pot-es pot-estis, iii pot-est pos-sunt).

\booksection{Pēnsum A}

Mārcus ad lūdum ven- nec iānuam puls-. Magister: ``Cūr tū iānuam nōn puls-, cum ad lūdum ven-?'' Mārcus: ``Ego iānuam nōn puls- cum ad lūdum ven-, quod nec Sextus nec Titus id facit. Audīte, Sexte et Tite: vōs iānuam nōn puls- cum ad lūdum ven-!'' Sextus et Titus: ``Nōs iānuam puls- cum ad lūdum ven-!'' Magister: ``Tacēte! Aperite librōs!'' Titus: ``Ego librum nōn hab-.'' Magister: ``Cūr librum nōn hab-5 Tite?'' Titus: ``Librum nōn hab-, quod Mārcus meum librum hab-!'' Mārcus: ``Sed vōs meās rēs hab-!'' Titus et Sextus: ``Nōs rēs tuās nōn hab-!''

Magister discipulōs dormīre vidēns exclāmat: ``Ō pueri! Dorm-! Ego recit-, vōs nōn aud-!'' Mārcus: ``Ego tē recitāre aud-. Nōn dorm-.'' Titus et Sextus: ``Nec nōs dorm-. Te recitāre aud-. Bene recit-, magister.'' Magister: ``Ego bene recit-, at vōs male recit-! Malī discipulī es-!'' Discipulī: ``Ve )rum nōn dīc-, magister. Bonī discipulī s-: in lūdō nec clāmnec rid-, et tē aud-!'' i )

\booksection{Pēnsum B}

Māne pueri in --- eunt. Pueri quī in lūdum eunt --- sunt. Quī lūdum habet --- est. Mārcus magistrum metuit, nam Diodōinfcrior rus magister --- est quī puerōs improbōs --- verberat. i

Intrat magister. Sextus dē --- surgit. Cēteri discipulī --- adsunt. Magister---: --- ``Ō, discipulōs improbōs!'' Sextus: ``Num \contr{} ego improbus ---?'' ---?``Magister: ''--- discipulus improbus nōn es, --- [= sed] cēteri discipulī improbī sunt!`` I

Post Sextum venit Titus, --- [= deinde] Mārcus. Mārcus --- [= ōstium] nōn pulsat, --- lūdum intrat, nec magistrum! salūtat. Magister: ``Discipulus improbus ---, Mārce! --- ad mē veni!'' Magister --- Mārcī verberat. Tergum est --- pars corpo \$ris. Magister puerum verberāre ---. Mārcus ad sellam suam --- neque ---. Magister: u--- [= cūr] nōn cōnsldis?`` Mārcus: ''Sedēre nōn ---, quod pars tergī --- mihi dolet!`` Y I

\booksection{Pēnsum C}

Quō pueri māne eunt? Quis est Diodōrus? ē Cūr pueri magistrum metuunt? Quis discipulus primus ad lūdum advenit? Quid facit Titus antequam lūdum intrat? Cūr Titus librum suum nōn habet? Quis est discipulus improbissimus? Cūr Quīntus in lūdum īre nōn potest? Cūr magister recitāre dēsinit? Tūne magister an discipulus es? Num tū iānuam pulsās antequam cubiculum tuum intrās?

\booksection{Vocābula nova}
\noindent\small pos-sum, pot-es, pot-est, pos-sumus, pot-estis, pos-sunt, persōna (pers) ego, nōs, tū, vōs, rīdeō, rīdēmus, rīdēs, ridetis, ridet, rident, [1] clāmāre clāmō, clāmāmus, clāmō, clāmā tis, clāmāls, clāmar, rideō, ridēes, ridē, tis, ridet rident, dlcimus, āīcō, dīc itis, dīc, dlc, dīc unt, facere, faciō, faci mus, facim, faci*, facir, faci unt, audiō, audīmus, audis, audītis, audit audiunt, esse, sum, sumus, es, estis, est, sunt, -amus, -ātis, -ās, -ēmus, -ēs, -ētis, ent, -et, -itis, -is, -it, -io, -imus, -iunt, -tmus, -ttis, -īs, īre, eo, īmus, īs, ītis, it, eunt, posse, pos-sunus, pos-sum --- pos-sumus, Vocābtdanova, lūdus, magister, discipulus, virga, seUa, iānua, vērum, tergum, mālum, lectulus, sevērus, tacitus, vērus, postcrior, prior, pūnīre, cōnsīdere, exclāmāre, dēsinere, redlre, reddere, recitāre, licēre, ego, nōndum, statim, tum, quid?, domī, antequam, at, sī, nisi.

\bookchapter{Cap. XVI}{Tempestās}

\navnote{nāvis Rōmāna}\navnote{inter-esse quōrum: ex quibus situs -a -um: situm est sīve = vel appellāre = nōmināre}Italia inter duo maria interest, quōrum alterum, quod / suprā Italiam situm est, 'mare Superum' sīve Hadriāaticum' appellātur, alterum, īnfrā Italiam situm, 4mare īnferum' sīve ``Tūscum'. Tōtum illud mare longum et lātum quod inter Eurōpam et Āfricam interest mare 5 nostrum' appellātur ā Rōmānīs.

Urbs Rōma nōn ad mare, sed ad Tiberim flūmen sita \navnote{Tiberis -is m (acc paulum --- \contr{} multum}est vīgintī milia passuum ā mari. Quod autem paulum \navnote{nāvigāre \contr{} nāvis}aquae est in Tiberī, magnae nāvēs in eō flumine nāvigāre nōn possunt. Itaque parvae tantum nāvēs Rōmam 10 adeunt.

\navnote{mare situs}Ōstiam omnēs nāvēs adire possunt, id enim est oppidum maritimum quod magnum portum habet. Ad ōstium Tiberis sita est Ōstia. (*Ōstium' sīve os' flūminis \navnote{īn-fluere}dīcitur is locus quō flūmen in mare īnfluit; Ōstia sita est 15 eō locō quō Tiberis in mare Inferum īnfluit,)

Alia oppida maritima quae magnōs portūs habent sunt Brundisium, Ariminum, Genua, Puteolī. Haec \navnote{Puteoli -ōrum m pi Ora -ae f}omnia oppida in ōrā maritimā sita sunt. (Ōra maritima est fīnis terrae, unde mare incipit. Portus est locus in \navnote{ōra = ōra maritima}ōrā maritimā quō nāvēs ad terram adire possunt.) In ōrā Italiae multī portūs sunt. Ex omnibus terrīs in portūs Italiae veniunt nāvēs, quae mercēs in Italiam vehunt. \navnote{merx}(Mercēs sunt rēs quās mercātōrēs emunt ac vēndunt.)

\navnote{mercis f Ostiensis -e \contr{} \contr{} Ōstia semper = omni tempore}Nōn modo mercēs, sed etiam hominēs nāvibus vehuntur. Portus Ōstiēnsis semper plēnus est hominum quī in aliās terrās nāvigāre volunt. Is quī nāvigāre vult adit nautam quī bonam nāvem habet. Si āēr tranquillus est, necesse est ventum opperiri. (Ventus est āēr quī \navnote{opperiri = exspectāre flāre: ventus flat (= movētur) turbidus -a -um: (mare) t.um \contr{} \contr{} tranquillum tempestās -ātis f turbāre = turbidum facere}movētur.) Cum nūllus ventus super mare flat, tranquil-lum est mare; cum magnus ventus flat, mare turbidum est. Tempestās est magnus ventus quī mare turbat ac flūctūs facit quī altiōrēs sunt quam nāvēs. Nautae tem\navnote{im-plēre = plēnum facere}pestātēs metuunt, nam magnī flūctūs nāvēs aquā implēre possunt. Tum nāvēs et nautae in mare merguntur. T, r flūctus

Nautae nec mari turbidō nec mari tranquillō nāvigāre. \navnote{marīturbidō (abl) = dum mareturbidumest opperiuntur --- exspectant ventō secundō (abl) = dum ventus secundus est \&-gredī = exīre; ēgrediuntur = exeunt puppis -is f (acc -im, medium mare = media maris pars gubernātor -ōris m = is quī nāvem gubernat oriri: sōl oritur = sōl in caelum ascendit oriēns -entis m e occidēns -entis m sōl EN P}volunt; itaque in portū ventum secundum opperiuntur (id est ventus quī ā tergō flat). Ventō secundō nāvēs ē ponū ēgrediuntur: vēla ventō implentur ac nāvēs plēnīs vēlis per mare vehuntur.

Pars nāvis posterior puppis dīcitur. In puppī sedet 77 nauta quī nāvem gubernat. Quōmodo nāvis in mediō mari gubernārī potest, cum terra nūlla vidētur? Gubernātor caelum spectat: in altō mari sōl aut stēllae ei ducēs sunt. Ea pars caelī unde sōl oritur dīcitur oriēns. Par- 45 tēs caelī sunt quattuor: oriēns et occidēns, merīdiēs et SEUTENTRIONES

\navnote{oriēns}\navnote{sōl occidēns (Y USD L1}OCCIDĒNS e ORIĒNS MERIDIES A]

\navnote{septem r. pem m. triōnēs'}\navnote{septen-triōnēs -um mpl \contr{} merīdiēs occidere --- oriri}septentriōnēs. Occidēns est pars caelī quō sōl occidit. Meridiēs dīcitur ea caelī pars ubi sōl merīdiē vidētur; pars contrāria septentriōnēs appellātur ā septem stēllīs \navnote{\contr{} contrāā}quae semper in eā caelī parte stant. Iīs quī ad septentriōnēs nāvigant, oriēns ā dextrā est, ā sinistrā occidēns, merīdiēs ā tergō. Oriēns et occidēns partēs contrāriae sunt, ut meridiēs et septentriōnēs. ---

Hodiē caelum serēnum et ventus secundus est. Nāvēs \navnote{manus d. sinistrā -ae f --- manus s. serēnus -a -um: (caelum) s.um = sine nūbibus simul = ūnō tempore}multae simul ē portū Ōstiēnsī ēgrediuntur. Inter eōs 55 hominēs quī nāvēs conscendunt est Mēdus, quī ex Italiā \navnote{cōn-scendere = ascendere proficīscī = abīre}proficīscitur cum amicā suā Lydiā. Mēdus, quī Graecus est, in patriam suam redīre vult. Graecia nōn modo ipsīus patria est, sed etiam Lydiae.

Mēdus et Lydia ex Italiā proficīscentēs omnēs rēs suās sēcum ferunt: pauca vestīmenta, paulum cibī nec \navnote{multum pecūniae = magna pecūnia praeter-eā = praeter eās rēs sōle oriente = dum sōl oritur multīs spectantibus = dum multī spectant altum -In - altum mare sequī = venīre post; eam sequuntur = post eam veniunt}multum pecūniae. Praetereā Lydia parvum librum fert, quem sub vestimentis occultat.

Sōle oriente nāvis eōrum ē portū ēgreditur multīs hominibus spectantibus. Nāvis plēnīsvēlīsaltum petit. Aliae nāvēs eam sequuntur.

Mēdus in puppim ascendit. Lydia eum sequitur. Ex altā puppī sōlem orientem spectant. Iam procul abest \navnote{vix = prope nōn cernere = vidēre (id quod procul abest)}Ōstia, hominēs quī in portū sunt vix oculīs cernl possunt. Mēdus montem Albānum, quī prope vīllam Iūliī situs est, cernit et ``Valē, Italia!'' inquit, ``Valēte montēs et vallēs, campī silvaeque! Ego in terram eō multō pulchriōrem, in patriam meam Graeciam!'' Mēdus laetātur neque iam dominum suum sevērum verētur.

Lydia collēs in quibus Rōma sita est procul cernit et ``Valē, Rōma!'' inquit, ``Nōn sine lacrimīs tē relinquō, nam tu altera patria es mihi.'' Lydia vix lacrimās tenēre potest.

Mēdus faciem Lydiae intuētur et ``Nōnne gaudēs'' \navnote{intuerī = spectāre}inquit, ``mea Lydia, quod nōs simul in patriam nostram redīmus?''

Lydia Mēdum intuēns ``Gaudeō'' inquit ``quod mihi licet tēcum venīre, At nōn possum laetāri quod omnēs amīcās meās Rōmānās relinquō.

Sine lacrimīs Rōmā \navnote{labi z- cadere}proficīscī nōn possum.`` Dē oculīs Lydiae lacrimae lā- 85 buntur.

Mēdus eam complectitur et ``Tergē oculōs!'' inquit, \navnote{complectī: eam complectitur = bracchia circum corpus eius pōnit}``Ego, amīcus tuus, quī tē amō, tēcum sum. In patriam nostram īmus, ubi multī amīcī nōs opperiuntur.''

Hīs verbīs Mēdus amīcam suam tristem cōnsōlātur.

\navnote{paulō post = post paulum (temporis) quam in partem? = in quam partem?}Paulō post nihil ā nāve cernitur praeter mare et caelum. Mēdus gubernātōrem interrogat: ``Quam in partem nāvigāmus?'' Ille respondet: ``In meridiem. Ecce \navnote{sōle duce = dum sōl dux est, sōle}sōl oriēns mihi ā sinistrā est. Sōle duce nāvem gubernō. Bene nāvigāmus ventō secundō atque caelō serēnō.``

\navnote{dūcente loquī = verba facere ūter -tra -trum = niger fit = esse incipit m n I gut p p I, IA}Dum ille loquitur, Mēdus occidentem spectat et nūbēs ātrās procul suprā mare oriri videt; simul mare tran quillum fit. ``Nōn serēnum est caelum'' inquit, ``Ecce

Gubernātor statim loquī dēsinit et nūbēs spectat; tum

īoo vēla aspiciēns exclāmat: ``Quid (mālum!) hoc est? Nūbēs ātrae ab occidente oriuntur et ventus simul cadit! Ō \navnote{fulgur in-vocāre}Neptūne! Dēfende nōs ā tempestāte!`` Nauta Neptūnum, deum maris, verētur.

``Cūr Neptūnum

invocās?`` inquit Mēdus, ''Prope 105 tranquillum est mare.``..

Gubernātor: ``Adhūc tranquillum est, sed exspectā \navnote{paulum = paulum temporis tonitrus -ūs m = quod audītur post fulgur}paulum: simul cum illīs nūbibus ātris tempestās orīri solet cum tonitrū et fulguribus. Neptūnum invocō, quod ille dominus maris ac tempestātum est.`` Nauta perterritus tempestātem venientem opperītur.

Paulō post tōtum caelum ātrum fit, ac fulgur ūnum et alterum, tum multa fulgura caelum et mare illūstrant. Statim sequitur tonitrus cum imbre, et simul magnus ventus flāre incipit. Mare tempestāte turbātur, ac nāvis, \navnote{iactāre = iacere atque iacere haurīre propter prp+acc}quae et hominēs et mercēs multās vehit, flūctibus iactātur et vix gubernāri potest. Nautae multum aquae ē nāve hauriunt, sed nāvis nimis gravis est propter mercēs. Hoc vidēns gubernātor ``Iacite mercēs!'' inquit nautīs, qul statim mercēs gravēs in mare iacere incipiunt, spectante-mercātōre, quī ipse quoque nāve vehitur. Ille \navnote{spectant mercātōre = dum mercātor spectat...}trīstis mercēs suās dē nāve labi et in mare mergī videt. Nēmō eum cōnsōlātur! Nāvis paulō levior fit, simul vērō \navnote{vērō = sed simul fit fiunt servāre}tempestās multō turbidior et flūctūs multō altiōrēs ftunt.

Mēdus perterritus

exclāmat: ``Ō

Neptūne!

Servā mē!`` sed vōx eius vix audītur propter tonitrum. Nāvis aquā implērī incipit, neque enim nautae satis multum aquae haurīre possunt.

Cēteris perterritīs, Lydia caelum intuētur et clāmat: \navnote{cēteris perterritīs = dum cēteri perterritī sunt}``Servā nōs, domine!''

Mēdus: ``Quis est ille dominus quem tū invocās?''

Lydia: ``Est dominus noster Iēsūs Chrīstus, quī nōn modo hominibus, sed etiam ventīs et mari imperāre potest.''

Mēdus: ``Meus dominus nōn est ille! Ego iam nūllius \navnote{nūllus -a -um, gen -ius}dominī servus sum. Nēmō mihi imperāre potest!``

Hic magnus flūctus nāvem pulsat. Mēdus Lydiam lābentem complectitur ac sustinēre cōnātur, nec vērō \navnote{nec vērō = sed nōn}ipse pedibus stāre potest. Mēdus et Lydia simul lābuntur. Mēdus surgere cōnātur, nec vērō sē locō movēre 140 \navnote{locō (abl) movēre = c locō movēre}potest, quod Lydia perterrita corpus eius complectitur.

\navnote{iterum = rūrsus fit fiunt, īnf fieri}Lydia iterum magnā vōce Christum invocat: ``Ō Chrīste! Iubē mare tranquillum fieri! Servā nōs, domine!''

Mēdus ōs aperit ac Neptūnum iterum invocāre vult, sed magnus flūctus ōs eius aquā implet. Mēdus loquī 145 cōnātur neque potest.

Tum vērō ventus cadere incipit! Iam flūctūs nōn tantī sunt quantī paulō ante. Nautae fessī aquam haurire dēsinunt ac laetantēs mare iterum tranquillum fieri aspiciunt.

\booksection{Grammatica Latīna}

\navnote{dēpōnēns -entis (dēp)}Verba deponentia

Mēdus laetātur (7 gaudet) nec dominum verētur (= timet). Nauta nōn proficīscitur (7 abit), sed ventum secundum opperītur (= exspectat).

'Lāetorf, Verērf, proficīscr,'opperfrT verba dēpōnentia sunt. Verbum dēpōnēns est verbum quod semper fōrmam verbī passīvī habet (praeter participium: laetāns,

verēns,prō- ficiscēns, opperiēns) atque in locō verbī āctlvī pōnitur.

Alia exempla verbōrum dēpōnentium: cōnsolāri, cōnāri, 160 intuērī, sequī, loquī, lābi, complectī, ēgredī, oriri.

Mēdus ex Italiā proficiscens sōlem orientem intuētur.

Nāvis ē portū ēgreditur; aliae nāvēs eam sequuntur. Dum Lydia loquitur, lacrimae dē oculīs eius lābuntur. Mēdus eam complecttrtir et consolāri cōnātur.

\booksection{Pēnsum A}

Nautae, quī Neptūnum ver-, tempestātem opper-. Mēdus et Lydia ex Italiā proficīsc-. Lydia Mēdum in puppim ascen tdentem sequ-. Dum nauta loqu-, Mēdus occidentem intu-, unde nūbēs ātrae ori-. Simul tempestās or-. Mercātor tristis: mercēs suās in mare lāb- videt; nēmō eum cōnsōl- potest. i

Quī ab amīcīs proficīsc- laet- nōn potest. t

\booksection{Pēnsum B}

Brundisium est oppidum --- (in --- maritimā situm) quod magnum --- habet. Portusest --- quō--- ad terram adire possunt.?

Cum nūllus --- flat, mare --- est. --- est magnus ventus qul īnfluere mare---etaltōs---facit.Gubemātorest---quīin---nāvissedēns flāre nāvem---.Partēscaellsunt---et---,---etmeridiēs. Oriēnsestea \contr{} caelī pars unde sōl ---, occidēns est ea pars quō sōl ---.

Mēdus et Lydia nāvem --- atque ex Italiā ---. Ventus --- est, gubernāre nāvis plēnīs --- ex portū ---; sed in altō mari tempestās ---: c magnus ventus --- incipit et tōtum mare turbidum ---. Nautae c aquam ē nāve--- --- et --- in mare iaciunt. Mēdus Neptūnum ---: sequī UŌ Neptūne! --- mē!`` Nāvis flūctibus --- nec --- [= sed nōn] laetāri mergitur.

\booksection{Pēnsum C}

Quae nāvēs Rōmam adire possunt? Quid est 'ōstium' flūminis? Num Tiberis in mare Superum Influit? Quandō nāvēs ē portū ēgrediuntur? Ubi sedet gubernātor et quid agit? Quae sunt quattuor partēs caelī? Quās rēs Mēdus et Lydia sēcum ferunt? Quō Mēdus cum amicā suī Īre vult? Cūr trīstis est Lydia? Quem deum invocant nautae? Cūr mercēs in mare iaciuntur? Num nāvis eōrum mergitur?

\booksection{Vocābula nova}
\noindent\small īre, eo, Imus, īs, ītis, it, eunt, laetārī= --- laetus esse, gaudēre, verēri = timēre, egredi, ēgrediuntur, oriri oritur Vocābula nova, nāvis, portus, locus, ōra, merx, nauta, ventus, tempestās, flūctus, vēlum, puppis, gubernātor, oriēns, occidēns, septentriōnēs, altnm, tonilms, fulgur, situs, superus, īnferus, maritimus, tranquillus, turbidus, contrārius, serēnus, āter, interesse, appellāre.

\bookchapter{Cap. XVII}{Numerī difficilēs}

\navnote{discipulus digitīs computat}--- / VMMESMNEE QISSMEPE MEEMESEEN Sus ERR NVMERI DIFFICILES

In lūdō pueri numerōs et litterās discunt. Magister pue- / \navnote{discere: discipulus discit docēre: magister docet}rōs numerōs et litterās docet. Magister puerōs multās rēs docēre potest, nam is multās rēs scit quās pueri \navnote{ne-scīre --- nōn scīre in-doctus = nōn doctus VN v! /}nesciunt. Magister est vir doctus. Pueri adhūc indoctī sunt. Quī paulum aut nihil discere potest, stultus esse dīcitur. Quī discere nōn vult atque in lūdō dormit, piger esse dīcitur. Discipulus quī nec stultus nec piger est, sed prūdēns atque industrius, multās rēs ā magistrō \navnote{prūdēns -entis --- \contr{} stultus industrius -a -um \contr{} \contr{} piger -gra -grum}discere potest. --- Magister Diodōrus recitāre dēsinit et puerōs aspicit, \navnote{quis-que \contr{} et ūnus et alter et tertius...; acc quem-que quisque ante suam sellam nēmō eōrum = nēmō ex iīs}quī taciti stant ante suam quisque sellam; nēmō eōrum dormit. Magister discipulum quemque cōnsīdere iubet, primum Sextum, deinde Titum, postrēmō Mārcum. . Magister: ``Nunc tempus est numerōs discere. Prī- i5 mum dic numerōs ā decem ūsque ad centum!'' \navnote{dīcere, imp dic! dīcite! tollere}Quisque puer manum tollit, primum Sextus, tum Ti tus, postrēmō Mārcus. Magister Titum interrogat.

Titus: ``Decem, ūndecim, duodecim, trēdecim, quattuordecim, quīndecim, sēdecim, septendecim, duodē-

Hīc magister eum interpellat: ``Necesse nōn est omnēs numerōs dīcere; ā vīgintī dīc decimum quemque numerum tantum, Tite!''

Titus: ``Vīgintī, trīgintā, quadrāgintā, quīnquāgintā, sexāgintā, septuāgintā, octōgintā, nōnāgintā, centum.''

Magister: ``Bene numerās, Tite. Iam tū, Sexte, dic numerōs ā centum ad mīlle!''

Sextus: ``Longum est tot numerōs dīcere!'' 30

Magister: ``At satis est decem numerōs dīcere, id est centēsimum quemque.''

Sextus: ``Centum, ducentī, trecentī, quadringentī, quīngentī, sescentī, septingentī, octingentī, nōngentī, mille.''

Magister: ``Quot sunt trīgintā et decem?''

Omnēs pueri simul respondent, nēmō eōrum manum tollit. Titus et Sextus ūnō ōre dīcunt: ``Quadrāgintā.'' Mārcus vērō dīcit: ``Quīnquāgintā.'' Quod Mārcus dīcit rēctum nōn est. Respōnsum Titi et Sextī rēctum est: iī \navnote{respōnsum -In = quod respondētur}rēctē respondent.

Respōnsum Mārcī est prāvum: is prāvē respondet. \navnote{laudāre}Il

Magister Titum et Sextum laudat: ``Vōs rēctē respondētis, Tite et Sexte. Discipulī prūdentēs atque industriī estis.'' Deinde Mārcum interrogat: ``Quot sunt trigintā et septem?'' et septem sunt trigintā septem.``

\navnote{trigintā septem z xxxvii (37)}``Rēctē respondēs, Mārce'' inquit magister, ``Quot sunt trīgintā et octō?''

Mārcus: ``Facile est ad hoc respondēre.''

Magister: ``Nōn tam facile quam tū putās!''

Mārcus: Trigintā et octō sunt trigintā octō.``

\navnote{duo-dē-quadrāgintā}Magister: ``Prāvē dīcis! Trigintā et octō sunt duodēquadrāgintā. Quot sunt trigintā et novem?''

Mārcus: ``Trigintā novem.''

Magister: ``Ō Mārce! Id quoque prāvum est. Trigintā \navnote{ūn-dē-quadrāginiā cōgitāre}et novem sunt undēquadrāgintā. Cūr nōn cōgitās antequam respondēs?``

Mārcus: ``Semper cōgitō antequam respondeō.''

Magister: ``Ergō puer stultus es, Mārce! Cōgitāre nōn 60 \navnote{stultē respondēs reprehendere \contr{} laudāre numquam = nūllō tempore numquam \contr{} \contr{} semper}potes! Nam stultē et prāvē respondēs!`` Magister Mārcum nōn laudat, sed reprehendit.

Mārcus: ``Cūr ego semper ā tē reprehendor, numquam laudor? Titus et Sextus semper laudantur, numquam reprehenduntur.''

Magister: ``Tū ā mē nōn laudāris, Mārce, quia numquam rēctē respondēs. Semper prāvē respondēs, ergō reprehenderis!''

\navnote{dif-ficilis -e \contr{} facilis}Mārcus: ``At id quod ego interrogor nimis difficile est. Atque ego semper interrogor!''

Magister: ``Tu nōn semper interrogāris, Mārce. Titus \navnote{saepe = multīs temporibus}et Sextus saepe interrogantur.``

Mārcus: ``Nōn tam saepe quam ego.''

Titus: ``Nōs quoque saepe interrogāmur, nec vērō prāvē respondēmus. Itaque nōs ā magistrō laudāmur, nōn reprehendimur.''

Mārcus: ``Et cūr vōs semper laudāmini? Quia id quod vōs interrogāminī facile est --- ego quoque ad id rēctē respondēre possum. Praetereā magister amīcus est patribus vestrīs, patrī meō inimīcus. Itaque vōs numquam reprehendiminī, quamquam saepe prāvē respondētis; \navnote{quamquam \contr{} \contr{} quia}ego vērō numquam laudor, quamquam saepe rēctē respondeō!``

Sextus: ``At frāter tuus saepe laudātur, Mārce.''

Magister: ``Rēctē dīcis, Sexte: Quīntus bonus discipulus est, et industrius et prūdēns.'' Quīntus ā magistrō laudātur, quamquam abest. HI

Sextus: ``Nōnne tū laetāris, Mārce, quod frāter tuus etiam absēns ā magistrō laudātur?'' 90

Mārcus: ``Ego nōn laetor cum frāter meus laudātur! Et cūr ille laudātur hodiē? Quia discipulum absentem reprehendere nōn convenit!''

Magister īrātus virgam tollēns ``Tacēte, puerī!'' in-quit, ``Nōnne virgam meam verēminl?''

Mārcus: ``Nōs nec tē nec virgam tuam verēmur!''

\navnote{dēnārius NSvasN c prōmere = sūmere (ē locō)}Magister, quī verba Mārcī nōn audit, ex sacculō suō duōs nummōs prōmit, assem et dēnārium, et ``Ecce'' inquit nummōs ostendēns ``as et dēnārius. Ut scītis, \navnote{i sēstertius = iv assēs I dēnārius = IV sēstertil}ūnus sēstertius est quattuor assēs et ūnus dēnārius quattuor sēstertiī. Quot assēs sunt quattuor sēstertiī, Tite?``

Magister: ``Et decem dēnāriī quot sestertii sunt?''

Titus: ``Quadrāgintā.''

Magister:

``Quot

dēnāriī sunt

duodēquīnquāgintā \navnote{duo-dē-quīnquāgintā}sēstertiī?``

Titus magistrō interrogantī nihil respondet.

Magister: ``Respōnsum tuum opperior, Tite. Cūr mihi nōn respondēs?''

Titus: ``Nōndum tibi respondeō, quod primum cōgi- 110 tāre oportet.

\navnote{oportēre: oportet = neccsse est, convenit}Duodēquīnquāgintā est numerus difficilis!``

Sextus manum tollēns ``Ego respōnsum sciō'' inquit.

\navnote{Sextus: ``ligo respōnsum sciō'' Sextus dīcit 'se respōnsum scīre'}Magister: ``Audī, Tite: Sextus dīcit 'sē respōnsum scīre.' Sed exspectā, Sexte! Nōn oportet respondēre antequam interrogāris.''

quattuordecim? an quīndecim? Nōn certus sum!`` Titus \navnote{in-certus = nōn certus}nōn certus, sed incertus est, et respōnsum incertum dat magistrō.

Magister: ``Respōnsum certum magistrō dare oportet. Respōnsum incertum nūllum respōnsum est. Nunc tibi licet respondēre, Sexte.''

Quamquam difficilis est numerus, Sextus statim rēctē respondet:

``Duodēquīnquāgintā

sēstertiī sunt decim dēnāriī.``

Magister iterum Sextum laudat: ``Bene computās, Sexte! Semper rēctē respondēs. Ecce tibi assem dō.'' Diodōrus assem dat Sextō, et dēnārium in sacculō repōnit.

Mārcus: ``Discipulō tam industriō dēnārium dare oportet. Largus nōn es, magister. Pecūniam tuam nōn \navnote{largus -a -um = quī multum dat largīri z multum dare}largīris.``

Diodōrus: ``Nōn oportet pecūniam largīri discipulis.'' 135

Mārcus:

``Quārē igitur tot numerōs

difficilēs discimus?``

\navnote{dē-mōnstrāre = mōnstrāre (verbīs)}Magister: ``Quia necesse est computāre scīre, ut multis exemplis tibi dēmōnstrāre possum. Si tria māla ūnō asse cōnstant et tū decem assēs habēs, quot māla emere

Mārcus: ``Tot assēs nōn habeō --- nōn tam largus est pater meus; nec mihi licet tot māla emere.''

Magister: ``Id nōn dīcō. Ecce aliud exemplum: Si tū sex assēs cum frātre tuō partīris, quot sunt tibi?''

\navnote{partīrī= --- dividere (in partēs)}Mārcus: ``Quīnque!''

Magister: ``In aequās partēs pecūniam partiri oportet. S1 vōs sex assēs aequē partīminī, quot tibi sunt?''

\navnote{aequē partīri = in aequās partēs dividere}Mārcus: ``Pecūniam meam exiguam nōn largior nec partior cum aliīs!''

Postrēmō magister ``Ō Mārce!'' inquit, hōc modō nihil tē docēre possum! Tam stultus ac tam piger es quam asinus!``

\booksection{Grammatica Latīna}

Persōnae verbī

Ego ā magistrō laudor = magister mē laudat. 155

Tu ā magistrō laudāns = magister tē laudat.

Nōs ā magistrō laudāmur = magister nōs laudat.

Vōs ā magistrō laudamini = magister vōs laudat.

\navnote{sing. plūr. I -or -mur II -ns -minī HI -tur -ntur III -tur}Persōna prīma laudor ------laudāmur 160 Persōna secunda

laudāris

laudāmini Persōna tertia

laudātur

Exempla:

[1] iactor iactāmur; [2] videor vidēmur; [3] mergor mergimur; [4] audior audimur.

Nauta, quī flūctibus iactatur ac mergitur: ``Flūctibus iactor 165 \navnote{iactā mur iactor iactā ris iactā mmf iactā tur iactā uittr videlor vidēmur vidēris vidē minī vidētur videntur}ac mergor/ Servā mē, Neptūne, quī prope ades, quamquam nec videris nec audfm ab hominibus.`` Neptūnus: ''Flūctibus iactans, nauta, nec vērō mergeris/ Nam ego adsum, quamquam nec videor nec audior.`` Deus ab hominibus nec vidētur nec audītur.

\navnote{mergor merg mur merg ens merg vninī merg itur merg untur audior audimur audiris audī minī audi unmr audī tur}Nautae, quī flūctibus iactantur ac merguntur: ``Flūctibus iactāmur ac mergimur/ Servāte nōs, ō deī bonī, quī prope adestis, quamquam

nec vidēminī nec audimini ab hominibus.`` Del: ''Flūctibus iactāmini, nautae, nec vērō mergtmim/ Nam nōs adsumus, quamquam nec vidēmur nec audīmur.`` 175 Deī ab hominibus nec videntur nec audiuntur. [1] Ut 'laudor laudāmur* dēclīnantur verba passīva et dēponentia quōrum infinitivus dēsinit in -ān. [2] Ut videor vidēmur' dēclīnantur verba passīva et dēpōnentia quōrum infinitivus dēsinit in -m. [3] Ut 'mergor mergimur' dēclīnantur verba passīva et dēponentia quōrum infinitivus dēsinit in -t (praeter ea quae -ihabent ante -or et -untur, ut 'capior capiuntur, ēgredior ēgrediuntur*). [4] Ut 'audior audīmur' dēcllnantur verba passīva et dēpōnentia quōrum infinitivus dēsinit in -ir? (sed oriri oritur').

Exempla verbōrum

dēpōnentium:

[1] laetor laetālmur; [2] verelor verēmur; [3] sequor sequimur; [4] partior parūmur.

Lydia: ``Ego Neptūnum nōn vereor. Tūne eum verēris, nauta?'' Nauta: ``Omnēs nautae deum maris verēmur.'' Lydia: * ``Cūr Neptūnum verēmtm?''

Nautae Neptūnum verentur, Lydia eum nōn verētur.

\booksection{Pēnsum A}

Mārcus ā magistrō nōn laud-, sed reprehend-. Mārcus: ``Cūr ego semper reprehend-, numquam laud-? Titus et Sextus semper laud-, numquam reprehend-.'' Magister: ``Tu nōn laud-, sed reprehend-, quia prāāvē respondēs.'' Mārcus: ``Sed: ego semper interrog-!'' Sextus: ``Tu nōn semper interrog-;! nōs saepe interrog- nec prāvē respondēmus: itaque ā magistrō laud-, nōn reprehend-.'' Mārcus: ``At facile est id quod vōs interrog-: itaque rēctē respondētis ac laud-!'' Titus: ``Nōs 1 magistrum ver-. Nōnne tū eum ver-?'' Mārcus: ``Ego magistrum nōn ver-. Cūr vōs eum ver-?'' 1 i

\booksection{Pēnsum B}

Discipulus est puer quī ---. Vir quī puerōs --- magister est. certus Magister est vir --- quī multās rēs --- quās pueri ---. Sextus largus nec stultus nec piger, sed --- atque --- est. Magister interrocentēsimus gat, discipulus quī respondēre potest manum ---. Nōn --- quattuordecim respondēre antequam magister interrogat. Magister: ``--- nuquīndecim merōs ā decem --- ad centum!'' --- puer manum tollit, priseptendecim mum Sextus, deinde Titus, --- Mārcus. Sextus numerōs dlcit: duodēvīgintī ``X, XI, XII, XIII, XIV, XV, XVI, XVII, XVIII, XIX, XX, XXX, XL, quadrāgintā L, lx, lxx, lxxx, xc, c.'' Sextus bonus discipulus est, quī quīnquāgintā semper --- respondet; itaque magister eum ---. Mārcus, quī c semper --- respondet, ā magistrō ---: ``--- rēctē respondēs, c Mārce! Cūr nōn --- antequam respondēs?'' Mārcus: ``Cūr Tiseptingentī tum nōn interrogās? Is nōn tam --- interrogātur quam ego.'' octingentī Magister: ``Iam eum interrogō: quot sunt xxi et lxviii?'' Tidiscere tus: ``lxxxvii? an lxxxviii?'' Titus nōn ---, sed --- est ac! d magistrō --- incertum dat. Sextus vērō rēctē respondet, --- difficile est id quod interrogātur. Quod Mārcus interrogātur nōn difficile, sed --- est.

Ūnus --- est IV sēstertii, ūnus sēstertius iv ---.

\booksection{Pēnsum C}

Quid pueri discunt in lūdō? Cūr Sextus ā magistrō laudātur? Cūr Mārcus reprehenditur? Estne facile id quod Titus interrogātur? Cūr Titus nōn statim respondet? Uter rēctē respondet, Titusne an Sextus? Quot sunt ūndētrigintā et novem? Trēs sēstertiī quot sunt assēs? Quot dēnāriī sunt octōgintā sēstertiī? Num facilēs sunt numeri Rōmānī? Tune ā magistrō tuō laudāris?

\booksection{Vocābula nova}
\noindent\small x, decem, xi, ūn-decim, xii, duo-decim, xiii, trē-decim, xiv, quattuor-decim, XV, quīn-decim, xvi, sē-decim, XVII, septen-decim, XVIII. duo-de-vīgintī, xix, ūn-dē-vīginiī, xx, vīgintī, xxx, trigintā, XL, quadrāgintā, L, quīnquāgintā, lx, sexāgintā, lxx, septuāgintā, lxxx octōgintā, xc, nōnāgintā, c, centum, tot (indēcl)= tam multī, CC, ccc, cccc quadringentī -ae -a, d, DC, DCC, dccc octin-gentī, DCCCCnÓn-gent -ae -a M/CIO mille, dare, dō, damus, dās, datis, dant, re-pōnere = rūrsus, pōnere, -or, -ris, -mur, -minī, -ans, -(tmini, -ātur, -untur, -ior, -imur, -eris, -iminī, itur, -iuntur, -īntuv, -īmini, -itur, vereor, veremi/r, verēns, verēminf, yerētur, verentur --- timet, respōnsum, as, dēnārius, doctus, indoctus, piger, prūdēns, industrius, rēctus, prāvus, facilis, scīre, nescīre, tollere, interpellāre, laudāre, cōgitāre, reprehendere, prōmere, oportēre, computāre, repōnere, largīri, dēmōnstrāre, partiri, quisque, tot, postrēmō, rēctē, prāvē, aequē, ūsque, numquam, saepe, quamquam, quārē.

\bookchapter{Cap. XVIII}{Litterae Latīnae}

ee Ae dk LITTERAE LATINAE

\navnote{litterae dīcuntur}I Discipulī nōn modo numerōs, sed etiam litterās dis-cunt. Parvī discipulī, ut Mārcus et Titus et Sextus, litterās Latīnās discunt. Magnī discipulī litterās Graecās et linguam Graecam discunt. Lingua Graeca difficilis est.

Ecce omnēs litterae Latīnae, quārum numerus est vīginti trēs, ab A ūsque ad z: A, B, c, D, E, f, g, h, i, k, l, m, N, o, P, Q,r, s, T, v, x, Y, Z. Litterae sunt aut vōcālēs aut cōnsonantēs: vōcālēs sunt A, E, i, o, v, y; cēterae \navnote{(littera) vōcālis -is f (littera) cōnsonāns}sunt cōnsonantēs. (Etiam i et v cōnsonantēs sunt in VOCābulīs IAM, VEL, VŌS, QVAM, Cēt.)

\navnote{cālis (cōnsonāns V v)]}Litterae Y et z in vocābulīs Graecis modo reperiuntur (ut in hōc vocābulō zephyrus, id est nōmen ventī quī ab occidente flat). v et z igitur litterae rārae sunt in linguā \navnote{frequēns -entis \contr{} \contr{} rārus}Latīnā, in linguā Graecā frequentēs. K littera, quae frequēns est in linguā Graecā, littera Latīna rārissima est, nam K in ūnō vocābulō Latino tantum reperitur, id est \navnote{Kaesō -ōnis m}kalendae (itemque in praenōmine Kaesō, quod praenōmen Rōmānōrum rārissimum est).

Hoc vocābulum amica quīnque litterās habet et trēs \navnote{quae-que syllaba:: et prima et secunda et tertia syllaba idem ea-dem idem ( \contr{} is merus (m), eadem littera (f), idem vocābulum per sē:: sōla}syllabās: a-mī-ca. Quaeque syllaba vōcālem habet, ergō 20 numerus syllabārum et vōcālium īdem est. In primā et in postrēmā syllabā huius vocabuli eadem vōcālis est: a. Vōcālis est littera quae per sē syllabam facere potest, ut a syllabam prīmam facit in vocābulō amīca. Sine vōcālī \navnote{facere/fierī (āct/pass): vōcālis syllabam facere potest; sine vōcālī syllaba fieri nōn potest iungere \contr{} dividere}syllaba fieri nōn potest. Cōnsonāns per sē syllabam nōn 25 facit, sed semper cum vōcālī in eādem syllabā iungitur. In exemplō nostrō m cum f iungitur in syllabā secundā mī, et c cum a in syllabā tertiā ca.

Cum syllabae iunguntur, vocābula fiunt. Cum vocābula coniunguntur, sententiae fīunt. Ecce duae senten- 30 \navnote{con-iungere = īungere}tiae: Lingua in ōre inest et Lingua Latīna difficilis est. \navnote{uter- utra- utrum-que, gen utrius-que}Vocābulum primum utriusque sententiae idem est, sed hoc idem vocābulum duās rēs variās significat. Item. \navnote{idem ea-dem idem, acc eun-dem ean-dem idem erus -Im = dominus}varia vocābula eandem rem vel eundem hominem significāre possunt, ut ōstium et iāmia, dominus et erus (sed 35 erus est vocābulum multō rārius quam dominus).

Quī litterās nescit legere nōn potest. Magister, quī puerōs legere docet, ipse et librōs Latīnōs et Graecōs \navnote{utraque lingua:: lingua Graeca et Latīna quis- quae- quod-que, acc quem- quam- quodque, gen cuius-que sīc= - hōc modō}legere potest, nam is utramque linguam scit. Quōmodo parvus discipulus hanc sententiam legit: Mīles Rōmānus 40 fortiter pugnat? Discipulus quamque litteram cuiusque vocāābuli sīc legit: M-1 ``m-i ml, l-e-s les: mī-les, R-o Rō, m-a \navnote{ita = eō modō}cābulum cuiusque sententiae ā discipulō legitur. (In hāc sententiā vocābulum fortiter novum est, sed quī vocābulum fortis scit, hoc vocābulum quoque intellegit, nam \navnote{intellegere fortis fortiter: mīles fortis fortiter pugnat}mīles fortis est mīles quī fortiter pugnat.)

In lūdō puerī nōn modo legere, sed etiam scribere discunt. Quisque discipulus in tabulā suā scribit eās sententiās quās magister ei dictat. Ita pueri scrībere dis-cunt. U ki Vti ANS NA Vu 2. pret: JA `` h d '' I [ /A: n

\navnote{discipulī scribunt magistrō dictante}Magister discipulis imperat: ``Prōmite rēgulās vestrās \navnote{līnea rēcta}et lineas rēctās dūcite in tabulīs. Tum scribite hanc sententiam: Z7omo oculōs et nāsum habet'' 55

Quisque puer stilum et rēgulam prōmit et dūcit līneam rēctam in tabulā suā; tum scribere incipit. Discipulī eandem sententiam nōn eōdem modō, sed variis modīs \navnote{ūnus:: sōlus trēs tria, abl tribus}scrībunt. Sextus ūnus ex tribus pueris rēctē scrībit: HOMŌ OCVLOS ET NASVM HABET. TitUS sīc scribit: HOMŌ HOCVLOS ET NASVM HABET. Mārcus VērŌ sīc: OMO OCLOS ET NASV ABET. Il

Magister: ``Date mihi tabulās, pueri!'' cuiusque pueri in manus sūmit litterāsque eōrum aspicit. Quālēs sunt litterae Sextī? Pulchrae sunt. Quālēs 65 sunt litterae Mārcī et Titī? Litterae eōrum foedae sunt ac vix legī possunt. Magister suam cuique discipulō tabulam reddit, primum Sextō, tum Titō, postrēmō Mārcō, atque ``Pulchrē et rēctē scribis, Sexte'' inquit, ``Facile est tālēs litterās legere. At litterae vestrae, Tite \navnote{tālis -e (= quī ita est)}et Mārce, legi nōn possunt! Foedē scrībitis, pigrī discipulī!`` Magister Titum et Mārcum sevērē reprehendit.

Titus: ``Certē pulcherrimae sunt litterae Sextī, sed \navnote{pulcher, comp pulchrior, sup pulcherrimus -a -um}meae litterae pulchriōrēs sunt quam Mārcī.``

Magister: ``Reddite mihi tabulās, Tite et Mārce!''

Titus et Mārcus tabulās suās reddunt magistrō, qul \navnote{simul = eōdem tempore}eās simul aspicit. Magister litterās Titī comparat cum litteris Mārcī, et ``Litterae vestrae inquit ''aequē foedae sunt: tū, Tite, neque pulchrius neque foedius scrībis quam Mārcus.``

Titus: At certē rēctius scribo quam Mārcus.``

Magister: ``Facile est rēctius quam Mārcus scrībere, nēmō enim prāvius scribit quam ille! Nōn oportet sē comparāre cum discipulō pigerrimō ac stultissimō! \navnote{piger, comp pigrior, sup pigerrimus -a -um}Comparā tē cum Sextō, quī rēctissimē et pulcherrimē scrībit.`` Tum sē vertēns ad Mārcum: ''Tu nōn modo foedissimē, sed etiam prāvissimē scrībis, Mārce! Nescīs Latīnē scribere! Puer pigerrimus es atque stultissimus!`` Iam Mārcus multō sevērius reprehenditur quam Titus.

Mārcus (parvā vōce ad Titum): ``Magister dīcit 'mē prāvē scribere': ergō litterās meās legere potest.''

\navnote{ex-audire = audire (quod parvā vōce dīcitur) turpis -e = foedus de-esse (: ab-esse) de-est dē-sunt}At magister, quī verba Mārcī exaudit, ``Litterās tuās turpēs'' inquit ``legere nōn possum, sed numerāre possum: quattuor litterās deesse cernō. Aspice: in vocābulō primō et in vocābulō postrēmō eadem littera H deest.''

Mārcus: ``At semper dīcō omo abet''

Magister: ``Nōn semper idem dīcimus atque scribimus. In vocābulō secundō V deest, in quārtō m. Quid significant haec vocābula turpia oclōs et nāsu? Talia verba Latīna nōn sunt! Nūllum rēctum est vocābulum praeter ūnum et, atque id vocābulum est frequentissimum \navnote{facilis, comp facilior, sup facillimus -a -um mendum -īn: m. facere = prāvē scribere}et facillimum! Quattuor menda in quīnque vocābulīs! Nēmō alter in tam brevī sententiā tot menda facit!``

Magister stilō suō addit litterās quae dēsunt; ita menda corrigit. Tum vērō ``Nec sōlum'' inquit ``prāvē et \navnote{corrigere- rēctum facere leviter scribere = levī manū scribere E}turpiter, sed etiam nimis leviter scribis. Hās lineas tenuēs vix cernere possum. Necesse est tē stilum gravius in cēram premere.`` (Discipulī in cērā scribunt, nam tabulae eōrum cērā operiuntur. Cēra est māteria mollis \navnote{apis -is f bēstiola -ae f --- parva bēstia}no quam apēs, bēstiolae industriae, faciunt.)

Mārcus: ``Stilum graviter premō, sed cēra nimis dūra est. Aliam tabulam dā mihi! Haec cēra prope tam dūra est quam ferrum.'' (Ferrum est māteria dūra ex quā \navnote{ef-ficere -iō = facere}cultrī, gladiī, stili aliaeque rēs multae efficiuntur.)

Magister Mārcō eandem tabulam reddit et ``Cēra tuaM inquit ''tam mollis est quam Sextī, et facile est eius litterās legere. Sūme tabulam tuam et scribe h litteram deciēs!``

Mārcus deciēs H scribit HHHHHHHHH H. \navnote{quotiēs? totis:}magister eum v sexiēs scrībere iubet, et Mārcus, quī 120 eandem litteram totiēs scribere nōn vult, v v v v scrībit. \navnote{p}Quotiēs Mārcus v scribit? Mārcus, ut piger discipulus, quater tantum v scribit. Deinde magister eum tōtum vocābulum nasvm quīnquiēs scribere iubet, et Mārcus scribit nasvm nasvm nasvm nasv nasv. Mārcus ter 125 rēctē et bis prāvē scrībit.

Tum Titus, quī duās litterās deesse videt, sīc incipit:

rum in partem corporis eius mollissimam premit! Titus tacet nec fīnem sententiae facere audet. Magister vērō 130 hoc nōn animadvertit.

\navnote{animadvertere = vidēre vel audire}Magister: ``Iam tōtam sententiam rēctē scribe!''

Mārcus interrogat: ``Quotiēs?''

Magister breviter respondet: ``Semel.''

Mārcus tōtam sententiam iterum ab initiō scrībere 135

Magister: ``Quid significat hoculōs? Illud vocābulum turpe nōn intellegō!''

Mārcus: ``Num hic quoque littera deest?''

Magister: ``Immō vērō nōn deest, at superest H litte- 140 \navnote{super-esse \contr{} \contr{} de-esse}ra! Num tū hoculōs dīcis?``

Mārcus, ut puer improbus, magistrō verba sua red-dit: Nōn semper dīcimus idem atque scrībimus!``

\navnote{dēlēre --- addere TS n}Magister:

``Tacē, improbe!

Dēlē illam litteram!`` Mārcus stilum vertit et litteram H dēlet. Simul Titus idem mendum eōdem modō corrigit in suā tabulā, ne-que vērō magister hoc animadvertit.

Postrēmō Mārcus ``Quārē'' inquit ``nōs scrībere docēs, magister? Mihi necesse nōn est scribere posse. Numquam

domī scrībō.``

Magister: ``Num pater tuus semper domi est?''

Mārcus: ``Nōn semper. Saepe abest pater meus.''

Magister: ``Cum pater tuus abest, oportet tē epistulās ad eum scrībere.''

\navnote{Zēnō frequentēs:: multās im-piger = industrius}Sextus: ``Ego frequentēs epistulās ad patrem meum absentem scribō.'' Sextus puer probus est ac tam impiger quam apis.

\navnote{Zēnō -ōnis m Latīnē scīre = linguam Latīnam scīre īdem servus: is quoque servus}Mārcus: ``Ego ipse nōn scrībō, sed Zēnōnī dictō. Zēnō est servus doctus quī et Latīnē et Graecē scit. īdem servus mihi recitāre solet.''

Magister calamum et chartam prōmit et ipse scribere incipit; is enim calamō in chartā scrībit, nōn stilō in cērā ut discipulī. (Charta ex papyrō efficitur, id est ex altā \navnote{charta/ -ae}herbā quae in Aegyptō apud Nīlum flūmen reperitur Charta et papyrus vocābula Graeca sunt.)

Mārcus magistrum scribere animadvertit eumque interrogat: ``Quid tū scribis, magister?''

``Epistulam'' inquit ille ad patrem tuum scribō. Breviter scribo 'tē esse discipulum improbum'.`` 170

Magistrō scribente, Mārcus ``Prāvē scribis'' inquit, ``syllaba im superest. Dēlē illam syllabam et scribe 'discipulum probum''!

Magister: ``Tacē, puer improbissime! Nihil dēleō, immō vērō vocābulum addō: te discipulum improbum atque pigrum esse' scribō!''

``Scrībe 'probum atque impigrum' '' inquit Mārcus, nec vērō haec verba ā magistrō audiuntur.

Diodōrus, quī diem epistulae addere vult, discipulōs \navnote{epistulae (dat) addere = ad epistulam addere}interrogat: ``Quī diēs est hodiē?''

Titus: ``Hodiē kalendae Iūniae sunt.''

\navnote{kalendae Iūniae = diēs primus mēnsis Iūniī}Magister: ``Rēctē dīcis. Kalendae sunt hodiē. Ergō \navnote{mercēs -ēdis f.}date mihi mercēdem!`` (Mercēs est pecūnia quam magister quōque mēnse ā patribus discipulōrum accipit. Discipulī kalendis cuiusque mēnsis mercēdem magistrō suō afferre solent.)

Sextus et Titus statim magistrō mercēdem dant. Mārcus vērō mercēdem sēcum nōn fert. Magister, antequam epistulam signat, pauca verba addit.

\navnote{signāre (epistulam) = cērā claudere}``Quid nunc scribis?'' interrogat Mārcus.

``Scribō te mercēdem ad diem nōn afferre''' respondet magister, atque epistulam signat ānulum suum in cēram imprimēns. d

\booksection{Grammatica Latīna}

Adverbium

\navnote{stultē \contr{} stultus -a -um fortiter \contr{} fortis -e adverbium -īn (adv)}Discipulus stultus est quī stultē respondet.

Mīles fortis est quī fortiter pugnat.

``Stultus! et 'fortis' adiectīva sunt. 'Stultē' et 'fortiter* sunt adverbia.

Adiectīvum: -us -a -um (dēcl. i/ll). Adverbium: -ē. 200

\navnote{-iter}Adiectīvum: -is -e (dēcl. m). Adverbium: -iter.

Exempla:

Sextus rēctē et pulchrē scribit, Titus et Mārcus prāvē et turpiter scrībunt ac sevērē reprehenduntur. Magister Latīnē et Graecē scit. Sōl clārē lūcet. Magister breviter respondet. Puer crassus nōn leviter, sed graviter cadit.

Certē Titus rēctius et pulchrius scribit quam Mārcus, nēmō prāvius aut turpius scribit quam ille; Mārcus pxavissimē et turpissimē scrībit, Sextus recussime et pulcherrimē. Nēmō foruus pugnat quam Rōmānī: Rōmānī fortissimē pugnant.

--- 'Rēctus', 'forūus' est comparātīvus adverbil, 'rēctissimē9, 'fortissimē9 superlātīvus est. Comparātlvus:

-ius. Superlātīvus: -issimē.

\booksection{Pēnsum A}

Sextus rēct- respondet, Mārcus prāv- respondet et sevērreprehenditur. Nēmō rēct- aut pulchr- scribit quam Sextus; is rēctissim- et pulcherrim- scrībit. Lūna plēna clār- lūcet, sed sōl clār- lūcet quam lūna et cēterae stēllae; sōl clārissimlūcet.

Hostēs nōn tam fort- pugnant quam Rōmānī. Mllitēs nōstrī fort- pugnant quam hostēs. Nostrī fortissim- pugnant. Mārcus nimis lev- scrībit. Magister brev- scribit 'Mārcum improbum esse.'

\booksection{Pēnsum B}

Ex xxiiJ litteris Latinis VI --- sunt: AEIO V Y; cēterae sunt Cōnsonāns semper cum vōcālī---. K, Y, z litterae --- sunt in --- Latīnā, in linguā Graecā ---. Vocābula ōstium et iānua --- rem significant; lingua duās rēs --- significat. Magister librōs Latīnōs et Graecōs --- potest.

Discipulī hanc --- in tabulīs suīs---: Hamō oculōs et nāsum habet. Magister tabulam --- pueri aspicit, et Sextī et Titī et Mārcī. In tabulā Mārcī iv litterae ---. Magister litterās quae dēsunt ---; --- [= eo modō] magister menda Mārcī---. ---. Tum ``Ō Mārce! Hīc nōn ---, immō --- littera h!'' Mārcus stilum vertēns H litteram ---. Tum magister calamum et --- sūmit et --- ad patrem Mārcī scribit.

Cēra est --- mollis. Ferrum est māteria ---.

\booksection{Pēnsum C}

Quot sunt litterae Latīnae? Estne D vōcālis an cōnsonāns? Quot syllabās habet vocābulum 'apis'? Quōmodo parvus discipulus legit? Quis discipulus rēctissimē scribit? Quālēs sunt litterae quās Titus scribit? Quālis discipulus est Mārcus? Cūr Mārcus nōn homō, sed omo scribit? Quotiēs Mārcus litteram H scribit? Stilōne in tabulā scrībit magister? Quid magister scribit? Estne facilis grammatica Latīna? Tune rēctius scrībis quam Mārcus?

\booksection{Vocābula nova}
\noindent\small P, pē*, Q, R, er', e, s, t, v, y, 'ypsllon', z, 'zēta', magister, epistulam, signat, rēctē, rēctius, rēctissimē, fortiter, fortius, fortissimē, -ius, -issimē, nova, vōcālis, cōnsonāns, zcphyrus, sententia, crus, mendum, cēra, māteria, apis, ferrum, epistula, calamus, charta, papyrus, mercēs, rārus, frequēns, varius, turpis, mollis, dūrus, quālis, tālis, impiger, iungere, coniungere, significāre, legere, intellegere, scrībere, dictāre, comparāre, exaudīre, deesse, addere, corrigere, premere, efficere, animadvertere, superesse, delēre, signāre, imprimere, īdem, eadem, idem, quaeque, quodque, sīc, ita, quotiēs, totiēs, semel, bis, ter, quater, quīnquiēs, sexiēs, deciēs, adverbium.

\bookchapter{Cap. XIX}{Marītus et uxor}

\navnote{tēctum}\navnote{columnae}\navnote{signa}FL LA JULI IL LLL LLL OD LL D DLL LA MARITVS ET VXOR

I lūlius cum uxōre suā in peristylō ambulat. Uxor eius est \navnote{uxor -ōris f marītus Aemiliae = vir Aemiliae}Aemilia. Iūlius maritus Aemiliae est. Maritus et uxor inter columnās et signa ambulant. Tēctum peristylī altīs columnis sustinētur, inter quās tria pulchra signa stant: ūnum Iūnōnis, alterum Cupidinis, tertium Veneris. \navnote{5 Iūnō -ōnis f Cupīdō -inis m Venus -eris f con-iūnx -iugis f/m alter -era -erum, gen -erius}Iūnō et Venus deae sunt. Cupīdō est filius Veneris et Mārtis, quamquam Venus Mārtis coniūnx nōn est, sed alterius deī, cui nōmen est Vulcānus. Venus enim māla uxor est, quae aliōs deōs amat praeter coniugem suum Vulcānum. Mārs et Vulcānus sunt filii Iānonis et Iovis; \navnote{Iuppiter Iovis m}Iūnō enim coniūnx Iovis est. Sed Iuppiter malus maritus est, qul multās aliās deās amat praeter Iūnōnem, coniugem suam. Nēmō deōrum pēior marītus est quam \navnote{malus, comp pēior -ius, sup pessimus -a -um ūllus -a -um: neque ūllus}Iuppiter, neque ūlla dea pēior uxor est quam Venus. Inter omnēs deōs deāsque Iuppiter pessimus maritus est ac Venus pessima uxor.

Mātrēs familiās Iūnōnem invocant, ea enim dea est \navnote{māter familiās = māter (domina) familiae mātrōna -ae f --- femina quae marītum habet amor -ōris m \contr{} amāre}mātrōnārum. Venerem et Cupīdinem invocant amantēs, nam Venus et Cupīdō amōrem in cordibus hominum excitāre possunt. Cupīdō enim est deus amōris, et 20 Venus, pulcherrima omnium deārum, dea amōris ac \navnote{pulchritūdō -inis f. \contr{} pulcher omnēs}pulchritūdinis est. Pulchritūdō Veneris ab omnibus laudātur.

Iūlius est bonus maritus quī uxōrem suam neque ūllam aliam feminam amat. Certē Iūlius maritus melior 25 \navnote{bonus, comp melior -ius, sup optimus -a -um Aemilia: ``Iūlius bonus vir est, melior quam ūllus alius vir: Iulius est optimus omnium virōrum!''}quam Iuppiter est! Item Aemilia bona uxor est quae marītum suum neque ūllum alium virum amat. Certē Aemilia uxor melior est quam Venus! Aemilia Iūlium virum optimum' appellat, Item Iūlius uxōrem suam 'optimam omnium feminārum' vocat.

\navnote{vocāre = appellāre trēs tria, gen trium duo -ae -o, gen duorum -ārum -ōrum pater ūnīusfiliī/fHiae; duōrum ūliōTum/dnārum filiārum; trium filiōrum /filiārum magnus, comp māior -ius, sup māximus -a -um parvus, comp minor -us, sup minimus -a -um pater familiās --- pater (dominus) familiae adulēscēns -entis m = vir quī nōndum xxx annōs habet virgō -inis f --- femina quae maritum nōndum habet ab-erat domus}Iūlius et Aemilia sunt parentēs trium līberōrum: duōrum filiōrum et ūnīus fīliae. Liberi adhūc parvī sunt. Mārcus octō annōs habet. Quīntus est puer septem annōrum. Iūlia quīnque annōs habet. Quīntus nōn tantus est quantus Mārcus nec tam parvus quam Iulia. Quīn- 35 tus māior est quam Iūlia et minor quam Mārcus. Māximus līberōrum est Mārcus, minima est Iūlia.

Ante decem annōs Iūlius pater familiās nōn erat, tunc enim nec uxōrem nec līberōs habēbat. Iūlius adulēscēns vīgintī duōrum annōrum erat. Aemilia mātrōna nōn 40 erat, sed virgō septendecim annōrum, quae Rōmae apud parentēs suōs habitābat. Domus eōrum nōn prōcul aberat ab aliā domō, in quā Iūlius cum parentibus \navnote{S oS Capitoljūm z P, sS Forum `` Rōmānum SCN Qnm uu, z Palatium AJ 2m 3}suīs habitābat. Iūlius et Aemilia in eādem urbe habitābant, nōn in eādem domō.

In urbe Rōmā multae domūs sunt et multa templa \navnote{templum -īn = domus deī forum -i n---locus in urbe apertus quō multī conveniunt con-venīre = in eundem locum venīre Optimus Māximus: Iovis cognōmen magnus et pulcher multī, comp plūrēs -a, sup plūrimī -ae -a templum, LS f CM a ( Ü (j/ S dM il: ALPE forum}deōrum. In mediā urbe inter collēs Capitōlium et Palātium est forum Rōmānum, quō hominēs ex tōtā Italiā atque ex omnibus prōvinciīs Rōmānīs conveniunt. In Capitōliō est templum Iovis Optimī Māximī; circum forum Rōmānum sunt alia multa templa deōrum. In Palātiō sunt domūs magnificae. Rōmae plūrēs hominēs habitant quam in ūllā aliā urbe imperiī Rōmānī. Urbs Rōma plūrimōs hominēs et plūrimās domōs habet. Ariss minum, Capua, Ōstia magna sunt oppida, sed minōra quam Rōma. Ōstia est oppidum minus quam Rōma, sed māius quam Tusculum. Rōma urbs māxima atque pulcherrima est totius imperiī Rōmānī. \navnote{tōtus -a -um, gen -ius}ll

Pater Iūliī, quī iam mortuus est, magnam pecūniam habēbat multāsque villās magnificās possidēbat praeter \navnote{possidēre = habēre dīves -itis = pecūniōsus pauper -eris \contr{} dīves tamen:: quamquam ea pauper erat tamen \contr{} r+itaque}domum Rōmānam: is homō dīves erat. Pater Aemiliae erat homō pauper, quī ūnam domum exiguam possidēbat. Aemilia igitur virgō pauper erat, sed tamen Iūlius eam amābat. Cūr Iūlius, adulēscēns dīves, virginem pauperem amābat? Quia ea virgō proba ac fōrmōsā erat. Nec Iūlius amōrem suum occultābat, nam Aemiliam \navnote{dōnum-īn = itaque dōnum -i 7 --- rēs quae datur}'amīcam' appellābat et multa dōna ei dabat. Sed tamen \navnote{miser -era -erum}Aemilia nōn Iūlium amābat, sed alium virum Rōmānum. Ergō Iūlius miser erat et nocte male dormiēbat.

Vir quī ab Aemiliā amābātur Crassus Dīves nōminābātur. Aemilia Crassum amābat, neque vērō ab eō amābātur, quod parentēs Aemiliae pauperēs erant. Itaque ea quoque misera erat.

Aemilia numquam Iūlium salūtābat, cum eum in forō Rōmānō vidēbat, quamquam ipsa ā Iūliō salūtābātur. 75 \navnote{cotīdiē = omni diē, quōque diē}Iūlius cotīdiē epistulās ad Aemiliam scribēbat, in quībus pulchritūdō eius laudābātur verbīs magnificīs, ac simul cum epistulīs rosās aliōsque flōrēs pulcherrimōs \navnote{flōs flōris m: rosa flōs est; rosae et līlia flōrēs sunt}ad eam mittere solēbat. Initiō Aemilia epistulās Iūliī nōn legēbat nec dōna eius accipiēbat, sed omnia ad eum 80 \navnote{re-mittere}remittēbat; sed post paucōs diēs neque epistulae neque

\navnote{annō (abl) post --- post annum marītus et uxor}Annō post Iūlius et Aemilia coniugēs erant sub eōdem tēctō habitantēs. Iūlius uxōrem suam amābat et ab eā amābātur.

Etiam nunc, decem annīs post, beātīsunt coniugēs. Iūlius Aemiliam amat et ab eā amātur, neque amor coniugum hodiē minor est quam tunc. z \contr{} m Zu AU

Iūlius uxōrem ōsculātur et ``Ō Aemilia'' inquit, ``mea IJ \navnote{ōsculārī (eam) = ōsculum dare (ei)}optima uxor! Decem annī longum est tempus, sed amor meus tempore nōn minuitur. Ut tunc tē amābam, ita \navnote{minuere = minōrem facere}etiam nunc tē amō.``

Tum Aemilia, quae verbīs Iūliī dēlectātur, ``Ō Iūli? \navnote{Iūlius, meus, voc Iūlī! mi!: O Iūli, mi vir!`` ergā prp -acc}inquit, mi optimē vir! Meus amor ergā tē multō māior est hodiē quam tunc! Tempus amōrem meum nōn minuit, immō vērō auget!``

\navnote{augēre = māiōrem facere; --- \contr{} minuere}Ad hoc Iūlius rīdēns ``Ita est ut dīcis'' inquit, ``nam \navnote{miser, comp miserior, sup miserrimus}tunc ego tē amābam, tū mē nōn amābās! Ego miserri mus eram omnium adulēscentium, quod tū numquam mē salūtābās, cum mē vidēbās, quamquam ego semper tē salūtābam, cum tē vidēbam. Neque epistulās, quās cotīdiē tibi scrībēbam, legēbās, neque flōrēs, quōs tibi multōs mittēbam, accipiēbās, sed omnēs ad mē remittēbās! Propter amōrem nocte vix dormiēbam --- semper

Aemilia item dē illō tempore cōgitāns ``Ego quoque'' inquit ``tunc miserrima eram. Amābam enim alium virum Rōmānum, qul mē nōn amābat, quod virgō pauper eram.'' no

Iūlius: ``Ille vir pessimus tē dignus nōn erat!''

\navnote{tē/amōre tuō d.)}Aemilia: ``Rēctē dīcis, mī Iūlī. Tu sōlus amōre meō dignus erās, tū enim mē amābās, mihi epistulās scribēbās ac flōrēs mittēbās, quamquam pecūniam nūllam habēbam. Certē nōn propter pecūniam mē amābās, Iūlī! At nunc tē sōlum nec ūllum alium virum amō. Fēmina beāta sum, quod marītum bonum habeō, et quia nōs trēs līberōs habēmus et cum magnā familiā in hāc vīllā magnifica habitāmus, quae familiā nostrā digna est. Hoc peristylum magnificum, hae columnae, haec signa, hī flōrēs --- haec omnia cotīdiē mē dēlectant! Ō, quam beātī hīc sumus, Iūlī! Quantus est deōrum ergā nōs amor!''

Iūlius: ``Est ut dīcis, Aemilia. Decem annīs ante apud \navnote{x annīs (abl) ante = ante x annōs}parentēs nostrōs habitābāmus, ipsī parentēs nōndum erāmus nec familiam habēbāmus.

Tu et parentēs tuī 125 pauperēs erātis, in exiguā domō habitābātis nec ūllam

\navnote{pauper, corp -erior, sup -errimus dīves, comp divitior, sup dīvitissimus}Aemilia: ``Quārē ego, virgō pauperrima, ā tē, adulēscente dīvitissimō, amābar?''

Iūlius iterum Aemiliam ōsculāns ``Tu ā mē amābāris'' inquit, ``quod pulcherrima erās omnium virginum Rōmānārum, prope tam pulchra quam ipsa Venus!''

Aemilia: ``Num hodiē minus pulchra sum quam tunc \navnote{chrior) = nōn tam pulcher}eram?``

Iūlius faciem uxōris intuēns ``Certē'' inquit ``mātrōna tam pulchra es quam virgō erās, mea Aemilia. Omnēs pulchritudinem tuam laudant.'' Tum vērō fōrmam eius spectāns: ``At minus gracilis es quam tunc: eō enim \navnote{gracilis -e = tenuis}tempore gracilior erās quam hoc signum Veneris.``

Aemilia signum Veneris aspicit, cuius corpus gracilius ac minus est quam ipsīus. ``Certē tam gracilis hodiē nōn sum'' inquit, ``sed quārē mē crassiōrem fieri putās?''

Iūlius ridēns respondet: Quia nunc cibum meliōrem ēs quam tunc edēbās!``

Aemilia: ``Id quod nunc edō nec melius nec pēius est quam quod apud parentēs meōs edēbam.''

Iūlius: ``Ergō plūs ēs quam tunc, Aemilia.''

\navnote{multum, comp plūs plāris, sup plūrimum -ī paulum, comp minus -ōris, sup minimum -ī opus est = necesse est /oportet}Postrēmō Aemilia ``Certē plūs edō quam solēbam'' inquit, ``sed nec plūs nec minus quam opus est. Nimis stultus vir es, Iūli! Nōnne intellegis nōn modo amōrem nostrum, sed etiam familiam tempore augēri? Num opus est mē plūs dīcere?''

Amor dōnum est Veneris. Liberi dōna Iūnōnis sunt.

\booksection{Grammatica Latīna}

Verbī tempora Tempus praesēns et praeteritum

\navnote{praesēns -entis (praes)}Tempus praesēns

(nunc): Māne est. Sōl lūcet. Omnēs vigilant, nēmō dormit. Avēs canunt.

Tempus praeteritum

(tunc): Ante trēs hōrās nox erat. Stēllae lūcēbant \contr{} sōl nōn lūcēbat. Nēmō vigilābat, omnēs dormiēbant. Avēs nōn canebant.

Praeteritum (pers. in): singulāris -bat -ebat, plūrālis -bant -ēbant. Praeteritum [A] Āctīvum.

\navnote{bam ba computā computā bā mus computā bā tis bam bās pārc bat pārē pārē bāmus pārē bātis pārēbant}Exempla: [1] computāre: computābat; [2] pārēre: pārēbat; [3] scribere: scribēbat; [4] dormīre: dormi6ar.

Ante vīgintī annōs Iūlius discipulus tam improbus erat quam Mārcus nunc est: in lūdō dormiēbat, male computābat, foedē scribēbat, neque magistrō suō parebat.

Mārcus: ``Num tū melior discipulus erās, pater, quam ego sum? Nōnne tū in lūdō dormiēbas? Num semper magistrō pārēbas, bene computābas et pulchrē scribebas?'' Iūlius: ``Ego numquam in lūdō dormiēbam, semper magistrō parēbam, \navnote{scrīb ēbā s scrīb ēba t scrīb ēbā mus scrib ēbā tis scrīb ēba nt dormi ēbā s dormi ēba t dormi ēbā mus dormi ebā tis dormi eba nt esse era m erā mus erā tis erani -bam -bāmus}bene computābam et pulchrē scrībēbam! Certē ego discipulus multō melior eram quam tū nunc es!``

Quīntus: ``Quid? Num vōs meliōrēs discipulī erātis quam nōs sumus? Nōnne vōs in lūdō dormiēbātis? Num semper magistrō parebatis, bene computābatis et pulchrē scrībēbātis?'' Iūlius: ``Nōs numquam in lūdō dormiēbāmus, semper magistrō pārebamus, bene computābamus et pulchrē scrībēbāmus! Certē nōs discipulī multō meliōrēs erāmus quam vōs nunc estis!''

Iūlius vērum nōn dīcit: ante vīgintī annōs discipulī tam malī erant quam nunc sunt: in lūdō dormiēbant,

male computabant, foedē scrībēbant, neque magistrō suō pārēbant! Persōna prīma

-bam.

-bāmus

-ēbam.

-ēbāmus Persona secunda

-bās

-bātis Persōna tertia

-bat

-bant \navnote{bāris laudā laudā bamur laudā laudā laudā ba ntur timē bā ris lur timē bā timē limē ba ntur [3] reprehend ēba r reprehend ēbā ris reprehend ebā tur ēbā mur reprehend reprehend ēbā minī reprehend ēbantur ēbar pūni ēbā ris ēbā tur pūni pūni ēbā mur pūni ēbā minī pūni ēba niur -bāris -ebālmur -ēbā ris -ēbāminī -ēbā tur -ebantur}[B] Passīvum.

Iūlius magistrum suum nōn timēbat, sed ipse ā magistrō timēbātur! Itaque nōn reprehend \contr{} 6ātor nec pūniēbātur, sed laudābātur ā magistrō.

Mārcus: ``Num tu semper ā magistrō tuō laudābāris? Nōnne reprehendēbāris ac pūniēbāris?'' Iūlius: ``Ego ā magistrō meō saepe laudāftar, numquam reprehendar aut pūniēbar. Ego magistrum nōn timēbam, sed ipse ab eō limēbar!'' Mārcus: ``Tune ā magistrō limēbāris?'' Iūlius: ``Nōs omnēs ā magistrō ūmēbāmur. Itaque nōn reprehendēbaāmur nec pūniēbāmur, sed semper laudābāmur!'' Mārcus: Nōn vērum dīcis! 200 Vōs nōn timebamini ā magistrō nec semper laudāfcāmmt, sed reprehendēfcāmmf ac pūniebāminī!``

In illō lūdō discipulī ā magistrō timēbantur! Itaque nōn reprehendēbantur

nec pūniēbantur, sed semper laudābantur.

Persōna prīma

-bar

-bāmur

-ēbar ---

-ēbāmur Persōna secunda

-bāris

-bāminl

-ēbāris

-ēbāminī Persōna tertia

-bātur

-bantur

-ēbātur -ēbantur

\booksection{Pēnsum A}

Ante x annōs Iūlius apud parentēs suōs habit- nec uxōrem i hab-. Iūlius et Aemilia Rōmae habit-. Iūlius Aemiliam amcolumna nec ab eā am-. Ea numquam Iūlium salūt- cum eum vid-, quamquam ipsa ab eō salūt-. Aemilia epistulās Iūliī nōn leg \contr{} nec dōna eius accipi-, sed omnia remitt-.

Iūlius: ``Tunc ego tē am- nec ā tē am-, nam tū alium virum: am- nec ab eō am-. Ego miser er- et tū misera er-. Tunc apud parentēs nostrōs habit-. Tu et parentēs tuī in parvā '' domō habit-. Vōs ā mē salūt-, quamquam pauperēs er--`` 1

\booksection{Pēnsum B}

Iūlius, quī --- Aemiliae est, --- suam nec --- aliam feminam! amat. Qulntus --- quam Mārcus et --- quam Iūlia est.

--- --- peristylī Xv1 altīs --- sustinētur. Inter columnās III --- stant. Iūnō, --- [= uxor] Iovis, --- mātrōnārum est. Nēmō 1 deōrum --- maritus est quam Iūppiter. Inter omnēs deōs Iūp ipiter --- marītus est. 1

Ante x annōs Aemilia nōn mātrōna, sed --- erat. Tunc Iūlius --- xxii annōrum erat. Pater Iūliī vir --- [= pecūniōsus] \contr{} erat. Quī nōn dīves, sed --- est, in parvā --- habitat. Rōmae! multī hominēs habitant, --- quam in ūllā aliā urbe. Rōma est! urbs --- imperiī Rōmānī. [

\booksection{Pēnsum C}

Quae est coniūnx Iovis? Num Iuppiter bonus maritus est? Cuius fīlius est Cupīdō? Estne Quīntus māior quam Mārcus? Ubi parentēs Iūliī habitābant? Cūr Iūlius tunc miser erat? Cūr Aemilia Iūlium nōn amābat? Num Iūlius Aemiliam propter pecūniam amābat? Cūr Aemilia nunc beāta est? Hodiēne Aemilia minus pulchra est quam tunc?

\booksection{Vocābula nova}
\noindent\small ēsse, edō, edimus, ēs, ēstis, ēst, edunt, tempus praeteritum, -bat, -bant, -ēbant, -ēbam -ibāmus, -ēbās, -ēbātis, -ēbat, uxor, maritus, signum, tēctum, dea, coniunx, mātrōna, amor, pulchritūdō, adulēscēns, virgō, domus, templum, forum, dōnum, flōs, melior, pēior, optimus, pessimus, māior, minor, māximus, minimus, magnificus, plūrēs, plūriml, dīves, pauper, miser, bcātus, dignus, gracilis, convenlre, possidēre, mittere, remittere, ōsculārī, minuere, augēre, mī, ūllus, tamen, cotīdiē.

\bookchapter{Cap. XX}{Parentēs}

\navnote{vīcēsimus -a-um}\navnote{cūnae parvulus -a -um- parvus fari = loquī In-fans -antis m/f}Puerī parvulī quī nōndum fari possunt īnfantēs dīcun- / tur. Parvulus īnfans in cūnīs cubāre solet. Cūnae sunt lectulus īnfantis. īnfans multās hōrās dormit nōn sōlum nocte, sed etiam diē; nam longus somnus īnfantī tam necessārius est quam cibus. īnfans neque somnō neque \navnote{\contr{} necesse carēre + abl: cibō c. = sine cibō esse postulāre = poscere vāgīre = 'vā' facere, plōrāre (ut īnfans) lac lactis n}cibō carēre potest. Quōmodo īnfans, quī fāri nōn potest, cibum postulat? Infans quī cibō caret magnā vōce vāgit. Ita parvulus infāns cibum postulat. Tum māter accurrit atque īnfantem ad pectus suum appōnit. īnfans lac mātris bibit. Parvulus īnfans, cui dentēs nūllī sunt, nōn 10 pāne, sed lacte vīvit.

\navnote{alere (īnfantem) = cibum dare (īnfantī) sī-ve = vel sī mulier -eris f --- femina aliēnus -a -um = nōn suus, alterius nūtrix -īcis f}Sī māter Infantem suum ipsa alere nōn potest sīve nōn vult, Infans ab aliā muliere alitur, quae eī in locō mātris est. Mulier quae aliēnum īnfantem alit nūtrix vocātur. Multī īnfantēs Rōmānī nōn ā mātribus

suīs, 15 sed ā nūtrīcibus aluntur. Multae mātrēs īnfantēs suōs ipsae alere nōlunt.

Ante quīnque annōs Iulia parvula īnfans erat. Tunc īnfantem in cūnīs habēbat Aemilia. Nunc ea parvulō Infante caret: cūnae vacuae sunt.

Sed post paucōs mēnsēs novus īnfans in cūnīs erit. Aemilia rūrsus parvulum Infantem habēbit neque cūnae vacuae erunt. Tum Iūlius et Aemilia quattuor līberōs habēbunt. Aemilia laeta cūnās movēbit et parvā vōce cantābit: ``Lalla''. Pater Infantem suum in manibus portābit eumque nōn minus amābit quam māter. Pater et māter īnfantem suum aequē amābunt.

Annō post pater et māter ab īnfante suō appellābuntur. Aemilia autem 'mamma' appellābitur, nōn 'māter', neque Iūlius 'pater', sed 'tata' appellābitur; neque enim īnfans ipsa nōmina 'patris' et 'mātris' dīcere poterit. Infans igitur parentibus suīs dīcet: ``Mamma! Tata!'' Neque sōlum prīma verba, sed etiam primōs gradūs faciet īnfans. Initiō pater eum sustinēbit ac manū dūcet, \navnote{mox = brevī tempore post}mox vērō īnfans sōlus ambulāre incipiet neque ā parentibus sustinēbitur neque manū dūcētur. Infans ambulāns ā parentibus laetīs laudābitur. Simul īnfans plūra verba discet et mox rēctē loquī sciet.

Utra īnfantem Aemiliae alet, māterne an nūtrīx? īnfans ā mātre alētur. Aemilia ipsa īnfantem vāgientem ad \navnote{Gc M SE NNNNGZA, [Sy MY primi gradūs īnfantis vesperi (adv) māne}pectus suum appōnet. īnfans lac mātris bibet, nōn nūtrīcis. Tum māter īnfantem in cunās impōnet. Vesperī duae ancillae cūnās in cubiculum parentum portābunt atque ante lectum eōrum ponent. Si infans bene dormiet nec vāgiet, parentēs quoque bene dormient neque ab īnfante vāgiente ē somnō excitābuntur. ---

\navnote{quī) = inter sē loquī}Iūlius adhūc in peristylō cum uxōre colloquitur. Ae-milia fessa in sellā cōnsīdit. Maritus et uxor iam nōn dē \navnote{futūrus -a -um: (tempus) futūrum \contr{} praeteritum sermō -ōnis m = id quod dīcitur, verba}tempore praeteritō colloquuntur, sed dē tempore futūrō. Sermō eōrum est dē rēbus futūris.

Iūlius, quī iam intellegit Aemiliam novum īnfantem exspectāre, ``Ō Aemilia!'' inquit, ``Mox parvulum filium habēbimus.''

Aemilia: ``Fīlium? Iam duōs filiōs habēmus. Ego alteram filiam habēre volō, plūrēs quam duōs fīliōs nōlō! 55 \navnote{volō filiola -ae f --- parva filia}Cūr tū filium habēre vīs, Iūlī? Nōnne laetus eris, sī filiolam habēbis? Num parvulam filiam minus amābis quam filium?``

Iūlius: ``Profectō laetus erō, si alteram filiam habēbō. \navnote{profectō = certē nēmō, acc nēminem}Nēminem magis amābō quam parvulam fīliam.``

Aemilia: uIam filiōs tuōs magis amās quam tuam Iūliam filiolam: Mārcum et Quīntum saepe laudās, sed Iūliam rārō laudās, quamquam proba est puella. Vōs virī \navnote{velle: volō volumus vīs vultis vult volunt adversus prp--acc manēre (= nōn abīre)}filiōs modo habēre vultis, filiās nōn amāātis!`` Aemilia surgit atque gradum adversus ōstium facit.

Iūlius: ``Manē hīc apud mē, Aemilia!''

Aemilia alterum gradum ad ōstium versus facit, tum \navnote{\contr{} f ōstium li ad ōstium versus = adversus ōstium nōlī abīre! = manē!}incerta cōnsistit.

Iūlius: ``Nōlī abīre! Tēcum colloquī volō.''

Aemilia nōn abit, sed apud maritum manet. Coniu- 70 \navnote{pergere (= nōn dēsinere) colloquium -i 2 \contr{} colloquī}gēs colloquī pergunt. Ecce colloquium eōrum:

Iūlius: ``Nōs viri nōn filios tantum, sed etiam filiās habēre volumus, nec filiās minus amāmus quam fīliōs.

\navnote{rārō adv, comp rārius, sup rārissimē}Aemilia: ``Multae mātrēs īnfantēs suōs apud nūtricēs relinquunt,

ego vērō manēbō apud īnfantem meum: numquam ab eō discēdam! Si aeger erit, ipsa eum cūrābō tōtamque noctem apud eum vigilābo; nēmō īnfantem aegrum tam bene cūrāre potest quam māter ipsa.`` s0

Iūlius: ``Nōnne ab īnfante sānō discēdēs?''

\navnote{minimē = certē nōn (discēdam) velsī dēbēre. officium -I n: mātris officium est = māter dēbet}Aemilia: ``Minimē! Bona māter semper apud īnfantem suum manēre dēbet. Sīve infāns valet sīve aegrōtat, māter ipsa eum cūrāre et alere dēbet --- hoc est mātris officium!'' 8\$

Iūlius: ``Tune ipsa īnfantem tuum lacte tuō alēs?''

Aemilia: ``Profectō īnfantem meum ipsa alam. Ego faciam officium meum! Neque sōlum diē, sed etiam nocte apud īnfantem erō: semper cum eō dormiam.''

\navnote{pot-erimus}Iūlius: ``Quid? Nōs et infans in eōdem cubiculō dormiēmus? Vix ūnam hōram dormīre poterimus, si infāns vāgiet. S1 tū et infans tuus in cubiculō nostrō dormiētis, ego profectō in aliō cubiculō dormiam, ubi ab īnfante vāgiente nōn excitābor!''

Aemilia: ``Ō

Iūlī! Ita loquitur homō quī officium suum nescit!``

Iūlius: ``Meum officium est pecūniam facere ac mag-nam familiam alere, nōn cum Infante vāgiente cubāre! Somnus virō industriō necessārius est!''

\navnote{pergere:: loquī pergere plūra:: plūra verba}Aemilia īrāta ``Nōlī pergere!'' inquit, ``Plūra ā tē audīre nōlō!'' atque iterum ad ōstium versus īre incipit.

``Manē, Aemilia!'' inquit Iūlius, ``Nōlī ita mē relinquere!'' sed illa adversus ōstium īre pergit.

Tum vērō Syra, quae eō ipsō tempore peristylum in- 1H trat ūnā cum Iūliā, dominae in ōstiō occurrit.

\navnote{ūnā cum = simul cum}Aemilia ante Syram et Iūliam cōnsistēns ``Quid vul- 105 tis?'' inquit, ``Cūr nōn manētis in hortō?''

Syra: ``Quia mox imbrem habēbimus: ecce caelum nūbibus ātris operitur. Si in hortō manēbimus, Qmidae erimus. Vōs quoque ūmidī eritis, dominī, si hic in peristylō manēbitis.''

Iūlius caelum spectāns ``Rēctē dīcis'' inquit, ``Illae nūbēs imbrem afferent. Venīte mēcum in ātrium! Mox sōl rūrsus lūcēbit.''

\navnote{Syra Aemiliae in ōstiō occurrit}Iūlius ātrium intrat; Aemilia eum sequitur ūnā cum \navnote{silēre = tacēre}Iūliā et Syrā. In ātriō Iūlius et Aemilia silentēs imbrem \navnote{\contr{} sermō}in impluvium cadentem aspiciunt. Iūlia silentium parentum animadvertit, et ``Quid silētis?'' inquit, ``Estisne tristēs? Ego vōs cōnsōlābor!''

ficium est nāvigāre, sīve mare tranquillum

sīve turbidum est. Multī nautae nunc in mare merguntur, dum \navnote{ad-vehere domum = ad domum suam re-vertī = red-Īre}rēs necessāriās ex terris alienis in Italiam advehere cōnantur. Ō, miserōs nautās, quī numquam domum revertentur! Ō, miserōs līberōs nautārum, quī post hanc tempestātem patrēs suōs nōn vidēbunt!``

Iūlia: ``Ego laetor quod pater meus nauta nōn est et \navnote{domī = in domō suā crās = diē post hunc diem}domī apud nōs manēre potest.``

Iūlius, quī crās Rōmam ībit, ``Nōn semper'' inquit ``mihi licet apud vōs manēre, Iūlia. Necesse est mihi crās rūrsus ā vōbīs discēdere, nec vērō in terrās aliēnās ībō, ut nauta.''

Iūlia: ``Quō ībis, tata? Quandō revertēris? Ego et iriamma tē sequēmur!''

Iūlius: ``Quārē mē sequēminl? Rōmam proficīscar, unde tertiō quōque diē revertar, s1 potero.''

\navnote{pot-erō nōs, abl nōbīs: ā nōbīs domō = ā domō suā}Iūlia: Nōlī ā nōbīs discēdere! Vel sī necesse erit domō abīre, nōn modo tertiō quōque diē, sed cotīdiē ad nōs revertī dēbēs. Hoc postulō ā tē! Nōlō tē carēre, tata. Cotīdiē tibi occurram.``

Iulius: ``Audīsne, Aemilia? Iūlia dīcit 'sē patre suō \navnote{= nōn velle dīligere = amāre (ut parentēs līberōs suōs) aequē (ac) = nec magis nec minus (quam)}carēre nōlle', ergō mē nōn minus dīligit quam tē. Atque ego profectō filiam meam aequē dīligō ac fīliōs. Nec alteram filiam minus dīligam.``

Iūlia: ``Quam 'alteram fīliam' dīcis? Mihi soror nōn

\navnote{Iūliō silente = dum Iūlius silet Iūliola = parva Iūlia}Iūliō silente, Aemilia ``Nōnne gaudēbis, Iūliola'' in-quit, si parvulam sorōrem habēbis?``

Iūlia: ``Sorōrem habēre nōlō! Nam si sorōrem habēbō, ea sōla ā vōbīs amābitur, ego nōn amābor!'' 150

Aemilia: ``Certē tū nōn minus ā nōbīs amāberis: tū et parvula soror aequē amābiminī.''

Iūlia: Si aequē amābimur, laeta erō. Sed multō magis laetābor, si frātrem habēbō, mamma! Nōnne tū quōque laetāberis, tata, si filiolum habēbis?``

\navnote{filio/us -īm = parvus filius}Iulius silet. Aemilia vērō, antequam Iūlia silentium 155 patris animadvertit, Nōlī dīcere tatam' et mammam', Iūliola!`` inquit, ''Ea nōmina ā tē audire nōlumus. Ita \navnote{nōlumus ( \contr{} ne-volumus) = nōn volumus decēre = dignus esse; eam decet = eā dignus est, ei convenit}loquuntur parvulī īnfantēs, nec sermō īnfantium tē decēt. Tatrem' et mātrem' dīcere oportet.``

Iulia: ``Sī īnfans nōn sum, nōlīte mē 'Iūliolam' cāre! Id nōmen mē nōn decet. Mihi nōmen est 'Iūlia'.''

Aemilia: ``Rēctē dīcis, Iūlia. Tu igitur ā nōbīs 'Iūlia' vocāberis, et nōs ā tē 'pater' et 'māter' vocābimur.''

Iūlia: ``Ita semper ā mē vocābiminī, tata et mamma!''

\booksection{Grammatica Latīna}

Verbī tempora Tempus futūrum

Tempus praesēns: Diēs est, Sōl lūcet. Omnēs vigilant. Nēmō dormit. Avēs canunt.

Tempus futūrum: Mox nox erit. Sōl nōn lūcēbit, sed stēllae 170 lūcēbunt. Omnēs dormienti. Nēmō vigilabit. Nūlla avis canet.

Futūrum (pers. Ill): singulāris -bit -et, plūrālis -bunt -ent. [A] Activum.

Exempla: [1] computāre: computā\&t'r; [2] pārēre: pārēbit; [3] scriblere: scribet; [4] dormīre: dormi.

Malus discipulus: ``Ab hōc diē bonus discipulus erō, magister: numquam in lūdō dormiam, semper tibi pārēbō, bene computābō et pulchrē scrībam!'' Magister: ``Quid? Tune bo-nus discipulus eris? Id fieri nōn potest! Crās rūrsus in lūdō dormi \contr{} s, male computābis, foedē scribes, nec mihi pārē\&ts/'' 180 Crās discipulus tam malus erit quam hodiē est: in lūdō dormier, male computābit, foedē scribet, nec magistrō pārebit.

Malī discipulī: ``Ab hōc diē bonī discipulī erimus, magister: numquam in lūdō dormiēmus, semper tibi pārēbintns, bene compmabimus et pulchrē scribemus/'' Magister: ``Quid? Vōsne bonī discipulī eritis? Id fieri nōn potest! Crās rūrsus in lūdō dormiētis, male computāābitis, foedē scribēm, nec mihi pārēbitis/'' Crās discipulī tam improbī erunt quam hodiē sunt: in lūdō dormient, male computābunt, foedē scribent, nec magistrō parēbunt.

Sing.

Plūr. Persōna prīma

-bō

-bimus

-am

-ēmus Persōna secunda

-bis ---

-bitis Persōna tertia

-bit ---

-bunt [B] Passīvum.

Fīlius: ``Ab hōc diē bonus discipulus erō. Ergō ā magistrō laudāfor, nōn reprehendar.'' Pater: Tū nōn laudāberis, sed reprehendēro ā magistrō!`` Crās discipulus rūrsus ā magistrō reprehendetur, nōn laudābitur.

Filii: ``Ab hōc diē bonī discipulī erimus. Ergō ā magistrō laudāfttmir, nōn reprehendēmur.'' Pater: ``Vōs nōn laudabiminī, sed reprehendemini ā magistrō!'' Crās discipulī rūrsus ā magistrō reprehenderztar, nōn laudabuntur.

Plūr. Persōna prima ---

-bor

-bimur

-ar

-ēmur Persona secunda

-beris

-biminī

-eris

--- -ēminī Persōna tertia

-bitur

-buntur

-ētur

-entur

\booksection{Pēnsum A}

Mox novus infāns in cūnīs Aemiliae er-. Aemilia rūrsus īnfantem hab-. Tum quattuor līberi in familiā er-. Iūlius et Aemilia quattuor līberōs hab-. Aemilia Infantem suum am-. Iūlius et Aemilia īnfantem suum aequē am-. Annō post īnfans prima verba disc- et primōs gradūs faci-. īnfans ambulāns ā parentibus laud-.

Aemilia: ``Ego īnfantem meum bene cūr-: semper apud eum man-, numquam ab eō discēd-.'' Iūlius: Certē bona māter er-, Aemilia: īnfantem tuum ipsa cūr- nec eum apud nūtricem relinqu-.`` Aemilia: ''Etiam nocte apud īnfantem er, semper cum eō dorm-. Nōs et Infans in eōdem cubiculō dorm-.`` Iūlius: ''Nōn dormiēmus, sed vigil-! Nam certē ab īnfante vāgiente excit-!`` Aemilia: ''Ego excit-, tū bene dormi- nec excit-!``

\booksection{Pēnsum B}

Parvulus puer quī in --- iacet --- appellātur. Infans quī cibō --- magnā vōce ---. Nōn pānis, sed --- cibus īnfantium est. Nūuix est --- [= femina] quae nōn suum, sed --- īnfantem ---. Multāe mātrēs īnfantēs suōs ipsae alere ---.

Marītus et uxor iam nōn dē tempore praeteritō ---, sed dē tempore ---. Aemilia: Cūr tū fīlium habēre ---, Iūlī? Ego alteram filiam habēre ---, plūrēs quam duōs filiōs --- [= nōn volō]. Vōs viri filiōs tantum ---! Fīliōs --- dīligitis quam fīliās!`` Iūlius --- [= tacet]. Aemilia īrāta --- adversus ōstium facit, sed Iūlius u--- hlc apud mē!'' inquit, ``--- discēdere!'' Aemilia manet ac loquī---: ---: ``Bona māter apud Infantem suum manēre ---, hoc mātris --- est. Nēmō enim īnfantem melius --- potest quam māter ipsa.''

\booksection{Pēnsum C}

Ubi parvulus īnfans cubat? Quid facit Infans qul cibō caret? Quid est cibus īnfantium parvulōrum? Num omnēs īnfantēs ā mātribus suīs aluntur? Quae sunt prīma verba īnfantis? Quandō Aemilia novum īnfantem habēbit? Cūr Aemilia ā marītō suō discēdere vult? Cūr Iūlius in aliō cubiculō dormīre vult? Quid Aemilia officium mātris esse dīcit? Cūr Syra et Iūlia in hortō nōn manent? Cūr Iūlia sorōrem habēre nōn vult?

\booksection{Vocābula nova}
\noindent\small vōs, abl vōbīs: ā vōbīs, īre: ībō, ībimus, Ibis, ībitis, ībit, ībunt, bis, computā, b't, computā b imus, computā b itis, computā b unt, bō, pārē, bit, blimus, blitis, bunt, pārē i scriblet scribZmus scribetis, scrībēs, scrīber, scrīb ēlmus, scrībītis, scnbelnt, [4] domijr dormi a m, dormi5, ]ētis, dormilent, esse erō, erimus, eris, eritis, erit, -bimus, o, is, Hīs, unt, it, m, -ētis, nt, beris, laudā, laudā b itur, ar, reprehend, reprehend ētur, or, tmur, erts, uunī, itur, r, mur, ris, Ttiinī, tur, ntur, nova, īnfans, cūnae, somnus, lac, mulier, nūtrix, gradus, sermō, filiola, filiolus, colloquium, officium, silentium, parvulus, necessārius, aliēnus, futūrus, ūmidus, fari, carēre, postulāre, vāgīre, alere, colloquī, volō, vīs, volumus, vultis, manēre, pergere, cūrāre, dēbēre, occurrere, silēre, advehere, revertī, dīligere, decēre, nōlle, domō, mox, magis, rārō, crās, adversus, nōlī nōlīte, sīve... sīve, profectō, minimē, ūnā cum.

\bookchapter{Cap. XXI}{Pugna discipulōrum}

I Ecce puer quī ē lūdō domum revertitur. Quis est hic puer? Mārcus est, sed difficile est eum cognōscere, nam sordidus est et sanguis dē nāsō eius fluit. Hodiē māne vestīmenta Mārcī munda erant atque tam candida quam \navnote{mundus -a -um = pūrus candidus -a -um = albus (ut nix) vestis -is f --- vestīmenta Ne e d NS ``D 9l}nix nova, nunc vērō sordida et ūmida sunt. Cūr vestis Mārcī ūmida est? Vestis ūmida est, quod Mārcus per imbrem ambulāvit. Nec modo Mārcus, sed etiam Titus et Sextus per imbrem ambulāvērunt. Omnēs discipulī vestimentis ūmidīs domum revertuntur. 10

Sed cūr sanguis dē nāsō fluit Mārcō? Sanguis eī dē nāsō fluit, quod Mārcus ā Sextō pulsātus est. Nōnne \navnote{Mārcus ā Sextō pulsātus est (pass) = Sextus Mārcum pulsāvit (āci)}Mārcus Sextum pulsāvit? Prīmum Sextus ōs Mārcī pugnō pulsāvit, deinde Mārcus et Titus Sextum pulsāvērunt. Sextus, quī māior est quam cēterī discipulī, cum Mārcō et Titō pugnāvit et ab iīs pulsātus est. \navnote{pugnāre cum = pugnāre contrā}Puerī pugnāvērunt in viā angustā quae Tusculō ad vīllam Iūliī fert.

Nōn modo vestis, sed etiam faciēs et manus et genua \navnote{genū -ūs n, p! genua -uum}Maārci sordida sunt. Cūr tam sordidus est puer? Sordidus est quod humī iacuit; humus enim propter imbrem 20 ūmida et sordida est. Et Mārcus et Sextus humī iacuērunt. Prīmum Mārcus iacuit sub Sextō. Titus vērō Mārcum vocāre audīvit ac Sextum oppugnāvit. Mox Sextus ipse humī iacēns ā duōbus puerīs pulsātus est; magnā vōce patrem et mātrem vocāvit, nec vērō parentēs eum 25 audīvērunt: vōx Sextī ā nūllō praeter puerōs audīta est.

Mārcus ātrium intrāns nōn statim ā patre suō cognōscitur, sed cum prīmum filius patrem salūtāvit, Iūlius \navnote{cruor -ōris m - sanguis cōn-spicere -iō = vidēre meus filius, voc mi fili! cornua}vōcem filii cognōscit. Tum cruōrem in faciē eius cōnspiciēns exclāmat pater: UŌ mī fili! Quis tē pulsāvit?``

Mārcus: ``Bōs īrātus cornū mē pulsāvit!''

Iūlius: \contr{} Id vērum nōn est! Is quī tē pulsāvit cornua \navnote{VV}nōn gerit. Ā quō pulsātus es?``

Mārcus: aĀ Sextō pulsātus sum.``

Iūlius: ``Intellegēbam tē nōn cornibus, sed pugnīs 35 pulsātum esse. At cūr tu pulsātus es? Certē nōn sine causā Sextus tē pulsāvit. Incipe ab initiō: ille prīmum ā tē pulsātus est!''

Mārcus: ``Minimē! Prīmum ille mē pugnō pulsāvit sine causā, deinde ego illum pulsavi!''

Iulius: ``Tune sōlus Sextum pulsāvistī?''

Mārcus: Ego et Titus eum pulsāvimus.``

Iūlius: ``Quid? Vōs duo ūnum pulsāvistis? Duo cum ūnō pugnāvistis?''

Mārcus sē et Titum excūsāre cōnātur: aAt pugnāvimus cum puerō māiōre. Sōlus Sextum vincere nōn possum, nam is multō māior ac validior est quam ego. Sextus tam validus est quam bōs!``

Iūlius: ``Et tū tam sordidus es quam porcus! Cūr sordida est vestis tua nova? Humīne iacuistī?''

\navnote{humī-ne}Mārcus: ``Humī iacuit: Sextus mē tenuit. Sed is quōque humī iacuit: nōs eum tenuimus!''

\navnote{pugna -ae f \contr{} pugnāre in-dignus -a -um = nōn dignus nōn decet = indignum est}Iūlius: ``Iam satis audīvī dē pugnā vestrā indignā. Nam certē nōn decet pulsāre puerum minōrem, sed duōs puerōs cum ūnō pugnāre indignissimum est --- hoc nūllō modō excūsāri potest! Nōlī mihi plūs nārrāre dē hāc rē! Age, I in cubiculum tuum ac mūtā vestīmenta! Dāvus tibi alia vestīmenta dabit.''

Mārcus Dāvum sequitur in cubiculum, ubi cruōrem et sordēs ē faciē, manibus, genibus lavat ac vestīmenta \navnote{sordēs -ium f pl (vestem) pōnere \contr{} \contr{} induere}mūtat. Puer vestem sordidam pōnit aliamque vestem mundam et candidam induit.

Interim Aemilia ātrium intrat. Māter familiās statim \navnote{interim = dum haec aguntur}sordēs et vestīgia Mārcī in solō cōnspicit et ``Sordidum est hoc solum!'' inquit, ``Aliquis pedibus sordidis in solō \navnote{ali-quis (: nesciō quis) ali-quid (: nesciō quid)}mundō ambulāvit! Quis per ōstium intrāvit? aliquis ē familiā nostrā?``

``Porcus intrāvit!'' ait Iūlius.

Aemilia: Quid ais?``

ūlius: Porcum intrāvisse' āiō.``

Aemilia: ``Ain' tū? Porcusne ātrium intrāvit?''

Iūlius: *'Porcus' quī intrāvit est tuus Mārcus filius!`` Iūlius 'Mārcum intrāvisse' dīcit, at nōn dīcit 'eum ā Sextō pulsātum esse et humī iacuisse.'

Aemilia, quae iam intellegit puerum sordidum ā pa- 75 tre 'porcum' appellātum esse, ``Ubi est Mārcus?'' in-quit, ``Cūr nōndum mē salūtāvit?''

Iūlius: ``Mārcus lavātur et vestem mutat.''

Mārcus, postquam vestem mūtāvit, mundus redit et \navnote{quam}mātrem salūtat.

``Salvē, mī filī!'' inquit māter, ``Bonus es puer, quod statim vestem mūtāvistī. In lūdōne quoque bonus puer fuistī?''

Mārcus: ``Profectō bonus puer fuī, māter. Laudātus sum ā magistrō.'' Mārcus dīcit sē bonum puerum fuisse 85 et laudātum esse', quamquam puer improbus fuit et ā magistrō verberātus est. Quod dīcit nōn vērum, sed falsum est: Mārcus mentītur.

\navnote{mentīri = falsum dīcere}Iūlius: ``Quid magister vōs docuit hodiē?''

\navnote{multa (npl) = multae rēs cētera (n pl) = cēterae rēs nōs, dat nōbīs}Mārcus: ``Multa nōs docuit: legere et scrībere et com- 90 putāre et cētera. Primum magister nōbīs aliquid recitāvit, nesciō quid: nihil enim audīvī praeter initium!'' Mārcus dīcit 'sē nihil audīvisse praeter initium',

id quod vērum est.

\navnote{quā dē causā? = dē quā omnia (n pl) = omnēs rēs}Aemilia: ``Quā dē causā nōn omnia audivisti?''

Iūlius: ``Hahae! Magistrō recitante, Mārcus dormīvit!'' Iūlius Mārcum dormīvisse' dīcit, nec vērō māter \navnote{crēdere = vērum esse putāre}id crēdit.

Aemilia: ``Audīsne, Mārce? Pater dīcit 'tē in lūdō dormīvisse'! Nōnne falsum est quod dīcit pater? Num tū vērē in lūdō dormīvistī?''

Mārcus: ``Ita est ut dīcit pater. Nec vērō ego sōlus dormīvī: omnēs dormīvimus!''

Aemilia: ``Ain* vērō? In lūdō dormīvistis? Malī discipulī fuistis! Nōnne pūnītī estis ā magistrō?''

Mārcus: ``Certē malī discipulī fuimus ac pūnītī sumus: omnēs verberātī sumus ā magistrō. Adhūc mihi dolet tergum. Sed paulō post magister litterās meās \navnote{tabella -ae f --- parva tabula vōs, dat vōbīs}laudāvit. Tabellam vōbīs ostendam. Ecce tabella mea.`` 110

Mārcus parentibus nōn suam, sed aliēnam tabellam ostendit. Cuius est ea tabella? Sextī est: Mārcus enim tabellās eōrum mūtāvit inter pugnam! Nōn Mārcus, sed Sextus scripsit litterās quae in eā tabellā leguntur. \navnote{scrips- \contr{} scribs-}Aemilia vērō, quae tabellam Mārcī esse crēdit, ā fīliō suō improbō fallitur. Nōn difficile est mātrem Mārcī fallere!

Aemilia: ``Tune ipse hās litterās pulcherrimās scripsistī?''

Mārcus: ``Ipse scripsī profectō. Mihi crēde!''

\navnote{crēdere + dat: mihi crēde! = crēde id quod dīcō! scrift- \contr{} scribt-}Mārcus mentītur; nam id quod Mārcus 'sē ipsum scripsisse5 dīcit, ā Sextō scriptum

est. Sed Aemilia, quae litterās ā Mārcō scriptās esse crēdit, ``Aspice, Iūlī!'' inquit, ``Mārcus ipse haec scripsit et ā magistrō laudātus est. Quid dīxit magister, Mārce? Nārrā nōbīs

Mārcus iterum mentītur: ``Dīxit 'mē pulcherrimē et rēctissimē scripsisse'.'' (Nōs vērō scīmus magistrum aliud dīxisse!)

Iūlius, quī Mārcum discipulum pigerrimum esse scit, iam dē verbīs eius dubitāre incipit. Aemilia vērō nihil \navnote{dubitāre = incertus esse; dere}dubitat, sed omnia crēdit!

Aemilia: ``Cēteri discipulī nōnne rēctē scripsērant?''

Mārcus: ``Minimē! Titus et Sextus prāvē scrīpsērunt et malī discipulī fuērunt, nec ā magistrō laudātī sunt. Ego sōlus laudātus sum!''

Iūlius: ``Vōsne omnēs eadem vocābula scripsistis?*'

Mārcus: ``Omnēs eadem scripsimus, sed ego sōlus rēctē scrīpsī--- --- ut iam vōbīs dM.''

Iam Iūlius, quī Sextum discipulum industrium ac prūdentem esse scit, Mārcō nōn crēdit. Mārcus patrem suum fallere nōn potest. Ergō Iulius, quī interim tabellam in manūs sūmpsit, Menūris, Mārce!`` inquit, ''Hoc tuā manū scriptum nōn est. Falsa sunt omnia quae nōbīs nārrāvistī!``

\navnote{falsa (n pI) = rēs falsae}Aemilia vērō ``Quā dē causā'' inquit ``eum falsa dīx- 145 isse putās? Cūr nōn crēdis filio tuō?''

Sed antequam Iūlius ad haec respondēre potest, ali-quis iānuam pulsat. Quis pulsat? Vidē capitulum quod sequitur!

\booksection{Grammatica Latīna}

Verbī tempora \navnote{perfectum (perf) im-perfectum (imperf)}Praeteritum perfectum et imperfectum

Nox obscūra erat: nūlla lūna caelum illūstrābat; caelum neque lūnā neque stēllīs illūstrabatur.

Tum fulgur caelum illūstrāvit; caelum obscūrum fulgure clārissimō illūstratum est.

Illūstrāabat, illūstrabatur' est praeteritum imperfectum. Illūstrāvit, illāstrātum est praeteritum perfectum

est.

Perfectum [A] Activum.

Mārcus dīcit 'sē Sextum pulsāvisse': ``Ego Sextum pulsāvi.'' [ūlius: ``Tune sōlus eum pulsāvistī?'' Mārcus: ``Ego et Titus eum pulsavimus.'' Iūlius: ``Vōsne ūnum pulsāvistis?'' Primum Mārcum pulsāvit Sextus, tum Mārcus et Titus eum pulsāaverunt. \navnote{pulsavērunt -is5e}Infinitivus perfectī: pulsāv isse. Singulāris Persōna prīma

pulsāvt Persōna secunda --- pulsāvistī

pulsāv istī Persōna tertia

pulsāvir pulsāvit pulsāverunt

Exempla: [1] recitāre: recitōz \contr{} isse; [2] pārere: pāāruisse; \navnote{pārē-: pāru-}[3] scribere: scrīpsisse; [4] audire: audīvisse.

Maārcus malus discipulus fuit: male recitāvit, foedē scripsit, \navnote{recixāvt Ttcilāvistt -āvistis recitōir āvēmnt pān/ imus pāru istī pāru istis i pāru ērunt scrlps imus sctlps istī scrīps istis scrīpsērunt dormīv; dormīvimus dois domi dormivit dormīvērunt}in lūdō dormm

nec magistrō pāmtr. Sed Mārcus dīcit 'sē bonum discipulum fuisse, bene recixāvisse, pulchrē scrīpsisse, magistrō pāruisse, nec in lūdō dormīvisse': ``Ego bonus discipulus fuī, bene recitāvi, pulchrē scrips?, magistrō pārut, nec in lūdō dormm/M Iūlius: ''Mentīris, Mārce! Tu malus discipulus fuistī, male recitāvistī, foedē scripsisti, in lūdō dormivistī, nec magistrō pāruistī!``

Mārcus et Titus malī discipulī fuērunt, male recitāvērunt, foedē scrīpsērunt, in lūdō dormīvērunt, nec magistrō pāruerunt. Mārcus et Titus: ``Nōs bonī discipulī fummus, bene recitdvimus, pulchrē scrīpsimus, magistrō pāruimus, nec in lūdō dormioimus!/'' Iūlius: ``Mentīminī, pueri! Vōs malī discipulī 185 fuistis, male Tecitāvistis, foedē scrīpsistis, in lūdō dormīvistis, nec magistrō pāvuistis!'' [B] Passīvum.

\navnote{laudāf us est laudāf us = quī laudātus est laudā:rt sunt laudār f = quī laudātī sunt}Pater fīlium/fīliam laudāvit = fīlius laudātus est/filia laudata est ā patre. Fīlius laudātus/filia laudāta gaudet. Pater filios/filiās laudāvit = filiī laudātī sunt/filiae laūdatae sunt ā patre. Fīlii laudār filiae laudatae

gaudent. Laudārus -a -um' est participium

perfectī. Participium perfectī est adiectīvum dēclīnātiōnis i/ii. Alia participia perfectī: pulsazus, appellātus, verberātus, scrīptus, auditus, pūnizus. Cum participium perfectī cum tempore praesentī \navnote{laudātus sum laudātus es laudātus est laudātī sumus laudat? estis laudātī sunt -tum esse mdātum esse -tus -tus sumus estis est sunt est -tasmt -tumest --- -tasunt}verbī 'esse' coniungitur, fit perfectum passīvī.

Discipulus

dīcit 'sē laudātum esse': UA magistrō laudātus sum.`` Pater: ''Tune laudātus es?``

Discipulī dīcunt 'sē laudātōs esse': ``Ā magistrō laudārt su- 200 mus'' Pater: ``Vōsne laudātī estis?'' īnfinītīvus: laudātum esse Singulāris

prima laudātus sum laudātī sumus Persōna ptītna Persōna secunda

laudātus es Persōna tertia

laudātus est laudātī sunt

\booksection{Pēnsum A}

Mārcus ūmidus est, quod per imbrem ambul-, ac sordidus, quod humī iac-. Discipulī ūmidl sunt, quod per imbrem ambul-, ac sordidī, quod humī iac-.

Mārcus: ``Ā Sextō puls- sum.'' Iūlius: ``Nōnne tū Sextum puls-?'' Mārcus: ``Prīmum Sextus mē puls-, tum ego illum puls-. Ego et Titus Sextum puls-.'' Iūlius: ``Quid? Vōs duo ūnum puls-? Et cūr sordida est vestis tua? Humīne iac-?'' Mārcus: Humī iac-: Sextus mē ten-.``

Aemilia: ``Quid magister vōs doc- hodiē?'' Mārcus: Magister aliquid recit-, sed ego initium tantum aud-.`` Aemilia: ''Cūr nōn omnia aud-?`` Iūlius: ''Mārcus in lūdō dorm-!`` Aemilia: ''Audīsne? Pater dīcit 'tē in lūdō dorm-.' Tune vērē dorm-?`` Mārcus: ''Certē dorm-, māter. Omnēs dorm-! Sed paulō post magister laud- litterās meās.`` Mārcus mātri litte rās quās Sextus scrips- ostendit. Aemilia: ''Tune ipse hoc! scrīps-?`` Mārcus: ''Ipse scrīps-. Cēteri puerī prāvē scrīps-.``!

\booksection{Pēnsum B}

Hodiē māne Mārcus --- erat, vestīmenta eius tam --- erant! quam nix. Nunc nōn modo --- [= vestīmenta], sed etiam faciēs et manūs et --- sordida sunt, atque --- [= sanguis] ei dē nāsō fluit. Nam Sextus, qul tam --- est quam bōs, Mārcum ji pulsāvit sine --- (ut ait Mārcus). Mārcus, quī --- iacuit, tam sordidus est quam ---. Difficile est eum ---.

In cubiculō suō Mārcus lavātur et vestem ---. --- [= dum haec aguntur] Aemilia ātrium intrat et vestīgia pueri ---. Ae 1milia: ``Quid hoc est? --- pedibus sordidis in --- mundō ambu 'lāvit!'' Mārcus, --- vestem mūtāvit, in ātrium redit. Puer mā \contr{} trī --- [= parvam tabulam] Sextī ostendit et 'sē ipsum eās j litterās ---' dīcit. Id quod Mārcus dīcit nōn vērum, sed --- est; Mārcus ---, sed Aemilia filio suō ---. Iūlius vērō dē verbīs eius: ---; Mārcus patrem suum --- nōn potest.; i

\booksection{Pēnsum C}

Cūr Mārcus rediēns ūmidus et sordidus est?! Ā quō pulsātus est? Cūr Mārcus sōlus Sextum vincere nōn potest? Quōmodo Mārcus sē excūsāre cōnātur? Quid agit Mārcus in cubiculō suō? Quid Aemilia in solō cōnspicit? Quid tum Iūlius dīcit uxōri? Quid Mārcus parentibus suīs ostendit? Quis litterās quae in eā tabulā sunt scripsit? Cūr Iūlius filium suum nōn laudat?

\booksection{Vocābula nova}
\noindent\small bōs bovis m/f, cornū -ūs n (dēcl IV): sing, plūr, sing, cornā, corntia, nōtn, cornua Ben cornūs cornuum dat corni cornibus abl, cornū, acc, cornās, cornuum, gen, dat, cornibus, cornu, abl, ait = dīcit, ais, ait, āiunt, lūdō-ne, sum, fuī, es, fuistī, est, fuit, esse, fuisse, pulsāvisse, pulsavi, pulsāāvimus, pulsāvistī, pulsāvistis, pulsāvit, i, -tmus, II, -istī, -istis, -ērunt, iisse, Juī, fuimus, firnstī, fiilistis, iit, fuērunt, nova, vestis, genū, humus, cruor, bōs, causa, porcus, pugna, sordēs, solum, tabella, mundus, candidus, angustus, validus, indignus, falsus, cognōscere, cōnspicere, excūsāre, vincere, nārrārc, mūtāre, mentīrī, crēdere, fallere, dubitāre, āiō ais ait, ain', aliquis, aliquid, humī, interim, postquam, perfectum, imperfectum.

\bookchapter{Cap. XXII}{Cavē canem}

Iānua vīllae ē duābus foribus cōnstat. Sub foribus est / \navnote{līmen cardō foris 4 2 cardō f) foris -is f līmen -inis n cardō -inis m vertī= --- circum movēri ] Y (}līmen, in quō salvē scrīptum est. Foris duōs cardinēs habet, in quibus vertī potest; cum foris in cardinibus vertitur, iānua aperitur aut clauditur. Servus cuius officium est forēs aperire et claudere ac vīllam dominī 5 cūstōdīre, ōstiārius vel iānitor appellātur.

\navnote{iānitor -ōris m = ōstiārius si}Sī quis vīllam intrāre vult, iānuam pulsat et extrā iānuam exspectat, dum iānitor forēs aperit eumque in \navnote{quis --- si aliquis extrā prp+acc ad-mittere: intrā prp-- acc \contr{} \contr{} extrā ferōx -ōcis = ferus \$ anteā=ante hoc tempus, ? antiquis temporibus ( catēna}vīllam admittit. Iānitor intrā iānuam sedet cum cane suō, quī prope tam ferōx est quam lupus; itaque necesse 10 est eum catēnā vincīre. Anteā dominī sevēri nōn sōlum canēs, sed etiam iānitōrēs suōs catēnis vinciebant.

Catēna quā canis vincītur ex ferrō facta est. Catēna \navnote{ferreus -a -um = ex ferrō factus}cōnstat ē multīs anulis ferreis quī inter sē coniunguntur. Ānulī quibus digitī ōrnantur nōn ex ferrō, sed ex aurō 15 facti sunt. Aurum est frūagni-pretiī sīcut gemmae. Ānulus aureus multō pulchrior est quam ānulus ferreus.

\navnote{aureus -a -um = ex aurō factus}Forēs ē lignō factae sunt sīcut tabulae. Lignum est māteria dūra, sed minus dūra quam ferrum. Quī rēs \navnote{ligneus -a -um = ex lignō factus faber -bri m claudere clausisse clausum (esse) lignum,}ferreās vel ligneās facit, faber appellātur. Deus fabrōrum est Vulcānus. ---

Iānua clausa est. Iānitor, quī forēs clausit postquam Mārcus intrāvit, iam rūrsus dormit! Iānitōre dormiente, canis vigilāns iānuam cūstōdit. Extrā forēs stat tabellārius (sīc appellātur servus qul epistulās fert, nam anteā in tabellīs scribēbantur epistulae). Is baculō ligneō forēs pulsat atque clāmat: ``Heus! Aperi hanc iānuam! Num quis hīc est? Num quis hanc aperit iānuam? Heus tū, \navnote{num quis = num aliquis quīn? = cūr nōn? quīn aperis? = cūr nōn aperis? aperī!}iānitor! Quīn aperis? Dormīsne?``

Cane lātrante iānitor ē somnō excitātur.

Tabellārius iterum forēs pulsat magnā vōce clāmāns: ``Heus, iānitor! Quīn mē admittis? Putāsne mē hostem esse? Ego nōn veniō vīllam oppugnātum sīcut hostis, \navnote{-tum: vīllam oppugnātum = quia vīllam oppugnāre volō tandem = postrēmō (post longum tempus)}nec pecūniam postulātum veniō.``

Tandem surgit iānitor. ``Quis forēs nostrās sīc pul-sat?'' inquit.

Tabellārius (extrā iānuam): ``Ego pulsō.''

Iānitor (intrā iānuam): ``Quis 'ego'? Quid est tibi nōmen? Unde venīs? Quid vis aut quem quaeris?'' 40

Tabellārius:

``Multa

simul

rogitās.

Admitte

mē! \navnote{rogitāre = interrogāre (multa) post-eā = deinde prius adv \contr{} \contr{} posteā}Posteā respondēbō ad omnia.``

Iānitor: ``Respondē prius! Posteā admittēris.''

\navnote{-tū nōmen facile dictZ est = facile est nōmen dīcere; vōx difficilis auditā est = difficile est vōcem audire}Tabellārius: ``Nōmen meum nōn est facile dictū: Tlēpolemus nōminor.''

Iānitor: ``Quid dīcis? Cleopolimus? Vōx tua difficilis est audītū, quod forēs intersunt.''

Tabellārius: ``Mihi nōmen est Tlēpolemus, sīcut iam \navnote{dīcere dīxisse dictum (esse) erus -īm = dominus}dictum est. Tusculō veniō. Erum tuum quaerō.``

Iānitor: ``Sī erum salūtātum venīs, melius est aliō \navnote{dormītum īre = ad somnum īre}tempore venīre, nam hāc hōrā erus meus dormītum īre 50 solet, post brevem somnum ambulātum exībit, deinde lavātum ībit.``

Tlēpolemus: S\$i quis per hunc imbrem ambulat, nōn opus est posteā lavātum īre! At nōn veniō salūtātum. Tabellārius sum.``

Tandem iānitor forēs aperit et Tlēpolemum foris in \navnote{foris adv = extrā forēs}imbre stantem videt. Canis Irātus dentēs ostendit ac

\navnote{fremere --- 'rrr' facere ut canis īrātus re-tinēre \contr{} re- t tenēre cavēre}quod catēnā retinētur.

Iānitor: ``Cavē! Canis tē mordēbit!'' Sīc iānitor hominem intrantem dē cane ferōcī monet..

\navnote{re-sistere = cōnsistere solvere \contr{} vincīre}Tlēpolemus in limine resistēns ``Retinē canem!'' in-quit, ``Nōlī eum solvere! Nec vērō opus est mē monēre dē cane, ego enim legere sciō.'' Tabellārius solum intrā līmen aspicit, ubi CAVE CANEM scriptum

\navnote{imāgō -inis/}est īnfrā imāgi\navnote{terrēre=timentem facere propius comp \contr{} \contr{} prope ac-cēdere = ad-īre monēre -uisse -itum}nem canis ferōcis. ``Neque haec imāgō neque canis vērus mē terret!'' inquit, et propius ad canem accēdit.

``Manē foris!'' inquit iānitor, NŌlī ad hunc canem accēdere! Iam tē monuī!``

Tabellārius vērō, quamquam sīc ā iānitōre monitus est, alterum gradum ad canem versus facit --- sed ecce canis in eum salit catēnam rumpēns! Homō territus ex \navnote{rumpere terrēre -uisse -itum. cēdere = Īre prehendere}ōstiō cēdere cōnātur, sed canis iratus pallium eius dentibus prehendit et tenet.

\navnote{canis saliēns catēnam rumpit et pallium Tlēpolemi dentibus prehendit}``Ei! Canis mē mordet!'' exclāmat tabellārius, quī iam neque recēdere neque prōcēdere audet: canis fremēns \navnote{prō-cēdere --- \contr{} re-cēdere sinere}eum locō sē movēre nōn sinit.

\navnote{aperire -uisse -rtum dē-rīdēre}Iānitor rīdēns ``Quīn prōcēdis?'' inquit, ``Nōlī resistere! Ego tē intrāre sinō. Iānuam aperuī. Prōcēde in so vīllam!'' Sīc iānitor virum territum dērīdet.

``Id facilius est dictā quam factū'' inquit tabellārius, atque alterum gradum facere audet, sed canis statim in pedēs posteriōrēs surgit atque pedēs priōrēs in pectore eius pōnit! Tabellārius, tōtō corpore tremēns, ex ōstiō cēdit: sīc canis eum ē vīllā pellit. ``Removē canem!'' \navnote{Vr re-movēre iste -a -ud: iste canis = ille canis apud te}inquit ille, ``Iste canis ferōx mē intrāre nōn sinit.''

Iānitor eum tremere animadvertit iterumque dērīdet: ``Quid tremis? Hicine canis tē terruit?''

Tlēpolemus: ``Nōlī putāre mē ab istō cane territum esse! Si tremō, nōn propter canem ferōcem, sed propter \navnote{sub + abl/acc: sub tēctō esse sub tēctum īre solvere -visse solūtum arbitrāri = putāre}imbrem frigidum tremō. Admitte mē sub tēctum, iānitor --- amābō tē! Vinci istum canem ferōcem! Cūr eum solvistī?`` Tabellārius enim canem ā iānitōre solūtum esse arbitrātur.

Iānitor catēnam manū prehendit canemque paulum ā 95 tabellāriō removet. ``Nōlī arbitrārī'' inquit ``mē canem solvisse. Canis ipse catēnam suam rūpit. Ecce catēna \navnote{rumpere rūpisse ruptum}rupta.``

Tlēpolemus: ``Num canis catēnam ferream rumpere potest? Id nōn crēdō. At certē vestem scindere potuit: \navnote{posse potuisse nūper = paulō ante emere ēmisse ēmptum scindere scidisse scissum}vidēsne pallium meum novum, quod nūper magnō pretiō ēmī, scissum esse ā cane tuō?``

Iānitor: ``Istud pallium nōn est magnī pretiī, neque id \navnote{venīre vēnisse num quid = num aliquid ferre: ferō ferimus fertis fert ferunt aureus -īm = nummus aureus (i aureus = xxv dēnāriī = c sēstertii) scī-licet= --- ut sciunt omnēs hic haec hoc -ne: hicine? haecine? hocine?}nūper ēmptum esse crēdō. Sed quid tū vēnistī? Num quid tēcum fers?``

Tlēpolemus: ``Stultē rogitās, iānitor, nam iam tibi dīxī 'tabellārium mē esse'. Quid tabellāriōs ferre arbitrāris? aureōsne iānitōribus? Profectō nōs aurum nōn ferimus.''

Iānitor: ``Vōs scīlicet epistulās fertis.''

Tlēpolemus: ``Rēctē dīcis. Epistulam afferō ad Lūcium Iūlium Balbum. Hocine erō tuō nōmen est?''

Iānitor. ``Est. Quīn mihi istam epistulam dās?''

Tlēpolemus: ``Prius vincī canem et sine mē intrāre! Nōlī iterum mē forās in imbrem pellere!''

Iānitor, postquam canem vīnxit, ``Nōn ego'' inquit, \navnote{vincīre vīnxisse vīnctum pellere pepulisse pulsum}``sed hic canis tē forās pepulit. Nōlbnarrare 'tē ā iānitōre forās pulsum esseM''

Cane vīnctō, tabellārius tandem intrat epistulamque \navnote{cane vīnctō= --- postquam canis vīnctus est}ostendit ianitori; quī statim epistulam prehendit et in ātrium ad dominum suum fert.

\booksection{Grammatica Latīna}

Supīnum \navnote{salūtātum -tum}[I] Amīcī salūtātum veniunt (= quia salūtāre volunt).

'Salūtātum' supīnum

vocātur. Supīnum in -tum dēsinēns significat id quod aliquis agere vult et pōnitur apud verba 'īre', 'venīre', 'mittere' et alia.

Exempla: Rōmānī cotīdiē lavātum eunt. Mīlitēs oppidum oppugnātum mittuntur. Vesperi omnēs dormītum eunt. Mēdus et Lydia ad tabernam eunt ānulum ēmptum. \navnote{dictā -tū}[II] Id est facile dicta = facile est id dīcere.

Dict4ā? est alterum supīnum in -tū dēsinēns, quod apud adiectīva 'facilis' et 'difficilis! et pauca alia reperitur.

Exempla: Multa sunt faciliōra dicta quam facta. Vōx difficilis auditā est.

\booksection{Pēnsum A}

Hōrā nōnā erus ambulā- īre solet. Tabellāriī nōn mittuntur pecūniam postulā-. Hostēs castra expugnā- veniunt. Multī barbari Rōmam veniunt habitā-.

Syra: ``Verba medicī difficilia sunt audi-.'' Ea rēs est facfacilis. Nōmen barbarum difficile dic- est.

Verba: terrēre -isse -um; claudere -isse -um; dīcere ---isse -um; solvere -isse -um; emere -isse -um; rumpere -isse -um; aperire -isse -um; vincīre -isse -um; pellere -isse Cum; scindere -isse -um; venīre -isse; posse -isse.

\booksection{Pēnsum B}

Iānua cōnstat ē duābus ---, quae in --- vertuntur. Ōstiārius vel --- dīcitur servus quī hominēs in vīllam---. ---. Canis eius prope tam --- quam lupus est; itaque necesse est eum catēnā --- [ \contr{} ferrum] ---. Ānulus Lydiae nōn ex ferrō, sed ex --- factus est. Servus qul epistulās fert --- dīcitur, nam --- epistulae in tabellīs scribebantur.

Tabellārius --- iānuam stat. Iānitor, quī --- iānuam sedet, tabellārium dē cane ferōcī ---: ``Cavē! Canis tē ---!'' In solō intrā --- scriptum est '--- *---canem' īnfrā--- --- canis. ``Nec haec --- nec canis vērus mē ---'' inquit tabellārius, ac propius ad canem ---. Canis catēnam --- et vestem eius dentibus ---. Tabellārius neque --- neque --- audet. Iānitor: ``--- prōcēdis? Ego tē intrāre ---!'' Sīc iānitor hominem territum ---. Ille alterum gradum facit, sed canis eum ex ōstiō ---. Tlēpolemus tōtō corpore --- ex ōstiō ---. --- [= postrēmō] iānitor canem vincit.

\booksection{Pēnsum C}

Quid est iānitōris officium? Cūr necesse est canem eius vincīre? Ex quā māteriā cōnstat iānua? Vulcānus quis est? Quid in solō intrā līmen vidētur? Quid tabellārius sēcum fert? Unde venit Tlēpolemus et quem quaerit? Quōmodo iānitor ē somnō excitātur? Quid agit Iūlius post meridiem? Cūr canis saliēns catēnā nōn retinētur? Quid facit iānitor antequam tabellārius intrat?

\booksection{Vocābula nova}
\noindent\small nova, foris, līmen, cardō, iānitor, catēna, aurum, lignum, faber, tabellārius, imāgō, pallium, aureus, ferōx, ferreus, ligneus, cūstōdīre, admittere, vincīre, rogitāre, fremere, mordēre, retinēre, cavēre, monēre, resistere, solvere, terrēre, accēdere, salīre, rumpere, cēdere, prehendere, recēdere, prōcēdere, sinere, dēridēre, tremere, pellere, removēre, arbitrārī, scindere, iste, scīlicet, anteā, posteā, prius, tandem, nūper, forīs, forās, sīcut, quīn, extrā, intrā, supīnum.

\bookchapter{Cap. XXIII}{Epistula magistrī}

Iūlius, quī canem lātrāre audīvit, iānitōrem ātrium intrantem interrogat: ``Quis advēnit?''

\navnote{ad-venīre -vēnisse}Iānitor: ``Tabellārius advēnit Tusculō. Ecce epistula quam illinc ad tē tulit.'' Hoc dīcēns iānitor epistulam \navnote{illinc = ab illō locō ferre tulisse lātum trā-dere = dare (dē manū in manum) quid-nam? = quid? quis-nam? = quis? mittere misisse missum 7L IT ī ))}dominō suō trādit.

Iūlius: ``Quidnam hoc est? Quisnam Tusculō epistulam ad mē mīsit?''

Iānitor: ``Nesciō. Tantum sciō epistulam Tusculō missam et ā tabellāriō ad tē lātam esse.''

lūlius: ``Nōn opus est idem bis dīcere. Ego id quod semel dictum est bene intellegō. Recēde hinc ad canem \navnote{hinc = ab hōc locō dī-mittere = ā se mittere, iubēre discēdere}tuum!`` Sīc Iūlius iānitōrem dimittit.

Dominus cēram aspiciēns signum magistrī cognōscit (est enim parva eius imāgō) et ``Missa est'' inquit ``ā \navnote{litterae -ārum f pl}magistrō Diodōrō. Nōlō hās litterās legere, nam certē magister poscit pecūniam quam eī dēbeo. Duōrum \navnote{dēbēre + dat: ei dēbeō : ei dare dēbeō}mēnsium mercēdem magistrō dēbeō.``

Aemilia: At fortasse epistula aliās rēs continet. Quis \navnote{con-tinēre = in sē habēre integer -gra -grum: (sigruptum signō integrō (abl) dum signum integrum est}scit? Fortasse magister aliquid scripsit dē Mārcō? Signō integrō, nēmō scit.``

Iūlius signum rumpit et epistulam aperit. Ecce ea quae in epistulā magistrī scrīpta sunt: Diodōrus Iūliō salūtem dīcit.

\navnote{salūs -ūtis/: f: (ei) salūtem dīcere = (eum) salūtāre}Discipulus improbus atque piger est tuus Mārcus filius. Male recitat, foedē et prāvē scrībit, computāre nūllō modō 25 \navnote{umquam = ūllō tempore (neque umquam = et numquam) nēmō magister = nūllus magister}potest, neque umquam rēctē respondet cum eum interrogāvī. Fīlium tuum nihil docēre possum quia ipse nihil discere vult. Nēmō magister pēiārem discipulum umquam docuit. Valē.

Scrībēbam Tusculīkalendīs Iūniīs. Hic diēs me monet de 30 \navnote{kalendis Iūniīs = diē primo mensis Iūniī}pecūniā quam mihi dēbēs. Quārē Mārcus hodiē mercēdem sēcum nōn tulit? Mercēs numquam mihi trāditur ad diem. Iterum valē.

\navnote{vultus -ūs m = faciēs pallidus -a -um = albus (dē colōre vultūs dīcitur) pallēre = pallidus esse ob prp--acc = propter}Interim Mārcus, cuius vultus ad nōmen magistri colōrem mūtāvit, pallidus et tremēns patrem legentem 35 spectat. Cūr pallet puer? Pallet ob timōrem. (Quī timet pallēre solet.)

Item Aemilia vultum Iūliī sevērum intuētur. Postquam ille epistulam lēgit ūsque ad finem, uxor eum \navnote{legere lēgisse lēctum}interrogat: ``Quid scripsit magister?''

Iūlius: ``Prior epistulae pars dē alia rē est; in parte posteriōre magister mē monet dē pecūniā quam ei debeo.''

Aemilia: ``Cūr nōn solvis pecūniam quae magistrō \navnote{solvere pecūniam = dare pecūniam quae dēbētur}dēbētur? Certē magister, quī filiōs nostrōs tam bene scrībere et legere docet, mercēdem suam meret. Sed \navnote{merēre -uisse -itum}quidnam scriptum est in priōre epistulae parte? Nōnne magister Mārcum laudat?``

\navnote{laus laudis f --- verba laudantia rēs epistulā continētur = epistula rem continet}Iūlius: ``Hāc epistulā nūlla laus continētur, nec enim puer piger atque improbus laudem meret! Tune putās tē hīs litteris laudārī, Mārce?''

\navnote{a-vertere}Mārcus vultum ā patre āvertit nec ūllum verbum re\navnote{plānus -a -um = quī facile intellegitur facile adv}spondet, at genua trementia et vultus pallidus respōnsum plānum est, quod pater facile intellegit. Saepe silentium est respōnsum plānissimum.

\navnote{facere -iō fecisse factum}Tacente Mārcō, Aemilia ``Quid fecit Mārcus?'' in-quit, ``Dīc mihi omnia!''

Iulius: ``Mārcus prope omnia fecit quae facere nōn \navnote{dēbēre -uisse -itum omnis -e = tōtus SEXTVS}dēbuit! Haec epistula omnem rem plānam facit.---Ō Mārce! Iam plānē intellegō falsa esse omnia quae nōbīs nārrāvistī: magister plānīs verbīs scribit 'tē discipulum improbissimum fuisse ac foedē et prāvē scrīpsisse'!``

\navnote{ostendere -disse superior -ius comp \contr{} inferior In-scrībere}Iūlius: ``Aspice hanc tabulam: vidēsne nōmen 'Sextī' litteris plānīs in parte superiōre īnscriptum? Tune solēs nōmina aliēna īnscribere in tabulā tuā? Haec nōn tua, sed Sextī tabula est. Hocine negāre audēs?''

\navnote{negāre = dīcere nōn vērum esse'}Mārcus, quī iam nōn audet mentiri, nihil negat, sed omnia fatētur: ``Rēctē dīcis, pater

Tabula Sextī est. \navnote{fatēri --- \contr{} negāre}Tabulās mūtāvi inter pugnam.``

Aemilia: ``Pugnam? Quam pugnam nārrās?''

Iūlius: ``Mārcus mihi iam nārrāvit sēcum Sextō pugnāvisse.' --- Nōnne tibi satis fuit vestem tuam novam

\navnote{per-dere --- \contr{} servāre per-dere -didisse}Mārcus: ``Tabulam Sextī nōn perdidī, pater. Vidē: integra est tabula.''

Iūlius: ``At certē pater Sextī putabit eum perdidisse tabulam suam. Fortasse Sextus ā patre suō pūniētur ob hanc rem. Intellegisne factum tuum indignum esse? \navnote{factum -īn = id quod factum est pudēre: mē pudet- intellegō mē indignē fecisse /indignē factum esse rubēre = ruber esse pudor -ōris m \contr{} pudēre pudēre + gen: facti suī : ob factum suum post-hāc = post hoc tempus futūrus sum = erō pāritūrus sum = pārēbō pugnātūrus sum = pugnābo dormītūrus sum = dormiam prō-mittere -mīsisse -missum}Nōnne tē pudet hoc fecisse? Profectō mē pudet hoc ā meō filiō factum esse!``

Mārcus, quī paulō ante ob timōrem pallēbat, iam ru-bet propter pudōrem. Puerum pudet factīsuī.(Is quem factōrum suōrum pudet rubēre solet.)

Mārcus: ``Certē malus puer fuī, sed posthāc bonus /H puer futūrus sum: semper vōbīs pāritūrus sum, num- 85 quam pugnātūrus sum in viā nec umquam in lūdō dormītūrus sum. Hoc vōbīs promitto, pater et māter! Mihi crēdite!''

Mārcus 'sē mālum puerum fuisse' fatētur ac simul prōmittit 'sē posthāc bonum puerum futūrum esse, 90 semper sē parentibus pāritūrum esse nec umquam in viā pugnātūrum nec in lūdō dormītūrum esse' --- id \navnote{pugnātūrum esse ante-hāc = ante hoc facere, imp fac! facite! prōmissum -n = quod prōmissum est}quod saepe antehāc prōmīsit!

Iūlius: ``Primum fac quod promisisti, tum tibi crēdēmus!'' Iulius nōn crēdit Mārcum prōmissum factūrum esse.

Mārcus: ``Omnia quae promisi factūrus sum. Nōlī mē \navnote{factūrus sum = faciam}verberāre! Iam bis verberātus sum ā magistrō.``

``Ergō verbera magistri nōn satis fuērunt!'' inquit Iūlius, ``Profectō verbera meruistī!'' Tum oculōs ā filiō \navnote{in-clūdere \contr{} in clāvis uz}āvertēns: ``Abī hinc ab oculīs meīs! Dūc eum in cubiculum eius, Dāve, atque inclūde eum! Posteā fer mihi clāvem cubiculī!''

Postquam Dāvus puerum ex ātriō dūxit, dominus \navnote{dūcere dūxisse ductum herī= --- diē ante hunc fugere -iō fūgisse īre, part iēns euntis comitāri: eum c. = cum comes -itis m = is quī comitātur in-clūdere -sisse -sum Marcoō...ductō/inclūsō = postquam Mārcus... ductus/inclūsus est}``Haec omnia'' inquit ``facta sunt, quod Mēdus heri domō fQgit nec hodiē Mārcum in lūdum euntem et illinc redeuntem comitāri potuit. Posthāc Mārcum sine comite ambulāre nōn sinam. Crās Dāvus eum comitābitur; is certē bonus comes erit.'' 1io

Mārcō in cubiculum ductō atque inclūsō, Dāvus redit et ``Mārcus'' inquit ``inclūsus est. Ecce clāvis cubiculī.''

Iūlius clāvem sūmit ac surgit. Aemilia, quae putat eum ad Mārcum īre, Quō Is, Iūli?`` inquit, ''Mārcumne verberātum 15?`` Aemilia Mārcum ā patre verberātum īri \navnote{māter putat filium ā patre verberātum īri= --- māter putat patrem filium verberātūrum esse}putat. ``Nōlī eum verberāre! Nōn putō eum posthāc in viā pugnātūrum neque in lūdō dormītūrum esse.''

Iūlius: ``Putāsne iam mūtātum esse istum puerum? Ego eum nec mūtātum esse nec posthāc mūtātum īri crēdō! Quamquam heri ā mē semel verberātus est atque hodiē bis ā magistrō, nec verbera patris nec magistrī eum meliōrem fecērant.''

Aemilia: ``Ergō nōn opus est iterum eum verberāre. Nec laudibus nec verberibus melior fieri potest.''

Iūlius: ``Nōlī timēre, Aemilia! Mārcum in cubiculō relinquam. Iam epistulam scriptūrus sum.'' Iūlius dīcit 125 \navnote{scriptūrus sum --- scrībam}'sē epistulam scriptūrum esse.'

Aemilia: ``Cuinam scriptūrus es?''

Iūlius: ``Magistrō scīlicet. Crās Dāvus, Mārcī comes, epistulam meam sēcum feret, quae ā Mārcō ipsō trādētur magistrō. Tabellārius, quī epistulam magistri trādi- 130 \navnote{trā-dere -didisse -ditum}dit, tempus perdit, si foris respōnsum meum opperitur. Dīmitte eum, Dāve! Dīc eī 'respōnsum meum crās ā \navnote{'respōnsum ā Mārcō trāditum ir? = 'Mārcum respōnsum trāditūrum esse}Mārcō trāditum ir?'.``

Aemilia: ``Nōnne Mārcus simul cum epistulā tuā mercēdem dēbitam trādet magistrō?''

\navnote{dēbitus -a -um = quī dēbētur ``mercēdem solvere nōlō''}Iūlius: ``Minimē! Ego enim plānē respondēbō 'mē mercēdem solvere nōlle'!''

Aemilia: ``Quid ais? Nōnne tē pudet pauperi magistrō mercēdem negāre? Quam ob rem mercēdem dēbitam \navnote{negāre - dat dare nōlle quam ob rem? = ob quam rem? quārē? af-ferre at-tulisse al-lātum}solvere nōn vīs? Causam afferre oportet.``

Iūlius: ``Magister ipse mihi causam attulit.''

Aemilia: ``Quōnam modō? Quaenam causa allāta est ā \navnote{quae-nam? = quae?}magistrō?``

Iūlius: ``In hīs litteris magister ipse fatētur 'sē fīlium meum nihil docēre posse': ergō mercēdem nōn meruit. Pecūniam quae merita nōn est nōn solvam. Nōlō pecūniam meam perdere!''

Epistulam sūmēns Aemilia ``Itane scrībit magister?'' inquit; tum, epistulā lēctā, Hoc tē nōn excūsat, nam \navnote{epistulā lectā--- postquam epistula lēcta est (ab eā)}plānē scribit 'Mārcum ipsum nihil discere velle', et quī 150 nihil discere vult, nihil discere potest. Nōn modo posse, sed etiam velle opus est: quod nōn vīs, nōn potes.``

Iūlius rīdēns ``Rēctē dīcis'' inquit, Ego enim pecūniam solvere nōlō: ergō solvere nōn possum!``

Hoc dīcēns Iūlius epistulam magistrī scindit.

\booksection{Grammatica Latīna}

\navnote{FAT C --- Iūlius epistulam scindit}Participium et infinitivus futūri [A] Āctīvum.

\navnote{pugnātūrus sum = pugnābō pugnātūrum esse}Mīles: ``Fortiter pugnātūrus sum.'' Mīles dīcit 'sē fortiter pugnātārum esse.``

``Pugnātūrus -a -um' est participium

futūrī. Participium futūri est adiectīvum dēclīnātiōnis V11. '``Pugnātūrum esse! \navnote{-tūrlumesse}est Infinitivus futūri, quī cōnstat ex participiō futūrī et esse

Exempla: pāritūrus, dormītūrus, factūrus, scrīptūrus, futūrus.

Mārcus: ``Posthāc bonus discipulus futārus

\navnote{[1] pugnātūrum esse [2] pugnātzrum esse (2] pāāritirum esse [3] factūrum esse [4] āormīlūrum esse futūrum esse}gistrō pārirūrum esse nec umquam in lūdō dormitāūrum esse.' Magister nōn crēdit eum prōmissum factārum esse. 170

Puerī dīcunt 'sē posthāc bonōs discipulōs futūros esse, semper magistrō pāritrōs esse nec umquam in lūdō dormitarōs esse. Magister nōn crēdit eōs prōmissum factūrōs esse. [B] Passīvum.

Iūlius: ``Sextus ā magistrō laudābitur.'' Iūlius dīcit 'Sextum \navnote{laudātum īri -tum īri}ā magistrō laudātum r1.`` Laudātum ir? est infinitivus futūri passīvī, quī ex supīnō et 'īrī' cōnstat.

\booksection{Pēnsum A}

Mārcus: ``Posthāc bonus puer fu- sum et vōbīs pāri- sum.'' Mārcus prōmittit sē bonum puerum fu- esse et sē pārentibus pāri- esse.' Aemilia putat Mārcum ā Iūliō verberā- īri, sed Iūlius se epistulam scrip- esse' dīcit.

Verba: dūcere -isse -um; legere -isse -um; mittere -isse -um; inclūdere -isse -um; facere -isse -um; ferre -isse -um; afferre ---isse -um; trādere ---isse Cum; perdere -isse; ostendere -isse; fugere -isse.

\booksection{Pēnsum B}

Iānitor dominō epistulam --- [= dat]. Iūlius --- rumpit; iam signum --- nōn est. Mārcus --- et tremēns --- patris legentis aspicit. Epistula magistri omnem rem --- facit. Iūlius: ``Sextī est haec tabula; num hoc --- audēs?'' Puer nihil negat, sed omnia ---. Iūlius: ``Indignum est --- tuum. Nōnne tē --- facti tuī?'' Mārcus, quī paulō ante --- timōrem pallēbat, iam --- ob ---. Mārcus: ``--- semper bonus puer erō, hoc vōbīs ---!''

Hodiē Mēdus Mārcum ad lūdum --- nōn potuit, nam --- domō fūgit; itaque Mārcus sine --- ambulāvit. Magister pecūniam quam Iūlius eī--- --- postulat. Iūlius: ``Magistrō pecūniam --- nōlō, neque enim is mercēdem ---. Pecūniam meam --- nōlō.''

\booksection{Pēnsum C}

Ā quō epistula missa est? Quōmodo Iūlius epistulam aperit? Quid magister scripsit dē Mārcō? Cuius nōmen in tabulā īnscriptum est? Quam ob rem rubet Mārcus? Negatne Mārcus sē mālum discipulum fuisse? Quid Mārcus parentibus suīs prōmittit? Quō Dāvus puerum dūcit? Quārē Iūlius surgit? Quid Iūlius magistrō respondēbit? Cūr mercēdem solvere nōn vult?

\booksection{Vocābula nova}
\noindent\small nova, signum, litterae, vultus, laus, factum, pudor, prōmissum, verbera, clāvis, comes, integer, pallidus, plānus, superior, trādere, dimittere, dēbēre, continēre, salūtem dīcere, pallēre, solvere, merēre, āvertere, īnscribere, negāre, fatēri, perdere, pudēre, rubēre, prōmittere, inclūdere, comitāri, illinc, hinc, quidnam?, quisnam?, fortasse, umquam, posthāc, antehāc, heri, ob.

\bookchapter{Cap. XXIV}{Puer aegrōtus}

I Interim Quīntus lectō tenētur. Puer aegrōtus iterum \navnote{aegrōtus -a -um = aeger con-vertere = vertere latus -eris n}iterumque super lectum sē convertit nec dormīre potest sīve in latere dextrō cubat sīve in latere sinistrō. Itaque ē lectō surgere cōnātur, sed pēs dēnuō dolēre incipit. Puer territus pedēs nūdōs aspicit et ``Quid hoc est?'' inquit, ``Pēs dexter multō māior est quam pēs laevus!'' \navnote{laevus -a -um = sinister pār paris = aequus}Qulntus mīrātur quod pedēs, quī heri parēs erant, ho\navnote{re-cumbere --- \contr{} surgere (ēlectō) subitus -a -um = nōn exspectātus subitō adv strepitus -ūs m = sonus magnus iūxtā prp--acc = apud}diē tam imparēs sunt. Tum puer aegrōtus in lectō recumbit. Vīlla quiēta est: nūllus sonus audītur ab ūllā parte; etiam avēs, tempestāte subitā territae, in hortō silent.

At subitō silentium clāmōre et strepitū māximō rumpitur, nam Mārcus in cubiculō suō, quod iūxtā cubiculum Quīntī est, magnā vōce clāmat et forem manibus pedibusque percutit. Quīntus, quī tantum strepitum \navnote{per-cutere -iō = pulsāre continuō = statim}mīrātur, Syram vocat; quae continuō accurrit.

Syra: ``Quid est, Quīnte? Putāvī tē dormīre.''

Quīntus: Meēne dormīre per tantum strepitum? Quid agitur, Syra? Quam ob rem Mārcus sīc clāmat ac \navnote{quid agitur? = quid fit? valdē \contr{} validē frangere = rumpere (rem dūram) aliter = aliō modō}forem tam valdē percutit?``

Syra: ``Nōlī hoc mirāri: Mārcus forem frangere cōnātur, quod aliter exīre nōn potest. Sed tū quid agis? Doletne tibi pēs adhūc?''

Puer 'pedem sibi dolēre' ait: ``Valdē mihi dolet pēs, \navnote{sē, dat sibi dolor -ōris m \contr{} dolēre cōnari -ātum esse}ob dolōrem ē lectō surgere nōn possum.``

Syra: ``Tune ē lectō surgere cōnātus es?''

Quīntus: ``Certē surgere cōnātus sum, at necesse fuit \navnote{dolēre -uisse intuērī, imp -re: intuēre/ = aspice! cadere cecidisse māāior pede laevōmāior pede laevd=māior quam pēs laevus}mē continuō recumbere, ita pēs doluit. Intuēre pedēs meōs, Syra! Comparā eōs! Antequam dē arbore cecidī, parēs erant, nunc pēs dexter māior est pede laevō.``

Syra: ``Ego nōn mīror pedem tuum aegrum esse, quod dē tam altā arbore cecidisti; at mīror tē crus nōn frēgisse. Facile os frangere potuistī.''

\navnote{frangere frēgisse frāctum}Quīntus: ``Quis scit? Fortasse os frāctum est, nam pedem vix movēre possum sine dolōribus.''

\navnote{flēre = lacrimāre patl -ior: dolōrem p. = dolōrem ferre et-sī= 2 quamquam dolor gravis--- d. magnus intus \contr{} foris cōnsōlāri, loquī, imp -re: Q. imperat: ``Cōnsōlāre mē! loquere mēcum!''}Syra: ``Ossa tua integra sunt omnia. Nōlī flēre! Puerum Rōmānum sine lacrimīs dolōrem patī decet.''

Quīntus: ``Nōn fleō, etsī dolōrem gravem patior. Trīstis sum, quod mihi necesse est intus cubāre, dum aliī pueri forīs lūdunt. Cōnsōlāre mē, Syra! Cōnsīde hīc 40 iūxtā lectum et loquere mēcum!''

Syra iūxtā lectum ad latus puerī laevum cōnsīdit eumque sīc cōnsōlātur: ``Nōlī tristis esse quod hic intus cubaās: immō laetāre tē nōn inclūsum esse in cubiculō ut \navnote{laetārī, imp -re: laetāre! = gaudē!}frātrem tuum! Nec Mārcō licet cum aliīs pueris lūdere.``

Quīntus: ``Is nōn aegrōtat nec dolōrēs patitur.''

Syra: ``Etsī valet, certē tergī dolōrēs passus est.''

\navnote{patī passum esse}Quīntus: ``Estne verberātus Mārcus? Cūr ei nōn licet \navnote{fieri factum esse. red-īre -iisse ignōrāre = nescīre cupere -iō = velle iani (+fut) = continuō īre iisse ( \contr{} īvisse) comitārī -ātum esse}exīre? Quid factum est postquam frāter meus domum rediit? Omnia ignōrō. Dum hīc sōlus cubō, nihil nōscere possum, etsī omnia scīre cupiō.`` IH

Syra: ``Iam nārrābō tibi omnia quae facta sunt: Frāter

Qulntus: ``Nōnne Mēdus eum comitātus est?'' 55

Syra: ``Mēdus heri domō fūgit --- putō quia amīcam suam, quae Rōmae habitat, vidēre cupīvit.''

\navnote{cupere -īvisse vix umquam = prope numquam certō adv ali-qua femina, acc aliquam feminam nōscere nōvisse loquī locūtum esse ut hominēs āiunt}Quīntus: ``Quōmodo Mēdus, quī vix umquam Rōmae fuit, puellam Rōmānam nōscere potuit?''

Syra: ``Nesciō quōmodo, sed certō sciō eum aliquam feminam nōvisse, nam saepe dē eā locūtus est. Nihil difficile est amantī, ut āiunt.

``Mēdus igitur hodiē Mārcum comitāri nōn potuit. Nūper Mārcus sōlus rediit, sed pater filium suum redeuntem vix cognōvit neque eum ōsculātus est ut solet, \navnote{cognōscere cognōvisse cognitum}nam Mārcus nōn modo ūmidus erat quod per imbrem ambulāverat, sed etiam sordidus atque cruentus quod \navnote{cruentus -a -um \contr{} cruor}humī iacuerat et ā Sextō pulsātus erat. Pueri enim in viā pugnāverant: primum Mārcum pulsāverat Sextus, tum Maārcus et Titus Sextum pulsāverant. Hōc audito, do-minus Mārcum sevērē reprehendit.``

\navnote{reprehendere -disse -ēnsum}Quīntus: ``Māter quid dīxit?''

Syra: ``Māter tua nōn aderat, sed paulō post intrāvit. \navnote{lavāre lāvisse lautum vidēre vidisse vīsum}Tunc Mārcus iam lautus erat et vestem mūtāverat, do-mina eum sordidum atque cruentum nōn vīdit. Mārcus vērō 'sē bonum discipulum fuisse' dīxit, etsī in lūdō 75 \navnote{cēterum = praetereā, sed pēior cēteris (abl) = pēior quam cēterī}dormīverat nec magistrum recitantem audīverat--- --- cēterum in hāc rē is nōn pēior fuerat cēterīs, nam omnēs dormīverant! Postrēmō litterās pulchrās quās Sextus scrīpserat mātri ostendit atque dīxit 'sē ipsum eās litterās scrīpsisse.' Tam turpiter frāter tuus mentītus est!``

\navnote{mentiri mentitum esse}Quīntus: ``Quōmodo scis Mārcum mentītum esse et eās litterās ā Sextō scriptās esse?''

Syra: ``Quia nōmen 'Sextī' in tabellā scriptum erat. Et

Qulntus: ``Nūper aliquem iānuam pulsāre et canem 85 \navnote{ali-quis, acc ali-quem}valdē lātrāre audīvī. Tum subitō magnā vōce clāmāvit \navnote{tumultus -ūs m = clāmor verēri veritum esse mordēre momordisse morsum lupō ferōcior = ferōcior quam lupus}aliquis. Quidnam tantus ille tumultus significāvit?``

Syra: ``Tabellārius clāmāvit quod canem veritus est, nec sine causā, nam canis eum momordit et vestem eius scidit. Is canis lupō ferōcior est!''

Quīntus: ``Ego canem iānitōris nōn vereor neque umquam ab eō morsus sum.''

Syra: ``Id nōn mīror, nam canis saepissimē tē vīdit. \navnote{saepe, comp saepius, sup saepissimē nōvisse \contr{} \contr{} ignōrāre}Canis tē nōvit, ignōrat illum.``

Quīntus: ``Canis mē nōn modo nōvit, sed etiam dīli- 95 \navnote{dare dedisse datum lūdere lūsisse}git, nam multa ei ossa dedī et saepe cum eō lūsi. Cēterum quid attulit tabellārius?``

Syra: ``Epistulam attulit in quā magister scripserat 'Mārcum discipulum pigerrimum fuisse atque foedē et \navnote{primo ado = initiō fatēri fassum esse}prāvē scrīpsisse.' Tum Mārcus, quī prīmō omnia negāverat, 'sē mentītum esse' fassus est.``

Qulntus: ``Profectō verbera meruit!''

Syra: ``Magister eum iam bis verberāverat, nec igitur pater eum dēnuō verberāre voluit, sed in cubiculum \navnote{velle voluisse}inclūsit. Cēterum facile tibi est frātrem tuum reprehendere, dum ipse hīc in molli lectulō cubās. Tune ipse semper bonus discipulus es?``

Quīntus: ``Melior sum frātre meō! Heri laudātus \navnote{melior frāure meō = m. quam frāter meus}sum, quia pulchrē scrīpseram et recitāveram.``

Syra: ``Tune sōlus pulchrē scrīpserās et recitāverās?''

Quīntus: ``Immō omnēs praeter Mārcum laudātī sumus, quod pulchrē scrīpserāmus et recitāverāmus.''

Syra: ``Sl vōs laudātī estis, quod pulchrius scripserātis et recitāverātis quam Mārcus, nōnne ille tam rēctē scripserat quam vōs? Ego Mārcum bene nōvī, nec putō eum vōbīs stultiōrem esse.''

\navnote{stultior vōbīs (abl) = stultior quam vōs}Quīntus: ``At certē pigrior est nōbīs!''

\booksection{Grammatica Latīna}

Verbī tempora Plūsquamperfectum [A] Āctīvum.

Puer ūmidus est quod per imbrem ambulāvit.

Puer ūmidus erat quod per imbrem ambulāverat

\navnote{ambulāverat}Ambulāvit' praeteritum perfectum est. 'Ambulāverat? est tempus praeteritum plūsquamperfectum.

\navnote{plūs-quam-perfectum -erat}Plūsquamperfectum dēsinit in -erat (pers. iii sing.)5 quod ad infinītīvum perfectī sine -isse additur.

Exempla: [1] rta\&verat;

[2] pāruerat; [3] scripserat; [4]

\navnote{[1] reciuro era m}\navnote{recitōp erā rmts recitōp erā tis rccitāv era nt pāru erā s pāru era t pāruerāmus pāruerant [3] setīps era m sctipserās scrips era t scnps erā mus scnps erā tis scrīps era nt [4] dormtu era m dornriu erā s donnīv era t dormu erā mus dormu erā tis dormrp era nt eram fuerāmus erās fiierā tis erant erat fu fu}dormīverat. Heri magister puerum laudāvit, quia bonus discipulus- fuerat: bene recitāverat et scripserat, magistrō pāāruerat,

nec in lūdō donnīverat.

Bonus discipulus malō: Heri magister mē laudāvit, quia bonus discipulus fueram: bene recitāveram et scrīpseram, magistrō pārueram, nec in lūdō dormīveram. At tē verberāvit magister, quia malus discipulus fuerās: male rtdlāverās et scrīpserās, in lūdō dormiverās, nec magistrō pāruerās.``

Heri magister puerōs laudāvit, quia bonī discipulī fuerant: bene recixāverant et scripsērant, magistrō pārnerant, nec in lūdō dormīverant.

Bonī discipulī malīs: ``Heri magister nōs laudāvit, quia bonī discipulī fuerāmus: bene recitaverdmus et scrīpserāmus, magistrō pāmerāmus, nec in lūdō dormiveramus. At vōs verberāvit magister, quia malī discipulī fuerātis: male recitāverātis et scrīpserātis, in lūdō dormīverātis, nec magistrō pāruerātis.''

Singulāris Plāūrālis Persōna prima Persōna prīma Persōna secunda -erās

-erās Persōna tertia

-erat [B] Passīvum.

Magister Mārcum verberāverat = Mārcus verberātus erat ā \navnote{verbeīātus erat verberō*z erant verberāt us eram verberār us erās verberār userat verberō* īerāmus verberāt īerātis verberōx īerant}magistrō. Magister puerōs verberāverat = pueri verberātī erant ā magistrō.

Fīlius: ``Pater mē verberāvit, etsī iam ā magistrō verberātus 155 eram.'' Māter: ``Cūr verberātus

erās?`` Filii: Pater nōs verberāvit, etsī iam ā magistrō verbera: erāmus.'' Māāter: ``Cūr verberātī erātis?''

Persōna prīma

verberātus eram

verberārf erāmus Persōna secunda

verberātus erās

--- verberātī erātis Persōna tertia

verberātus erat

--- verberātī erant

\booksection{Pēnsum A}

Mārcus rediēns sordidus erat quod humī iac- et cruentus quod Sextus eum puls-. Pueri in viā pugn-. Tergum Mārcō dolēbat dolor quod magister eum verber-. Mārcus ā magistrō verber quod in lūdō dorm- nec magistrum recitantem aud-. Mārcus: aegrōtus ``Magister mē verberāvit quod in lūdō dorm- nec eum recitan:tem aud-, sed mē laudāvit quod bene comput- et scrips-.'' impār

Verba: lavāre -isse -um; vidēre -isse -um; mordēre -isse -um; dare -isse -um; reprehendere -isse -um; frangere \contr{} -isse -um; cognōscere -isse -um; nōscere -isse; lūdere mīrārī -isse; cadere -isse; īre ---isse; cupere -isse; velle -isse; patī: -um esse; loquī -um esse; fatēri ---um esse. f

\booksection{Pēnsum B}

Puer --- [= aeger] dormīre nōn potest sīve in --- dextrō iacet sīve in --- [= sinistrō]. Itaque surgere cōnātur, sed in lectō i ---, nam pēs --- [= rūrsus] dolēre incipit. Mārcus clāmat et forem valdē --- [-pulsat],.ita puer inclūsus forem --- cōnātur. \contr{} Quīntus clāmōrem et --- audit et Syram vocat; quae --- accurrit. Puer dīcit 'pedem --- dolēre': Pēs mihi dolet, ob --- --- dormīre nōn possum. Fortasse --- crūris frāctum est.`` Syra: Nōlī ---!'' Quīntus: ``Nōn fleō, --- [= quamquam] dolōrem \contr{} gravem ---.''

\booksection{Pēnsum C}

Cūr Quīntus mīrātur pedēs suōs aspiciēns? Quam ob rem pēs eius aegrōtat? Estne frāctum os crūris? Quārē Mārcus forem frangere cōnātur? Quid Syra Quīntō nārrat dē Mārcō? Cūr Mārcus rediēns cruentus erat? Quārē Mēdus ē vlllā fūgit? Quid Syra dē cane iānitōris putat? Cūr canis Quīntum dīligit? Cūr Iūlius Mārcum nōn verberāvit? Mare Chios * Dclphi Aegaeum Corinthus Athene Samos loponnesus4 e lo A. a, Ae \contr{} * e? O

a? Sparta t \contr{} Q. y 25 Melose\$d SS e o oc o 3 5 a e Rhodus ' ( '

\booksection{Vocābula nova}
\noindent\small -eram, -erāmus, -erās, erātis, -erat, erant, -tus, erāmus, crcon, -tx, -ta, erās, -tae, erai, -taerant, -tumerat, sonus, strepitus, tumultus, laevus, pār, subitus, cruentus, convertere, recumbere, percutere, frangere, flēre, patī, ignōrāre, nōscere, cupere, iūxtā, dēnuō, subitō, continuō, certō, primō, valdē, aliter, intus, etsī, cēterum, plūsquamperfectum.

\bookchapter{Cap. XXV}{Thēseus et Mīnōtaurus}

I Syra, postquam facta Mārcī nārrāvit, abīre vult, sed Quīntus ``Nōlī'' inquit ``mē relinquere! Te hīc manēre volō. Nārrā mihi aliquam fabulam!''

\navnote{agnus -Im = parvula ovis forte = nesciō cūr, sine causā regere = gubernāre Sōl:: deus sōlis trahere trāxisse tractum Hector -oris m Trōiānus -a -um \contr{} Trōia -ae/, f, urbs Asiae inter-ficere -iō -fecisse -fectum = mortuum facere moenia -ium npl = mūri}Syra: ``Quam fabulam mē tibi nārrāre vis? Fabulamne dē lupō et agnō quī forte ad eundem rivum vēnērunt? an fabulam dē puerō quī cupīvit regere equōs quī currum Sōlis per caelum trahunt?''

Tacente Quīntō, Syra pergit: An cupis audīre fabulam dē Achille, duce Graecōrum, quī Hectorem, ducem Trōiānum, interfecit atque corpus eius mortuum post currum suum trāxit circum moenia urbis Trōiae? an fabulam dē Rōmulō, quī prīma moenia Rōmāna aedifi-

nia humilia

dērīdēbat! Omnēs istās fabulās antīquās saepe audīvī. Iam vērō nec dē hominibus nec dē bēstiīs audīre cupiō. Nārrā mihi fabulam dē aliquō ferōcī mōnstrō, cui caput bēstiae et corpus hominis est et quod hominēs vīvōs vorat! Talem fabulam audīre cupiō.``

\navnote{vorāre = tōtum ēsse}Syra: ``At tāle mōnstrum tē terrēbit, Quīnte.''

\navnote{timidus -a -um = quī timet}Quīntus: ``Nōlī putāre mē puerum timidum esse! Ti-mor mōnstrōrum puerum Rōmānum nōn decet!M

``Nārrābō tibi fabulam dē Thēseō et Mīnōtaurō'' in-quit Syra, et sīc nārrāre incipit:

\navnote{ōlim = aliquō tempore, tempore antīquō terribilis -e \contr{} terrēre taurus -I m = bōs masculus}``In īnsulā Crētā ōlim vīvēbat mōnstrum

terribile, nōmine Mīnōtaurus, cui caput taurī, corpus virī erat. Mīnōtaurus in magnō labyrinthō habitābat.`` EN OO Q: DO

\navnote{taurus}Quīntus: ``Quid est labyrinthus?''

Syra: ``Est magnum aedificium quod frequentibus \navnote{aedificium -īn \contr{} aedificāre}mūris in plūrimās partēs dīviditur. Nēmō quī tāle aedificium semel intrāvit rūrsus illinc exīre potest, etsī iānua \navnote{patēre = apertus esse Athēniēnsis -e \contr{} Athēnae -ārum / pl,urbs Graeciae quī:: Daedalus com-plūrēs -a = plūrēs quam duo mīrābilis -e \contr{} mīrārī quot-annīs = omni annō, quōque annō}patet. Labyrinthus ille in quō Mīnōtaurus inclūsus tenēbātur, ā Daedalō, virō Athēniēnsl, aedificātus erat. Quī iam antequam ex urbe Athēnīs in Crētam vēnit, complūrēs rēs mīrābilēs fecerat.

``Mīnōtaurus nihil praeter hominēs vīvōs edēbat. Ita-que

complūrēs

adulēscentēs virginēsque

quotannīs \navnote{Athēnīs (abl) = ab/ex urbe Athēnīs saevus -a -um = ferōcissimus illūc = ad illum locum mors mortis f}Athēnīs in Crētam mittēbantur, quī omnēs in labyrinthō ā mōnstrō illō saevō vorābantur. Nāvis quā Athēniēnsēs illūc vehēbantur vēla ātra gerēbat, nam eō colōre significātur mors.`` Il

Quīntus: ``Quam ob rem tot Athēniēnsēs ad mortem certam mittēbantur?''

\navnote{rēx rēgis m: vir qul terram/urbem regit Mīnōs -ōis m: rēx Crētae expugnātiō -ōnis f \contr{} expugnāre pere; auri cupidus = quī aururi cupit velle: volēbat (imperf) bene/male velle + dat = amlcus/inimīcus esse necāre = interficere patriae amāns = quī patriam amat f 2 magna laus Athēnīs- in urbe Athēnīs Athēnās = ad/in urbem Athēnās ibi = illīc}Syra: ``Rēx Mīnōs, quī tunc Crētam regēbat, paucīs annīs ante urbem Athēnās bellō expugnāverat. Post expugnātiōnem urbis Mīnōs, cupidus auri atque sanguinis, nōn modo magnam pecūniam, sed etiam hominēs vīvōs ab Athēniēnsibus postulāverat. Rēx enim Athēniēnsibus male volēbat, quod fīlius eius paulō ante ab iīs so necātus erat.

``Eō tempore Thēseus, vir patriae amāns atque glōriae cupidus, Athēnīs vīvēbat. Quī nūper Athēnās vēnerat neque ibi fuerat cum urbs ā rēge Minōe expugnata est. Thēseus, quī patrem Mīnōtauri, taurum album, iam necāverat, novam glōriam quaerēns Mīnōtaurum ipsum quoque interficere cōnstituit. Itaque ūnā cum cēteris \navnote{cōn-stituere -uisse -ūtum vēlis ōrnātus = vēla gerēns cōnscendere -disse proficīscī -fectum esse iubēre iussisse iussum eum ā mīlitibus dūcī iussit = mīlitēs eum dūcere iussit cōn-spicere -iō -spexisse -spectum incipere -iō coepisse coeptum}Athēniēnsibus nāvem velis ātris ōrnātam cōnscendit et in Crētam profectus est. Ibi continuō rēgem Mlnōem adiit, quī eum ā mīlitibus in labyrinthum dūcī iussit. 6(

``Mīnōs autem filiam virginem habēbat, cui nōmen erat Ariadna. Quae cum primum Thēseum cōnspexit, eum amāre coepit cōnstituitque eum servāre. *

\navnote{ac-cēdere -cessisse}``Ariadna igitur, antequam Thēseus labyrinthum intrāvit, ad eum accessit et sīc loquī coepit: ''Contrā Mīinquit Thēseus ``mihi auxilium ferent contrā illum. Hodiē certē Mīnōtaurum

occīdam atque

cīvēs meōs ā \navnote{occīderc-disse -sum = interficere (gladiō) cīvis-is m = quī/quae in eādem urbe habitat}mōnstrō illō terribili servābō. Bonum gladium gerō. Ad pugnam parātus sum.`` Tum Ariadna ''Hoc nōn dubitō'' \navnote{exitus -ūs m \contr{} exīre reperire repperisse repertum filum -īn = līnea tenuis}inquit, ``sed quōmodo exitum labyrinthī posteā reperiēs? Nēmō adhūc per sē viam ē labyrinthō ferentem repperit. Ego vērō tibi auxilium feram: ecce fīlum ā Daedalō factum quod tibi viam mōnstrābit. Auxiliō huius fili hūc ad mē redībis.'' Haec locūta, Ariadna \navnote{hūc = ad hunc locum haec locūta = postquam haec locūta est opperire/ = exspectā!}Thēseō fīlum longum dedit; atque ille ``Opperīre mē'' 75 inquit ``hīc ad iānuam! Nōlī timēre! Ego mortem nōn timeō. Sine timōre mortis contrā hostem eō. Brevī hūc \navnote{brevī ad - brevī tempore, mox}redībō, neque sine tē, Ariadna, in patriam revertar. Illūc tē mēcum dūcam neque umquam tē relinquam. \navnote{pollicēri = prōmittere}Hoc tibi polliceor.``

``Tum Thēseus, filum Ariadnae post sē trahēns, labyrinthum intrāvit ac sine morā Mīnōtaurum in mediō \navnote{mora -ae/: f: sine morā petere -īvisse -Itum Mīnōtaurō occīsō --- postquam M.us occisus est sequī secūtum esse filum secūtus:: postquam filum secūtus est, filum sequēns}labyrinthō exspectantem petīvit, quem post brevem pugnam gladiō occīdit. Mīnōtaurō occīsō, Thēseus fīlum Ariadnae secūtus exitum labyrinthī facile repperit. Ita Thēseus ob amōrem patriae cīvēs suōs ā mōnstrō saevissimō servāvit.

\navnote{nex necis f \contr{} necāre}``Haec sunt quae nārrantur dē nece Mīnōtaurī.''

Hīc Quīntus Perge'' inquit ``nārrāre dē Thēseō et 1H Ariadnā! Nōnne illa Thēseum secūta est?''

Syra: ``Thēseus ē labyrinthō exiēns ''Mīnōtaurus ne\navnote{laetāmini/ = gaudēte! intuēminī! = spectāte! sequiminī mē! = venīte mēcum! nāvigāre: ad nāvigandum}cātus est'' inquit, Laetāmini, cīvēs meī! Intuēminī gladium meum cruentum! Sequiminī mē ad portum! Ibi nāvis mea parāta est ad nāvigandum.`` Tum Ariadnam cōnspiciēns ''Et tū'' inquit ``sequere mē! Proficiscere mēcum Athēnās!'' Ariadna, quae nihil magis cupiēbat, ``Parāta sum ad fugiendum'' inquit, atque sine morā nāvem Thēseī cōnscendit. Thēseus nāvem solvit et cum \navnote{fugiendum Naxus -i f Naxum = ad īnsulam Naxum Naxō = ab insulā Naxō relinquere -liquisse -lictum litus -oris n quaerere -sīvisse -situm}filiā rēgis nāvigāvit Naxum; ibi vērō nocte silenti Ariadnam dormientem relīquit atque ipse Naxō profectus est. Māne Ariadna ē somnō excitāta amīcum in lītore quaesivit neque eum repperit. Puella misera ab humili lītore

\navnote{prō-spicere -iō = ante sē aspicere ascendere -disse Thēseus, voc Thēseu! revertere! = redī! red-dere -didisse -ditum cōnspectus -ūs m \contr{} cōnSpicere; ē cōnspectū =}in altum saxum ascendit, unde prōspiciēns nāvem Thēsel procul in mari vīdit. Tum, etsī vōx eius ā nūllō audiri poterat, Ariadna amīcum suum fugientem vocāvit: ``Thēseu! Thēseu! Revertere ad mē!'' neque ūllum respōnsum ei redditum est praeter vōcem ipsīus quam dūra saxa reddidērunt.

Brevī nāvis ē cōnspectū eius abiit neque iam ūllum vēlum in mari cernēbātur. Arino adna igitur in litus dēscendit atque hūc et illūc currēns multīs cum lacrimīs capillum et vestem scindēbat, ut hominēs quī maerent agere solent --- ita maerēbat virgō \navnote{maerēre \contr{} \contr{} gaudēre ante omnēs:: magis quam omnēs aliōs}miserrima, quae ā virō quem ante omnēs amābat sōla relicta erat inter ferās īnsulae sīcut agnus timidus inter saevōs lupōs.``

Quīntus: ``Cūr Thēseus amīcam suam ita dēseruit?''

\navnote{dēserere -uisse -rtum}Syra: ``Tālēs sunt viri, mī puer. Montēs aurī feminīs pollicentur, tum prōmissa obliviscuntur ac feminās sine \navnote{oblivisci -litum esse}nummō dēserunt! Quis tam facile prōmissum oblīvīscitur quam vir quī feminam amāvit? Ego quoque ōlim 120 dēserta sum ab amīcō pecūniōsō quī mihi magnās rēs \navnote{: nēemo...! pollicēri -itum esse cupiditās -ātis f \contr{} cupidus ob amōrem viri:: quia virum amō/amāvl}pollicitus erat. Nōlī vērō putāre mē ob cupiditātem pecūniae amāvisse eum, ego eum amābam quia eum prōbum virum esse crēdēbam. Etiam nunc maereō ob amōrem illīus virī.``

Quīntus: ``Oblīvīscere illīus virī improbī quī tē tam \navnote{oblīvīsci + genJacc: o. hominis, o. reī/rem}turpiter dēseruit!``

Syra: ``Nōn facile est amōris antīquī oblivisci. Sed \navnote{nārrātiō -ōnis f \contr{} \contr{} nārrāre}hoc tū nōndum intellegis, mī Quīnte. Redeō ad nārrātiōnem fabulae, quam prope oblīta sum, dum dē aliīs 130 rēbus loquor.

\navnote{Naxf = in insulā Nāxō}``Ariadnā Naxi relictā, Thēseus ad patriam suam nāvigābat. Interim pater eius Aegeus, rēx Athēniēnsium, ab altō saxō in mare prōspiciēbat. Brevī nāvis filii in cōnspectum vēnit, sed nāvis rediēns eadem vēla ātra 135 gerēbat quae abiēns gesserat: Thēseus enim post necem \navnote{gerere gessisse gestum arbitrātus = quī arbi- trātus est}Minotauri vēla mūtāre oblītus erat! Itaque Aegeus, arbitrātus mortem filii eo colōre significārī, sine morā dē saxō sē iēcit in mare, quod ā nōmine eius etiam nunc \navnote{iacere -iō iēcisse iactum}'mare Aegaeum' vocātur. ``Post mortem rēgis Aegeī fīlius eius Thēseus rēx Athēniēnsium factus est. Qul multōs annōs Athēnās magnā cum glōriā rēxit.'' \navnote{regere rēxisse rēctum finem nārrātiōnis facit = nārrāre dēsinit}Hīs verbīs Syra fīnem nārrātiōnis facit.

\booksection{Grammatica Latīna}

Verba dēpōnentia Imperātīvus. Laetāre/ = gaudē!

Laetāmini/ = gaudēte! \navnote{laetāre laetamini intuē]re intuēminī revertere revertimini [4] [4]partīri partīre partf)minī -re ----minī}Intuēre/ = spectā!

Intuēmtm/ = spectāte! 150 Revertere/ = redī!

Revertimini! = redīte! Partīre/ = dīvide! 'Laetāre laetamini! et cētera sunt imperāativi verbōrum dēpōnentium. Imperāātivus verbī deponentis: singulāris -re, plūrālis -minī.

\booksection{Pēnsum A}

Mēdus (ad Lydiam): ``In Graeciam ībō. Comit- mē! Proficīsc- mēcum! Hoc mihi pollic-! Oblīvīsc- Rōmae! Ōsculmē!'' Mēdus (ad nautās): ``Iam proficīsc-, nautae! Ventus secundus est: intu- caelum! Laet-!'' Verba: trahere -isse -um; petere -isse -um; quaerere -isse -um; occidere -isse -um; relinquere -isse -um; interficere -isse -um; cōnstituere ---isse -um; dēserere ---isse -um; cōnspicere ---isse -um; incipere ---isse -um; reddere -isse -um; reperire -isse -um; iubēre ---isse -um; gerere -isse -um; iacere -isse -um; regere ---isse -um; accēdere ---isse; ascendere ---isse; proficīscī -um esse; sequī ---um esse; oblivisci -um esse.

\booksection{Pēnsum B}

Quīntus --- dē lupō et --- audīre nōn vult nec fabulam dē Achille, quī Hectorem --- et corpus eius mortuum post --- suum circum --- [= mūrōs] Trōiae ---.

Labyrinthus est magnum --- unde nēmō exīre potest, etsī iānua ---. Thēseus, qul patrem Mīnōtauri, --- album, iam necāverat, Mīnōtaurum ipsum quoque --- [= interficere] ---. Antequam Thēseus, ad pugnam ---, labyrinthum intrāvit, Ariadna, fīlia---, ---, eī--- --- longum dedit. Ita Ariadna Thēseō --- tulit, nam ille filum sequēns --- labyrinthī repperit. Post --- Minōtauri Thēseus cum Ariadnā Naxum nāvigāvit atque --- [= illic] eam --- [= relīquit]. Ariadna ab altō --- prōspiciēbat, sed --- [= mox] nāvis Thēseīē --- eius abiit. Puella misera in --- dēscendit, ubi --- et --- currēns capillum scindēbat, ut faciunt eae quae ---. Thēseus post --- patris multōs annōs Athēnās ---.

\booksection{Pēnsum C}

Quid fecit Achillēs? Quid fecit Rōmulus? Quālem fabulam Qulntus audīre cupit? Ubi habitābat Mīnōtaurus? Quid est labyrinthus? Quid Mīnōtaurus edēbat? Cūr tot Athēniēnsēs ad eum mittēbantur? Quis Athēniēnsēs ā Mīnōtaurō servāvit? Quōmodo Thēseus exitum labyrinthī repperit? Sōlusne Thēseus ē Crētā profectus est? Ubi Thēseus Ariadnam relīquit? Ā quō mare Aegaeum nōmen habet?

\booksection{Vocābula nova}
\noindent\small nova, fabula, agnus, currus, moenia, mōnstrum, taurus, labyrinthus, aedificium, mors, rēx, expugnātiō, glōria, auxilium, cīvis, exitus, filum, mora, nex, litus, saxum, cōnspectus, cupiditās, nārrātiō, humilis, timidus, terribilis, mīrābilis, saevus, cupidus, parātus, regere, trahere, interficere, aedificāre, vorāre, patēre, necāre, cōnstituere, occīdere, pollicēri, prōspicere, dēscendere, maerēre, dēserere, oblivīscī, coepisse, complūrēs, forte, quotannīs, ōlim, ibi, illūc, hūc, brevī.

\bookchapter{Cap. XXVI}{Daedalus et Īcarus}

Quīntus: ``Nōnne rēx Mīnōs Thēseum cum Ariadnā fugientem persecūtus est?''

\navnote{per-sequī = (fugientem) sequī}Syra: ``Certē rēx eōs persequī coepit, sed nāvis Thēseī \navnote{celer -eris -ere: equus celer (m) nāvis celeris (f) pīlum celere (/) pflum celere (n) capere-iō cēpisse captum cōn-ficere -iō -fecisse -fectum = facere vērum = sed}nimis celeris fuit. Mīnōs, quamquam celeriter nāvigāvit, nōn tam celer fuit quam Thēseus neque eum cōnsequī potuit. Tum rēx īrātus cēpit Daedalum, quī fīlum cōnfecerat et Ariadnae dederat, eumque in labyrinthum inclūdī iussit unā cum eius [carō filiō; vērum pater et filius mīrābilī modō ē labyrinthō fūgērunt. Crās tibi \navnote{fuga -ae f \contr{} fugere ad mnandtm = ad nārrātiōnem cōn-sūmere in nārrandō = in nārrātiōne finis nārrandī = finis nārrātiōnis quoniam = quia reliquus -a -um = cēterus}nārrabō dē fugā eōrum, hodiē plūs temporis ad nārrandum nōn habeō: iam hōram cōnsūmpsi in nārrandō.``

Quīntus: ``Neque tempus melius cōnsūmere potuistī! Nōn oportet in mediā fabulā fīnem nārrandī facere. Quoniam māiōrem fabulae partem iam nārrāvistī, partem reliquam quoque nārrāre dēbēs. Ego parātus sum ad audiendum.''.

\navnote{cupidus audiendi --- quī audire cupit reliqua fabula = reliqua fabulae pars}Ad hoc Syra ``Ergō'' inquit, ``quoniam tam cupidus es audiendi, reliquam fabulam tibi nārrābō:

``Daedalus in labyrinthō inclūsus cum filio suō intrā mūrōs errābat nec exitum invenīre poterat, etsī ipse 20 \navnote{in-venīre = reperire}labyrinthum aedificāverat. Quoniam igitur aliae viae \navnote{audāx -ācis = audēns ef-fugere \contr{} ex + fugere cōnsilium -īn: c. eius = quod facere cōnstituit cōnsīdere -sēdisse carcer -eris m per nōs: sine auxiliō quis-quam = ūllus homō nēmō) iuvāre iūvisse: eum i. = el auxilium ferre haud = nōn paene = prope diī = deī}clausae erant, ille vir audāx per āera effugere cōnstituit. īcarus autem, quī cōnsilium patris ignōrābat, humī cōnsēdit et ``Fessus sum'' inquit ``ambulandō in hōc carcere, quem ipse nōbīs aedificāvistī, pater. Ipsī per nōs hinc effugere nōn possumus, neque quisquam nōs in fugiendō iuvāre poterit, ut Thēseum iūvit Ariadna. Haud longum tempus nōbīs reliquum est ad vīvendum, nam cibus noster paene cōnsūmptus est. Ego iam paene mortuus

sum. Nisi diī nōs iuvābunt, numquam vivi 30 hinc ēgrediēmur. Ō diī bonī, auxilium ferte nōbīs!``

``Daedalus vērō ''Quid iuvat deōs invocāre'' inquit, ``dum hīc quiētus sedēs? Quī ipse sē iuvāre nōn vult, auxilium deōrum nōn meret. At nōlī timēre! Ego cōnsilium fugae iam excōgitāvī. Etsī clausae sunt aliae viae, ūna via nōbīs patet ad fugiendum. Intuēre illam aquilam quae in magnum orbem circum carcerem nostrum volat! Quis est tam liber quam avis quae trāns montēs, \navnote{carcer ex-cōgitāre orbis -ism liber -era -erum trānsprp+acc trāns prp acc: - AM mōns trāns zt}vallēs, flūmina, maria volāre potest. Quīn avēs caelī imitāmur? Mīnōs, quī terrae marique imperat, dominus \navnote{imitārī: avēs i. ē-volāre quidem = certē cupidus discendī}āeris nōn est: per āera hinc ēvolābimus! Hoc est cōnsilium meum. Nēmō nōs volantēs persequī poterit.`` ''Ego quidem studiōsus sum volandī'' inquit īcarus, ``sed ālae necessāriae sunt ad volandum. Quoniam diī nōbīs ālās nōn dedērunt, volāre nōn possumus. Hominēs sumus, nōn avēs. Nēmō nisi deus nātūram suam mutāre potest. \navnote{nisi deus: praeter deum nōn item: nōn possunt}Avēs nāturā volāre possunt, hominēs nōn item.`` Tum \navnote{ars artis/: ars mea = id quod facere possum ignis `` -ism A}Daedalus ``Quid ego facere nōn possum?'' inquit, ``Prōfectō arte meā ipsa nātūra mūtārī potest. Multās rēs mīrābilēs iam cōnfecī, quae artem meam omnibus dēmōnstrant, nōn sōlum aedificia magnifica, ut hunc labyrinthum, vērum etiam signa quae sē ipsa movēre possunt tamquam hominēs vivi. Ālās quoque cōnficere \navnote{tam-quam = sīcut opus -eris n}possum, quamquam opus haud facile est.`` ''Audāx quīss dem est cōnsilium tuum'' inquit īcarus, ``sed omne cōn\navnote{cōnsilium fugiendī = cōnsilium fugae per-ficere -iō -fecisse}silium fugiendī mē dēlectat, ac tū id quod semel excōgitāvistī perficere solēs.`` ''Certē cōnsilium meum perficiam'' inquit pater, ``Ecce omnia habeō quae necessāria \navnote{penna}sunt ad hoc opus: pennās, cēram, ignem. Igne cēram \navnote{mollīre = mollem facere lacertus -īm = bracchium superius ingēns -entis = valdē magnus iungere iūnxisse iūnctum figere fixisse fixum}molliam, cērā molli pennās iungam et in lacertīs figam.``

``Daedalus

igitur arte mīrābilī sibi et fīliō suō ālās ingentēs cōnfecit ex pennīs, quās cērā iūnxit et in lacertis fixit. Postquam finem operis fecit, ``Opus iam perfectum est'' inquit, ``ecce exemplum artis meae novissimum. ÁAvēs quidem nōn sumus, sed avēs imitābimur in volandō. Ventō celerius trāns mare volābimus, nūlla avis nōs cōnsequī poterit.''

``īcarus studiōsus volandī ālās hūc illūc mōvit, nec sē \navnote{movēre movisse mōtum levāre ( \contr{} levis)}suprā humum levāre potuit. ``Ālae mē sustinēre nōn possunt'' inquit, ``Tu docē mē volāre!'' Statim Daedalus sē ālis levāvit et ``Nisi ālās rēctē movēs'' inquit, ``volāre nōn potes. Imitāre mē! Haud difficilis est ars volandī. Movē ālās sūrsum deorsum hoc modō!'' Ita pater fīlium \navnote{sūrsum[f]edeorsum[4]}suum artem volandī docuit tamquam avis pullōs suōs. Tum puerum ōsculātus ``Parātī sumus ad volandum'' \navnote{puerum ōsculātus = postquam puerum ōsculātus est}inquit, ``sed prius hoc tē moneō: volā post mē in mediō āere inter caelum et terram, nam si in infimo āere prope \navnote{ferus, comp inferior sīn = si autem}mare volābis, pennae

ūmidae

fīent, sīn volābis in summō āere prope caelum, ignis sōlis cēram molliet \navnote{īnfimus; comp superior ūrere ussisse ustum = igne cōnsūmere estō! estōte! = es! este!}atque pennās ūret. Nōlī nimis audāx esse in volandō! Cautus estō, mi fili! Iam sequere mē! Carcerem nostrum effugimus --- liberi sumus!``

``Haec verba locūtus Daedalus cum filiō sūrsum ē labyrinthō ēvolāvit, neque quisquam fugam eōrum animadvertit

nisi aliquī pāstor, quī forte suspiciēns eōs 85 \navnote{aliquis:: aliquī vir su-spicere -iō = sūrsum aspicere}tamquam magnās avēs volantēs vīdit ac deōs esse arbitrātus est. Mox pater et filius Crētam relīquērunt, ne-que vērō rēctā viā Athēnās in patriam suam volāvērunt, \navnote{lībertās -ātis f \contr{} liber dē-spicere (--- (*-*suspicere) = deorsum aspicere multitūdō -inis f/( \contr{} multī) = magnus numerus}sed novā lībertāte dēlectātī in magnum orbem suprā mare Aegaeum volāre coepērunt.

īcarus dēspiciēns 90 multitūdinem īnsulārum mirātus est: ``Ō, quot parvae īnsulae in marī ingentī sunt!'' Daedalus vērō ``Illae īnsulae'' inquit ``haud parvae sunt3 quamquam parvae esse videntur. Certē Mēlos īnsula, quae īnfrā nōs est, nōn \navnote{tibi vidētur: tū putās rēs magna esse mihi vidētur = rem magnam esse putō paen-insula -ae f \contr{} paene Insula}tam parva est quam tibi vidētur.`` īcarus: ''Sed illa In- 95 sula quae nōbīs ā sinistrā est multō māior esse mihi vidētur. Quae est illa īnsula?`` Daedalus: ''Peloponnēsus est, Graeciae pars, nec vērō īnsula est, sed paenīnsula, nam Peloponnēsus terrā angustā, quae Isthmus vocātur, cum reliquā Graeciā coniungitur. Prope Isthmum sita est Corinthus, urbs pulcherrima, nec procul absunt \navnote{orbis terrārum = omnēs terrae (Eurōpā, Asiā, Āfrica) nimis audāx}Athēnae, patria nostra.`` ''Sī altius volābimus, nōn sōlum Graeciam, sed paene tōtum orbem terrārum spectābimus'' inquit puer temerārius atque etiam altius sē levāvit. Illinc nōn sōlum magnās Eurōpae et Asiae partēs dēspiciēbat mīrāns, vērum etiam ōram Āfricae procul cernēbat, deinde suprā sē sōlem in caelō serēnō lūcen\navnote{-spicere -iō -spexisse -spectum}tem suspexit. Statim puer, cupidus sōlem propius aspiciendī, quamquam pater eum monuerat, in summum

Hīc Quīntus, quī exitum fabulae studiōsē exspectat, \navnote{exitus = finis accidere -disse = fierī:}interrogat: ``Quid tum accidit?'' EDT TU d P2 * Ss: N AN.

\navnote{orbis terrārum}Syra: ``Tum factum est id quod necesse erat accidere: ignis sōlis propinquī cēram, quā pennae iūnctae et fīxae \navnote{propinquus -a -um = qul prope est quatere -iō = celeriter hūc illūc movēre mergere mersisse mersum [carius -a -um \contr{} Icarus}erant, mollīvit et pennās ussit. Puer territus, lacertōs nūdōs quatiēns, in mare cecidit ac mersus est, neque pater ei auxilium ferre potuit. Ea maris Aegaeī pars in quā Icarus mersus est ā nōmine eius 'mare īcarium' \navnote{invenīre -vēnisse -ventum: omnis -e = tōtus tempus dormiendi est = tempus est dormīre}appellātur. Item īnsula propinqua, in cuius lītore cor-pus pueri inventum est, etiam nunc 'īcaria' vocātur.

``Ecce omnem fabulam habēs dē puerō temerāriō qul lībertātem quaerēns mortem invēnit. Iam tempus dormiendī est. Nōnne fessus es longās fabulās audiendō?''

Quīntus caput quatit et ``Nōn sum fessus, nec illa fabula longa esse mihi vidētur. Ex omnibus fabulīs haec 125 \navnote{cāsus -ūs m \contr{} cadere valdē, comp magis, sup māximē ab-errāre}dē cāsū īcari mē māximē dēlectat, etiam magis quam illa dē filiō Sōlis, quī currum patris regere cōnātus item dē summō caelō cecidit, quod ab orbe sōlis stultē aberrāverat. Semper valdē dēlector tālēs fabulās audiendō.``

Syra: ``Ego nōn minus dēlector nārrandō illās fabulās, nōn modo quod ipsae per sē pulcherrimae esse mihi videntur, sed etiam quia exitūs fabulārum hominēs temerāriōs optimē monent. Talis enim est hominum nātūra, et quidem māximē puerōrum. Nōn sōlum dēlectandī \navnote{quārē nārrātur fabula? monendī causā nārrātur rTe-vocāre}causā, vērum etiam monendī causā nārrātur fabula dē 135 fīliō Daedalī, nam quod illī puerō accidit, idem omnī puerō accidere poterit, nisi patri suō pāret. Nōlīīcarum imitārī, ml Quīnte! Semper cautus estō! Vērum haud necesse est tē ā mē monērī post id quod heri tibi accidit. Certē ille cāsus tuus melius tē monet quam ulla fabula!``

Hīs verbīs puerō monitō, Syra tandem nārrandī fīnem facit. Neque Quīntus eam abeuntem revocat, sed in lectō recumbit oculōsque claudit. Mox puer in somnīs sibi vidētur ālis ōrnātus trāns montēs et flūmina \navnote{volāre sibi vidētur = se volāre videt/putat}volāre.

\booksection{Grammatica Latīna}

Gerundium

\navnote{nārrandum nārrandī nārrazdō}Hōrā in nārrandō cōnsūmptā, Syra plūs temporis ad nārrandum nōn habet et fīnem facit nārrandī.

--- 'Nārrandum' est gerundium;

\navnote{-ndum -ndī-ndō amandum amanif [2]docērc doce ndum}quod pōnitur in locō īnfinītivi et sīc dēclīnātur: accūsātīvus -ndum, genetīvus -ndī, ablātīvus (et datīvus) -ndō.

Exempla:

Ovidius, semper parātus ad amandum, librum dē amandō scrīpsit, qul appellātur 'Ars amandi.'

Magister, quī artem docendi scit, parātus est ad docendum. Magister ipse discit aliōs docendō.

Industriī estōte in scribendō, discipulī! Tempus scribendī \navnote{scrib end ō dormi dormi dormi dīcere = loquī}est. Estisne parātī ad scrībendum? 160

Iam tempus dormiendi est, sed Qulntus nōn est fessus audiendō neque parātus ad dormiendum.

Scrībere scribendō, dicendo dīcere discēs.

\booksection{Pēnsum A}

Nāvēs necessāriae sunt ad nāvig-. Iūlia dēlectātur in hortō ambul- et flōrēs carp-. Mēdus cōnsilium fugi- excōgitāvit; Lydia eum iūvit in fugi-. Paulum satis est ad beātē vīv-.

Verba: iungere -isse -um; figere ---isse Cum; mergere -isse -um; ūrere -isse Cum; movēre -isse -um; capere -isse -um; invenīre -isse -um; cōnsīdere -isse; iuvāre -isse; accidere -isse.

\booksection{Pēnsum B}

Daedalus exitum labyrinthī --- [= reperire] nōn poterat nec --- [= ūllus homō] eum in fugiendō --- poterat. --- igitur aliae viae clausae erant, vir --- [= audēns] per āera--- --- cōnstituit. --- patris fīlium dēlectāvit. Tum Daedalus ālās cōnfecit ex --- quās cērā iūnxit et in lacertīs---. ---. Postquam hoc --- perfecit, Daedalus fīlium suum --- volandī docuit: ``--- mē! Movē ālās --- --- hōc modō! Ars volandī --- [= nōn] difficilis est. Sed antequam hinc ---, hoc tē moneō: nōlī volāre in --- āere prope mare nec in --- āere prope sōlem. Iam sequere mē! --- nōstrum relinquimus, --- sumus!''

īcarus in summum caelum ascendēns nōn sōlum Graeciam, --- [= sed] etiam Asiam ac --- [= prope] tōtum --- terrārum dēspiciēbat. Tum vērō id quod pater timuerat --- [= factum est]: --- sōlis propinquī cēram --- atque pennās ---.

Ecce fabula mīrābilis dē puerō --- qul --- quaerēns mortem invēnit.

\booksection{Pēnsum C}

Quis Daedalum in labyrinthum inclūdī iussit? Quōmodo Daedalus effugere cōnstituit? Cūr hominēs volāre nōn possunt? Ex quibus rēbus Daedalus ālās cōnfecit? Cūr ālae [cārum sustinēre nōn poterant? Quid pater et filius volantēs vīdērunt? Estne Peloponnēsus magna īnsula? Quārē īcarus in summum caelum ascendit? Quid tum puerō accidit? Ubi corpus pueri inventum est? Num haec fabula dēlectandī causā modo nārrātur?

\booksection{Vocābula nova}
\noindent\small nova, fuga, cōnsilium, carcer, orbis, nātūra, ars, opus, penna, ignis, lacertus, llbertās, multitūdō, paenīnsula, cāsus, celer, reliquus, audāx, liber, studiōsus, ingēns, infimus, summus, cautus, temerārius, propinquus, persequī, cōnsequī, cōnficere, cōnsūmere, invenīre, effugere, iuvāre, excōgitāre, imitārī, ēvolāre, perficere, mollire, figere, levāre, ūrere, suspicere, dēspicere, accidere, quatere, aberrāre, revocāre, vidēri, estō, quisquam, sūrsum, deorsum, haud, paene, quidem, tamquam, quoniam, vērum, sīn, trāns, gerundium.

\bookchapter{Cap. XXVII}{Rēs rūsticae}

\navnote{mmip}\navnote{quiēscere -ēvisse = nihil agere, quiētus esse, dormīre dēnique = postrēmō dēsinere -siisse}I Quid agit pater familiās post meridiem? Prīmum quiēscit, tum ambulat, dēnique lavātur. Iūlius igitur, postquam paulum quiēvit, ambulātum exit. Iam dēsiit imber, avēs rūrsus in hortō canunt. Dominus hūc illūc in \navnote{cher ager -gri m cingere: locum cingere = circum locum esse crēscere crēvisse = māior fierl}hortō suō amoenō ambulat, deinde exit in agrōs, quī hortum cingunt.

In agrīs frūmentum

crēscit vēre et aestāte. Mēnse Augustō frūmentum metitur et ex agris vehitur. Deinde agri arantur et novum frūmentum seritur. Quī agrōs \navnote{serere \contr{} metere}arant ac frūmentum serunt et metunt, agricolae appellantur. Agricola est vir cuius negōtium est agrōs colere.

\navnote{negōtium -īn = officium, opus}Agricola arāns post arātrum ambulat. Arātrum est īnstrūmentum quō agri arantur. Arātor duōs validōs \navnote{arātor -ōris m = quī arat prae prp-r abl --- ante agere = euntem facere spargere -sisse -sum}bovēs quī arātrum trahunt prae sē agit. Quōmodo frūmentum seritur? Agricola sēmen manū spargit. Ex parvīs sēminibus quae in agrōs sparsa sunt frūmentum crēscit. Mēnse Augustō frūmentum mātūrum est. Quōstrūmentum quō agricola metit. Quō īnstrūmentō serit agricola? Quī serit nūllō īnstrūmentō ūtitur praeter manum. Quī arat arātrō ūtitur; quī metit falce ūtitur; quī serit manū suā ūtitur.

Deus agricolārum est Sāturnus, quī ōlim rēx caelī fuit, sed ā filio suō Iove ē caelō pulsus in Italiam vēnit, \navnote{p ūti ūsum esse (4- abl) regiō -ōnis f --- pars terrae Latīnī -ōrum m = quī Latium incolunt rudis -e = indoctus}ubi eam regiōnem quae Latium appellātur optimē rēxit 25 Latīnōsque, hominēs ut tunc erant rudēs ac barbarōs, agrōs colere docuit. In forō Rōmānō est templum Sāturnī.

\navnote{-ve = vel: aliās-ve = vel aliāās frūgēs = quī frūgēs fert loca -ōrum npl = regiō pāscere = herbā alere pecus -oris n = ovēs, porcī, bovēs pābulum -īn = cibus pecoris et equōrum}Ager quī multum frūmentī aliāasve frūgēs ferre potest, fertilis esse dīcitur. Italia est terra fertilis, sed multa loca 30 Italiae nōn arantur nec ullas frūgēs ferunt praeter her-bam. Iīs locis ovēs, porcī, bovēs pāscuntur, nam herba est pecoris pābulum, et facilius est pecus pāscere quam agrōs colere. Praetereā lānā ovium ūtuntur hominēs, nam ē lānā vestēs efficiuntur. Itaque pecus māiōris pre- 35 til est quam frūmentum, et quī pecus pāscit plūs pecūniae facit quam quī agrōs colit.

\navnote{cōpia -ae/: f: magna cōpia = multum in-vehere solum z terra, humus parum = nōn satis}Frūmentum minōris pretiī est, quia magna cōpia frūmenti ex Āfricā in Italiam invehitur. Solum Āfricae fer-tile est, nisi aquā caret, sed multīs locis Āfricae parum aquae est. Ergō necesse est agrōs aquā flūminum rigāre. \navnote{bis ter-ve = bis vel ter}Agricolae qul agrōs prope Nīlum flūmen colunt bis terve in annō metere possunt, Aegyptus enim terra fertilissima est, quia solum eius aquā Nili rigāātur.

Agri Iūliī, quī sub monte Albānō sitī sunt, nōn sōlum frūmentum, sed etiam vītēs ferunt. Iī agri in quibus vītēs crēscunt vīneae dīcuntur. Frūgēs vīneārum sunt ūvae, quae mēnse Septembri mātūrae sunt. Ex ūvis mātūris vīnum efficitur. II

\navnote{La Tti, e lusculum villze S UI 3 2 z Me ----. Lacus z Albānus z ZAlbands 7, oN Do A) \$ ūva di -ae f V vītis -isf praedium -īn = vīlla cum agris circā prp-- acc = circum}Iūlius, quī nūper ex urbe in praedium suum Albānum vēnit, circā agrōs et vīneās suās ambulat. Suprā eum est mōns Albānus, post montem lacus Albānus, quī villis magnificīs cingitur. Nullā in parte Italiae tot et tantae vīllae sitae sunt quam in Latiō et māximē circā lacum illum amoenum. Nē in Campāniā quidem plūrēs \navnote{= etiam nōn}vīllae sunt, quamquam multī Rōmānī in ōrā maritimā eius regiōnis villas possident; nam plūrimī Rōmānī sub \navnote{sub urbe = prope urbem}urbe Rōmā habitāre volunt in villis suburbanis.

\navnote{sub urbe situs}Iūlius aspicit agricolās quī in agris et in vīneīs opus faciunt, gaudēns quod ipsī, ut dominō divitī, nōn \navnote{laborāre = opus facere; \contr{} \contr{} quiēscere labor -ōris m = opus exīstimāre = cēnsēre prae agricolīs beātus = beātior quam agricolae}necesse est in agris labōrare. Quamquam nūllō modō labōrem agricolārum sordidum indignumve esse existimat, tamen sē prae agricolīs beātum esse cēnset. Neque enim labōrat dominus, sed quiēscit, cum in praediō suō est. Rūs quiētum et amoenum eum dēlectat. In urbe \navnote{agrī, silvae, campī... rūri (loc) in urbe ōtium -i z e negōtium ōtium) urbānus -a -um \contr{} urbs}Iūlius semper in negōtiō est, sed rūrī in ōtiō cōgitat dē negōtiīs urbānīs. Itaque Iūlius, quī ōtium rūris valdē amat, cum prīmum cōnfecta sunt negōtia urbāna, in praedium suum suburbānum proficīscitur.

\navnote{colōnus -i m \contr{} colere prō prp-abl: prō dominō z in locō dominī}Agrī Iūliī nōn ā dominō ipsō coluntur, sed ā colōnīs. 70 Colōnus est agricola quī nōn suōs, sed aliēnōs agrōs prō dominō absentī colit et mercēdem dominō solvit prō frūgibus agrōrum.

Colōnī Iūliī sunt agricolae validī quī industriē labōrant omnēsque mercēdem ad diem solvere solent. At 75 hōc annō quīdam colōnus mercēdem nōndum solvit. \navnote{quī- quae- quod-dam: quī-dam colōnus = ūnus ex colōnīs dominus imperat ut colōnus accēdat = dominus colōnum accēdere iubet}Iulius eum colōnum in agrō cōnspicit et ``Hūc accēde, colōne!'' inquit. Dominus imperat ut colōnus accēdat, tum interrogat: Cūr nōndum solvistī mercēdem quam ter quaterve iam abs tē poposcī? Octingentōs sēstertiōs \navnote{poscere poposcisse v}mihi dēbēs. Solve eōs!`` Iūlius colōnō imperat ut mercēdem solvat.

Colōnus pallidus prae metū loquī nōn potest.

\navnote{prae metū = M}Iūlius: ``Audīsne? Imperō tibi ut mercēdem solvās. Quīn respondēs?''

Colōnus: ``Nūlla est mihi pecūnia. Nē assem quidem habeō.''

nu Iūlius: ``Nisi hīc et nunc solvis mercēdem dēbitam, servīs meīs imperābō ut tē agris meīs pellant. Iam trēs \navnote{agrīs:: ex agris}mēnsēs exspectō ut ea pecūnia mihi solvātur. Etsī vir \navnote{patiēns -entis \contr{} patī patientia -ae f \contr{} patiēns prō-icere (-iicere) -iō}patiēns sum, hic finis est patientiae meae!``

Colōnus ad pedēs dominī sē prōicit eumque ōrat ut pateenūam habeat: patientiam habeat: Tatientiam

habē, domine! Nōlī ā mē postulāre ut tantum pecūniae statim solvam! Intrā \navnote{tantum -īn = tam multum}duōs trēsve mēnsēs omnia accipiēs. Nōlī mē ē domō meā rapere! Octō līberi mihi sunt, quōs ipse cūrāre \navnote{rapere -iō -uisse raptum}dēbeō. Cūra īnfantium multum temporis magnamque patientiam postulat, itaque parum temporis habeō ad opus rūsticum.``

\navnote{rūsticus -a -um \contr{} rūs}Iūlius: ``Quid? Num uxor abs tē postulat ut tu prō \navnote{neglegere -exisse -ēctum \contr{} \contr{} cūrāre}mātre īnfantēs cūrēs? Itane īnfantēs suōs neglegit? Mātris officium est īnfantēs cūrāre. Tu vērō cūrā ut agri bene colantur et mercēs ad diem solvātur!``

Colōnus: ``Uxor mea officium suum nōn neglegit nec postulat ut ego īnfantēs cūrem; sed nunc nec īnfantēs cūrāre nec quidquam aliud agere potest, quia aegrōtat: \navnote{quid-quam = ūlla rēs}intrā paucōs diēs novum Infantem exspectat. Nōlī mē ab \navnote{gravidus -a -um: (femina) gravida = quae īnfantem exspectat precēs -um f pl = verba ōranlia, quod ōrātur}uxōre gravidā rapere! Per omnēs deōs tē ōrō!``

Hīs precibus dominus sevērus tandem movētur. Colōnō imperat ut taceat atque surgat, tum ``Quoniam'' inquit ``uxor tua gravida est, abī domum! Primum cūrā ut uxor et liberi valeant, tum vērō labōrā ut pecūniam omnem solvās intrā finem huius mēnsis, id est intrā tricēsimum diem!''

\navnote{dī-mittere -misisse}Colōnō dīmissō, Iūlius alium agricolam vocat eumque dē rēbus rūsticīs rogitat, ac primum dē vīneīs: ``Quōmodo vīneae sē habent hōc annō?''

``Optimē'' inquit agricola, ``Aspice hanc vītem: tot et tantae ūvae magnam vīnī copiam prōmittunt, ac vīnum bonum futūrum esse exīstimō, nam sōl duōs iam mēn\navnote{prōd-esse prō-fuisse calor -ōris m = āēr calidus, tempus calidum nocēre -uisse \contr{} \contr{} prōdesse frigus -oris n \contr{} +calor vīnum bonum est: calor efficit ut v. bonum sit satis dē vīneīs dictum est}sēs prope cotīdiē lūcet ūsque ā māne ad vesperum. Nihil enim vīneīs magis prōdest quam sōl et calor, nec quidquam iīs magis nocet quam imber et frīgus.``

Iūlius: ``Calor sōlis nōn ipse per sē efficit ut vīnum bonum sit. Vītēs probe cūrāre oportet. Itaque vōs moneō ut industriē in vlneīs labōrētis. Sed satis dē vīneīs. Frūmentum quāle erit?''

``NŌn ita bonum'' inquit alter agricola, ``Solum nimis siccum est nec rigāri possunt agri quod procul absunt ā rīvō. Imber brevis quem hodiē habuimus frūmentō prōfuit quidem, sed parum fuit. Item sicca est herba, pecus parum pābulī invenit. Sed scīsne ovem herī paene rap-tam esse ā lupō?''

Iūlius: ``Quid? Lupusne ovem ē grege rapuit?''

Agricola: ``Ovis ipsa ē grege aberrāverat.

Nec vērō 135 \navnote{grexgregism : MD --- (7 X Y Ld ' -. Qum grex gregis m z multitūdō bēstiārum}lupus ovi nocuit, nam pāstor eam in silva repperit atque ē dentibus lupī servāvit!``

Iūlius: ``Ō \contr{} pāstōrem pigerrimum, quī officium suum ita neglēxit! Pāstōris officium est cūrāre nē ovēs aberrent nēve silvam petant. Ego vērō cūrābō nē ille pāstor posthāc officium neglegat!''

Agricola: Nōlī nimis sevērus esse! Nōn cēnseō illum pāstōrem prae cēteris pigrum esse.``

\navnote{prae cēteris piger = cēteris pigrior neglegēns -entis = quī officium neglegit}Iūlius: ``Rēctē dīcis: nam pigri ac neglegentēs sunt omnēs! At ego faciam ut industriī sint!''

Agricola: ``Certē pāstōrēs minus labōrant quam agricolae. Nōbīs nūllum est ōtium, nec opus est nōs monēre ut industriī sīmus nēve quiēscāmus.''

Iūlius: ``Nōlī cēnsēre opus pāstōrum facilius esse. Cūra pecoris magnum est negōtium, nōn ōtium, ut pāstōrēs nēquam in molli herbā dormientēs putant. Ego \navnote{nē-quam indēcl = improbus}vērō cūrābō nē ille pāstor neglegēns sit nēve dormiat! Faciam ut tergum ei doleat! Arcesse eum!``

Sed eō ipsō tempore pāstor gregem prae sē agēns ē \navnote{optimē: optimo tempore}campīs revertitur. Cum primum is prope vēnit, ``Optimē advenīs'' inquit dominus Irātus baculum prae sē tenēns, ``nam verbera meruistī!''

Pāstor humī sē prōiciēns dominum ōrat nē sē verberet: ``Nōlī mē verberāre, ere! Nihil fecī!''

``At propter hoc ipsum'' inquit Iūlius ``tē verberābō, homō nēquam, quod nihil fecistī! Officium tuum est cūrāre nē ovēs aberrent nēve ā lupō rapiantur. Precēs tibi nōn prōsunt. Prehendite eum, agricolae, et tenēte!'' \navnote{prōd-est prō-sunt}Iūlius duōbus agricolīs imperat ut pāstōrem prehendant et teneant.

Tum vērō, dum pāstor territus verbera exspectat, ovēs sine pāstōre relictae dē viā in agrōs aberrant ac frūmentum immātūrum carpere incipiunt. Agricolae \navnote{prō-hibere = retinēre mittere \contr{} \contr{} prehendere}hoc videntēs clāmant: ``Prohibē ovēs tuās ab agrīs nōstris, pāstor!'' --- tum pāstōrem mittunt atque celeriter 170 ovēs in agrīs sparsās persequuntur.

\navnote{modo = nūper}Pāstor sōlus cum dominō relictus ``Modo dīxistī'' in-quit ``meum officium esse cavēre nē ovēs aberrent. Nōlī mē officiō meō prohibēre!''

\navnote{officio (abl) prohibēre = ab officio prohibēre}Iūlius: ``Ego tē nōn prohibēbō officium facere. Fac ut ovēs ex agrīs agantur! Age, curre, pāstor!''

Vix haec dīxerat Iūlius, cum pāstor quam celerrimē \navnote{quam celerrimē potest = tam celeriter quam maāxime fieri potest}potest ad ovēs suās currit. Dominus ridēns eum currentem aspicit, tum ad vīllam revertitur. Etsī dominus sevērus exīstimātur, tamen inhūmānus nōn est.

\booksection{Grammatica Latīna}

\navnote{sevērus esse}\navnote{coniūnctīvus -īm (coni)}Coniūnctīvus Tempus praesēns. [A] Āctīvum.

\navnote{[l]intre* [3]claudar}Dominus: ``Intrā, serve! Claude forem! Tacē et audī!'' Do-minus servō imperat ut intret, forem claudat, taceat et audiat. Servus intrat, forem claudit, tacet et audit.

'Intrat', 'tacet', 'claudit', 'audit' est indicātīvus. 'Intret', ``taceat, claudar', audia? coniūnctīvus

est. Coniūnctīvus

\navnote{et -at}Exempla: [1] cōgitāre: cōgiter; [2] respondēre: respondeat; [3] scribere: scrīb; [4] audīre: audiat.

\navnote{[1] [1] cogitlem cōgit2s cogite \contr{} cōgit em}Magister: Studiōsus estō, discipule! Primum audī quod interrogō, tum cōgitā, dēnique surge et respondē!`` Magister discipulum monet ut studiōsus sit: prīmum audiat, tum cōgi- 195 tet, dēnique surgat et respondeat. Discipulus silet. Magister: * ''Audīsne, puer? Moneō tē ut studiōsus sīs: primum audias, tum cogites, dēnique surgās et respondeas.``

Discipulus: ``Nōn opus est mē monēre ut audiam et cōgitem atque ut studiōsus sim. Sed respondēre nesciō. Nōlī igitur ā mē postulāre ut surgam et respondeam/''

Magister: ``Studiōsī estōte, discipulī! Primum audīte quod interrogō, tum cōgitāte, dēnique surgite et respondēte!'' Magister discipulōs monet ut studiōsī sit: primum audiant, tum cogitent, dēnique surgant et respondeant. Discipulī silent. Magister: ``Audītisne, pueri? Moneō vōs ut studiōsī sitis: primum audiātis, tum cōgitetis, dēnique surgātis et respondeātis.'' Discipulī: ``Nōn opus est nōs monēre ut audiamus et cogitemus atque ut studiōsī sīmus. Sed respondēre nescimus. Nōlī igitur ā nōbīs postulāre ut surgāmus et respondeāmus/''

Persōna prīma -em -ēmus -am -amus Persōna secunda -es -ētis -ūs Persōna secunda -ēs Persōna tertia -et

-et [B] Passīvum.

\navnote{verber ē ris verber ē tur verber ēmur verber ē minī verber c ntur morde ā ris morde ā tur morde ā mur morde ā mim morde a ntur inclūd āris inclūd ā tur inclūd mur inclūd ntur inclūd vinci ā ris vinci ā tur vinci ā mur vinci ā minī vinci a ntur}Dominus imperat ut servus improbus leneātur et verberētur, deinde ut vinciātur et inclūdātar.

Dominus imperat ut servī improbī teneantur et verberentur, deinde ut vinciantur et inclūdanntr.

Iānitor tabellārium monet ut caveat nē ā cane mordeatur: ``Cavē nē ā cane mordeāns/'' Tabellārius: ``Tuum negōtium est cūrāre nē ego mordear. Vincī canem! Ego cūrābō ut tū vinciam et inclūdans et verberms ab erō tuō!'' Iānitor: ``Num tuum negōtium est cūrāre ut ego vinciar et inclūdar et ver-

Iānitōrēs tabellāriōs monent ut caveant nē ā cane mordeantur: ``Cavēte nē ā cane mordeommf/'' Tabellāriī: ``Vestrum negōtium est cūrāre nē nōs mordeāmur. Vincīte canem! Nōs cūrābimus ut vōs vincidmim et indūdāminī et verberēmmt ab erō vestrō!'' Iānitōrēs: ``Num vestrum negōtium est cūrāre ut nōs vinciāmur et inclūdāmur et verberemur? Persōna prima

-er Persōna secunda

-ēris

-ēminī

-ātur

-antur Persōna tertia

-ētur

-entur -ātur -antur

\booksection{Pēnsum A}

Iūlius colōnō imperat ut mercēdem solv-. Ille dominum ōrat ut patientiam habe-: ``Nōlī postulāre ut tantam pecūniam statim solv-!'' Dominus colōnō imperat ut tacē- et surg-, tum ``Primum cūrā'' inquit ``ut uxor et liberi valē-, tum vērō cūrā ut agrōs bene col- et mercēdem solv-!''

Dominus colōnōs monet ut labōr- nēve quiēsc-: Moneō vōs ut labōr- nēve quiēsc-!``

Māter filiam monet ut cauta s-: ``Moneō tē ut cauta s-!'' Fābula nōs monet nē temerāriī s-.

Verba: spargere -isse -um; rapere -isse -um; neglegere -isse -um; dēsinere -isse; quiēscere -isse; crēscere -isse; poscere -isse; prōdesse ---isse.

\booksection{Pēnsum B}

Mēnse Augustō --- metitur, deinde --- arantur et novum frūmentum ---. Agricola qul --- post arātrum ambulat duōs bovēs --- [= ante] sē agēns; arātrum est --- quō agri arantur. Agricola quī serit nūllō īnstrūmentō --- et sēmen manū ---. Qul --- falce ūtitur. Agricola est vir cuius --- est agrōs ---. Ex Aegyptō, quae terra --- est, magna --- frūmentī in Italiam ---. Frūgēs vīneārum sunt ---, ex quibus --- efficitur. --- sōlis vīneīs ---, ---, --- [ \contr{} calor] vīneīs ---.

Iūlius in --- suō Albānō nōn labōrat, sed ---. --- quiētum et --- [= pulchrum] eum dēlectat. Colōnus aliēnōs agrōs --- dominō absentī colit.

Pāstor est vir quī pecus --- et cūrat. --- pecoris magnum negōtium est, nōn ---, ut pigri pāstōrēs--- [= [- cēnsent]. Pāstōrēs nōn tam industriē --- quam agricolae. Pāstōris officium est cūrāre --- ovēs aberrent --- [= et nē] ā lupō rapiantur.

\booksection{Pēnsum C}

Quid est negōtium agricolae? Quandō frūmentum metitur? (. -. I Num arātor ipse arātrum trahit? Quid est pābulum pecoris? Unde frūmentum in Italiam invehitur? s Quae regiō Āfricae fertilissima est? Cūr necesse est agrōs rigāre? Quae sunt frūgēs vīneārum? Num Iūlius ipse in agris labōrat? Omnēsne colōnī mercēdem solvērunt? Quot sēstertiōs colōnus Iūliō dēbet? Num uxor colōnī officium suum neglegit? Quid est officium pāstōris? Estne Iūlius dominus inhūmānus?

\booksection{Vocābula nova}
\noindent\small ēmus, cōgit, ētis, ent, responde ā mus, responde ā tis, [3] surglam iīs surg a m, surga, mus, surgā, tis, [4] audijam audias audi a m, audiā, aud, ātis, ant, esse, sīmus, sijm4, sītis, sint, stt, 'ēmur, -ēris, -ēminī, -QUur, -entur, āris, -ā tur -a ntur, ager, frūmentum, agricola, negōtium, arātrum, īnstrūmentum, sēmen, falx, regiō, frūgēs, pecus, pābulum, lāna, cōpia, vītis, vīnea, ūva, vīnum, praedium, labor, Tūs, ōtium, colōnus, patientia, cūra, precēs, calor, frigus, grex, amoenus, mātūrus, rudis, fertilis, suburbānus, urbānus, patiēns, rūsticus, gravidus, siccus, neglegēns, nēquam, immātūrus, inhūmānus, tricēsimus, cingere, crēscere, metere, serere, spargere, ūti, pāscere.

\bookchapter{Cap. XXVIII}{Perīcula maris}

Interim Mēdus et Lydia ventō secundō per mare īnfe- / rum nāvigāre pergunt ad fretum Siculum (id est fretum \navnote{Siculus -a -um = Siciliae iungere) = dīvidere ē-icere (-iicere) -iō -iēcisse -iectum}angustum quō Sicilia ab Italiā disiungitur). Gaudent omnēs quī eā nāve vehuntur praeter mercātōrem cuius mercēs necesse fuit ē nāve ēicere.

Mēdus vērō multum cōgitat dē verbīs Lydiae et dē \navnote{cessāre --- opus neglegere, minus agere; ventus cessat:: v. minuitur}tempestāte quae tam subitō cessāvit, postquam Lydia dominum suum invocāvit. Ut tempestās mare tranquil-\navnote{animus -īm \contr{} corpus}lum turbāvit, ita verba Lydiae animum Mēdī turbāverunt.

Lydia amīcum suum colōrem mūtāvisse animadvertit et ``Quid pallēs?'' inquit, ``Utrum aegrōtās an territus es?''

``Nōn aegrōtō'' inquit Mēdus, ``Corpus quidem sānum est mihi, animus vērō turbātus. Quis est ille domi- 15 nus tuus cui mare et ventī oboedīre videntur?''

Lydia: ``Nōn meus tantum, sed omnium hominum est dominus, et Rōmānōrum et Graecōrum et barbarōrum.''

Mēdus: ``Utrum homō an deus est?''

Lydia: ``Chrīstus est Deī filius quī horno factus est. In \navnote{indēcl Iūdaea -ae f Iūdaei -ōrum m pi eo adv = illūc}oppidō Bethlehem nātus est in Iūdaeā, patriā Iūdaeōrum, quae inter Syriam et Aegyptum sita est. Eō vēnērunt rēgēs, quī stēllam eius vīderant in oriente, et invēnērunt puerum cum Marīā, mātre eius, et adōrāvērunt \navnote{ad-ōrāre vel-ut = tamquam, sīcut}eum velut deum. Posteā Chrīstus ipse plānē dēmōnstrāvit sē esse filium Deī, nam discipulōs docēbat, quōrum

\navnote{turba -ae/= f - multitūdō hominum}Mēdus: ``Omnis medicus id facit.''

\navnote{caecus -a -um = quī vidēre nōn potest surdus -a -um = quī audire nōn potest lo- quī nōn potest claudus -a -um = quī ambulāre nōn potest}Lydia: ``Quī medicus verbīs solis potest facere ut hominēs caecī videant, surdī audiant, mūtī loquantur, claudī ambulent?''

Mēdus: ``Potestne dominus tuus haec facere?''

Lydia: ``Profectō potest. In Iūdaeā Iēsus nōn sōlum faciēbat ut caecī vidērent, surdī audīrent, mūtī loquerentur, vērum etiam verbīs efficiēbat ut mortuī surge-rent et ambulārent. Ex ūniversā Iūdaeā hominēs aegri, \navnote{ūniversus -a -um = tōtus fāma -ae f - quod nārrātur dē aliquō}quī famam dē factīs eius mīrābilibus audiverant,

ad eum conveniēbant. Postrēmō tamen Iēsus Chrīstus ab improbīs hominibus necātus est.``

Mēdus: ``Quid? Nōn vīvit dominus tuus?''

\navnote{surgere surrēxisse}Lydia: ``Immō vērō vīvit, nam tertiō diē Iēsūs surrēxit ā mortuīs et quadrāgesimo diē post in caelum ascendit.

Immortālis est fīlius Deī sīcut pater eius, \navnote{im-mortālis -e \contr{} mornāsci nātum esse --- \contr{} mori -ior mortuum esse}Deus vīvus. Hominēs mortālēs nāscuntur ac moriuntur, Deus immortālis semper vīvit. Sed ipsa male nārrō: ex hōc libellō recitābō tibi aliquid.``

\navnote{libellus -im --- parvus liber}Lydia libellum, quem adhūc intrā vestem occultāvit, \navnote{ex-tendere ap-prehendere = prehendere (manū) quī-dam, abl quō-dam vīvere vīxisse}prōmit et Mēdō ostendit. Qul manum extendēns libel-lum apprehendit et ``Qul liber est iste?'' inquit.

Lydia: ``Scriptus est ā quōdam Iūdaeō, nōmine Matthaeō, quī simul cum Christō vīxit et discipulus eius fuit. In hōc librō Matthaeus, quī suīs oculīs auribusque dominum nostrum vīderat et audīverat, dicta et facta \navnote{dictum est, verbum memorāre = nārrāre, dīcere discere didicisse rogāre = ōrāre ē-volvere: librum ē. = librum aperīre tollere = sūmere et sēcum ferre}eius memorat.``

Mēdus, quī legere nōn didicit, Lydiae librum reddit eamque rogat ut aliquid sibi legat; quae continuō librum ēvolvit et ``Legam tibi'' inquit ``dē virō claudō cui Iēsūs imperāvit ut surgeret et tolleret lectum suum et domum ambulāret.''

Mēdus: ``Modo dīxistī \contr{} Christum etiam mortuīs imperāvisse ut surgerent et ambulārent.' Plūra de eā rē audīre cupiō.''

\navnote{princeps -ipis m = vir inter aliōs primus (quī aliīs imperat) sus-citāre = excitāre ūnus = quīdam}Lydia: Audi igitur quod scriptum est dē Iaīrō, principe quōdam Iūdaeōrum, quī Iēsum rogāvit ut fīliam 65 suam mortuam suscitāret:..

Ecce princeps ūnus accessit, nōmine Iaīrus, et adōrābat eum dīcēns: Filia mea modo mortua est, sed venī, impōne manum tuam super illam, et vivet.`` Et surgēns Iesus sequēbātur eum cum discipulis suīs. --- Et veniēns Iēsūs in do- 70 \navnote{tumultuārī= --- tumultum facere}mum prīncipis, vidēns tibicines et turbam tumultuantem, dīcēbat: ``Discēdite! Nōn enim mortua est puella, sed dor-mit.'' Et dēridebant eum. Et, eiectā turbā, intrāvit et tenuit manum eius et dīxit: ``Puella, surge!'' Et surrexit puella. Et exiitfāma haec in ūniversam terram illam.

Mēdus: ``Per deōs immortālēs! Sī hoc vērum est, princeps omnium deōrum est deus tuus; neque enim \navnote{Īnferi -ōrumm = mortuī, quī loca īnfera (sub terrā) incolere dīcuntur potestās -ātis f. habērī = exīstimāri}ūllus deus Rōmānus hominem mortālem ab īnferis suscitāre potest --- nē Iuppiter quidem tantam potestātem habet, etsī ille deus māximus habētur.``

Lydia: ``Est ut dīcis; nec sōlum deus prīnceps, sed \navnote{mundus -Im = caelum et terra et mare}Deus ūnus et sōlus est ille. Tōtus mundus in potestāte Deī est, et caelum et terra et mare.``

Hīc gubernātor, qul sermōnem eōrum exaudivit, ``Tanta'' inquit, ``unīus deī potestās nōn est. Nam trēs diī, Neptūnus, Iuppiter, Plūtō, mundum ūniversum ita \navnote{Plūtō -ōnis m: deus, rēx Inferōrum dividere -vīsisse -vīsum rēgnāre = rēx esse anima: forma hominis mortuī ut ex animā facta versāri = hūc et illūc movēri, errāre}inter sē divisērunt, ut Iuppiter rēx caelī esset, rēx maris esset Neptūnus, Plūtō autem rēgnāret apud Īnferōs, ubi animae mortuōrum velut umbrae versāri dīcuntur.``

Mēdus: ``Num quis tam stultus est ut ista vēra esse crēdat? Perge legere ē libellō tuō, Lydia!''

Lydia iterum librum ēvolvit et ``Ecce'' inquit ``quod nārrātur dē Christō super mare ambulante:

\navnote{nāvicula -ae f --- parva nāvis (ventus) contrārius \contr{} 4-*.secundus vigilia -ae f: quārta v.= nōna hōra noctis (nox in iv vigiliās dīviditur)}Nāvicula autem in mediō mari iactābātur flūctibus, erat enim ventus contrārius. Quārtā autem vigiliā noctis vēnit ad eōs Iēsūs ambulāns super mare. Discipulī autem videntēs eum super mare ambulantem turbātī sunt dīcentēs: Phan\navnote{phantasma -ātis --- anima quae ambulāre vidētur eīs = iīs cōnstāns -antis \contr{} turbātus Petrus -Im: Christl discipulus ipse:: ille}tasma est!`` et prae timōre clāmāvērunt. Statimque Iēsūs locūtus est eīs dīcēns: ''Cōnstantēs estōte! Ego sum. Nōlīte timēre!`` Respondēns autem ā Petrus dīxit: ''Domine, si tā. 100 es, tube mē ad tē venīre super aquam!``

At ipse ait: Veni! Et descendens Petrus de nāviculā ambulabat super aquam, \navnote{ut venīret = quia venīre volēbat salvus -a -um = servātus; salvum facere = servāre}ut venīret ad Iēsum. Vidēns vērō ventum validum timuit, et incipiens mergi clāmāvit dīcēns: ``Domine! Salvum

mē fac!`` Et continuō Iēsūs extendēns manum apprehendit eum, 105 \navnote{dum ascendentibus eīs - dum ascendunt}et ait illī: ``Quārē dubitāstī?'' Et ascendentibus eīs in nāviculam, cessāvit ventus. Quī autem in nāviculā erant vēnērunt et adōrāvērunt eum dīcentēs: Vere fihus Deī es.``

Gubernātor, quī Mēdum attentum videt, ``Num tū'' \navnote{attentus -a -um = quī studiōsē audit. per-suādēre -suāsisse facere ut homō crēdat}inquit ``tam stultus es ut haec crēdās? Mihi nēmō per- no suādēbit hominem super mare ambulāre posse!''

Lydia: ``Chrīstus nōn est homō, sed fīlius Deī, quī omnia facere potest. Ipse dīxit: ''Data est mihi omnis potestās in caelō et in terrā.`` Modo nōs ē tempestāte servāvit, nōnne id tibi persuāsit eum habēre potestātem 115 maris et ventōrum? Audi igitur quod in eōdem librō nārrātur dē potestāte Chrīstī:

\navnote{= et eum in nāviculam ascendentem secūtī sunt discipulī eius}Et ascendente eō in nāviculam, secti sunt eum discipulī eins. Et ecce tempestās magna facta est in marī, ita ut nāvicula operiretur flūctibus --- ipse vērō dormiēbat! Et \navnote{salvare = salvum facere per-īre -eō -iisse = pertunc = tum vanquillitās -ātis f \contr{} tempestās}accessērunt ad eum discipulī eius et suscitāvērunt eum dīcentēs: ``Domine, salvā nōs! Perīmus!'' Ait illis Iesūs: ``Quid timidi estis?'' Tunc surgēns imperāvit ventīs et marī, et facta est tranquillitās magna. Hominēs autem mirāti sunt dīcentēs: ``Quālis est hic, quod ventī et mare oboediunt eī?''

Gubernātor: ``Mare et ventī nēminī oboediunt nisi \navnote{nēmō, dat nēminī}Neptūnō. Ille cūrāvit ut nōs ē tempestāte servārēmur nēve mergerēmur---vel potius nōs ipsī quī mercēs ēiēcimus. Nōlīte vērō cēnsēre nōs iam extrā perīculum esse. Tempestās quidem dēsiit, sed multa alia pericula nōbīs impendent, ut saxa quibus nāvēs franguntur, vorāagines \navnote{impendēre (4 dat) vorāgō -inis f --- locus quō nāvēs flūctibus vorantur (: merguntur) praedō -ōnis m rapere volunt}in quās nāvēs merguntur, praedōnēs maritimī quī nāvēs persequuntur, ut mercēs et pecūniam rapiant nautāsque occīdant. Semper in periculō versāmur.``

\navnote{tūtus -a -um = sine periculō}Mēdus: ``Sed hīc tūtī sumus ā praedōnibus.''

Gubernātor: ``Nūllum mare tūtum est ā praedōnibus, \navnote{per-venīre}nē mare īnferum quidem, quamquam rāri hūc perveniunt. Nec tūtī sumus ā cēteris periculis quae modo memoraāvi. Brevī nāvigabimus per fretum Siculum, ubi ab utrāque parte magnum perīculum impendet nautīs: ab ōrā Italiae saxa periculōsa quibus Scylla nōmen est, ab \navnote{\contr{} periculum Charybdis -is f (acc -im) (periculum) vītāre= --- nōn adire, effugere Īre, coni praes: eam eāmus eās eātis eant spērāre = rem bonam futūram esse putāre}ōrā Siciliae vorāgō terribilis quae Charybdis vocātur. Multae nāvēs quae Scyllam iam vītāverant, deinde in Charybdim mersae sunt. At bonum animum habēte! Ego, ut gubernātor cōnstāns, cūrābō ut omnia pericula vītēmus ac salvī in Graeciam eāmus.``

Mēdus: ``Omnēs id futūrum esse spērāmus. Quandō eō perveniēmus?''

Gubernātor: ``Intrā sex diēs, ut spērō, vel potius octō. \navnote{īre: eundi (gerundium); cupidus es eundī = īre cupis mālō = magis volō}Sed cūr tam cupidus es in Graeciam eundī? Ego Rōmae vīvere mālō quam in Graeciā.``

\navnote{ma-vis = magis vīs esse}Mēdus: ``Num Rōmae māvīs servīre quam līber esse in Graeciā?''

Gubernātor:

uNōs cīvēs Rōmānī mori

mālumus \navnote{mālumus = magis volumus}quam servīre!``

Mēdus: ``Nōlī putāre mē servīre mālle, nam ego quōque liber nātus sum, nec quisquam quī liber fuit lībertātem spērāre dēsinit. In Italiā dominō sevērō serviēbam, quīā mē postulābat ut opus sordidum facerem nec mihi \navnote{pecūnia dominiquaeservo datur si quid = si \$1aliquid}pecūlium dabat. Si quid prāvē feceram, dominus imperābat ut ego ab aliīs servīs tenērer et verberārer. Sed herī ē vīllā fūgī, ut verbera vītārem, atque ut amīcam meam vidērem ac semper cum eā essem.Multīs promis- Sīs el persuasi ut mēcum ex Italiā proficīscerētur, Lydia \navnote{persuādēre ut = verbīs mā-vult = magis vult}enim Rōmae vīvere māvult quam in Graeciā. Ōstiā igi- 165 tur hanc nāvem cōnscendimus, ut in Graeciam nāvigārēmus.``

Gubernātor Lydiam interrogat: ``Tune quoque dominō Rōmānō serviēbās?''

Lydia: ``Minimē vērō. Ego nēminī serviō nisi dominō \navnote{in caelīs = in caelō}nostrō quī est in caelīs. Nēmō potest duōbus dominīs servīre. Certē nōn laetō animō Rōmā profecta sum, et difficile fuit mihi persuādēre ut amīcās meās Rōmānās dēsererem. Nec prōmissīs solis Mēdus mihi persuāsit ut sēcum venirem, sed etiam dōnō pulcherrimō. Ecce ānu- 175 lus aureus gemmātus quem amīcus meus prope centum sēstertiīs emit mihi.`` Lydia manum extendēns digitum ānulō aureō ōrnātum gubernātōri ostendit.

Gubernātor ānulum tam pulchrum admīrātur, tum \navnote{ad-mirāri = mirāri (rem magnificam)}conversus ad Mēdum ``Profectō'' inquit ``dīves esse vidēris, ut servus! Num dominus ille sevērus, quī tibi imperābat ut opus sordidum facerēs, tantum pecūlium tibi dabat prō opere sordidō?''

Mēdus rubēns nescit quid respondeat, et velut homō surdus mūtusque ante eōs stat.

\navnote{prōmere -mpsisse -mptum}Lydia: ``Quīn respondēs? Sacculum prōmpsistī pecūniae plēnum --- nōnne tua erat ista pecūnia?''

Mēdus turbātus, dum oculōs Lydiae vītāre cōnātur, mercātōrem celeriter accēdere videt.

\booksection{Grammatica Latīna}

Coniūnctīvus Tempus imperfectum [A] Āctīvum.

Dominus servum monet ut sibi pāreat. 195

\navnote{pārēret}Dominus servum monēbat monuit ut sibi pārēret.

Pareat' est coniūnctīvus praesentis. Pārēret' coniūnctīvus imperfectī est. Coniūnctīvus imperfectī (pers. iii sing.) -ret.

Exempla: [1] recitāre: recitaret; [2] tacēre: tacēm; [3] scrībere: scriberet; [4] audīre: audiret.

Magister discipulum monuit ut tacēre: et audiret et studiōsus esset. Tum ei imperāvit ut scriberet et recitāret.

Pater: ``Nōnne magister tibi imperāvit, fili, ut scriberes et \navnote{tacc rē mus tacē rē tis}recitāres?`` Fīlius: ''Primum mē monuit ut tacērem et audirem et studiōsus essem, tum mihi imperāvit ut scriberem et recitārem.`` Pater: ''Num necesse erat tē monēre ut tacēres et audirēs \navnote{scrib erē s scrīb tremus scrib crēltis scrīb erent}et studiōsus essēs?``

Magister discipulōs monuit ut lacērent et audīrent et studiOsI essent. Tum imperaāvit ut scriberent et recitārent.

Pater: ``Nōnne magister vōbīs imperāvit, filii, ut scrīberētis et recitāarētis?'' Filii: ``Primum nōs monuit ut taceremus et 210 \navnote{servareris}audirēmus et studiōsī essēmus, tum nōbīs imperāvit ut scrīberēmus et recitāremus.`` Pater: Num necesse erat vōs monēre ut tacērēm et audīr\&is et studiōsī essētis?''

Persōna secunda

-ret Persōna tertia -ret -rent [B] Passīvum.

Dominus imperāvit ut servus tenērētur et verberāretur, deinde ut vincīrētur et inclūderēetur.

Dominus imperāvit ut servī tenērentur et verberāarentur, de-inde ut vincirentur et inclūderentur.

\navnote{servā rē ris servā re tur servā remur servā revminī servā rentur [3] Vocābula nova:}Mēdus: ``Salvus sum. Neptūnus cūrāvit ut ego ē perīculō servārer nēve in mare mergerer'' Lydia: Nōlī putāre Neptūnum cūrāvisse ut tū servārms nēve mergereris. Nēmō nisi 225 Chrīstus cūrāvit ut nōs ē periculō servārēmur nēve in mare mergerēmur.`` Gubernātor: ''Ego bene gubernando cūrāvi ut vōs servāremini nēve mergeremini!``

Plūrālis Persōna prīma

-rer Persōna secunda ---

-rēris Persōna tertia

-rētur

\booksection{Pēnsum A}

Servus dominum ōrābat nē sē verberā-, sed dominus imperāvit ut tacē- et surg-, tum aliīs servīs imperāvit ut eum prehend- et tenē-.

Mēdus ā dominō fūgit, ut amīcam suam vidē- et semper cum eā es-. Mēdus: *À dominō fūgi, ut amīcam meam vidēet semper cum eā es-.``

Mīnōs imperāvit ut Daedalus et īcarus in labyrinthum inclūd-. īcarus: ``Quis imperāvit ut nōs inclūd-?'' Daedalus: Mīnōs imperāvit ut ego inclūd- et ut tū mēcum inclūd-.''

Verba: dīvidere -isse -um; ēicere -isse -um; prōmere! -isse -um; vīvere ---isse; discere -isse; persuādēre -isse; survorāgō gere -isse.

\booksection{Pēnsum B}

Lydia Rōmae vīvere--- --- quam in Graeciā, sed Mēdus multīs mūtus prōmissīs eI --- ut sēcum proficīscerētur. In Italiā Mēdus ūniversus dominō sevērō ---. I

Nāvis nōndum extrā --- est. Multa pericula nautīs---. ---. --- --- maritimī nāvēs persequuntur; nūllum mare --- est ā praedō;nibus.

Hominēs mortālēs nāscuntur et ---, diī vērō --- sunt. Nūl;lus deus Rōmānus hominem mortuum ab īnferīs--- --- [= excitāre] potest, nē Iuppiter quidem tantam --- habet, etsī ille deus māximus --- [= exīstimātur]. Trēs diī ūniversum --- inter sē divīsērunt. T

Chrīstus in oppidō Bethlehem --- est. Ille verbīs solis efficiēbat ut hominēs --- vidērent, --- audīrent, --- loquerentur, --- --- ambulārent.: I

\booksection{Pēnsum C}

Fretum Siculum quid est? Ubi nātus est Christus? I Quid Iaīrus Chrīstum rogāvit? ) Quae pericula nautīs impendent? Cūr nautae praedōnēs maritimōs metuunt? Quali dominō Mēdus in Italiā serviēbat? I Quārē Mēdus ā dominō suō fūgit? Num Lydia laetō animō Rōmā profecta est? s Quōmodo Mēdus ei persuāsit ut sēcum venīret? Cūr Mēdus nescit quid respondeat? I Quem Mēdus accēdere videt? Tune in Graeciā vīvere māvīs quam in patriā tuā?

\booksection{Vocābula nova}
\noindent\small mālle = magis velle, mālō, mālumus, māvīs, māuultis, mālunt, recitā, fretum!, animus, turba, fama, libellus, dictum, princeps, tībīcen, potestās, mundus, nāvicula, vigilia, phantasma, tranquillitās, periculum, praedō, pecūlium, mortālis, immortālis, cōnstāns, salvus, anentus, tūtus, periculōsus, quadrāgēsimus, disiungere, ēicere, cessāre, oboedlre, adōrāre, nāscī, mori, extendere, apprehendere, memorāre, rogāre, ēvolvere, suscitāre, tumultuāri, habērī, rēgnārc, versāri, persuādēre, salvāre, perīre, impendēre, pervenlre, vītāre, spērāre, servīre, mālle, admīrārī, poiius, utrum, velut.

\bookchapter{Cap. XXIX}{Nāvigāre necesse est}

x boc qq V VĒNĒ 0 NĀVIGĀRE NECESSE EST

Multae nāvēs multīque nautae quotannīs in mari pereunt. In fundō maris plūrimae nāvēs mersae iacent. Nec tamen ūllīs periculīs ā nāvigandō dēterrentur nautae. ``Nāvigāre necesse est'' āiunt, et mercātōrēs, quī ipsī \navnote{ad-icere (-iicere) -iō aestimāre: magnī ae. = magnī pretiī esse cēnsēre vīvere re-manēre = manēre nōn-nūllī -ae -a = haud paucī, complūrēs magnī pretiī}pericula maris adire nōn audent, haec adiciunt: ``Vīvere 5 nōn est necesse!'' Mercātōrēs mercēs suās magnī aestimant, vītam nautārum parvī aestimant!

Nec vērō omnēs mercātōrēs domī remanent, cum mercēs eōrum nāvibus vehuntur. Nōnnūllī in terrās aliēnās nāvigant, quod mercēs pretiōsās nautīs crēdere nōlunt, sīcut mercātor ille Rōmānus qul eādem nāve vehitur quā Mēdus et Lydia. Is laetus Ōstiā profectus est cum mercibus pretiōsīs quās omnī pecūniā suā in Italiā ēmerat eō cōnsiliō ut eās māiōre pretiō in Graeciā vēnderet, Ita spērābat sē magnum lucrum factūrum esse. 15 \navnote{pecūniam suam augēre spēs -el f --- id quod spērātur}Iam vērō ea spēs omnis periit, nam flūctibus mersae sunt mercēs in quibus omnem spem posuerat. Subitō \navnote{pōnere posuisse positum}\navnote{quod mīrantem facit rēns, tristis}mercātor ē dīvitissimō pauperrimus factus est. Nōn mīrum est eum maestum esse.

Mercātor ad gubernātōrem accēdēns multīs cum la\navnote{heu! (dolōrem animī significat)}crimis queritur: ``Heu, mē miserum! Omnia quae possidēbam in fundō maris sunt. Quid faciam? Quid spērem? Quōmodo uxōrem et līberōs alam? Nē as quidem \navnote{ā-mittere -mīsisse -missum \contr{} \contr{} accipere}mihi reliquus est: omnia āmisi. Heu!`` 25

``Dēsine querī!'' inquit gubernātor, ``Nōn enim omnia amisisti sI uxor et līberi tuī salvī sunt. Nōnne līberōs \navnote{plūris aestimāre = māiōris pretiī esse cēnsēre dīvitiae -ārum f ( \contr{} dīves) = magna pecūnia}plūris aestimās quam mercēs istās? Dīvitiās āmittere mi-serum est, at multō miserius līberōs āmittere.``

Talibus verbīs nauta mercātōrem maestum cōnsōlāri \navnote{prōtinus = statim}cōnātur, sed frūstrā, nam ille prōtinus ``Nōlī tū mē cōnsōlārr

inquit, quī ipse imperāvistī ut mercēs meae iacerentur!``

Gubernātor: ``Iactūrā mercium nāvis servāta est.''

\navnote{iactūra -ae f \contr{} iacere}Mercātor: ``Rēctē dīcis: meae mercēs ēiectae sunt, ut nāvis tua salva esset!''

Gubernātor: ``Mercēs iēcimus ut nōs omnēs salvī essēmus. Iactūrā mercium nōn modo nāvis, sed etiam vīta \navnote{omnēs nōs, gen omnium nostrum nēmoō nostrum = nēmō nōbīs-cum = cum nōbīs felix -īcis= --- cui rēs bona accidit nēmō vestrum = nēmō}omnium nostrum servāta est. Mercēs quidem periērunt, sed nēmō nostrum periit. Ergō bonum animum habē! Laetāre ūnā nōbīscum tē vītam nōn āmīsisse simul cum mercibus! Hominēs felīcēs sumus.``

Mercātor:

``Vōs quidem

felīcēs estis, nēmō enim vestrum assem āmīsit. Mīrum nōn est vōs laetāri. At nōlīte mē monēre ut laetus sim, postquam omnia mihi \navnote{ē-ripere -iō -ripuisse -reptum \contr{} ē + rapere laetitia -ae f \contr{} laetus af-ficere -iō -fecisse -fectum tristitia -ae f \contr{} tristis; t.ā afficere = tristem facere; t.ā affici --- tristis fieri nāvigātiō -ōnis f \contr{} nāvigāre}ēripuistis! Laetitia vestra mē nōn afficit.``

Gubernātor: ``Nec quisquam nostrum tristitiā tuā afficitur. Semper gaudeō cum dē līberis meīs cōgitō, quī māximā laetitiā afficiuntur cum patrem suum ē nāvigātone periculōsā salvum redīre vident.''

Mercātor: ``Ego quoque līberōs meōs amō nec eōs 50 \navnote{dolor (animī) \contr{} \contr{} laetitia dolōre afficere = dolentem (: maestum) facere}dolōre afficere volō. Sed quōmodo vīvāmus sine pecūniā? Quōmodo cibum et vestem emam Infantibus meīs? Ō dī immortālēs! Reddite mihi mercēs!``

Gubernātor: ``Quid iuvat deōs precāri ut rēs āmissae \navnote{precāri ( \contr{} precēs)}tibi reddantur? Frūstrā hoc precāris.``

Mercātor: ``Quid ergō faciam? Ipse dē nāve saliam, \navnote{vōbīs-cum = cum vōbīs per-turbāre = valdē turbāre delphinus -īm a. Ariōn -ōnis m fidēs -ijum fpi}an in eādem nāve maneam vōbīscum?`` Vir ita perturbātus est ut sē interroget, utrum in mare saliat an in nāve remaneat.

``Sali modo!'' inquit gubernātor, ``Nēmō nostrum tē prohibēbit. At certē nōn tam felix eris quam Ariōn, quī delphlnō servātus est.''

\navnote{nōtus -a -um: is mihi nōtus est = eum nōvi ignārus -a -um = ignōrāns, rudis nōbilis -e = multīs nōtus Lesbos 1f Orpheus -17m}Mercātor, quī Arionem ignōrat, ``Quisnam est Ariōn?'' inquit, ``Nē nōmen quidem mihi nōtum est.''

Gubernātor: ``Ignārus quidem es sī illum ignōrās. 65 Ariōn, vir nōbilis Lesbi nātus, tam pulchrē fidibus canēbat ut alter Orpheus appellārētur. An tam ignārus es ut etiam Orpheus tibi ignōtus sit?''

Mercātor: ``Minimē vērō. Orpheus quidem omnibus \navnote{fidicen -inis m qul fidibus canit}nōtus est. Is fidicen nōbilissimus fuit quī tam pulchrē 70 \navnote{oblītae:: cum oblivīscerentur}canēbat ut bēstiae ferae, nātūram suam oblītae, accēderent, ut eum canentem audīrent, ac rapidī fluviī consisterent, nē strepitū cantum eius turbārent. Orpheus \navnote{nōlēbant) cantus -ūs 7 \contr{} canere}etiam ad īnferōs dēscendit ut uxōrem suam mortuam

\navnote{inde = illinc re-dūcere cum videat = vidēns, quia videt}Gubernātor, cum omnēs attentōs videat, hanc fabulam nārrat:

``Cum Arīōn, nōbilissimus suī temporis fidicen, ex \navnote{cum... nāvigāret = dum nāvigat}Italiā in Graeciam nāvigāret magnāsque dīvitiās sēcum habēret, nautae pauperēs, quī hominī diviti invidebant, \navnote{imīcus esse ob bonum aliēnum}eum necāre cōnstituērunt. Ille vērō, cōnsiliō eōrum cognitō, pecūniam cēteraque sua nautīs dedit, hoc sōlum ōrāns ut sibi ipsī parcerent. ``Ecce'' inquit ``omnia quae \navnote{necāre, salvum esse sinere per-mittere (eī)= --- sinere per-movēre: permōtī sunt:: animī eōrum}possideō iam vestra sunt. Dīvitiās meās habēte, parcite vītae! Permittite mihi in patriam revertī! Hoc sōlum precor.`` Nautae precibus eius ita permōtī sunt ut manūs quidem ab eō abstinērent; sed tamen imperāvērunt ut statim in mare dēsilīret! Ibi homō territus, cum iam \navnote{abs -- tenēre de-silire -uisse \contr{} dē dē-spērāre \contr{} spērāre carmen -inis n = verba quae canuntur}vītam dēspērāret, id ūnum ōrāvit ut sibi licēret vestem ōrnātam induere et fidēs capere et ante mortem carmen canere. Id nautae, studiōsī cantum eius audiendi, eī permīsērunt. Ille igitur, pulchrē vestītus et ōrnātus, in \navnote{celsus -a -um = altus canere cecinisse al-licere -iō -lēxisse -lectum}celsā puppī stāns carmen clārā vōce ad fidēs cecinit. Ut Orpheus cantū suō ferās ad sē alliciēbat, ita tunc Arīōn canendō piscēs allēxit ad nāvem. Postrēmō autem cum fidibus ōrnāmentīsque, sīcut stābat canēbatque, in mare dēsiluit.

*' ``Tum vērō nova et mīra rēs accidit: delphīnus, cantū \navnote{repente = subitō sub-īre eum = sub eum īre dorsum -In = tergum bēstiae vehere vēxisse vectum ex-pōnere Periander -dri m parum:: haud crēdere -didisse quasi = tamquam fallāx -ācis = quī fallit}allectus, repente hominem natantem subiit eumque in dorsō suō sedentem vēxit et in lītore Graeciae salvum exposuit. Inde Ariōn prōtinus Corinthum petīvit, ubi rēgem Periandrum, amīcum suum, adiit eīque rem sīcut acciderat nārrāvit. Rēx haec parum crēdidit, et Arīonem quasi virum fallācem cūstōdīri iussit. Sed postquam nautae Corinthum vēnērunt, rēx eōs interrogāvit 105 'num scirent ubi esset Ariōn et quid faceret?' Respon\navnote{inde:: ex Italiā}dērunt hominem, cum inde abīrent, in terrā Italiā fuisse eumque illīc bene vīvere, aurēs animōsque hominum cantu suō dēlectāre atque magnum lucrum facere/ Cum haec falsa nārrārent, Ariōn repente cum fidibus 110 ōrnāmentīsque cum quibus sē in mare iēcerat appāruit. \navnote{apparēre -uisse = in cōnspectum venīre stupēre = valdē mīrārī maleficium -īn = mālum factum cōn-fitēri -fessum esse}Nautae stupentēs, cum eum quem mersum esse putābant ita vīvum appārēre vidērent, prōtinus maleficium suum confessi sunt.``

Hic Mēdus Etsī nōtum est'' inquit ``nōnnūllōs homi- 115 \navnote{dubitō num fabula vēra sit = nōn crēdō fabulam vēram esse}nēs delphīnīs vectōs esse, tamen dubitō num haec fabula vēra sit.``

Gubernātor: Sīve vēra sīve falsa est, valdē mē dēlectat fabula dē felici salūte Arionis, nam sīcut ille mīrum \navnote{si-ve salūs -ūtis f \contr{} salvus in modum = modō (al) dē salūte dēspērāre nōn-numquam = haud rārō, satis saepe bonō animō (abl) esse = bonum animum habēre. anima:: vīta}in modum servātus est, cum iam dē salūte dēspērāret, 120 ita hominēs nōnnumquam contrā spem ē māximisperi- culis ēripiuntur. Hāc fabulā monēmur ut semper bonō animō sīmus nēve umquam dē salūte dēspērēmus. Dum anima est, spēs est.``

Haec verba tandem mercātōrem perturbātum aliquid cōnsōlāri videntur. HI

Tum vērō Lydia ad Mēdum versa ``Modo tē interrogāvī'' inquit ``tuane esset pecūnia quā hunc ānulum ēmistī. Cūr nōndum mihi respondistī?''

Ita repente interrogātus Mēdus 'sē pecūniam ē sacculō dominī surripuisse' confitetur.

\navnote{sur-ripere -iō -ripuisse fur fūris m nēquam, comp nēquior, sup nēquissimus}``Ō Mēde!'' exclāmat Lydia, ``Fūr es! Iam mē pudet tē, fūrem nēquissimum, amāvisse!''

At Mēdus ``Nōlī'' inquit ``mē fūrem appellāre, mea Lydia! Dominus enim aliquid pecūlii mihi dēbēbat. Pecūlium dēbitum sūmere fūrtum nōn est.''

\navnote{fūrtum -īn = maleficium furis accūsāre + gen: fūrti accūsāre --- dē fūrto a.}Sed Lydia pergit eum fūrti accūsāre: ``Fūrtum fecistī, Mēde! Frūstrā tē excūsāre cōnāris.''

Mēdus: ``Sī fūrtum fecī, tuā causā id fecī. Eō enim cōnsiliō nummōs surripuī ut dōnum pretiōsum tibi \navnote{bene-ficium -īn++ maleficium}emerem. Nōnne hoc beneficium potius quam maleficium esse tibi vidētur?``

Lydia:

``Facile est aliēnā pecūniā dōna pretiōsa emere. Tale dōnum mē nōn dēlectat. Hunc ānulum iam \navnote{ab-icere (-iicere) -iō -jēcisse -iectum \contr{} ab dē-trahere}gerere nōlō: in mare eum abiciam!`` Hoc dīcēns Lydia ānulum dē digitō dētrahit, sed gubernātor prōtinus bracchium eius prehendit.

Simul Mēdus ānulum ē manū Lydiae lāpsum capit.

\navnote{lābi lāpsum esse}Lydia irāta exclāmat: ``Abstinē manum, nauta!'' at ille ``Nōlī stultē agere!'' inquit,

``Nēmō tibi ānulum ita abiectum reddet --- nisi forte tam felix eris quam Polycratēs, tyrannus Samius, cuius ānulus, quem ipse in \navnote{sevērissimus}mare abiēcerat, mīrum in modum inventus est nōn in fundō maris, sed in ventre piscis!``

Lydia: ``Cūr ille tyrannus ānulum suum abiēcit?''

Gubernātor: ``Ānulum

abiēcit, cum sēsē nimis felīcem esse cēnsēret. Nihil malī umquam ei acciderat ac tanta erat potestās eius, tanta glōria tantaeque dīvitiae, ut nōn sōlum aliī tyrannī, sed etiam dī immortālēs ei \navnote{felīcitās -āātis f \contr{} felix suādēre -sisse (+ dat) = persuādēre cōnārī iactūram facere reī= --- rem abicere/āmittere invidia -ae f \contr{} invidēre ā-vertere:: prohibēre}invidērent. Tum amīcus eius, rēx Aegyptī, cum felīcitā- 160 tem atque glōriam eius ingentem vidēret, tyrannō suāsit, ut iactūram faceret eius reī quā māximē omnium dēlectābātur: ita deōrum invidiam āvertī posse spērābat. Polycratēs igitur nāvem cōnscendit et ānulum quem pretiōsissimum habēbat in mare abiēcit.

``Paucīs post diēbus aliquī piscātor in eōdem mari \navnote{piscātor -ōris m = quī piscēs capit}piscem cēpit quī tam formōsus erat ut piscātor eum nōn \navnote{dōnāre( \contr{} dōnum) dare secāre -uisse sectum = cultrō dīvidere}vēnderet, sed tyrannō dōnāret. Vērum antequam piscis ad mēnsam tyrannī allātus est, servus quī piscem secābat eī ānulum attulit 'quem in ventre piscis inventum \navnote{re-cognōscere}esse' dīxit. Polycratēs, cum ānulum suum recognōsceret, māximā laetitiā affectus est.``

Mēdus: ``Nēmō umquam eō tyrannō beātior fuit!''

Gubernātor: ``Nōlī quemquam ante mortem beātum \navnote{quis-quam, acc quemquam = ūllum hominem fortūna -ae/: f. f. hominis = quod forte homini accidit}dīcere! Hoc nōs docet fortūna illīus tyrannī. Polycratēs enim paulō post ā quōdam virō fallācī, qul eum falsīs prōmissīs Samō in Asiam allēxerat, terribilem in mo-dum necātus est. Ita nōnnumquam vita beāta morte miserrimā fīnītur. Varia quidem est hominum fortūna, \navnote{= finem facere reī varius = qul mūtātur aequus animus = animus cōnstāns ferre = patī}sed homō prūdēns bonam et malam fortūnam aequō animō fert nec alterius fortūnae invidet.``

Dum gubernātor loquitur, altera nāvis procul in marī appāret. Mēdus eum apprehendit et ``Dēsine loquī!'' inquit, ``Cūrā negōtium tuum! Quīn prospicis? Vidēsne \navnote{vēlōx -ōcis = celer ap-propinquāre (+ dat) = prope venīre}nāvem illam vēlōcem quae ā septentriōnibus nōbīs appropinquat?``

``Per deōs immortālēs!'' inquit gubernātor, cum primum nāvem appropinquantem prōspexit, ``Illa nāvis vēlōx nōs persequitur. Certē nāvis praedōnum est. Omnia vēla date, nautae!''

Nāvis autem vēlīs sōlīs nōn tam vēlōciter vehitur quam ante tempestātem, nam vēla ventō rapidō scissa sunt. Itaque gubernātor imperat ut nāvis rēmīs agātur. \navnote{rēmus}Mox rēmīs vēlisque vehitur nāvis quam vēlōcissimē po-test, sed tamen altera nāvis, cuius rēmi quasi ālae ingentēs sūrsum deorsum moventur, magis magisque appropinquat.

Gubernātor perterritus exclāmat: ``Ō di bonī! Quid faciāmus? Brevī praedōnēs hīc erunt.''

Tum mercātor, cum gubernātōrem pallidum videat, ``Bonō animō es!'' inquit, ``Nōlī dēspērāre! Spēs est, dum anima est.''

\booksection{Grammatica Latīna}

Ut', 'nē* cum coniūnctīvō [A] Tempus praesēns.

Iūlius Dāvō imperat ut puerum excitet. Aemilia Syram mo-net nē puellam excitet.

\navnote{= D. clāmat quia puerum excitāre vult = S. tacet quia puellam excitāre nōn vult}Dāvus clāmat, ut puerum excitet. Syra tacet, nē puellam excitet.

Dāvus ita (tam clārē, tantā vōce) clāmat ut puerum excitet. Syra tam quiēta est ut puellam nōn excitet. [B] Tempus praeteritum.

Iūlius Dāvō imperāvit ut puerum excitāret. Aemilia Syram monuit nē puellam excitaret.

Dāvus clāmāvit, ut puerum excitāret. Syra tacēbat, nē puel- 215 lam excitāret.

Dāvus ita clāmāvit ut puerum excitāret. Syra tam quiēta erat ut puellam nōn excitāret.

\booksection{Pēnsum A}

Magister puerōs monet --- pulchrē scrib-. Sextus tam pulchrē scribit --- magister eum laud-. Magister ipse calamum sūmit, --- litterās scrib-.

Daedalus ālās cōnfecit --- ē labyrinthō ēvol-. Icarus tam altē volāvit --- soli appropinqu-, quamquam pater eum monuerat --- temerārius es-.

Herī Quīntus arborem ascendit, --- nīdum quaer-, etsī pater eum monuerat --- cautus es-. Medicus Quīntō imperāvit --- oculōs claud-, --- cultrum medicī vid-. Quīntus tam pallidus erat --- Syra eum mortuum esse put-.

Sōl ita lūcēbat --- pāstor umbram pet-, --- in sōle ambul-.

Tantus atque tālis deus est Iuppiter --- Optimus Māximus appell---.

Verba: vehere -isse -um; pōnere -isse -um; āmittere -isse -um; allicere ---isse Cum; ēripere ---isse -um; secāre -isse -um; suādēre -isse; dēsilīre -īsse; canere -isse; crēdere -isse; cōnfitēri -um esse; lābi -um esse.

\booksection{Pēnsum B}

Orpheus, fidicen ---, tam pulchrē canēbat ut ferae --- [=; prope venīrent] ac fluviī--- --- cōnsisterent. Etiam ad --- dēscendit, ut --- [= illinc] uxōrem suam ---. Nēmō tam --- est ut i Orpheum ignōret.

Ariōn quoque omnibus --- est. Cum ille magnās --- sēcum celsus in nāve habēret, nautae pauperēs homini dīvitī--- --- eumque necāre cōnstituērunt. Arīōn, cum --- suam in perīculō esse vēlōx sentīret, pecūniam nautīs --- [= dedit] eōsque ōrāvit ut sibi adicere ---. Precibus --- nautae ei permīsērunt ut ante mortem --- --- \navnote{remanēre āmittere}caneret. Hōc factō, Ariōn in mare ---; sed delphinus eum in queri --- sedentem ad lītus vēxit. Ita ille servātus est, cum iam ēripere salūtem ---. Nautae, cum Arionem --- [= in cōnspectum veafficere \navnote{precārī}nīre] vidērent, --- [= statim] --- suum confessi

\booksection{Pēnsum C}

Quōmodo mercātōrēs lucrum faciunt? Cūr mercātor Rōmānus tristis est? Quārē mercēs ēiectae sunt? Quid mercātor deōs precātur? Quārē ad īnferōs dēscendit Orpheus? Num nautae Arionem gladiīs Quōmodo Ariōn servātus est? Quid nōs monet haec fabula? Cūr Polycratēs ānulum suum abiēcit? Ubi ānulus eius inventus est?

\booksection{Vocābula nova}
\noindent\small nova, fundus, vīta, lucrum, spēs, dīvitiae, iactūra, laetitia, tristitia, nāvigātiō, delphīnus, fidēs, fidicen, cantus, carmen, dorsum, maleficium, salūs, fur, fūrtum, beneficium, tyrannus, felīcitās, invidia, piscātor, fortūna, rēmus, pretiōsus, mīrus, maestus, fēlīx, nōtus, ignārus, nōbilis, ignōtus, rapidus, fallāx, dēterrēre, parcere, permittere, permovēre.

\bookchapter{Cap. XXX}{Convīvium}

\navnote{LECTVS SVMMVS}Ex agris reversus Iūlius continuō balneum petit, atque \navnote{H 3 B i] MENSAKS z VV z z LECTVsIMvs B triclīnium dn revertī -tisse/-sum esse balneum -īn = z locus ubi corpus lavātur}primum aquā calidā, tum frigidā lavātur. Dum ille post balneum vestem novam induit, Cornēlius et Orontēs, amīcī et hospitēs eius, cum uxōribus Fabiā et Paulā \navnote{hospes -itis m}adveniunt. (Hospitēs sunt amīcī quōrum alter alterum \navnote{re-cipere -io -10= accipere, admittere = nōn exspectāātus}semper bene recipit domum suam, etiam si inexspectātus venit.)

Hodiē autem hospitēs Iūliī exspectātī veniunt, nam Iūlius eōs vocāvit ad cēnam. (Cēna est cibus quem Rōmānī circiter hōrā nōnā vel decimā sūmunt.)

\navnote{circiter adv \contr{} circum: c. IX = plūs minus ix}Aemilia ātrium intrāns hospitēs salūtat et maritum suum tardum excūsat: ``Iūlius tardē ex agris revertit \navnote{tardus -a -um = quī ad * tempus nōn venit diū = per longum tempus induere -uisse -ūtum; indūtus = vestitus}hodiē, quod nimis diū ambulāvit. Ergō nōndum exiit ē balneō. Sed brevī lautus erit.``

Tum Iūlius lautus et novā veste indūtus intrat et amī- 15 \navnote{salvēre iubēre = salūtāre}cōs salvēre iubet: ``Salvēte, amīcī! Gaudeō vōs omnēs iam adesse. Quam ob rem tam rārō tē videō, mi Cornēlī?''

\navnote{vīsere = vīsum īre}Cornēlius: ``Nōnnumquam tē vīsere volui, nec prius urbem relinquere potuī prae multīs et magnīs negōtiīs \navnote{dēmum = dēnique, tandem \contr{} Tusculum re-quiēscere}meīs. Nunc dēmum, postquam heri ad vīllam Tusculānam rediī, paulum requiēscere possum et amīcōs vīsere. Post tanta negōtia magis quam umquam ōtiō fruor.``

Iūlius: ``Tune quoque Rōmā venīs, Orontēs?''

Orontēs: ``Nuper longum iter fecī in Graeciam. īdibus Māiis dēmum ex itinere Rōmam revertī, unde hodiē veniō.''

Iūlius: ``Ergō vōs mihi aliquid dē rēbus urbānīs novissimīsnūntiābitis.''

\navnote{= (verba) afferre}Cornēlius: ``Et tū nōs docēbis dē rēbus rūsticīs, ut agricola studiōsus et dīligēns.''

\navnote{dīligēns -entis \contr{} \contr{} neglegēns con-trahere minus esse, imperāre -ēns -entis, adv -enter ad dīligenter; prūdēns, ado prūdenter}Iūlius frontem contrahit et ``Agricola'' inquit ``ipse nōn sum, sed multīs agricolīs praesum ac dīligenter cūrō ut colōnī agrōs meōs bene colant.''

Orontēs, quī vitā rūsticā nōn fruitur, ``Prūdenter facis'' inquit ``quod agrōs ipse nōn colis. Si necesse est in agrīs labōrāre, vīta rūstica nōn iūcunda, sed molesta \navnote{iūcundus -a -um = quī dēlectat}est. Ego numquam īnstrūmentō rūsticō ūsus sum.`` Il

\navnote{jūcundus ūti ūsum esse cēnāre = cēnam sūmere}lūlius: ``Dē rēbus rūsticīs et urbānīs colloquēmur inter cēnam. Prīmum omnium cēnābimus. Sex hōrae iam sunt cum cibum nōn sūmpsī. Venter mihi contrahitur propter famem.''

Cornēlius: ``Sex hōrae nihil est. Homō sex diēs cibō carēre potest, nec tamen fame moritur.''

Iūlius: ``Dubitō num ego tam diū famem ferre pos- 45 \navnote{pos-sim carēre -uisse}sim. Sex hōrās cibō caruisse iam molestum est. Magnum mālum est famēs.``

\navnote{acc -im, abl ā paulis-per e \contr{} diū}Cornēlius: Id nōn negō, sed multō molestior est sitis. Sine cibō diū vīvere possumus, sine aquā paulisper tantum.`` 50

Orontēs ridēns ``Equidem'' inquit ``sine aquā iucundē \navnote{equidem = ego quidem}vīvere possum, sine vīnō nōn item! Magnum bonum est \navnote{= bona rēs}vīnum.``

Cornēlius: ``Nēmō negat vīnum aquā iūcundius esse, sed tamen aquam bibere mālō quam sitim patī. Num tu 55 sitim perferre māvīs quam aquam bibere?''

\navnote{per-ferre = ferre (ūsque ad finem), diū patī ē-ligere -lēgisse -lēctum PN S. a---- ---À fi 1 T. cocus; vamsriti:.}Orontēs: ``Melius quidem est aquam bibere quam sitī perire. Ex malīs minimum ēligere oportet. Nec vērō iūcundē vīvō nisi cotīdiē bonō vīnō fruor. Vīnum vīta est.''

Cornēlius: ``Nōn vīvimus, ut bibāmus, sed bibimus, ut vīvāmus.''

Hīc Aemilia ``Necesse est'' inquit ``paulisper famem \navnote{coquere coxisse coctum ex-ōrnāre = ōrnāre}et sitim ferre, dum cibus coquitur et triclīnium exōrnātur.``

\navnote{parāre = parātum facere culina -ae f - locus ubi cibus coquitur}(Servus cuius negōtium est cibum coquere atque cēnam parāre in culīnā, cocus appellātur. Aliī servī, ministrī quī vocantur, cibum parātum ē cūlīnā in triclīnium portant. In triclīniō sunt trēs lecti, lectus summus, me\navnote{īmus -a -um = īnfimus medium -In = medius locus convīvium -īn = cēna quae amīcīs datur sternerestrāvisse stratum convīva -ae mlf--- qul/ quae in convivio adest ac-cubāre = ad mēnsam cubāre}nium flōribus exōrnātur et vestis pretiōsa super lectōs sternitur. Neque enim sedentēs cēnant Rōmānī, sed in lectīs cubantēs. Quot convīvae in singulīs lectīs accubant? In singulīs lectīs aut singulī aut bīnī aut ternī convīvae accubāre solent. Cum igitur paucissimi sunt convīvae, nōn pauciōrēs sunt quam trēs, cum plūrimī, nōn plūrēs quam novem --- nam ter ternī sunt novem.)

\navnote{pridem = multō ante tardus --- \contr{} celer}Iūlius: ``Hōra decima est. Cēnam iam prīdem parātam esse oportuit! Nimis tardus est iste cocus!''

Aemilia: ``Tuumne hoc negōtium est an meum? Uter nostrum in culina praeest? Nōndum hōra decima est. \navnote{lectumsternere = vestem super lectum stemerc}Patienter exspectā, dum servī lectōs sternunt. Cēnābimus cum primum cocus cēnam parāverit et servī triclīnium ornaāverint. Brevī cēna parāta et triclīnium ōrnātum erit.`` HI

\navnote{vāsa SES ER, argentea CĒS C EJ DS puer = servus intrēmus! = quīn intrāmus?}Tandem puer 'cēnam parātam esse' nūntiat. ``Triclīnium intrēmus!'' inquit Iūlius, atque convīvae laeti triclīnium flōribus exōrnātum et veste pulcherrimā strātum intrant. Rosae et līlia et alia multa flōrum genera in \navnote{genus -eris n vās vāsis n pl vāsa -ōrum argentō factus argentum -īn: ex argentō fiunt dēnāriī}mēnsā sparsa sunt inter vāsa et pōcula argentea; nec enim quidquam nisi argentum mēnsam decet viri nōbilis. (Argentum quidem minōris pretiī est quam aurum, nec vērō quisquam ex vāsis aureis cēnat nisi hominēs divitissimi atque glōriōsī, ut rēgēs Orientis.)

\navnote{glōriōsus -a -um = nimis glōriae cupidus}Iūlius, dominus convīviī, cum Aemiliā in lectō mediō \navnote{ac-cumbere = accubantem sē pōnere}accumbit;

in aliīs duōbus lectīs bīnī convīvae accumbunt: Cornēlius et Fabia in lectō summō ad sinistram Iūliī, Orontēs et Paula ad dextram Aemiliae in lectō īmō. Tum dēmum incipit cēna.

Primum ōva convivis appōnuntur; deinde piscēs cum \navnote{holus -eris n = herba quam edunt hominēs nux}holeribus; sequitur caput cēnae: porcus quem Iūlius ipse ē grege ēlēgit; postrēmō mēnsa secunda: nucēs, ūvae, varia genera mālōrum. Cibus optimus est atque \navnote{dus esse carō carnis f acūtus -a -um: (culter) a. = quī bene secat}convivis placet, māximē vērō laudātur carō porcī, quam minister cultrō acutō secat convivis spectantibus.

\navnote{sānē = certē}``Haec carō valdē mihi placet'' inquit Fabia cum prīmum carnem gustāvit, ``Cocus iste sānē negōtium suum scit.''

``Ego cocum nōn laudō'' inquit Orontēs et salem carnī \navnote{sal salis m carnī:: in carnem gere cibum sale aspergere = cibō salem aspergere}aspergit, quī sale nōn ūtitur! Optima quidem est carō, sed sale caret.`` Orontēs cibum sale aspergere solet, ut sitim augeat! (Sāl est māteria alba quae in marī et sub terrā invenītur.)

\navnote{calida -ae/f --- aqua calida funderc fūdisse fūsum merum = vīnum sine aquā miscēre -uisse mixtum pōtāre = bibere līberāre = līberum facere}Iam ministrī vīnum et calidam in pōcula fundunt. Rōmānī vīnum cum aquā miscent neque vīnum merum bibere solent. Sōlus Orontēs, cui nōn placet vīnum mix-tum, merum pōtat, sed is Graecus est atque lībertīnus. (Lībertīnus est quī servus fuit et līberātus est; in lectō īmō accubant lībertīnī.)

Iūlius pōculum tollēns ``Ergō bibāmus!'' inquit, ``Hoc 120 \navnote{bibāmus! = quīn bibimus?}vīnum factum est ex optimīs ūvis mearum vīneārum. Nec vīnum meum pēius esse mihi vidētur quam vīnum illud Falernum quod vīnum Italiae optimum habētur.`` \navnote{Falernum -In = vīnum ager = terra, regiō y UC X J, c N 7 Ager Ve N /'alernus CĒN Via Appia. Capua}(Falernum est vīnum ex agrō Falernō, regiōne Campā-

Statim Cornēlius ``Sānē optimum'' inquit ``vīnum est tuum, etiam melius quam Falernum'', itemque Fabia Sānē ita est'' inquit, nam ea omnibus dē rēbus idem sentit quod maritus. 130

\navnote{Campania}At Paula vīnum gustāns ``Hoc vīnum'' inquit ``nimis acerbum est: ōs mihi contrahitur. Ego vīnum dulce \navnote{mel mellis n vīnō (dat): in vīnum ap-portāre \contr{} ad-portāre}amō; semper mel vīnō misceō.`` Statim minister mel apportat, quod Paula in pōculum suum fundit. (Mel est quod apēs ex flōribus quaerunt; nihil melle dulcius est.)

Iūlius: ``Idem nōn omnibus placet. Sed quidnam tū sentīs, Orontēs? Utrum vīnī genus melius esse tibi vidē tur, Falernum an Albānum?''

``Equidem'' inquit Orontēs ``sententiam meam nōn \navnote{sententia mea = id quod sentiō quam}ante dicam quam utrumque gustāverō.`` 140

Ad hoc Iūlius ``Rēctē mē monēs'' inquit ``ūnum vīnī genus parum esse in bonā mēnsā. Profectō utrumque gustābis. Age, puer, prōfer Falernum quod optimum \navnote{prō-ferre = prōmere}habeō! Tum dēmum hoc vīnum cum illō comparāre poterimus, cum utrumque gustāverimus. Ergō pōcula \navnote{ex-haurire -sisse -stum \contr{} \contr{} +implēre}exhaurite, amīcī! Cum primum meum vīnum potāveritis, Falernum pōtābitis!``

Pōculum Orontis primum Falernō complētur, nam is \navnote{com-plēre = implēre}iam prīdem pōculum suum exhausit. Deinde ministrī Falernum in cētera pōcula fundunt. Omnēs, postquam vīnum gustāvērunt, idem sentiunt: vīnum Falernum multō melius esse vīnō Albānō! Cornēlius et Orontēs \navnote{inter sē aspiciunt = alter alterum aspicit}inter sē aspiciunt. Neuter eōrum sententiam suam apertē dīcere audet.

Tum Orontēs sīc incipit: ``Nesciō equidem utrum melius sit. Dulcius quidem est Falernum,

nec vērō 155

At Cornēlius prūdenter ``Utrumque'' inquit ``aequē bonum est. Neutrum melius esse mihi vidētur.''

\booksection{Grammatica Latīna}

Verbī tempora Futūrum perfectum [A] Āctīvum.

Dux mīlitem laudābit, si fortiter pugnāverit.

\navnote{pugnāverit -ent}'Laudābit' tempus futūrum est. 'Pugnāverit est tempus futūrum

perfectum.

Futūrum perfectum dēsinit in -erit 165 \navnote{ītaxāv eri s recitāu eri mus reciufo eri ti\$}(pers. iii sing.), quod ad infinītīvum perfectī sine -isse adicitur.

Exempla: [1] recitāverit; [2] pārulerit; [3] scripserit; [4] audīv erit.

\navnote{pāru eris pāru ent pāru eri mus pāru eri tis pāru eri nt [3] scrīps erō scrīps eri s scfīps eri t scrlps erilmus scrīps eriltis scrlps erint}Discipulus laudābitur sl magistrō pāruerit et industrius fuerit: s1 rēctē scrīpserit, bene recitāverit et attentē audīverit.

Magister: Tē laudāābō si mihi pārueris et industrius fueris.`` Discipulus: ''Quid mihi faciēs si piger fuerō nec tibi pāruerō?`` Magister: ''S1 prāvē scripseris et male recitāveris nec attentē audīveri, tē verberābō!`` Discipulus:

``Ergō mē laudābis sī 175 rēctē scrīpserō, bene recilāverō et attentē audtveroō.''

Discipulī laudābuntur si magistrō pāruerint et industriī fuerint: si rēctē scripserint, bene rediāverint et attentē audīverint.

\navnote{audifc \contr{} eri t auāīv erilmus audtv eriuis audīiymnt}Magister: Vōs laudābō si mihi pārueritis et industriī fueritis.`` Discipulī: ''Quid nōbīs faciēs sī pigri fuerimus nec tibi pliuerimus?`` Magister: ''SI prāvē scripseritis et male recitāveritis nec attentē mdtuerins, vōs verberābo!`` Discipulī: ''Ergō nōs laudābis si rēctē scripserimus, bene recitaverimus et attentē zūdīveritnus.``

Singulāris

Plūrālis Persōna prīma

-erō Persōna secunda

-eris Persōna tertia

-erit [B] Passīvum. Pater gaudēbit sī fīlius ā magistrō laudātus erit (= si magister fīlium laudāverit). \navnote{erunt -tus erō tus -tī -ta -tae eris eritis erit erunt -Ium erit laudātus 145erō laudōr us eris laudātus erit laudat]i erimus laudat i eritis laudat erunt}Pater gaudēbit si filii ā magistrō teudātī erunt (= si magister fīliōs laudāverit). Pater: ``Gaudēbō, fili mi, sī laudātus eris.'' Fīlius: ``Quid mihi dabis si laudātus erō?'' Pater: ``Gaudēbō, filii meī, sī laudātī eritis.'' Filii: ``Quid nōbīs dabis si laudātī erimus?''

Singulāris

Plūrālis Persōna prīma

laudātus erō

laudārf erimus Persōna secunda

laudātos m\$

laudatf eritis Persōna tertia

laudātus erit

laudātf erunt

\booksection{Pēnsum A}

Syra: ``Iam dormi, Quīnte! Cum bene dormīv-, valēbis.'' Quīntus: ``Nōn dormiam antequam tū mihi fabulam nārrāv-. Cum fabulam audīv-, bene dormiam. Cum bene dormīv-, brevī sānus erō, nisi medicus mē necāv-!''

Patria salva erit sī mīlitēs nostrī fortiter pugnāv-.

Dux: ``Nisi vōs fortiter pugnāv-, mīlitēs, hostēs castra nostra expugnābunt.''

Mīlitēs: ``Num quid nōbīs dabitur, sī fortiter pugnāv-?''

Verba: induere -isse -um; ēligere -isse -um; coquere -isse -um; sternere -isse -um; fundere -isse -um; miscēre -isse -um; exhaurire -isse -um; revertī ---isse/---um esse; ūtī -um esse.

\booksection{Pēnsum B}

Iūlius Cornēlium et Orontem, amīcōs et --- suōs, cum uxōribus ad --- vocāvit. Cum hospitēs veniunt, Iūlius in --- lavātur. Aemilia eōs --- iubet [= salūtat] et marītum suum --- excūsat. Cornēlius ad vīllam suam reversus ōtiō ---. Orontēs, qul ex longō --- revertit, vītam rūsticam nōn ---, sed --- esse cēnset. Hospitēs in ātriō exspectant, dum cibus ---. Servus quī in --- cibum coquit, --- appellātur. In --- sunt trēs lecti: lectus sum-mus, medius, ---; in singulīs lectīs --- aut ---. aut ternī --- accubant. Rōmānī in lectīs cubantēs ---.

Tandem puer 'cēnam parātam esse' ---. In mēnsā sunt --- et pōcula ex --- facta. Cibus omnibus ---, māximē autem --- laudātur. --- [= servī] vīnum in pōcula ---. Orontēs vīnum merum --- [- [=bibit], cēteri convīvae aquam vīnō ---.

Sine cibō homō --- vīvere potest, sine aquā --- tantum. --- māla rēs est, sed multō pēior est ---. Ex malīs minimum --- oportet.

\booksection{Pēnsum C}

Quī sunt Cornēlius et Orontēs? Ubi est Iūlius cum hospitēs adveniunt? Quid est balneum? Nōnne iūcunda est vita rūstica? Quid est coci negōtium? Num Rōmānī in sellis sedentēs cēnant? Quot lecti sunt in triclīniō? Quot convīvae in singulīs lectīs accubant? Ex quā māteriā pōcula et vāsa facta sunt? Tūne vīnum aquā calidā mixtum bibis?

\booksection{Vocābula nova}
\noindent\small erō, fuerimu\$, Jmerutis, en, t, fuerint, -erō, -erimus, eris, -eritis, -erit, laudari, nova, balneum, hospes, cēna, iter, famēs, sitis, bonum, triclīnium, culīna, cocus, minister, medium, convīvium, convīva, genus, vās, argentum, holus, nux, carō, sāl, calida, mcrum, lībertīnus, mel, inexspectātus, tardus, dīligēns, iūcundus, molestus, īmus, argenteus, glōriōsus, acūtus, merus, acerbus, dulcis, recipere, salvēre iubēre, vīsere, requiēscere, fruī, nūntiāre, contrahere, praeesse, cēnāre, perferre, ēligere, coquere, exōrnāre, parāre, sternere, accubāre, accumbere, placēre, gustāre, aspergere, fundere, miscēre, pōtāre, līberāre, apponāre, prōferre, exhaurire, complēre, singulī, bīnī, ternī, circiter, diū, paulisper, dēmum, pridem, equidem, sānē.

\bookchapter{Cap. XXXI}{Inter pōcula}

Nōn sōlum dē cibō et pōtiōne est sermō convīvārum. \navnote{pōtiō -ōnis f - quod pōtātur}Iūlius hospitēs suōs dē rēbus urbānīs interrogat: ``Quid novi ex urbe? Octō diēs iam sunt cum Rōmae nōn fuī, nec quisquam interim mihi litterās inde mīsit. Quam ob rem nec ipse praesēns nec absēns per litterās quidquam \navnote{praesēns -entis \contr{} \contr{} absēns}cognōvi dē eō quod nuper Rōmae factum est.``

Aemilia: ``Nēmō tibi quidquam scribet dē rēbus urbānīs, nisi prius ipse epistulam scripseris.''

Orontēs: ``Opus nōn est epistulās exspectāre, nam facile aliquid novī per nūntiōs cognōscere potes. Cūr \navnote{nūntius -īm = is quī nūntiat}nōn servum aliquem Rōmam mittis?``

Iūlius: ``Servī sunt malī nūntiī: saepe falsōs rūmōrēs \navnote{rūmor -ōris m = rēs ex aliīs audīta quae nārrātur quī-dam, acc quen-dam}nūntiant. Numquam servōs meōs Rōmam mittō.``

Cornēlius: Quid? Heri quendam servum tuum vīdī in viā Latīnā. Faciem recognōvī, saepe eum hic vīdī.``

Iūlius ā Cornēliō quaerit 'quod nōmen et sit?

\navnote{ab aliquō quaerere = aliquem interrogāre ali-quī -qua -quod}Cornēlius: ``Aliquod nōmen Graecum, putō. ' 'Midās' fortasse, nec vērō certus sum. Semper nōmina oblīvīscor, nam māla memoria mihi est.''

Orontēs: ``Midās est nōmen rēgis, dē quō haec fabula \navnote{quī- quae- quod-dam, abl quō- quā- quō-dam, gcn cuius-dam avārus -a -um = cupidus pecūniae optāre = cupere}nārrātur: In quādam urbe Asiae ōlim vīvēbat rēx quīdam avārus, nōmine Midās, quī nihil magis optābat

Iūlius, quī fabulam audīre nōn vult, Orontem interpellat: ``Nōn Midās'' inquit, ``sed Mēdus est nōmen

Orontēs vērō, minimē turbātus,

nārrāre pergit: \navnote{minimē = nūllō modō Bacchus -i zr: deus vīnī quid-quid = omnis rēs quae}ūTum Bacchus deus, quī ob quoddam beneficium rēgī bene volēbat, Dabō tibi'' inquit ``quidquid optāveris.'' Statim Midās ``Ergō dā mihi'' inquit ``potestātem quid-quid tetigerō in aurum mūtandi. Hoc sōlum mihi optō.'' \navnote{tangere tetigisse tāctum mūnus -eris n = dōnum}Bacchus, etsī rēgem avārum mūnus pessimum optāvisse cēnsēbat, tamen prōmissum solvit.``

Iūlius impatiēns ``Tacē, Orontēs!'' inquit, ``Omnēs \navnote{im-pauens -entis \contr{} inpatiēns}illam fabulam nōvimus.``

At Aemilia, quae fabulam ignōrat, ab Oronte quaerit 'quamobrem id mūnus pessimum sit?'

\navnote{quam-ob-rem = cūr cui:: Aemiliae}Cui Orontēs ``Stultē id quaeris'' inquit, ``Midās enim, quamquam terram, lignum, ferrum manū tangendō in aurum mūtāre poterat, fame et sitī moriēbātur, cum cibus quoque et pōtiō, simul atque ā rēge tācta erat, \navnote{simul atque = eōdem tempore quō, cum primum}aurum fieret. Postrēmō rēx miser deum ōrāvit ut mūnus illud īnfelīx revocāret. Bacchus igitur ei suāsit ut in \navnote{īn-fēlīx -Īcis e}quōdam flūmine lavārētur; cuius flūminis aqua, simul atque corpore rēgis tācta est, colōrem aureum accēpit.``

\navnote{felix ac-cipere -iō -cēpisse -ceptum}Iūlius: ``Hicine finis fabulae est?''

Orontēs: ``Huius quidem fabulae finis est, sed aliam fabulam dē eōdem rēge nōvī. Deus Apollō effecerat ut

\navnote{Apollō -inis m asinīnus -a -um \contr{} asinus}Iūlius: ``Satis est! Fābulās tuās Graecās audīre nōlu\navnote{ferre ali-quantum --- haud paulum (nesciō quantum) quantum=quam multum au-ferre abs-tulisse ablātum fidere (+ dat/abl) cōn-fidere = fidere fidēs -ei f - animus fidēns fidus -a -um īn-fidus}mus. Redeāmus ad meum Mēdum servum, quī heri aufūgit aliquantum pecūniae sēcum auferēns.``

Cornēlius: ``Quantum pecūniae abstulit?''

Iūlius: ``Centum circiter sēstertiōs. Atque ego illī servō praeter cēterōs fīdēbam! Posthāc servō Graecō nūllī cōnfidam, neque enim fidē meā dignī sunt: īnfidī et nēquam sunt omnēs! In familiā meā ūnum sōlum servum fīdum esse crēdō.'' Il

Hic Aemilia maritum interpellat et ``St, Iūli!'' inquit, ``Nōlī servum praesentem laudāre!''

Iūlius Dāvum cōnspiciēns ``Sed is servus adest'' in-quit, ``Nōlō eum laudāre praesentem. Mēdus vērō plānē \navnote{cruciāre = dolōribus afficere prius-quam = antequam crux -ucis f. latere = se occultāre, occultāri fugitīvus -a -um = quī aufūgit tot... quot = tam multī}īnfīdissimus omnium est. Profectō eum verberābō at-que omnibus modīs cruciābō, si eum invēnerō priusquam Italiam relīquerit. Nisi pecūniam mihi reddiderit, in cruce fīgētur!``

Cornēlius: ``Etiam si adhūc Rōmae latet, difficile erit servum fugitīvum in tantā urbe reperīre. Rōmae enim tot servī sunt quot hominēs līberī.''

Aemilia: ``Fortasse Rōmam abiit ob amōrem alicuius 70 mulieris. Iuvenis est Mēdus: quid nōn faciunt iuvenēs \navnote{īuvenis -1s m = vir circiter xxx annōrum crēdere = putāre}amōris causā? Crēdō eum apud puellam Rōmānam latēre.``

Orontēs: ``Ergō numquam reperiētur, nam vērum est \navnote{Ovidius}quod scripsit Ovidius in librō quī vocātur 'Ars amandf:

Quot caelum stēllās, tot habet tua Rōma puellās.?

\navnote{praemium -īn = mūnus quod datur prō beneficiō re-trahere}Iūlius: ``Profectō magnum praemium dabō ei quī ser-vum meum fugitīvum hūc retrāxerit.''

\navnote{statuere -uisse -ūtum tantum quantum = tam multum quam nimius -a -um = nimis magnum}Cornēlius: ``Quantum pecūniae dabis? Certum praemium statuere oportet.''

Iūlius: ``Tantum quantum ille surripuit.''

Orontēs: ``Centum tantum sēstertiōs? Sānē nōn nimium praemium prōmittis!''

Aemilia autem marītō suō suādet ut clēmēns sit: \navnote{clēmēns -entis --- \contr{} sevērus}``Nōlī Mēdum cruciāre, si eum invēneris. Clēmēns estō, mi Iūli! Centum sestertii haud magna pecūnia est, ut ait Orontēs, nec aliud quidquam surripuit Mēdus.''

Iūlius: ``An cēnsēs eum praemium meruisse quod manūs abstinuit ā gemmīs tuīs? Nimis clēmentēs sunt mulierēs: quam facile viris nēquissimīs ignōscunt! At 90 \navnote{ignōscere (+ dat): alicui i. = maleficia alicuius oblīvīsc/nōn pūnire}nostra melior est memoria!``

Aemilia: ``Nōvistīne hoc dictum: Dominō sevērō tot esse hostēs quot servōs'? Servī enim dominum clēmentem amant, sevērum ōdērunt.''

\navnote{ōdisse (perf) \& \contr{} amāre}Iūlius: Servī mē metuunt quidem, nec vērō ōdērunt. Nec enim umquam sine causā servum pūnīvī. Sum do-minus iūstus. Servus dominum iniūstum ōdit, iūstum \navnote{iūstus -a -um \contr{} in-iūstus in-iūria -ae f --- factum iniūstum nec enim quidquam = nihil enim pūnire \contr{} poena -ae f eius generis servī = tālēs servī}et sevērum metuit, nōn ōdit. Nē servō quidem iniūriam facere oportet, sed necesse est servōs infidos aut fugitīvōs sevērē pūnīre, nec enim quidquam nisi poena sevēra eius generis servōs ā maleficiīs dēterrēre atque in officiō tenēre potest. Neque quisquam mē accūsābit sī servum meum cruciāverō aut interfecerō, id enim est iūs dominī Rōmānī. Servum aliēnum necāre nōn licet, \navnote{quod licet et iūstum est lēx legis f vetāre ++ permittere}ut scriptum est in lēgibus, nec vērō ūlla lēx dominum vetat servum suum improbum interficere.``

Cornēlius: ``Nec ūlla lēx id permittit. Nōn idem est \navnote{Solō -ōnis m sapiēns -entis = prūdēns}permittere ac nōn vetāre. Solō, vir sapiēns et iūstus, quī Athēniēnsibus lēgēs scrīpsit, nūllam poenam statuit in \navnote{parricīda -ae m = qul patrem suum occīdit id-eō = ob eam rem}parricīdās. Num ideō cēnsēs civi Atheniensi licuisse patrem suum necāre?``

Iūlius: uIta sānē nōn cēnseō. At quamobrem Solō nūllam poenam in parricīdās statuit? Quia nēmō Athēniēnsis umquam post hominum memoriam patrem suum occīderat, nec ille vir sapientissimus arbitrābātur quemquam posteā tam inhūmānum scelus factūrum \navnote{scelus -eris n = grave maleficium scelestus -a -um \contr{} scelus capite punīre = pūniendī causā interficere supplicium -īn = poena sevērissima, poena capitis (: vītae)}esse. At profectō aliud est patrem suum necāre, longē aliud servum scelestum capite pūnīre, illud enim turpissimum scelus, hoc supplicium iūstum est. Ōlim iūs erat patri familiās nōn modo servōs, sed etiam līberōs suōs interficere. Eius reī exemplum memorātur Titus Mānlius Torquātus, quī fīlium suum cōram exercitū necārī \navnote{cōram frp-r abl: c. aliquō = ante oculōs alicuius}iussit quia contrā imperium patris cum hoste pugnāverat! Sānē pater crudēlis fuit Mānlius, sed illō suppliciō \navnote{crūdēlis -e = saevus atque inhūmānus}sevērissimō cēteri mīlitēs dēterrēbantur nē officium dē- 125 sererent.``

Aemilia: ``Nōtum est veterēs Rōmānōs etiam ergā \navnote{vetus -eris = antiquus}līberōs suōs crūdēlēs fuisse, nec vērō quisquam hodiē exemplum sūmit ab illō patre crūdēlissimō.``

Orontēs: ``At etiam nunc patrī licet īnfantem suum invalidum in montibus expōnere.''

Aemilia: ``Pater quī infantem exposuit ipse necandus \navnote{-um:virnecandus/pūniendus est = virum necārī/pūnīrīoportet}est! Nōnne tālis pater tibi vidētur cruce dignus esse?``

Iūlius: ``Certē pater tam inhūmānus sevērē pūniendus est, namque īnfantēs invalidōs expōnere est mōs 135 \navnote{nam-que = nam mōs mōris m = id quod fieri solet crucī:: in cruce}antīquus atque crūdēlis. Aliī nunc sunt mōrēs. Vērum hominem līberum crucī fīgere nōn est mōs Rōmānōrum; id supplicium in servōs statūtum est.``

Aemilia: ``Ergō quī Infantem suum dēbilem ad ferās \navnote{dēbilis -e --- \contr{} validus}expōnī iussit, ipse ad bēstiās mittendus

est cum aliīs 140 hominibus scelestīs!``

\navnote{\contr{} Christus; m pi quī Chrīstum adōrant}minem Iūdaeum tamquam novum deum adōrant, deōs veterēs Rōmānōs dērident. In convīviīs suīs sanguinem hūmānum bibere solent, ut rūmor est.``

Aemilia: ``Nōn omnēs vēri sunt rūmōrēs qul afferuntur super Christiānīs.''

Fabia: ``Nec omnēs īnfantēs expositī pereunt. Aliī in silvīs ab ipsīs ferīs aluntur, aliī inveniuntur

\navnote{\&ducāre-(līberōs) alere, cūrāre, docēre}ā pāstōribus, quī eōs cum liberis suīs ēducant.``

\navnote{Paris -idis m Priamus -Im: rēx Trōiae}Orontēs: ``Sīcut Paris, rēgis Priamī filius dēbilis, quī ā servō rēgis fīdō in quōdam monte prope urbem

At Cornēlius ``Opus nōn est'' inquit ``vetus exemplum \navnote{liber vetus fabula vetus vīnum vetus}Graecum afferre, cum complures fabulae nārrentur dē Rōmānīs puerīs quī ita servātī sunt. Cēterum fabulam male intellēxistī, nec enim dēbilis fuit Paris nec fīdus \navnote{intellegere -lēxisse -lēctum}servus Priamī, nam rēx ei imperaverat ut Paridem interficeret, et quidquid dominus imperāvit, servō facien-

\navnote{servō (dat):: ā servō}Orontēs: ``Ille servus nōn puniendus,

sed potius laudandus fuit: namque ita Paridem servāvit --- eum \navnote{Menelāus -I m: rēx Spartae ab-dūcere}quī posteā Helenam, feminam omnium pulcherrimam, ā marītō Menelāō abdūxit.``

Paula:

``Num

tantam

iniūriam

laudandam

esse cēnsēs?``

Orontēs: ``Quod Venus suādet iniūria nōn est! Sānē laudandus est ille iuvenis qul nōn modo feminam illam \navnote{audēre ausum esse}pulcherrimam abdūcere ausus est, sed etiam mīles fortissimus fuit quī et multōs aliōs hostēs et ipsum Achillem occīdit.`` Hīc pōculum tollit Orontēs et exclāmat: ''Vīvat fortissimus

quisque!

Vīvant omnēs feminae \navnote{fortissimus quisque = omnēs viri fortēs senex senis m= vir annōrum plūs quam LX Nestor -oris m}amandae! Gaudeāmus atque amēmus! Iuvenēs sumus ut Paris, nōn senēs ut Priamus, rēx Trōiānōrum, aut Nestor, dux Graecōrum senex, quī ad nōnāgēsimum annum vīxit. Quisquis feminās amat, pōculum tollat et \navnote{quis-quis = omnis homō quī}bibat mēcum! Nunc merum bibendum est!``

\navnote{nimium = nimis multum bibere bibisse alterum tantum = bis tantum}Cornēlius: ``Tacendum est, nōn bibendum! Iam nimium bibistī. Cēnseō tē ūnum

tantum

vīnī bibisse quantum nōs omnēs, vel potius alterum tantum!``

Orontēs: ``Vōs igitur parum bibistis. Numquam nimium huius vīnī bibere possum. Valeat quisquis vīnum bonum amat! Vīvat Bacchus, deus vīnī! Vīvāmus omnēs et bibāmus! Pōcula funditus exhauriāmus!''

\navnote{funditus (adv) = ā fundō ab ōvō ūsque ad māla = ab initiō cēnae ūsque fabulārī = loquī i}Paula: ``Iam tacē! Satis est. Nōnne tē pudet ita ab ōvō 185 ūsque ad māla fabulārī? Sānē pudendum est!''

Orontēs autem, simul atque pōculum suum funditus exhausit, ā Paulā ad Aemiliam versus ``Omnēs m-mē interpellant'' inquit, ``praeter t-tē Aemilia. Tu t-tam p-

rudēs pāstōrēs ēducātus erat! Numquam mōrēs urbānōs didicistī, rūstice! Nimium pōtāvistl, ēbrius es. Abstinē \navnote{mium vīnī}manum ā mē!``

Orontēs iterum pōculum tollēns haec cantat: ``Quisquis amat valeat! Pereat quī nescit amāre!

\navnote{pōtāvit per-īre, coni -eat bis tantō = bis tantō magis nūgac -ārum = pl --- rēs stultae negat sē esse* = dīcit sē nōn esse' falsus = quī fallit, fallāx}Bis tantō pereat quisquis amāre vetat!``

Aemilia: ``Nōlumus istās nūgās audīre. Ēbrius es!''

Orontēs negat 'sē esse ēbrium' atque in lectō surgēns aliud carmen super feminā falsā et infidā cantāre incipit, sed priusquam fīnem facit, sub mēnsam lābitur!

Duo servī eum ē triclīniō auferunt atque in cubiculō pōnunt. Tum vestem super eum iam dormientem sternunt.

\booksection{Grammatica Latīna}

\navnote{Haec Inscriptiō inventa est Pompēiisin oppidō Campāniae (item imāgō}\navnote{inscriptio -ōnis f/.--- quod īnscriptum est}Gerundīvum

Vir laudandus. Fēmina laudanda. Factum laudandum.

'Laudandus -a -um' gerundīvum

appellātur. Gerundīvum est adiectīvum dēclīnātiōnis i/ii. Cum verbō esse coniūnctum gerundīvum significat id quod fieri oportet; is ā quō aliquid fieri oportet apud gerundivum significātur datīvō.

Exempla:

\navnote{magistrō (dat)}Discipulus industrius magistrō laudandus est. Discipulus piger reprehendendus et pūniendus est. Tacendum est.

Lingua Latīna vōbīs discenda est. Vocābula dīligenter scribenda sunt. Omnia menda corrigenda sunt: zddendae sunt litterae quae dēsunt, quae supersunt stilō versō dēleruftze sunt. Quidquid magister imperāvit discipulo faciendum est.

Dominus dīcit *'ovēs bene cūrandds esse.'

\booksection{Pēnsum A}

In hīs exemplis syllabae quae dēsunt add- sunt:

Mercēs ad diem solv- est. Quī fūrtum fecit pūn- est. Quid-quid dux imperāvit mīlitibus faci- est. Quid magis opt- est quam vīta beāta? Ē malīs minimum ēlig- est. In perīculīs dēspēr- nōn est. Pater dīcit 'filium pūn- esse, nōn laud-.' Verba: tangere -isse -um; accipere -isse -um; auferre -isse -um; statuere ---isse -um; intellegere ---isse -um; bibere -isse; audēre ---um esse. In exemplis quae sequuntur vocābula add- sunt.

\booksection{Pēnsum B}

Herī Mēdus ā dominō --- aliquantum pecūniae sēcum ---. Mēdus dominum suum nōn amat, sed ---. Iūlius, qul eum Rōmae --- [= occultāri] putat, magnum praemium dabit ei quī eum invēnerit --- [= antequam] Italiam relīquerit. Iūlius dīcit 'mulierēs nimis --- esse ac facile viris nēquissimīs ---.' --- dominus imperāvit servō faciendum est. Solō, vir---, ---, Athēniēnsibus --- optimās scrīpsit. Patrem suum necāre --- inhūmānum est. Iūlius nōndum --- est ut Nestor, sed adhūc --- ut Paris ille qul Helenam ā maritō ---. ``--- amat valeat!'' cantat Orontēs, qul --- est quod --- [= nimis multum] vīnī pōtāvit.

\booksection{Pēnsum C}

Quis fuit Midās? Quamobrem Midās fame et sitī cruciabāatur? Cūr Iūlius illam fabulam audire nōn vult? Ubi Cornēlius servum Iūliī vīdit? Quantum pecūniae Mēdus sēcum abstulit? Estne Mēdus adhūc Rōmae? Quid faciet Iūlius si Mēdum invēnerit? Quōmodo hominēs ā maleficiīs dēterrentur? Quam feminam Paris abdūxit? Quī hominēs ad bēstiās mittuntur? Cūr Orontēs pedibus stāre nōn potest?

\booksection{Vocābula nova}
\noindent\small nova, pōtiō, rūmor, memoria, mūnus, fidēs, crux, iuvenis, praemium, poena, iūs, lēx, parricīda, scelus, supplicium, mōs, iniūria, senex, nūgae, praesēns, avārus, impatiēns, īnfelīx, asinīnus, fidus, īnfidus, fugitīvus, nimius, clēmēns, iūstus, iniūstus, sapiēns, scelestus, crūdēlis, vetus, invalidus, dēbilis, ēbrius, nōnāgēsimus, optāre, interpellāre, aufugere, auferre, fidere, cōnfidere, cruciāre, latēre, retrahere, statuere, ignōscere, ōdisse, vetāre, ēducāre, abdūcere, fabulāri, quidquid, quisquis, quantum, aliquantum, nimium, quamobrem, ideō, funditus, priusquam, namque, cōram.

\bookchapter{Cap. XXXII}{Classis Rōmāna}

\navnote{nāvēs longae}\navnote{cūnctus -a -um = omnis, Infestus -a -um: (locus) īnfestus e \contr{} tūtus \contr{} lībertās sub-mergere = mergere (sub aquam)}I Ōlim cūncta maria tam īnfesta erant praedōnibus ut nēmō nāvigāret sine māximō perīculō mortis aut servītūtis. Multī nautae et mercātōrēs, mercibus ēreptīs nāvibusque submersīs, ā praedōnibus aut interficiēbantur aut in servitutem abdūcēbantur.

Iī sōlī quī magnam \navnote{līberāre + abl = 1.ah/ex}pecūniam solvere potuerant servitūte līberābantur. Ipse Gāius Iūlius Caesar, cum adulēscēns ex Italiā Rhodum nāvigāret, ā praedōnibus captus est nec prius līberātus quam ingēns pretium solvit.

Nec sōlum nautae, sed etiam incolae ōrae maritimae \navnote{incola -ac m = quī/quae incolit}īnsulārumque in metū erant. Nōnnūllae īnsulae ab incolis dēsertae erant, multa oppida maritima ā praedōnibus \navnote{(oppidum) capere = expugnāre vis f, acc vim, abl vi audacia -ae f \contr{} audāx contemnere = parvī aestimāre, nōn timēre}capta. Tanta enim erat vīs et audācia eōrum, ut vim Rōmānōrum contemnentēs etiam portūs Italiae oppugnārent.

Quoniam igitur propter vim atque multitūdinem \navnote{mare Tūscum --- mare TGscum = īnferum}praedōnum nē mare Tuscum quidem tutum erat, pa-rum frūmentī ex Sicilia et ex Āfricā Rōmam advehēbātur. Ita māxima inopia frūmentī facta est, quam ob rem pretium frūmentī semper crēscēbat. Postrēmō, cum 20 \navnote{crēscere = augērī cārus -a -um = magnī pretiī populus -īm = cūnctī cīvēs classis -is f --- nāvium numerus adversus = contrā ēgregius -a -um = melior cēteris, optimus prae-pōnere (+ dat) proximus -a -um sup (comp propior -ius) Rōmae}iam tam cārum esset frūmentum ac pānis ut multī pauperēs inopiā cibī necessarii perīrent, populus Rōmānus ūnō ōre postulāvit ut ūniversa classis Rōmāna ad-versus

hostēs

illōs audācissimōs

mitterētur.

Ergō Gnaeus Pompēius, dux ēgregius, classi praepositus est. Quī prīmum ē mari Tuscō, quod mare proximum Rōmae est, et ex Siciliā, īnsulā Italiae proximā, praedōnēs pepulit, tum eōs in Āfricam persecūtus est. Dēnique, aliquot nāvibus in Hispāniam missīs, ipse cum classe in \navnote{ali-quot indēcl= 2 complūrēs (nesciō quot) vincere vicisse victum victoria -ae f \contr{} vincere gēns gentis f Aegyptiī -ōrum m pl commūnis -e = nōn ūnius sed omnium grātus -a -um = quī placet, optātus nūntius = quodnūntiātur minul coeptum est = minul coepit}Asiam profectus praedōnēs quōs ibi invēnit brevī tem- 30 pore omnēs vīcit. Hāc victōriā ēgregiā omnēs gentēs, ab Hispanis ūsque ad Aegyptiōs Iūdaeōsque, commūnī perīculō līberātae sunt. Simul atque grātus nūntius dē eā victōriā grātissimā Rōmam pervēnit, pretium frūmentī minuī coeptum est; victīs enim praedōnibus,

nautae sine metū per maria, quae omnium gentium commūnia \navnote{summus:: māximus}sunt, nāvigābant. Rōmae igitur ex summā inopiā repente māxima frūmentī cōpia facta est ac pānis tam vīlis fuit quam anteā fuerat --- id quod populō Rōmānō grātissimum fuit. Pompēium victōrem cūnctus populus \navnote{victor -ōris m = quī vīcit}Rōmānus summīs laudibus affecit.

\navnote{fit = accidit}Ex eō tempore rārō fit ut nāvis praedōnum in mari \navnote{internus -a -um \contr{} intrā; mare I.um (intrā fretum Ōceanī)- mare nostrum per-currere \contr{} mercātor tuērī= --- tūtum facere, cūstōdīre iūre:: rēctē, vērē ubī-que = omnī locō super-esse = reliquus esse tantā audāciā esse = tam audāx esse}Internō appāreat, nam classēs Rōmānae, quae cuncta maria percurrunt, nāvēs mercātoriās atque ōram maritimam dīligenter tuentur. Mare Internum iterum 'nōstrum mare' iūre appellātur ā Rōmānīs. Neque tamen classis Rōmāna omnēs nautās quī ubīque nāvigant tuēri potest. Adhūc supersunt aliquot praedōnēs maritimī, quī tantā audāciā sunt ut nē armīs quidem Rōmānōrum dēterreantur.

Amīcī nostri in mari Tuscō nāvigantēs tālēs praedōnēs audācissimōs nāvem suam persequī arbitrantur. Cūnctī perturbātī sunt. Etsī nautae omnibus vīribus \navnote{vīrēs -ium f pl \contr{} vīs rēmigāre = nāvem rēmīs agere ad-iuvāre = iuvāre}rēmigant, tamen illa nāvis, ventō secundō adiuvante, magis magisque appropinquat.

\navnote{rēnus adversus -a -um: (venagere ēgisse āctum}Caelum nūbilum suspiciēns gubernātor optat ut ventus in adversum vertātur. Is enim nautās suōs tam validōs esse crēdit ut nulla alia nāvis rēmīs solis āācta nāvem suam cōnsequī possit. I

\navnote{inter-eā = interim flectere -xisse -xum voluntās -ātis f \contr{} velle; v. tua = quod tū vīs}Intereā Lydia genua flectit et Deum precātur ut sē adiuvet: ``Pater noster qul es in caelīs! Fīat voluntās tua! Sed līberā nōs ā malō!''

Mēdus autem gladium brevem, quem adhūc sub \navnote{\contr{} dūcere in-ermis -e --- \contr{} armātus}veste occultāvit, ēdūcit et ``Equidem'' inquit ``nōn inermis occīdar. Si praedōnēs mē armīs petīverint, omnibus vīribus repugnābō! Fortēs fortūna adiuvat, ut āiunt.''

\navnote{re-pugnāre = contrā pugnāre}Tum vērō Lydia ``Converte gladium tuum'' inquit ``in locum suum! Omnēs enim quī cēperint gladium, gladiō peribunt, ut ait Christus.''

``Sed tē quoque'' inquit Mēdus ``gladiō meō dēfen- 70 dam. Nōlō tē ā praedōnibus occidi spectāre inermis. Dōnec ego vīvam, nēmō tibi nocēbit!''

\navnote{dōnec = dum, tam diū quam}Sed Lydia, quae Mēdum ut furem contemnit, ``Nōn ā tē'' inquit, ``sed ā Deō auxilium petō. Is sōlus nōs tuēri \navnote{ab aliquō petere = aliquem rogāre dis-suādēre \contr{} \contr{} persuādēre neu = nē-ve pugnāre opus esse + abl}potest.``

Item gubernātor multīs verbīs Mēdō dissuādēre cōnātur nē gladium ēdūcat neu praedōnibus vi et armīs resistat: ``Quid opus est armīs? Tanta est vis praedōnum ut nūllō modō iīs resistere possīmus. Neque praedōnēs nautās inermēs occīdunt, cum eōs magnō pretiō servōs 80 vēndere possint.''

Mēdus: ``Iamne oblītus es quid modo dīxeris? Dīxistī \navnote{prae-ferre: mortem servitūtī p. = mori mālle quam servīre haud sciō an dīxerim = fortasse dīxī cānis = qu! dīligitur}enim 'tē mortem servitūtī praeferre.``

Gubernātor: ``Haud sciō an ego ita dīxerim, sed prōfectō lībertās mihi vītā cārior est. Nihil lībertātī praeferō. Quam ob rem omnem pecūniam meam praedōnibus dabō, sī lībertātem mihi reddent. Hanc grātiam \navnote{grātia -ae f --- rēs grāta, beneficium}sōlam ab iīs petam.``

Mēdus: ``Certē praedōnēs pecūniam tibi ēripient, sed \navnote{pecūniae grātiā = ob grātiam pecūniae, pecūniae causā sēstertium = sēstertiorum of-ferre ob-tulisse oblātum = sē datūrum esse dlcere/ostendere}felix eris s1 pecūniae grātiā vītae tuae parcent.``

Gubernātor: Si necesse erit, decem milia sēstertium praedōnibus offerre possum. Quod ipse nōn possideō amīcī mel prō mē solvent.``

Mēdus: ``Ergō nūlla spēs est mihi, quī nec ipse pecūniam habeō nec amīcum tam pecūniōsum, ut mē ē ser- 95 vitūte redimere possit aut velit.''

Tum mercātor ``Mihi quidem'' inquit ``multī sunt amīcī pecūniōsl, sed valdē dubitō num pecūniā suā mē redimere velint. Fortūnā adversā amīcīs fidendum nōn est! Namque amīcī, quōs in rēbus secundīs multōs habēre vidēmur, temporibus adversīs nōbīs dēsunt. Duōs \navnote{lium nōn ferre re-minīscī \contr{} \contr{} oblīvīscī poēta -ae m/f --- quī/quae carmina scrībit (multōs) numerāre: habēre nūbila: adversa}versūs reminīscor ē carmine quod dē hāc rē scrīpsit poēta quīdam:

Dōnec eris felīx, multōs numerābis amīcōs.

Tempora si fuerint nūbila, sōlus eris!``

Gubernātor: ``Nesciō quī poēta ista scripserit. Tune nōmen eius meministī?''

\navnote{meminisse (perf) = memoriā tenēre (}Mercātor: ``Illōs versūs scripsit Ovidius, poēta ēgregius, nisi memoria mē fallit. Quī ipse, cum fortūnā adversā premerētur, ab amīcīs suīs dēsertus erat.''

\navnote{premere:: malā rē afficere}Gubernātor: ``Nōn vērum est quod dīxit Ovidius. Nam etsī rāra est vēra amicitia ac fidēs, nōn omnēs \navnote{amicitia -ae f \contr{}: amīcus fidēs = animus fidus seu = sīve (amīcus) certus: fidus}amīcī sunt falsī seu infidi. Multō melius Ennius poēta:

Amīcus certus in rē incertā cernitur. Certus ac vērus amīcus est quī numquam

amīcō suō deest seu secunda seu adversa fortūna est. Mihi vērō multī sunt tālēs amīcī, quī semper mihi aderunt in rēbus adversīs, seu pecūniā seu aliā rē mihi opus erit. Ipse \navnote{ad-esse af-fuisse ( \contr{} adfuisse) mihi grātus = quī mihi bene vult prō beneficiō; mum grātum) habēre grāliam re-ferre = g.am prō grātiā reddere grātiās agere = dīcere sē grātiam habēre*}enim saepe amīcīs meīs affuī, nēmō amīcus umquam frūstrā auxilium ā mē petīvit. Ergō omnēs mihi grātī sunt prō beneficiīs.``

Mercātor: ``Aliud est grātiam habēre, aliud grātiam referre. Nōn omnēs quī tibi prō beneficiīs grātiās agunt, ipsī posteā, si opus fuerit, grātiam tibi referent. Facile est grātiās agere prō beneficiīs, nec vērō quidquam dif- 125 ficilius esse vidētur quam beneficiōrum meminisse.''

\navnote{meminisse + gen/acc: m. hominis, m. ret/rem}Gubernātor: ``Sed ego ipse soleō amīcīs meīs grātiam referre. Numquam beneficii oblitus sum, semper pecūniam acceptam reddidī.''

Hic Mēdus ``Ergō'' inquit ``melior es amīcus quam ille quem ego aliquandō ē servitūte redēmī.''

\navnote{ali-quandō = aliquō tempore (nesciō}Gubernātor: ``Mīror unde pecūniam sūmpseris ut aliōs redimerēs, cum tē ipse redimere nōn possīs.''

Lydia: ``Ego mīror cūr id mihi nōn nārrāveris.''

\navnote{quandō) nē quis = nē aliquis,}Mēdus: ``Nihil cuiquam nārrāvī dē eā rē, nē quis mē 135 glōriōsum exīstimāret. Sed quoniam omnēs mē quasi servum scelestum contemnitis, nārrābō vōbīs breviter quōmodo amīcum ē servitūte redēmerim atque ipse ob eam grātiam servus factus sim:

``Cum homō līber Athēnīs vīverem, ā quōdam amīcō 140 \navnote{pīrāta -ae 5 = praedō maritimus}epistulam accēpī quā ille mihi nūntiāvit sē ā pīrātīs cap-tum esse' ac mē per amīcitiam nostram ōrāvit ut sē ē servitūte redimerem magnum pretium solvendō. Cum \navnote{mūtuus -a -um: pecūnia mūtua = pecūnia quae reddenda est condiciō -ōnis f --- lēx inter duōs statūta diēs/= f --- tempus statūtum}autem tantum pecūniae nōn habērem, necesse fuit pecūniam mūtuam sūmere. Ergō virum dīvitem adiī, quī 145 mihi omnem pecūniam mūtuam dedit hāc condiciōne, ut annō post ad certam diem omnia sibi redderentur. Pecūniā solūtā, amīcus meus ā pīrātīs līberātus grātiās mihi ēgit prō beneficiō, ac simul mihi prōmīsit 'sē intrā annum omnem pecūniam redditūrum esse' --- sed annō 150 post nē assem quidem ab eō accēperam! Diē ad solvendum cōnstitūtā, cum pecūniam dēbitam solvere nōn \navnote{cōn-stituere = statuere}possem, homō ille dīves mē in carcerem mīsit et aliquot diēbus post servum vēndidit. --- Sed nesciō cūr hoc vōbīs nārrāverim, nec enim sine māximō dolōre eius \navnote{reminīscī + gen/acc: r. hominis, r. Tet/rem ut}temporis reminīscor cum in patriā līber inter cīvēs Itiberōs versārer. Utinam aliquandō līber patriam videam! Sed frūstrā hoc optō, nam iam illī pīrātae eam spem mihi ēripient, idque eōdem diē quō ab amica meā dēsertus sum!`` Hoc dīcēns Mēdus ānulum quem Lydia abicere voluit prae sē fert.

\navnote{prae sē ferre = ante sē tenēre, ostendere nē dēspērāveris! = nōlī dēspērāre!}Gubernātor: ``Nē dēspērāveris! Fortasse ānulō istō aureō tē redimere poteris. Namque avāri atque auri cupidī sunt omnēs pīrātae. Magna est vis auri.'' 165

Mēdus ānulum parvum aspiciēns ``Putāsne'' inquit ``mē tam parvī aestimārī ā pīrātīs?''

Gubernātor: ``Nōn omnēs hominēs parī pretiō aestimantur. Scīsne quantum pīrātae ā Iūliō Caesare captō \navnote{talentum -īn (pecūnia Graeca) = xxiv mīlia sēstertium superbus -a -um = qul aliōs contemnit}postulāverint?

Vīgintī talenta postulāvērunt,

id est prope quīngenta mīlia sēstertium. At Caesar, vir super-bus, cum id parum esse cēnsēret, quīnquāgintā talenta \navnote{mināri(+iar) = poenam prōmittere}pīrātīs obtulit, simul vērō supplicium iīs minātus est! Tum praedōnibus quasi servīs suīs imperāvit ut tacērent neu somnum suum turbārent --- ita Caesar praedōnēs contemnēbat, cum in eōrum potestāte esset. Ubi \navnote{ubi primum+ - perf = cum primum armāre = armīs parāre}prīmum redēmptus est, ipse nāvēs armāvit et captōs praedōnēs in crucem tolli iussit.``

Mēdus: ``Nōn sum tam superbus ut mē cum Caesare comparandum esse putem. Utinam nē pīrātae mē ut \navnote{utinam nē = optō nē}servum fugitīvum occīdant! Vērum hōc ānulō sī quis servāri potest, nōn ego, sed amīca mea servanda est. \navnote{nē abiēceris! = nōlī abicere!}Ecce ānulum reddō tibi, Lydia. Nē eum abiēceris! Utinam ille ānulus vītam tuam servet!``

Lydia ānulum oblātum accipit et ``Grātiās tibi agō'' inquit, ``sed quōmodo tua vīta servābitur?''

``Id nōn cūrō'' inquit Mēdus, ``nec enim mortem metuō, si tē salvam esse sciō.''

Tum Lydia ``Ō Mēde!'' inquit, ``Nunc dēmum intellegō mē tibi vītā cāriōrem esse. Ignōsce mihi quod tē \navnote{meā grātiā = meā causā}accūsāvī! Omnia meā grātiā fecistī. Quōmodo tibi grātiam referam?``

Mēdus: ``Nihil rogō, nisi ut mē amēs ita ut mē amābās. Hōc nihil grātius mihi fieri potest.''

\navnote{hōc nihil grātius = nihil. grātius quam hoc}Lydia nihil respondet, sed Mēdum complectitur at-que ōsculātur. Quid verbīs opus est?

Intereā gubernātor in mare prōspicit et Quid hoc?`` inquit, ''Aliae nāvēs illam sequuntur. Tot nāvēs praedō\navnote{nāvis longa = nāvis ar māta (ē classe Rōmānā) nē timueritis! = nōlīte timēre!}nēs nōn habent. Nāvēs longae sunt, quae mare percur-runt, ut nōs ā pīrātīs tueantur. Nē timueritis, amīcī!``

Mercātor: ``Sed cūr illae nōs persequuntur?''

Gubernātor: ``Quia tamquam praedōnēs ab iīs fugimus. Rēmōs tollite, nautae!''

\navnote{dē-sistere = dēsinere tollere sus-tulisse sub-lātum}Nautae statim rēmigāre dēsistunt ac rēmīs sublātīs nāvēs longās salūtant. Classis celeriter appropinquat. Iam mīlitēs armātī in nāve proximā cernuntur.

Lydia magnā cum laetitiā classem appropinquantem spectat, \navnote{etiam-nunc = adhūc}sed Mēdus etiamnunc perterritus esse vidētur.

``Nōnne laetāris, Mēde'' inquit Lydia, ``quod nōs omnēs ē commūnī perīculō servātī sumus?''

Mēdus: ``Laetor sānē quod vōs servātī estis; sed ego \navnote{ne oblīta sīs! = nōlī oblivisci!}mīlitēs aequē timeō atque pīrātās. Nē oblīta sīs mē ser-vum fugitīvum esse. Timeō nē mīlitēs mē captum Rōmam abdūcant, ut cōram populō ad bēstiās mittar in amphitheātrō. Hoc dominus mihi mināri solēbat.``

``Nē timueris!'' inquit Lydia, ``Mīlitēs ignōrant quī \navnote{timēre!}homō sīs et quid anteā feceris. Iam nēmō nōs prohibēbit simul in patriam nostram commūnem redīre.``

Intereā nāvēs longae tam prope vēnērunt ut mīlitēs cognōscant nāvem mercātōriam esse. Itaque persequī \navnote{cursus -ūs m \contr{} currere}dēsistunt atque cursum ad orientem flectunt. Brevī cūncta classis ē cōnspectū abit.

--- Hīc amīcōs nostrōs in mediō cursū relinquimus. Utinam salvī in Graeciam perveniant! Omnia bona iīs optāmus. '

\booksection{Grammatica Latīna}

Coniūnctīvus Tempus perfectum [A] Āctīvum.

Iūlius dubitat num magister Mārcum laudāverit.

\navnote{laudāverit}Laudāverit' est coniūnctīvus temporis praeteritī perfectī. 230 \navnote{-erit}Coniūnctīvus perfectī dēsinit in -erit (pers. iii sing.), quod ad īnfīnltīvum perfectī sine -isse adicitur.

Exempla: [1] recitāverit; [2] pārulerit;

[3] scrips]erit; [4] audīveri.

Pater fīlium interrogat num bonus discipulus fuerit: num magistrō pāruerit, attentē mdtverit, rēctē scripserit et pulchrē \navnote{itcāv eri mus reciufo eri tis nt recitāz? eri [2] pāru erilm pāru erils pāru erii pāru eri mus pāru eri tis pāru eri ni scrīpi \contr{} n}recitaverit.'

Pater: ``Audīsne? Interrogō tē 'num bonus discipulus fueris, num magistrō parueris, attentē audīveris, rēctē scripseris et pulchrē recitāveris'.'' Fīlius: ``Iam tibi dīxī me industrium fuisse.' Quārē igitur mē interrogās (num bonus discipulus fuerim, num magistrō pāruirim, attentē audīverim, rēctē scrīpserim et pulchrē leGitāverim'? Crēde mihi! Nē dubitaveris dē verbīs meīs!''

Parentēs filiōs interrogant 4num bonī discipulī fuerint: num magistrō pāruerint, attentē audīverint, rēctē scripserint et pulchrē recilāverint.'

Parentēs: ``Audītisne? Interrogāmus vōs num bonī discipulī fueritis, num magistrō pzrueritis, attentē mdīveritis, rēctē sctipseritis et pulchrē recitaveritis'.'' Filii: ``Iam vōbīs dīximus 'nōs industriōs fuisse.' Quārē igitur nōs interrogātis 'num bonī discipulī fuerimus, num magistrō pāruerimus, attentē audīverimitSyrēctē scripserimus et pulchrē reciiāverimus'? Crēdite nōbīs! Nē dubitāveritis dē verbīs nostrīs!''

Persōna prima ---

-erim Persōna secunda ---

-eris Persōna tertia

-erit [B] Passīvum. 260

Pater dubitat num filius ā magistrō laudātus sit.

Pater: ``Tune ā magistrō laudātus es?'' Filius: ``Nesciō num \navnote{laudātus sim laudāt us sīs laudār us sit teuāātlī sīmus laudār līsītis l}laudātus sim!`` Pater: ''Quōmodo nescīs num laudātus sīs ā magistrō?`` Fīlius: ''Nesciō quid ab eō dictum sit, nam in lūdō dormīvī!`` 265

Pater dubitat num filii ā magistrō laudatf sint.

Pater: ``Vōsne ā magistrō laudātī estis?'' Filii: ``Nescīmus num laudātī sīmus!'' Pater: ``Quōmodo nescitis num laudātī sītis?'' Filii: ``Nescīmus quid magister dixerit, nam in lūdō dormīvimus!''

Singulāris Persōna prīma

laudātas sim Persona secunda

laudātas sīs Persōna tertia

laudātus sit

\booksection{Pēnsum A}

Dominus dubitat num pāstor ovēs bene cūrāv-. Dominus: ``DIc mihi, pāstor, utrum in campō dormīv- an vigilāv-.'' Pāstor: ``Mīror cūr mē interrogēs utrum dormīv- an vigilāv-. Semper officium meum faciō.'' Dominus: ``Ergō dic mihi cūr herī ovis ē grege aberrāv- ac paene ā lupō capta s-.'' Pāstor mīrātur unde dominus hoc audīv-.

Iūlius servōs interrogat num Mēdum vīd-. Servī: Nescimus quō fug-, ut tibi dīximus. Cūr nōs interrogās num eum vid?`` Iūlius: ''Id interrogō, quia dubitō vērumne dīx-!``

Verba: vincere -isse -um; agere -isse -um; flectere ---isse -um; offerre ---isse -um; redimere -isse -um; tollere -isse -um; adesse -isse.

\booksection{Pēnsum B}

Priusquam Pompēius, dux --- [= optimus], --- Rōmānae praepositus est, --- [= omnia] maria --- erant praedōnibus, quī Rōmānōs ita --- ut etiam portūs Italiae oppugnārent. Tan-ta erat --- praedōnum ut nēmō iīs--- --- posset. Nēmō sine metū mortis aut --- nāvigābat. --- [= complūrēs] īnsulae ab --- relinquēbantur. Rōmae frūmentum tam --- erat ut multī pauperēs fame (= --- cibī)perirent.

Pīrātae, quī--- --- [= aliquō tempore] Caesarem cēperant, xx talenta --- [= sīve] D milia sēstertium ab eō postulāvērunt, sed Caesar, vir ---, l talenta iīs---. ---. Tantō pretiō Caesar --- est.

Omnēs amīcī prō beneficiīs --- agunt, sed paucī posteā --- referre volunt. Vēra--- --- rāra est. Difficile est beneficiōrum ---. Fortēs fortūna ---.

\booksection{Pēnsum C}

Cūr Pompēius classi Rōmānae praepositus est? Cūr amīcī nostri ab alterā nāve fugiunt? Utrum ventus secundus an adversus est? Quantum gubernātor pīrātīs offerre potest? Num ipse tantum pecūniae possidet? Quandō vērus amīcus cognoscitur? Quōmodo Mēdus servus factus est? Cūr Mēdus adhūc nihil nārrāvit de eā rē? Suntne praedōnēs quī eōs persequuntur? Cūr Mēdus etiamnunc perterritus est? Quō nāvēs longae cursum flectunt?

\booksection{Vocābula nova}
\noindent\small velle, coni praes, velim, velimus, velīs, velītis, velit, velint, adversa fortūna = māla, fortuna, audiven! audi?erimus sie, mus, tis, duāīv eri, audfz, eri, ni, fueri, m, nōll, nē dubitāzeris! 2 nōlī, dubitāre!, nē dubitāveritis/ 2 nōlīte dubitāre!, -erim, erimus, -eris, -eritis, erit, -erint, -tus, sim, slmus, -tī, -ta, -tae, sītis, sīs, sit, sint, nova, servitūs, incola, vīs, vīrēs, audācia, inopia, populus, classis, victōria, gēns, victor, voluntās, grātia, poēta, amīcitia, pīrāta, condiciō, talentum, amphitheātrum, cursus, cūnctus, īnfestus, cārus, ēgregius, proximus, commūnis, grātus, vīlis, internus, mercātōrius, nūbilus, adversus, inermis, mūtuus, superbus, submergere, contemnere, praepōnere, percurrere, tuēri, rēmigāre, adiuvāre, flectere, ēdūcere, repugnāre, dissuādēre, praeferre, offerre, redimere, reminīscī, meminisse, referre, mināri, armāre, dēsistere, aliquot, ubīque, aliquandō, intereā, etiamnunc, dōnec, neu, seu, utinam.

\bookchapter{Cap. XXXIII}{Exercitus Rōmānus}

\navnote{agmen -inis}I Exercitus Rōmānus ūniversus cōnstat ex legiōnibus \navnote{n legiō cohors -rtis}duodētrīgintā, quae in dēnās cohortēs dīviduntur. In singulīs legiōnibus sunt sēna vel quīna vel quaterna mīlia mīlitum, quī omnēs cīvēs Rōmānī sunt. Praetereā \navnote{auxilia -oōrum n pl ad-iungere = addere legiōnārius -a -um \contr{} le-}magna auxilia exercitui adiunguntur. Auxilia sunt peditēs equitēsque ex prōvinciīs, quī arma leviōra, sīcut arcūs sagittāsque, ferunt. Legiōnāriī sunt peditēs scūtīs, gladiīs pīlīsque armātī.

Signum legiōnis est aquila argentea, quae in itinere ante agmen fertur. In itinere cohortēs alia post aliam in \navnote{ōrdō -inis m prō-gredī -ior = prōcēdere veniēns īnstruere -ūxisse -Qctum proelium -īn = pugna exercituum hortāri = monēre prō-currere -currissc}longō ōrdine prōgrediuntur. Talis ōrdō mīlitum prōgredientium dīcitur agmen. Cum agmen ad hostēs pervēnit, si tempus et locus idōneus est ad pugnandum, cohortēs in trēs ōrdinēs Instruuntur. Exercitus ita īnstrūctus aciēs appellātur. Ante proelium dux exercitūs mīlitēs suōs hortātur ut fortiter pugnent. Tum peditēs prō\navnote{mittere = iacere caedere = pulsāre; (gladiō) c. = occīdere imperātor -ōris m = quī imperat (titulus ducis) circum-dare = cingere (locum) mūnīre = vāllō /mūrō circumdare suprā (ado): in cap. xii morāre aetās -ātis f --- annī quōs aliquis vīxit mīlitāris -e \contr{} mīles puer vii annōs nātus = puer vii annōrum velle, coni imperf: vellem vellēmus vellēs vellētis vellet vellent studēre (+ dat) litteris s. = studiōsus esse litterārum (: legendi) studium -īn \contr{} studēre}gladiīs caedunt. Hostibus proeliō victīs, dux ā mīlitibus 'imperātor' salūtātur. Vesperi exercitus locō ad dēfendendum idōneō castra pōnit, quae vāllō et fossā circum- 20 dantur. Ita mūniuntur castra Rōmāna.

Aemilius, frāter Aemiliae minor, quem suprā commemorāvimus, ā primā aetāte studiōsus fuit reī mīlitāris. Iam puer septem annōs nātus gladiōs ligneōs et arcūs sagittāsque sibi faciēbat, ut cum aliīs puerīs eius- 25 dem aetātis proelia luderet. Septendecim annōs nātus ā patre interrogātus 'quid tum discere vellet?' filius statim respondit 'sē nihil nisi rem mīlitārem discere velle.' Voluntās filii patri haud placēbat, cum ipse litteris studēret neque ūllum aliud studium filiō suō dignum esse arbi- 30 trārētur. Cum vērō filius nūllō modō contrā voluntātem ad studium litterārum cōgi posset, pater eum ūnā cum \navnote{cōgere co-ēgisse -āctum persuādēre}Pūbliō Valeriō, adulēscente eiusdem aetātis, in Germāniam ad exercitum Rōmānum mīsit, ut apud ducem \navnote{stipendium -īn = mercēs militis; s. merēre = mīlitāre}quendam ēgregium stipendia merēret. Ibi Aemilius prō 35 patriā pugnāns iam magnam glōriam mīlitārem sibi \navnote{virtūs -ūtis f \contr{} vir; (militis) v. = animus fortis}quaesīvit. Virtūs eius ēgregia ab omnibus laudātur.

Et pater et soror Aemiliī frequentēs epistulās ad filium et frātrem suum mittunt. Pater māximē dē rēbus pūblicīs scribit, sed Aemilia dē rēbus prīvātīs, ut dē 40 \navnote{pūblicus -a-um( \contr{} popu-}liberis suīs et dē convīviīs, scribere solet. Epistulae quās \navnote{(epistulam) dare = mittere praecipuē = prae cēteris rēbus, māximē}Aemilius ad patrem dat praecipuē dē glōriā et virtūte mīlitārī sunt, sed ex iīs litteris privātīs quās Aemilia nuper ā frātre suō accepit plānē appāret eum iam vītā \navnote{appāret = plānum est, intellegitur fatīgāre = fessum facere pri-diē = diē ante}mīlitārī fatīgātum esse. Ecce litterae novissimae, quās Aemilia pridie kalendās Iūniās ā frātre accēpit et posterō diē in convivio recitāvit: Il ``Aemilius sorōri suae cārissimae s. d.

\navnote{s. d. = salūtem dīcit ante diem septimum kalendās Māiās}Hodiē dēmum mihi allāta est epistula tua quae a. d. vii kal. Māi. scripta est, id est ante vīgintī diēs. Quam tardus est iste tabellārius! Celsī sānē et arduī montēs \navnote{arduus}Germāniam ab Italiā disiungunt, ac difficillimae sunt viae quae trāns Alpēs ferunt, sed tamen celer tabellārius \navnote{NEN Alpēs -ium fel: montēs Eurōpae altissimi ferē = circiter}idem iter quīndecim ferē diēbus cōnficere potest, ut tabellāriī pūblici quibus ūtitur dux noster. Ego vērō istum tabellārium properāre docēbō, cum meās litterās ad tē referet.

Sed quamquam tardē advēnit, grātissima mihi fuit tua epistula magnōque cum gaudiō ex eā cognovi tē et \navnote{gaudium -īn ( \contr{} gaudēre) valētūdō -inis f \contr{} valēre Albānum -īn = praedium Albānum}Iūlium et līberōs vestrōs, quōs valdē dīligō, bonā valētūdine fruī. Cum epistulās tuās legō, apud vōs in Albānō esse mihi videor, neque in hāc terrā frigidā inter hominēs barbarōs. Tum, nesciō quōmodo, ita permoveor ut lacrimās vix teneam --- ita patriam et amīcōs meōs dēsīderō.

\navnote{desiderāre: amīcōs d.; --- dolēre sē abesse ab}Ō, quam longē absum ab Italiā et ab iīs omnibus quōs praecipuē dīligō! Utinam ego Rōmae essem aut tu apud mē essēs! Cum Dānuvium flūmen aspiciō, quod praeter \navnote{praeter castra}castra nostra fluit, dē Tiberī cōgitō atque dē Rōmā. \navnote{amnis -is m = flūmen}Quandō tē aspiciam, urbs pulcherrima? Utinam hic amnis Tiberis esset et haec castra essent Rōma!

Sed frūstrā haec optō, cum nēmō nisi deus ita subitō \navnote{trāns-ferre}in alium locum trānsferrī possit. Sī Mercurius essem alāsque habērem, ventō celerius trāns montēs amnēsque in Italiam volārem, ubi Venus, soror mea pulcherrima, \navnote{Venus est Mercurii soror (pater utriusque Iuppiter)}frātrem rīdēns reciperet.

\navnote{ef-fundere \contr{} ex + fundere cupidus patriae videndas = cupidus patriam videndī prae-stāre: officium p. z officium facere}Rīdēbis certē, mea soror, nec sine causā, nam rīdiculum est tālēs rēs optāre, nec lacrimās effundere mīlitem decet, cuius officium est sanguinem effundere prō patriā. Ego vērō, etsī cupidus sum patriae videndae, offi- 80 cium meum praestābō sīcut cēteri mīlitēs Rōmānī, quōrum magnus numerus in Germāniā est. Nisi nōs hīc \navnote{trāns-īre -iisse}essēmus finēsque imperiī dēfenderēmus, hostēs celeriter Dānuvium et Alpēs transirent atque ūsque in Italiam pervenīrent, nec vōs in Latiō tūtī essētis. Nē hoc fīat, 85 \navnote{tam-diū... quam-diū citrā e ultrā prp-racc: ultrā flūmen}legiōnēs Rōmānae hīc sunt ac tamdiū remanēbunt quamdiū hostis armātus seu citrā seu ultrā Dānuvium flūmen reliquus erit.

Quoniam igitur ipse ad tē properāre nōn possum, litterās ad tē scrībere properō. Quaeris ā mē cūr tibi \navnote{ūnī -ae -a: ūnae litterae = ūna epistula trini -ae -a: trinae litterae = trēs epistulae quaternae litterae = quattuor epistulae}ūnās tantum litterās scripserim, cum interim trinās quaternāsve litterās ā tē accēperim. Haud difficile est mē excūsāre, quod neglegēns fuerim in scribendō. S1 mihi \navnote{in epistulīs scribendis}tantum esset otii quantum est tibi, in epistulīs scrībendīs nōn minus dīligēns essem quam tū. Sed cum per complūrēs mēnsēs vix tempus habuerim ad dormiendum, facile intellegēs nūllum mihi ōtium fuisse ad epistulam scrībendam. Prope cotīdiē aut Germānī castra nostra oppugnāvērunt aut nōs impetūs fecimus in illōs. 100

Hodiē vērō nūllum hostem armātum citrā flūmen vidēmus. Magnus numerus eōrum aut caesus aut captus \navnote{caedere cecidisse caesum}est ā nostrīs, reliquī ultrā flūmen in magnīs silvīs latent. Heri enim exercitum Germānōrum proeliō vīcimus. Quod sīc factum est: Ill

Mediā nocte in castra nūntiātum est 'magnum hōstium numerum nāviculīs ratibusque cōpulātīs flūmen \navnote{secundum prp+acc: s. flūmen = flūmen sequēns con-vocāre = in eundem locum vocāre cum... cēpissent= --- postquam cēpērunt flūmen s}trānsiisse celeriterque secundum flūmen adversus cāstra nostra prōgredī.' Hōc nuntiō allātō, mīlitēs statim convocātī sunt. Quī cum arma cēpissent et vāllum ascendissent, primō mīrābantur quamobrem mediā nocte ē somnō excitāti essent, cum extrā vāllum omnia tranquilla esse vidērentur. Ego quoque dubitāre coeperam \navnote{secundum flūmen}num nūntius vērum dīxisset, cum subitō paulō ante \navnote{ante lūcem = ante diem ex-currere}lūcem magnus numerus Germānōrum ē silvīs proximis excurrēns castra nostra oppugnāvit. Nostrī, cum parātī essent ad castra dēfendenda, illum prīmum impetum \navnote{ad castra dēfendenda dē-sistere -stitisse etiam atque etiam = iterum atque iterum cōnstāns, adv cōnstanter pugnātur ā nostrīs = nostri pugnant ē-rumpere = subitō excurrere plēri- plērae- plēra-que = prope omnēs convertere -tisse -sum}facile sustinuērunt. Nec tamen hostēs castra oppugnāre dēstitērunt, sed aliī ex aliīs partibus etiam atque etiam sub vāllum prōcurrērunt. Cum complurēs hōrās ita fortissimē ā nostrīs, ab hostibus cōnstanter ac nōn timidē pugnātum esset, equitātus noster repente portā dextrā ērumpēns impetum in latus hostium apertum fecit. Paulō post, cum plērīque hostēs sē ad equitēs convertissent, peditēs nostri portā sinistrā ērūpērunt. Hostēs hāc rē perturbātī, cum iam longā pugnā fatigati essent, im\navnote{diū, comp diūtius, sup diūtissime tergum vertere:: fugere ripa -ae f --- lītus flūminis}petum Rōmānōrum ab utrāque parte venientium diūtius sustinēre nōn potuērunt, ac post brevem pugnam terga vertērunt. Cum ad ripam flūminis fugientēs pervēnissent, alil nāviculīs ratibusque sē servāvērunt, aliī armīs abiectīs in flūmen sē prōiēcērunt, ut natandō ad 130 \navnote{ulterior -ius comp \contr{} ultrā; sup ultimus citerior -ius comp \contr{} citrā; sup citimus caedēs -is f \contr{} caedere}ripam ulteriōrem pervenīrent, reliquī omnēs ab equitibus, quī ad eōs persequendōs missī erant, in ripā citeriōre caesi aut capti sunt. Tanta ibi caedēs hostium facta est ut meminisse horream.

Duo ferē milia hominum magnamque armōrum cō- 135 piam hostēs eō proeliō āmīsērunt. Ē nostrīs haud multī \navnote{dēsīderantur:: āmissi sunt vulnerāre \contr{} vulnus}dēsīderantur. Ipse sagittā in bracchiō laevō vulnerātus sum, sed vulnus meum leve est, multī graviōra vulnera accēpērunt. Nec vērō ūllus mīles legiōnārius ā tergō vulnerātus est. Plērique autem mīlitēs nostrī ex tantō proeliō incolumēs sunt.

\navnote{incolumis -e = salvus, integer}Hōc proeliō factō, dux victor, cum ā mīlitibus imperātor salūtātus esset, virtūtem nostram laudāvit 'quod \navnote{numerō superiōrēs}contrā hostēs numerō superiōrēs fortissimē pugnāvissēmus'; (tot hominibus ūnō proeliō āmissis, hostēs brevī 145 arma positūrōs esse' dīxit. Hīs verbīs māximō gaudiō affecti sumus, nam post longum bellum omnēs pācem \navnote{pāx pācis/ \contr{} f --- bellum lēgātus -Im = nūntius pūblicus}dēsideramus. Hodiē lēgātī ā Germānīs missī ad castra vēnērunt, ut cum imperātōre colloquerentur. Nesciō an lēgātī pācem petītum vēnerint, sed certō sciō im- 150 perātōrem nostrum cum hoste armātō colloquī nōlle.

Ego cum cēteris hāc victōriā glōriōsā gaudeō quidem, \navnote{glōriōsus: (rēs) glōriōsa = quae glōriam affert}sed multō magis gaudēerem si amīcus meus Pūblius Valerius, quōcum primum stipendium meruī, incolumis \navnote{quō-cum = cum quō}esset et mēcum gaudēre posset. Qul cum mihi in hostēs prōgressō auxilium ferre vellet, ipse ex aciē excurrēns \navnote{prō-gredī -ior -gressum esse. percutere -iō -cussisse -cussum ex vulnere = propter vulnus}pilo percussus cecidit. Graviter vulnerātus in castra portātus est, ubi in manibus meīs ex vulnere mortuus est, postquam mē ōrāvit ut per litterās parentēs suōs dē morte fili consolarer. Sed quōmodo aliōs cōnsōler, cum \navnote{ef-fundere -fudisse -fusum}ipse mē consolari nōn possim? Fateor mē lacrimās effūdisse cum oculōs eius clausissem, sed illae lacrimae et \navnote{et-enim = namque}mīlitem et amīcum deceēbant, etenim malus amīcus fuissem, nisi lacrimās effudissem super corpus amīcī mortuī, cum ille sanguinem suum prō mē effūdisset.

Utinam patrem audīvissem, cum mē ad studium litte\navnote{tum = tunc dēspicere = contemnere ōtiōsus -a -um \contr{} ōtium ōderam \contr{} amābam}rārum hortārētur! Sed tum litterās et studiōsōs litterārum dēspiciēbam. Poētās ut hominēs ōtiōsōs ōderam, praecipuē Tibullum, quī vītam rūsticam atque ōtiōsam laudābat, vītam mīlitārem dēspiciēbat. Mīrābar cūr ille \navnote{dīrus -a -um = terribilis ēnsis -is m = gladius dīrus}poēta mortem glōriōsam prō patriā 'dīram' vocāret et ēnsēs quibus patria dēfenditur 'horrendōs', ut in hīs versibus, quōs senex magister etiam atque etiam nōbīs recitābat:

\navnote{prō-ferre = in lūcem p., prīmum facere genus hominum = cūnctī hominēs proelia nāta sunt : brevior via ad diram mortem aperta est}Quis fuit horrendōs primus quī protulit ēnsēs?

Quam ferus et verē ferreus ille fuit!

Tum caedēs hominum generi \contr{} tum proelia nāta,

tum brevior dirae mortis aperta via est!

Ridiculi mihi puerō videbantur ii versūs, cum nōndum caedem vīdissem. At hodiē Tibullum vērum dīx- 180 isse intellegō. Si iam tum hoc intellēxissem, certē patrem audīvissem nec ad bellum profectus essem, ut tot caedēs et tot vulnera vidērem.

At iam satis dē caede scrīpsī. Nōlō tē fatīgāre nārrandō dē bellō cruentō, cum brevī pācem fore spērē- 185 \navnote{fore (īnf fut) = futūrum}mus. Nisi ea spēs mē fallit, posthāc plūrēs epistulās ā \navnote{exspect\&o/ = exspectā (posthāc)! scribirō/ = scribe! monētō/ = monē! nārrātōte/ = nārrāte!}mē exspectātō, atque plūrēs etiam ipsa scrībitō! Etiam aliōs monētō ut ad mē litterās dent. Et dē rē pūblicā et dē rē prīvātā nārrātōte mihi! Scītōte mē omnia quae apud vōs fiunt cognōscere velle.

\navnote{īdūs Maāiae: diēs xv mēnsis Māiī}Valētūdinem tuam cūrā dīligenter! Datum

pridiē īdūs Māiās ex castris.``

\booksection{Grammatica Latīna}

Coniūnctīvus Tempus plūsquamperfectum [A] Activum.

\navnote{laudāvisset}Iūlius dubitābat num magister Mārcum laudāvisset.

Laudāvisset'! coniūnctīvus est temporis praeteritī plūsquamperfectl. Coniūnctīvus plūsquamperfectī dēsinit in -isset \navnote{-isset}(pers. iii sing.) \contr{} quod ad Infinītīvum perfectī sine -isse adicitur.

\navnote{[3] scrīps isset}Exempla: [1] recitāv isset; [2] paruisset; [3] scripsisset; [4] audīvisset.

Pater fīlium interrogāvit 'num bonus discipulus fuisset: num magistrō pāruisset, attentē audīvisset, rēctē scripsisset et \navnote{retilāvissēmus rtciiāvlissēuis rtciiāv issent}pulchrē recilāvisset.' Fīlius nōn respondit. Tum pater ab eō quaesīvit 'nurn ā magistrō laudātus esset.' Cum filius id negavisset, ``Ergō'' inquit pater ``malus discipulus fuistī; nam sī \navnote{pārtt issē s pāru isse t}bonus discipulus fuissēs, magister tē laudāvisset. S1 magistrō pāruissēs, attentē audīvissēs, rēctē scripsisses et pulchrē recitāvissēs, ā magistrō laudātus essēs.`` Fīlius: ''Etiam s! \navnote{pāru [3] scrīps isselm scrīps issēls scrīps isset scrīps issē mus scnps issētis}industrius fuissem, magister mē nōn laudāvisset.`` Pater: ''Quid?`` Fīlius: ''Etiam si magistrō pāāruissem, attentē audīvissem, rēctē scrīpsissem et pulchrē recitavissem, laudātus nōn essem ab illō magistrō iniūstō!``

Pater filios interrogāvit 'num bonī discipulī fuissent: num \navnote{issem aud?}magistrō pztuissent, attentē audīvissent, rēctē scripsissent et pulchrē recitāvissent.' Filii nōn respondērunt. Tum pater ab iīs quaesīvit 'num ā magistrō laudātī essent.`` Cum filii id negāvissent, ''Ergō'' inquit pater ``malī discipulī fuistis; nam si bonī discipulī fuissētis, magister vōs lauddvisset. S1 magistrō pāruissētis, attentē zudīvissētis, rēctē scrīpsissētis et pulchrē recilāvissētis, ā magistrō laudātī essētis.'' Filii: Etiam si magistrō pztuissēmusy attentē audīvissēmus, rēctē scripsissemus et pulchrē reciiāvissēmus, laudātī nōn essēmus ab illō magistrō iniūstō!``

Persōna prima

-issem Persōna secunda ---

-issēs Persōna tertia

-isset [B] Passīvum.

Exempla: vidē suprā!

Persōna prīma

laudātas essem

laudātī essēmus Persōna secunda --- laudātus essēs laudātī essētis Persōna tertia

laudātus esset

laudārf essent

Imperātīvus

futuri Pugnāto, mīles! Pugnātōte, mīlitēs!

\navnote{pugnāto pugnātōte}Pugnātō, pugnātōte' imperāativus futūrī est. Imperātīvus \navnote{-tā -tōte}futūri dēsinit in -tō (sing.) et -tōte (plūr.); idem ferē significat atque imperātīvus praesenus ('pugnā, pugnāte').

Exempla:

Magister: ``Posthāc bonus discipulus estō, puer! Semper mihi pāretō/ Dīligenter auditō/ Pulchrē recitātō et rēctē scrībitō!''

Magister: ``Posthāc bonī discipulī estōte, pueri! Semper mihi parētōte! Dīligenter andītōte! Pulchrē recilātōte et rēctē sctībitōte!''

\booksection{Pēnsum A}

Magister epistulam ad Iūlium scripsit, cum Mārcus in lūdō dormīv- nec magistrō pāru-. Sī Mārcus bonus discipulus fu-5 magister eum laudāv- nec epistulam scrips-.

Mārcus: Heri in cubiculō inclūsus sum, cum pater epistulam tuam lēg-. Nisi tū eam epistulam scrips-, ā patre laudātus es-.`` Magister: ''Epistulam scripsi cum in lūdō dormīvnec mihi pāru-. Si industrius fu-, tē laudāv- nec epistulam scrips-.`` Mārcus: ''Etiam si industrius fu- et tibi pāru-, mē nōn laudāv-!``

Clēmēns es-, domine! Patientiam habē-! Servōs probōs laudā-, sed improbōs pūnī-!

Industriī es-, servī! Cum dominus loquitur, tacē- et audī-! Semper officium faci-!

Verba: Instruere -isse -um; cōgere -isse -um; caedere ---jsse ---um; convertere -isse -um; convertere-isse -um; percutere -isse -um; prōcurrere -isse; dēsistere -isse; prōgredī -um esse.

\booksection{Pēnsum B}

Ūna --- cōnstat ex v vel VI mīlibus hominum, quī in x --- dīviduntur. Exercitus prōcēdēns --- dīcitur. Exercitus ad --- [= pugnam] īnstrūctus --- appellātur. Post victōriam dux ā mīlitibus --- nomināātur. Si mīlitēs fortiter pugnāvērunt, --- eōrum ab imperātōre laudātur. Officium militis est sanguinem --- prō patriā.

Aemilius, quī in capitulō Xil --- est, ūnā cum Valeriō, adulēscente eiusdem ---, in Germāniā --- meruit. Tabellārius pūblicus xv --- [= circiter] diēbus Rōmā in Germāniam --- [= trīnī celeriter īre] potest. Difficile est Alpēs ---. Mīlitēs Rōmānī adiungere Dānuvium sunt, hostēs sunt --- Dānuvium.

\booksection{Pēnsum C}

Quae arma gerunt auxilia? Quid est signum legiōnis? Quōmodo mīlitēs in aciem īnstruuntur? Ad quod studium pater Aemilium hortābātur? Quō Aemilius adulēscēns missus est? Cūr Aemilius epistulās legēns permovētur? Quamobrem ipse paucās epistulās scripsit? \navnote{trānslre}Qul nūntius nocte in castra allātus est? Cūr hostēs castra Rōmāna nōn expugnāvērunt? Num Tibullus vītam mīlitārem laudat?

\booksection{Vocābula nova}
\noindent\small -celeriter, properāre \contr{} celeriter, īre/agere, re-ferre, issēs, issent, isselt, fu, -issem, -issēmus, -issēs, -issētis, -isset, -issent, essēmus, -tī, os, ta, essēs -tae essētis esset, essent, esset, -tum esset ---, p, pugnātōte, scribitōte, anāitōte, estō, estōte, nova, legiō, cohors, agmen, ōrdō, aciēs, proelium, imperātor, aetās, studium, stipendium, virtūs, gaudium, valētūdō, amnis, ratis, ripa, caedēs, vulnus, pāx, lēgātus, ēnsis, legiōnārius, idōneus, mīlitāris, pūblicus, privātus, posterus, arduus, rīdiculus, ulterior, citerior, incolumis, ōtiōsus, dīrus, horrendus, dēnī, sēnī, quīnī, quaternī, vulnerāre, fore, plērique, pridiē, praecipuē, tamdiū, quamdiū, diūtius, ferē, etenim, citrā, ultrā, secundum.

\bookchapter{Cap. XXXIV}{Dē arte poēticā}

Epistulā Aemiliī convivis recitāta, ``Ergō'' inquit Aemilia ``nōn modo filiō, sed etiam frātri meō bracchium vulnerātum est.''

Fabia: ``Quisnam vulnerāvit filium tuum?''

\navnote{pēs turgidus scalpellum -īn = parvus culter medicī}Aemilia: ``Medicus

bracchium

eius scalpellō suō s acūtō secuit! Quī, cum Quīntus heri dē altā arbore cecidisset, pedem eius turgidum vix tetigit, sed misellō \navnote{misellus -a -um = miser sanguinem mittere = vēnam aperire opera -ae/: f: o. alicuius = quod aliquis agit diīs = deīs (abl/dar)}puerō sanguinem mīsit! Profectō filius meus nōn operā medicī sānābitur, sed diīs iuvantibus spērō eum brevī sānum fore. Līberīne tuī bonā valētūdine ūtuntur?``

Fabia: ``Sextus quidem hodiē nāsō turgidō atque cru\navnote{certāre = pugnāre laedere -sisse -sum = pulsandō vulnerāre, inter-esse lūdus -Im \contr{} lūdere certāmen -inis n ( \contr{} certāre) = pugna}entō ē lūdō rediit, cum certāvisset cum quibusdam pueris quī eum laeserant.``

Aemilia: ``Cum pueri lūdunt, haud multum interest inter lūdum et certāmen. Mārcus vērō, si quis eum lae- 15 dere vult, ipse sē dēfendere potest.''

Fabia: ``Item Sextus scit sē dēfendere, dummodo cum singulīs certet. Nēmō sōlus eum vincere potest. Sed filius meus studiōsior est legendi quam pugnandī.'' 20 X Aemilia: ``Id dē Mārcō dīci nōn potest. Is nōn tam litteris studet quam lūdis et certāminibus!'' Hīc Cornēlius mulierēs interpellat: ``Ego quoque lūdīs et certāminibus studeō, dummodo aliōs certantēs spectem! Modo in amphitheātrō certāmen magnificum spectāvl: plūs trecentī gladiātōrēs certābant. Plērīque \navnote{plūs c = plūs quam c gladiātor -ōris m = vir quī populō spectante cum alterō pugnat \contr{} gladiātor im-plicāre -āvisse/-uisse -ātum/-itum}gladiīs et scūtīs armātī erant, aliī retia gerēbant.`` Aemilia, quae certāmen gladiātōrium nōn spectāvit, ā Cornēliō quaerit 'quōmodo gladiātōrēs rētibus certent?' Cornēlius: ''Alter alterum in rēte implicāre cōnātur, nam quī rētī implicitus est nōn potest sē dēfendere et \navnote{spectātor -ōris m = quī spectat plērum-que = prope semper}sine morā interficitur, nisi tam fortiter pugnāvit ut spectātōrēs eum vīvere velint. Sed plērumque is quī victus est occīditur, victor vērō palmam accipit, dum spectātō-

\navnote{plaudere -sisse -sum}Iūlius: ``Mihi nōn libet spectāre lūdōs istōs ferōcēs. \navnote{libēre: libet grātum est, placet}Mālō cursūs equōrum spectāre in circō.``

\navnote{circēnsis -c \contr{} circus; m pl lūdi circēnsēs iuvāre = dēlcctāre tum auriga -ae m = qul currum regit velle lūgēre = maerēre (ob alicuius mortem) scaenicus -a -um \contr{} scaena cōmoedia -ae f --- fabula scaenica dē rēbus levibus et ridiculis Plautus -īm: poēta quī cōmoediās scripsit Amphitryōn -ōnis m ācer ācris ācre = cupidus agendī, impiger gerere = agere}Cornēlius: ``Lūdī circensēs mē nōn minus iuvant quam gladiatorit: modo in amphitheātrum, modo in circum eō. Sed ex novissimīs circēnsibus maestus abiī, cum ille aurīga cui plērīque spectātōrēs favēbant ex 40 currū lāpsus equis prōcurrentibus occisus esset. Illō occlsō spectātōrēs plaudere dēsiērunt ac lūgēre coepērunt.''

Fabia: ``Ego lūdōs scaenicōs praeferō: mālō fabulās spectāre in theātrō. Nūper spectāvī cōmoediam Plautī 45 dē Amphitryōne, duce Graecōrum, cuius uxor Alcmēna ab ipsō Iove amābātur. Dum dux ille ācer et fortis prōcul ā domō bellum gerit, Iuppiter sē in fōrmam eius \navnote{geminus -a -um: (filiī) geminī= --- duo eōdem diē nāti parere peperisse partum}mūtāvit, ut Alcmēnam vīseret; quae, cum putāret con-

peperit, quōrum alter fuit Herculēs. Illa fabula omnibus nōta est. Sed scitisne cūr fēminis libeat in theātrum \navnote{ingenium -īn = nātūra animī ratiō -ōnis f --- causa quae rem plānam facit reddere:: dare}īre? Ovidius poēta, quī ingenium mulierum tam bene nōverat quam ipsae mulierēs, ratiōnem reddit hōc 55 versū:

\navnote{ut ipsae spectentur bellus -a -um = pulcher poēticus -a -um \contr{} poēta prīncipium -In = prima pars, initium}Spectātum veniunt, veniunt spectentur ut ipsae!``

Fabia: ``Et viri veniunt ut bellās feminās spectent!''

Tum Iūlius, quī artis poēticae studiōsus est, ``Ovidius IJ ipse'' inquit ``id dīcit. Ecce principium carminis quod scripsit ad amīcam sēcum in circō sedentem:

\navnote{nōbilium equōrum tamen precor ut vincat ille auriga cui tū ipsa favēs}Nōn ego nōbilium sedeō studiōsus equōrum;

cui tamen ipsa favēs vincat ut ille precor. Ut loquerer tēcum vēnī tecumque sederem,

\navnote{nōn nōtus:: ignōtus amor quem facis (: excitās)}nē tibi nōn nōtus quem facis esset amor.

Tū cursūs spectās, ego tē --- spectemus uterque

quod iuvat, atque oculōs pāāscat uterque suōs!`` \navnote{pāscere:: dēlectāre}\navnote{bellum canere = dē bello c. (: carmen scrlbere) fatum -īn- quod necesse est accidere, mors meae puellae dīxī gremium -īn = genua sedentis; in gremiō meō sedēre sēdisse retinēre -uisse -tentum}Cornēlius: ``Ego memoriā teneō versūs Ovidii de puellā quae poētam industrium prohibēbat bellum Trōiānum canere et fatum rēgis Priamī: Saepe meae ''Tandem'' dīxī ``discēde!'' puellae

--- in gremiō sēdit prōtinus illa meō! Saepe Pudet'' dīxī. Lacrimīs vix illa retentīs

``Mē miseram! Iam tē'' dīxit ``amāre pudet?'' Implicuitque suōs circum mea colla lacertōs

\navnote{suōs lacertōs colla:: collum mihi dedit ingenium meum ab sūmptīs armīs mea bella: meōs amōrēs SM Z}et, quae me perdunt, ōscula mille dedit! Vincor, et ingenium sūmptīs revocātur ab armīs,

rēsque domi gestās et mea bella canō.`` Fabia: ''Iste poēta viris sōlīs placet. Mē vērō magis iuvant carmina bella quae Catullus scripsit ad Lesbiam amīcam. Si tibi libet, Iūlī, recitā nōbīs ē librō Catullī.`` \navnote{libenter (adv \contr{} libēns) = cum gaudiō tenebrae -ārum f \contr{} lūx accendere -disse -ēnsum}``Libenter faciam'' inquit Iūlius, ``sed iam nimis obscūrum est hoc triclīnium; in tenebris legere nōn possum. Lucernās accendite, servī!'' 85 Lucernīs accēnsīs, Iūlius librum Catullī prōferri iubet, tum ``Incipiam'' inquit ``ā carmine dē morte passeris quem Lesbia in deliciis habuerat: \navnote{passer VVS -eris m 4 dēliciae -ārum f - id quod dēlectat Venerēs, Cupidines pi : Venus, Cupīdō quantum est hominum = quot sunt hominēs venustus -a -um = bellus, amandus}Lūgēte, ō Venerēs Cupīdinēsque et quantum est hominum venustiōrum! Passer mortuus est meae puellae, passer, dēliciae meae puellae, \navnote{[Catullus3]}quem plūs illa oculīs suīs amābat; nam mellītus erat suamqne nōrat \navnote{= dulcis (ut mel) nōrat = nōuerat ipse m = dominus; ipsa/ f = domina: suam ipsam}ipsam tam bene quam puella mātrem, nec sēsē ā gremiō illīus movēbat, sed circumsiliens modo hūc modo illūc \navnote{circum-sillre \contr{} -salīre ūsque = semper plpiāre = 'pipi' facere}ad sōlam dominam ūsque pīpiābat. Quī nunc it per iter tenebricōsum \navnote{nebrae) = obscūrus negant redīre quemquam = dīcunt 'nēminem redīre' Orcus -i m- Plūtō: īnferi dē-vorāre = vorāre}illūc unde negant redīre quemquam. At vōbīs male sit, malae tenebrae Ora, quae omnia bella dēvorātis: tam bellum mihi passerem abstulistis. Ō factum male! Ō miselle passer! \navnote{male factum = maleficium tuā operā = tuā causā turgidulus = turgidus ocellus -Im = (parvus) oculus ultimus -a -um (sup \contr{} ultrā) = postrēmus}Tuā nunc operā meae puellae flendō turgidult rubent ocellī! ``Hīs versibus ultimīs poēta vēram ratiōnem dolōris suī reddit: quod oculī Lesbiae lacrimīs turgidī erant ac rubentēs! Tunc enim Catullus Lesbiam sōlam amābat \navnote{perpetuus -a -um = quī numquam finiētur,sine. fine mēns mentis f --- animus (cōgitāns) [Catullus}atque amōrem suum perpetuum fore crēdēbat. Ecce aliud carmen quō mēns poētae amōre accēnsa dēmōns- 110 trātur: Vīvāmus, mea Lesbia, atque amemus, rūmōrēsque senum sevēriōrum omnēs ūnīus aestimēmus assis! \navnote{ūnius assis aestimāre: minimi aestimāre, nihil cūrāre}Sōlēs occidere et redīre possunt --- nōbīs, cum semel occidit brevis lūx, \navnote{: nōbīs dormiendum est ūnam noctem perpetuam}nox est perpetua ūna dormienda.

Dā mī bāsia mīlle, deinde centum,

\navnote{bāsium -īn = ōsculum dein = deinde ūsque = sine fine}dēin mille altera, dein secunda centum \contr{}

deinde ūsque altera mille, deinde centum!

Dēin, cum mīlia multa fecerimus,

\navnote{con-turbāre = turbāre nē numerum sciāmus nōbīs invidēre tantum bāsiōrum = tot bāsia uxōrem dūcere = uxōrem suam facere affirmāre = certō dīcere nūbere -psisse (+ dat): alicui n. = alicuius uxor fieri}conturbābimus illa, ne sciāmus,

aut nē quis malus invidēre possit,

cum tantum sciat esse bāsiōrum.

``Catullus Lesbiam uxōrem dūcere cupiēbat, nec vērō illa Catullō nūpsit, etsī affirmābat 'sē nūllī aliī virō nūbere mālle.' Mox vērō poēta dē verbīs eius dubitāre coepit:

'Nūulli s? dīcit mulier mea 'nūbere mālle

\navnote{uxōrem petat id quod mulier dīcit amantī}quam mihi, nōn si sē Iuppiter ipse petat!

Dīcit. Sed mulier cupīdō quod dīcit amantī

in ventō et rapidā scribere oportet aquā!

``Postrēmō poēta intellēxit Lesbiam īnfidam et amōre suō indignam esse, neque tamen dēsiit eam amāre. Ecce \navnote{cupīdō dolēre (animō) \contr{} gaudēre certus, dubitāns odium -i \& amor}duo versūs quī mentem poētae dolentem ac dubiam inter amōrem et odium dēmōnstrant:

Ōdī et amō! Quārē id faciam, fortasse requīris?

\navnote{re-quīrere = quaerere ex-cruciāre = valdē cruciāre}Nesciō; sed fieri sentiō --- et excrucior!``

Hīs versibus recitātīs convīvae diū plaudunt.

Tum Paula ``Iam satis'' inquit ``audīvimus dē amandō \navnote{risus -ūs m \contr{} ridēre iocōsus -a -um = quī risum excitat, ridiculus}et dē dolendō. Ego ridēre malo, neque iste poēta risum excitat. Quīn versūs iocōsōs recitās nōbīs?``

Cui Iūlius ``At Catullus'' inquit ``nōn tantum carmina sēria, sed etiam iocōsa scripsit. Ecce versūs quibus po\navnote{[Catullus 13.1-8 \contr{} v MESES Mises ``uf arānea SA}ēta pauper quendam amīcum dīvitem, nōmine Fabullum, ad cēnam vocāvit:

Cēnābis bene, mi Fabulle, apud mē

paucis, si tibi di favent, diēbus

--- si tēcum attuleris bonam atque magnam

\navnote{nōn sine:: cum candidā:: pulchrā sāl:: sermō iocōsus et urbānus cachinnus -īm = risus inquam (pers I) = dīcō venuste noster:: mi amīce}cēnam, nōn sine candidā puellā

et vīnō et sale et omnibus cachinnis.

Haec si, inquam, attuleris, venuste noster,

cēnābis bene --- nam tuī Catullī

plēnus sacculus e\$t arāneārum!``

\navnote{(risum) movēre:: excitāre}Hi versūs magnum risum movent. Tum vērō Cornēlius ``Bene quidem'' inquit ``et iocōsē scrīpsit Catullus, \navnote{epigramma -ātis n (pl abl -atīs) Mārtiālis -is m: poēta Rōmānus quī XII librōs epigrammatum scripsit ``ÁS Ee Y sinus}nec tamen versūs eius comparandī sunt cum epigrammatīs sale plēnīs quae Mārtiālis in inimīcōs scripsit. Semper librōs Mārtiālis mēcum in sinu ferō.``

Ab omnibus rogātus ut epigrammata recitet, Cornē- 160 lius libellum ēvolvit et ``Incipiam'' inquit ``ā versibus quōs poēta de suīs libellis scrīpsit:

Laudat, amat \contr{} cantat nostrōs mea Rōma libellōs,

\navnote{nostrōs (: meōs) libellōs}meque sinūs omnēs, me manus omnis habet.

Ecce rubet quīdam, pallet, stupet, ōscitat, ōdit.

\navnote{ōscitāre = ōs aperire (ob dormiendi cupiditātem) nōbīs, nostra: mihi, mea}Hoc volō: nunc nōbīs carmina nostra placent''

Post hoc prīncipium Cornēlius aliquot epigrammata Mārtiālis recitat, in iīs haec quae scripta sunt in aliōs poētās:

\navnote{poētās: Pontiliānum, Cinnam, Māmercum}Cūr nōn mitto meōs tibi, Pontiliāne, libellōs?

Nē mihi tā mittāsy Pontiliāne, tuōs! --- Versiculōs in me nārrātur scribere Cinna. \navnote{versiculus -īm = versus (parvus) nōn scribit is nīl = nihil}Nōn scribit, cuius carmina nēmō legit! --- Nīl recitās --- et vis, Māmerce, poēta vidērī. Quidquid vis, estō --- dummodo nīl recitēs! Ecce alia epigrammata Mārtiālis quae Cornēlius convivis attentis atque dēlectātīs recitat: Nōn amō tē, Sabidī, nec possum dīcere quāāre. Hoc tantum possum dīcere: nōn amō tē! --- --- \navnote{tēdatūrum esse fata:: fatum (: mortem) Marō -ōnis m ā mē petis (= mē rogās)}Nīl mihi dās vīvus, dīcis ``post fata datūrum.' Sī nōn es stultus, scis, Maro, quid cupiam! --- 'Esse nihil dīcis, quidquid petis, improbe Cinna. Sī nīl, Cinna \contr{} petis \contr{} nīl tibi, Cinna, negō! --- Nesciō tam multīs quid scribās, Fauste, puellīs. \navnote{tam multīs puellīs(= (7 tot puellīs) E}Hoc sciō: quod scrībit nūlla puella tibī! Sequuntur epigrammata quibus dēridentur feminae, praecipuē anūs, ut Laecānia et Paula: \navnote{--- femina vetus Thāis -idis f = candidus ut nix haec:: Laecānia illa:: Thāis sapere -iō -iisse = sapiēns esse; -istī = -iistī dūcere = uxōrem dūcere nōbīs:: mihi magis anus (: brevī moritūra!)}Thāis habet nigrōs, niveōs Laecānia dentēs. Quae ratiō est? Émptōs haec habet, illa suōs! --- Nūbere vis Prisco; nōn mīror, Paula,

sapīstī. Dūcere tē nōn vult Prīscus: et ille sapit! --- Nūbere Paula cupit nōbīs, ego dūcere Paulam nōlō: anus est; vellem si magis esset anus! Cēterīs rīdentibus ``Quid

ridētis?`` inquit

Paula, ``Num haec in mē, uxōrem fōrmōsam atque puellam, \navnote{puella anus}scrīpta esse putātis?`` Cornēlius: ''Minimē, Paula. Nec scīlicet in tē, sed in Bassam scriptum est hoc:

Dīcis 'formosam'), dīcis te, Bassa,

'puellam.

\navnote{dīcis tē esse fōrmōsam' ea quae fōrmōsā nōn est ē-rubēscere = rubēre incipere, rubēns fieri}Istud quae nōn est, dīcere, Bassa, solet!``

Hoc audiēns ērubēscit Paula atque cēteri convīvae risum vix tenent. Cornēlius vērō prūdenter ``Nōn omnia'' inquit ``iocōsa sunt carmina Mārtiālis. Ecce duo versūs dē fatō virī pauperis, et quattuor in Gelliam, quae cōram testibus lacrimās effundit super patrem 205 \navnote{testis -is m = quī adest}mortuum:

Semper pauper eris, si pauper es, Aemiliāne.

Dantur opēs nūllīs nunc nisi dīvitibus. ---

\navnote{opēs -um / pl = dīvitiae G. patrem āmissum (: dē patre āmissō/mortuō) nōn flet prō-silire \contr{} -salīre quisquis:: is quī quaerit:: cupit}Āmissum nōn flet, cum sōla est, Gellia patrem;

si quis adest, iussae prōsiliunt lacrimae!

Nōn lūget quisquis laudāri, Gellia, quaerit:

ille dolet vērē quī sine teste dolet.``

Centum ferē epigrammatīs recitātīs, Cornēlius cum hōc fīnem facit recitandī:

Cui lēgisse satis nōn est epigrammata centum,

nīl illī satis est, Caediciāne, malī!

\navnote{nīl malī plaudere + dat}Rīdent omnēs et Cornēliō valdē et diū plaudunt.

\booksection{Grammatica Latīna}

\booksection{Dē versibus}

[I] Syllabae brevēs et longae. Nōn ego nōbilium sedeō studiōsus equorum.

Hic versus cōnstat ex hīs syllabīs: nō-:i e-go- nō-bi-li-um- \$ede-ō- stu-di-ō-su-\$ e-quō-rum.

Syllaba brevis est quae in vōcālem brevem \contr{} a \contr{} e, i, o, u, y) dēsinit; quae dēsinit in vōcālem longam (4, \contr{}, 1, 0, 4, y) aut in diphthongum (ae, oe, au, eu, ei, ui) aut in cōnsonantem \navnote{diphthongus -i f 2 duae vōcālēs in ūnā syllabā coniūnctae}c, d,f,g,l, m,n, p, v s, t, x) syllaba longa est. Syllabae brevēs: ne, goy bi, li, se, de, stu, di, su, se; syllabae longae: nō, um, 0, quō, rum.

\navnote{notae: - syllaba brevis --- syllaba longa inter-dum = nōnnumquam}Haec nota H syllabam brevem significat, haec [---] sylla-

Cōnsonantēs

brygry cr, try quae initium syllabae facere so-lent (ut li-bri), interdum dīviduntur: nig-rōs, pat-rem.

Vōcālis -ō ultima interdum fit brevis: vo-lō, ne-mō. [II] Syllabae coniūnctae.

Litterae vocābulōrum ultimae cum vōcālibus sequentibus coniunguntur hīs modīs: [A] Cōnsonāns ultima cum vōcālī prīmā (vel h-) vocābuli sequentis ita coniungitur ut initium syllabae faciat: a-nu-\$ est; \navnote{haec nota [o] significat litterās in syllabā coniungendās}vin-ca-Pu-Pil-le. [B] Vōcālis ultima (item -am, -em, -um, -im) ante vōcālem primam (vel h-) vocābull sequentis ēlīditur: Vīvāmus, mea Lesbia, atque amemus: Lesbi'atqu'amēmus In est et es ēlīditur e sōla est: sōla'sty vērum e\$t: venmC\%ty bella es: bella's. [IIT] Pedēs.

Singulī versūs dīviduntur in pedēs, quī bīnās aut ternās syllabās continent. Pedēs frequentissimī sunt trochaeī, iambī, dactylī, spondet. Trochaeus cōnstat ex syllabā longā et brevī, ut /z-na, iambus ex brevī et longā, ut vi-ri, dactylus ex longā et duābus brevibus, ut fe-mi-na, spondēus ex duābus longīs, ut ne-mō. [IV] Versus hexameter. Nōnego nōbilium sedeō studiōsus equōrum.

\navnote{haec nota [] inter pedēs pōnitur}Hic versus hexameter vocātur ā numerō pedum, nam sex Graecē dīcitur hex. Hexameter cōnstat ex quīnque pedibus dactylīs et ūnō spondēō (vel trochaeō); prō dactylīs saepe spondēī inveniuntur, sed pēs quīntus semper dactylus est: uu I uul u I uu I uul Done eris felix multōs numerābis amīcōs. \navnote{[Horātius: Ars poetica333,}Aut prōdesse volunt aut dēlectāre poētae. Caelum, nor animum māūtant quī trāns mare currunt. [V] Versus pentameter.

cui tamelrī ipsa favēs vincaffille precor.

Hic versus dīviditur in duās partēs, quārum utraque bīnōs pedēs et dīmidium continet. Ita versus tōtus ex quīnque pedibus cōnstat et vocātur pentameter: quīnque enim Graecē dīcitur pente. Utraque pars est ut initium hexametri, sed pars posterior spondēōs nōn admittit:

--- nābila

sōlus?ris.

In gremiō sēdit protinu\$illa meō.

Versus pentameter semper hexametrum sequitur: Nō amō tā, Sabidī, nec possum dīcere quārē. Hoc tantum possum dīcere norī amō tē. [VI] Versus hendecasyllabus.

Passer mortuufest meae puellae.

Hic versus, qul ūndecim syllabās habet, vocātur hendecasyllabus, nam ūndecim Graecē hendeca dīcitur. Versus hendecasyllabus in quīnque pedēs dividi potest: spondēum, dactylum, duōs trochaeōs, spondēum aut trochaeum (pēs primus rārius iambus aut trochaeus est).

E Vīvāmus mea Lesbi'atqu'a mēmus. Cēnābis bene mi Fabull'apud mē

paucis si tibi di favent diēbus.

\booksection{Pēnsum A}

Hī versūs in syllabās brevēs et longās et in pedēs dividendi sunt notīs appositīs: Scrībere mē quereris, Vēlōx, epigrammata longa.

Ipse nihil scribis: tū, breviōra facis! [Mārtiālis 1.110] Dās numquam, semper prōmittis, Galla, rogant.

Si semper fallis, iam rogō, Galla, negā! [11.25] Quem recitās meus est, ō Fīdentīne, libellus.

Sed male cum recitās, incipit esse tuus! [1.38]

Bella es, nōvimus, et puella, vērum est,

et dīves, quis enim potest negāre?

Sed cum tē nimium, Fabulla, laudās,

nec dīves neque bella nec puella es! [1.64]

\booksection{Pēnsum B}

Dum gladiātōrēs --- [= pugnant], --- dēlectātī manibus ---. Fēminis nōn --- gladiātōrēs spectāre.

Ovidius in --- sedēns ōrābat ut vinceret ille cui amica eius iocōsus ---. Hoc nārrātur in --- --- [=initiō] carminis. Cum poēta --- [= niveus fortūnam] Priamī canere vellet, puella in --- eius sēdit eīque c mīlle--- --- dedit. Ovidius --- mulierum bene nōverat.

---, Iūlius recitat carmen --- [= pulchrum] dē --- plaudere mīs] dēmōnstrātur --- [= causa] dolōris. Catullus Lesbiam lūgēre uxōrem --- cupiēbat, at illa Catullō --- nōluit. --- poētae inter circumsilīre amōrem et --- dīvidēbātur.

\booksection{Pēnsum C}

Quid Rōmaāni in amphitheātrō spectant? Quid Fabia in theātrō spectāvit? Quis fuit Ovidius? Quārē Ovidius in circum vēnerat? Quae carmina recitat Iūlius? \navnote{libenter plērumque}Cūr ocellī Lesbiae turgidī rubēbant? l Cūr poēta Lesbiam et amābat et ōderat? Num sacculus Catullī plēnus erat nummōrum? Quid scripsit Mārtiālis? Cūr librōs suōs nōn mīsit Pontiliānō? t Tune Cinnam bonum poētam fuisse putās? Cūr Laecāniae dentēs niveī erant? s Ex quibus syllabīs cōnstat pēs dactylus? Ex quibus pedibus cōnstat hexameter?

\booksection{Vocābula nova}
\noindent\small Vocābvlanava, scalpellum, opera, lūdus, certāmen, gladiātor, rēte, spectātor, palma, circus, aurīga, theātrum, cōmoedia, ingenium, ratiō, prīncipium, fatum, gremium, tenebrae, lucerna, passer, dēliciae, ocellus, mēns, bāsium, odium, risus, cachinnus, arānea, epigramma, sinus, versiculus, anus, testis, opēs, diphthongus, nll, trochaeus, iambus, spondēus, hexameter, pentameter, hendecasyllabus.

\bookchapter{Cap. XXXV}{Ars grammatica}

[Ex Dōnātī 'Arte grammaticā minōre']

\booksection{Dē partibus ōrātiōnis}

\navnote{grammaticus -a -um: ars}\navnote{culō iv post Christum nātum ōrāātiō -ōnis f --- sermō}[Magister:] Partēs ōrātiōnis quot sunt? [Discipulus:] Octō.

[D.:] Nōmen, prōnōmen, verbum, adverbium, partici- 5 pium, coniūnctiō, praepositiō, interiectiō.

\booksection{Dē nōmine}

\navnote{coniūnctiō, interiectiō:}[M.:] Nōmen quid est? \navnote{cāsus (nōminis): v. īnfrā \contr{} commūnis nōmina propria, ut : Roma, Tiberis appellātīvus -a -um \contr{} appellāre; nōmina appellātīva, uturbs, flūmen}[D.:] Pars ōrātiōnis cum cāsū, corpus aut rem propriē

et 'proprium' dīcitur, aut multōrum et 'appellātīvum'.

[M.:] Genera nōminum quot sunt? [D.:] Quattuor.

[D.:] Masculīnum, ut hic magister, femininum, ut haec \navnote{vem deae quae singulīs artibus praesunt scamnum -n = sella sacerdōs -ōtis mJf---vir /femina cuius negōtium est diīs servīre}Mūsa, neutrum, ut hoc scamnum, commūne, ut hic et haec sacerdōs. [M.:] Numerī nōminum quot sunt? [D.:] Singulāris, ut hic magister, plūrālis, ut A? magistri.

[M.:] Cāsūs nōminum quot sunt?

[D.:] Nōminātīvus, genetīvus, datīvus, accūsātīvus, vocātīvus, ablātīvus. Per hōs omnium generum nōmina,

[M.:] Comparātiōnis gradūs quot sunt?

\navnote{comparātiō -ōnis f \contr{} comparāre}[D.:] Positīvus, ut doctus, comparātivus, ut doctior, superlātīvus, ut doctissimus. [M.:] Quae nōmina comparantur? \navnote{appellātīva: scīlicet adiectīva dumtaxat = tantum quālitās -ātis f \contr{} quālis quantitās -ātis f \contr{} quantus}[D.:] Appellātīva dumtaxat quālitātem aut quantitātem

\booksection{Dē prōnōmine}

[M.:] Prōnōmen quid est? [D.:] Pars ōrātiōnis quae prō nōmine posita tantundem

\navnote{tantun-dem ( \contr{} tantumdem) = idem}[M::] Genera prōnōminum quae sunt? [D.:] Eadem ferē quae et nōminum: mascullnum, ut quis, feminīnum, ut quae, neutrum, ut quod, commūne, ut quālis, tālis, trium generum, ut ego, tū. [M.:] Numerī prōnōminum quot sunt?

\navnote{Prōnōmina sunt [1] persōnālia: ego, tā, fis, VŌS, se [2] possessīva: meus, tuus, suus, noster, vester [3] dēmōnstrātīva: hky iste, ille, i\$, īdem, ipse [5] interrogātīva: quis/ quī, uter [6] indēfinīta: aliquis/ -quī, quis, quisquam, quisque, uterque, quīdam, nēmō, nihil, neuter}[D.:] Singulāris, ut hic, plūrālis, ut hī. [M.:] Persōnae prōnōminum quot sunt?

[D.:] Prima, ut ego, secunda \contr{} ut tū, tertia, ut ille. [M.:] Cāsūs item prōnōminum quot sunt? [D.:] Sex, quem ad modum et nōminum, per quōs om- 55

\booksection{Dē verbō}

\navnote{quem ad modum = sīcut in-flectere = dēclināre}[M.:] Verbum quid est? [Z).:] Pars ōrātiōnis cum tempore et persōnā, sine cāsū,

\navnote{neutrum: nec agere nec}[M.:] Modī verbōrum quī sunt? \navnote{pau optātīvus -a -um \contr{} optāre}[D.:] Indicātīvus, ut /egō, imperātīvus, ut lege, optātlvus, ut utinam legerem, coniūnctīvus, ut cum legam, īnfinītīvus, ut legere. [M.:] Genera verbōrum quot sunt? [D.:] Quattuor.

[D.:] Āctīva, passīva, neutra, dēpōnentia. [M.:] Āctīva quae sunt? [D.:] Quae in -ō dēsinunt et acceptā -r litterā faciunt ex 70 \navnote{accepta:: additā faciunt ex se passīva}sē passīva, ut legō: legor. [M.:] Passīva quae sunt? \navnote{dēmere -mpsisse -mptumaddere}[D.:] Quae in -r dēsinunt et eā dēmptā redeunt in āctīva, ut legor: legō. [M.:] Neutra quae sunt? [D.:] Quae in -ō dēsinunt ut āctiva, sed acceptā -r litterā \navnote{Latina: rēcta}Latīna nōn sunt, ut s \contr{} ō, currō (``stor, curror'' nōn dīcimus!). [M.:] Dēpōnentia quae sunt? [D.:] Quae in -r dēsinunt ut passīva, sed eā dēmptā Latīna nōn sunt, ut luctor, loquor. \navnote{luctārī= --- certāre complectendis corporibus}[M.:] Numeri verbōrum quot sunt?

[D.:] Singulāris, ut /ego, plūrālis, ut legimus. [M.:] Tempora verbōrum quot sunt?

[D.:] Praesēns, ut legō, praeteritum, ut \%ī, futūrum, ut legam. [M.:] Quot sunt tempora in dēclinātione verbōrum? [D.:] Quīnque.

[D.:] Praesēns, ut legō, praeteritum imperfectum, ut legēbam, praeteritum

perfectum,

ut legi, praeteritum plūsquamperfectum, ut lēgeram, futūrum, ut legam. [M.:] Persōnae verbōrum quot sunt?

[D.:] Prīma, ut /egō, secunda, ut legis, tertia, ut legit.

\booksection{Dē adverbiō}

[M.:] Adverbium quid est? \navnote{ad-icere -iēcisse -iectum significātiō -ōnis f \contr{} significāre ex-plānāre = plānum facere}[D.:] Pars ōrātiōnis quae adiecta verbō significātiōnem

[M.:] Significātiō adverbiōrum in quō est? [D.:] [ \contr{}.:]Sunt aut locī adverbia aut temporis aut numeri aut negandī aut affīrmandī aut dēmōnstrandī aut optandī aut hortandī aut ōrdinis aut interrogandī aut quālitātis

[M.:] Dā adverbia locī! [D.:] Ut hic vel ibi, intus vel foris, illūc vel inde. [M.:] Dā temporis! [D.:] Ut hodiē, nunc, nūper, crās, aliquandō; numeri, ut semel, bis, ter; negandī, ut nōn; affīrmandī, ut etiam, quidnī; dēmōnstrandī, ut ēn, ecce; optandī, ut utinam; \navnote{quid-nī = quīn, certē}hortandī, ut eia; ōrdinis, ut deinde; interrogandī, ut cūr, quārēy quamobrem; quālitātis, ut doctē, pulchrē, fortiter; \navnote{forsitan = fortasse}quantitātis, ut multum, parum; dubitandī, ut forsitan,

[M.:] Comparātiō adverbiōrum in quō est? [D.:] In tribus gradibus comparātiōnis: positīvō, comparātīvō, superlātīvō. [M.:] Dā adverbium positivi gradūs! [D.:] Ut doctē; comparātīvī, ut doctius; superlativi,

ut

\booksection{Dē participiō}

[M.:] Participium quid est? [D.:] Pars ōrātiōnis partem capiēns nōminis, partem \navnote{parti-cipium \contr{} pars}verbī: nōminis genera et cāsūs, verbī tempora et signifi-

[M.:] Genera participiōrum quot sunt? [D.:] Quattuor.

[D.:] Masculīnum,

ut hic lēctus, feminīnum, ut haec \navnote{lēctus -a -um part perf \contr{} legere}lēcta, neutrum, ut hoc lēctum, commūne tribus generibus, ut hic et haec et hoc legēns. [M.:] Cāsūs participiōrum quot sunt?

[D.:] Nōminātīvus, ut hic legēns, genetīvus, ut huius legentis, datīvus, ut huic legentī, accūsātīvus, ut hunc legentem, vocātīvus, ut ō legēns, ablātīvus, ut ab hōc legente. [M.:] Tempora participiōrum quot sunt?

[D.:] Praesēns, ut legēns, praeteritum,

ut Jectus, futū-

\navnote{legendus dīcitur 'gerundīvum' vel (participium futūri passīvī'}[M.:] Numeri participiōrum quot sunt?

\booksection{Dē coniūnctiōne}

\navnote{con-iūnctiō -ōnis f \contr{} coniungere ad-nectere = adiungere, coniungere ōrdināre = in ōrdine pōnere potestās = significātiō speciēs -ei f --- forma, genus}[M.:] Coniūnctiō quid est? [D.:] Pars ōrātiōnis adnectēns ōrdinānsque sententiam.

[M.:] Potestās coniūnctiōnum quot speciēs habet? [D.:] Quīnque.

\navnote{cōpulātīvās, cēt.,}[D.:] Cōpulātīvās, disiūnctivās, explētīvās, causālēs, ratiōnālēs. [M.:] Dā cōpulātīvās! \navnote{\contr{} cōpulāre}[D.:] Et, -que, atque, ac. [M.:] Dā disiūnctīvās! \navnote{\contr{} disiungere}[D.:] Aut, -ve, vel, nec, neque. [M.:] Dā explētīvās! \navnote{explētīvus -a -um \contr{} explēre = vacuum implēre}[D.:] Quidem, equidem, quoque, autem, tamen. \navnote{causālis -e \contr{} causa sī-quidem = quoniam quandō = quoniam}(M.:] Dā causālēs! [D.:] St, etst, sīquidem, quandō, nam, namque, etenim \contr{}

\navnote{ratiōnālis -e \contr{} ratiō quā-propter = quam obrem propter-eā = ideō}[M.:] Dā ratiōnālēs! [D.:] Itaque, enim, quia, quāpropter, quoniam, ergō, ideō,

\booksection{Dē praepositiōne}

\navnote{praepositiō -ōnis f \contr{} prae-pōnere}[M.:] Praepositiō quid est? [D.:] Pars ōrātiōnis quae praeposita aliīs partibus ōrātiōnis significātiōnem eārum aut complet aut mūtat aut minuit. [M.:] Dā praepositiōnēs cāsūs accūsātīvī! \navnote{adversum = adversus cis = citrā}[D.:] Ad, apud, ante, adversum, cis, citrā, circum, circā, contrā, ergā,

extrā \contr{} inter, intrā,

īnfrā, juxtā,

ob, per, prope, secundum, post, trāns,

ultrā,

praeter,

propter,

[M.:] Quō modō? [D.:] Dīcimus enim ad patrem, apud vīllam, ante do-mum, adversum inimīcōs, cis Rhenum, citrā forum,

circum oppidum, circā templum, contrā hostem, ergā parentēs, extrā vāllum, inter nāvēs, intrā moenia, īnfrā tēctum, iūxtā \navnote{Īra -ae f --- animus īrātus}viam, ob iram, per portam \contr{} prope fenestram,

secundum ripam, post tergum, trāns flūmen, ultrā finēs, praeter offi-

[M.:] Dā praepositiōnēs cāsus ablātīvī!

[M.:] Quō modō? [D.:] Dīcimus enim ā domō, ab homine, cum exercitū, \navnote{iūs = locus ubi iūs dīcitur}cōram testibus, d\& forō, \& ture, ex prōvinciāy prō patriā,

[M.:] Dā utriusque cāsūs praepositiōnēs! [D.:] In, sub, super. [M.:] In et sub quandō accūsātivōo cāsuī iunguntur? \navnote{quandō = cum significāre = verbīs ostendere}[D.:] Quandō 'nōs in locum īre/iisse/itūrōs esse significāmus. [M.:] Quandō ablātīvō? [D.:] Quandō 'nōs in locō esse/fuisse/fiiturōs esse' signi-

[M.:] Super quam vim habet? \navnote{vīs= --- potestās, significātiō magis A quam B = nōn BsedA mentiō -ōnis/: f: mentiō- nem facere alicuius = loquī de aliquō}[D.:] Ubi locum significat, magis accūsātīvō cāsul servit quam ablātīvō; ubi mentiōnem alicuius facimus, ablā- 210 tivo tantum, ut

multa super Priamō rogitāns, super Hectore multa*

\booksection{Dē interiectiōne}

\navnote{[Vergilius: Aenēis 1.750] inter-iectiō -ōnis f}[M.:] Interiectiō quid est? \navnote{affectus -ūs m=id quō afficitur animus (ut laetitia, īra, timor, cēt.) inconditus -a -um = nōn ōrdinātus, rudis, sine arte}[D.:] Pars ōrātiōnis significāns mentis affectum vōce

[M.:] Significātiō interiectiōnis in quō est? [D.:] Aut laetitiam significāmus, ut euax! aut dolōrem, \navnote{admirātiō -ōnis f \contr{} admīrāri similis -e = quī idem esse vidētur, eiusdem generis similia: ut ō/ ei! heus!}ut heu! aut admīrātiōnem, ut papae! aut metum, ut at- 220 tat! et si qua sunt similia.

Versus sumptus ē librō primō illīus carminis cui titulus est \navnote{Aenēis -idis f Aenēis -ae m Carthāgō -inis/* f: urbs Africae Dīdō -ōnis f rēgīnā -ae f --- femina rēgnāns}Aenēis. In hōc carmine Vergilius poēta nārrat dē Aenēā Trōiānō, quī ē patriā fugiēns Carthāginem vēnit, ubi ā Dīdōne rēgīnā receptus est. Ā rēgīnā interrogātus Aenēās nārrat dē bellō Trōiānō et dē fugā suā (v. cap. XXXVII et XXXVIII

in alterā LINGVAE LATINAE parte).

\booksection{Grammatica Latīna}

D ē dēclīnātiōne Nōmina, prōnōmina, verba dēclīnantur. Cēterae partēs ōrātiōnis sunt indēcllnābilēs.

Dēcllnātiōnēs nōminum sunt quīnque: \navnote{Dēclīnātiōnēs:}Dēclīnātiō prīma: gen. sing. -ae, ut terra -ae. Dēclīnātiō secunda: gen. sing. -x, ut annus -ī, verbum -t. Dēclīnātiō tertia: gen. sing. -is, ut sōl -is, urbs -is. Dēclīnātiō quārta: gen. sing. -4s, ut portū\$ -ūs, genū -ūs. Dēclīnātiō quīnta: gen. sing. -ēil-ei, ut diēs -ēi, rēs reī.

Declināatiōnes verbōrum sunt quattuor, 'coniugātiōnēs' \navnote{con-iugāliō-ōnis/ Coniugāliōnēs: 1. con-iugātiō -ōnis f Coniugātiōnēs: l. āre 2.}quae vocantur: Coniugātiō prīma: īnf. -ārel-ari, ut amārey -n. Coniugātiō secunda: īnf. -ērel-ēri, ut monēre, -n. Coniugāliō tertia: īnf. -erel-ī, ut legere, -ī. Coniugātiō quārta: īnf. -Irel-xn, ut audire3 -n.

\booksection{Pēnsum A}

Dēclīnā haec vocābula: [1] āla, nōmen femininum 1 dēclīnātiōnis: Singulāris: nōm. haec āl-, acc. h- al, gen. h- āl-, dat. hdl-, abl. h- āl-. Plūrālis: nōm. h- āl-, acc. h- āl-, gen. h-

[2] pēs, nōmen mascullnum III dēclīnātiōnis: Singulāris: nōm. A- pēs, acc. h- ped-, gen. h- ped-, dat. hped-, abl. h- ped-. Plūrālis: nōm. h- ped-, acc. h- ped-, gen. h--- ped-, dat. h- ped-, abl. h- ped-. [3] ōrāre, verbum āctīvum I coniugātiōnis (pers. I sing.): Indicātīvus: praes. ōr- \contr{} imperf. ōr- \contr{} fut. ōr- \contr{} perf. ōrāc-, plūsquamperf. ōrāv-, plūsquamperf. ōrāo- \contr{} fut. perf. ōrāv-. Coniūnctīvus: praes. ōr- \contr{} imperf. ōr-, perf. ōrāv- \contr{} plūsquamperf. ōrāv-. [4] [4] dīcere, verbum āctīvum iii coniugātiōnis (pers. I sing.): Indicātīvus: praes. dic-, imperf. dic-, fut. dic-, perf. dīx-, plūsquamperf.

dīx-, fut. perf. dīx-. Coniūnctīvus: praes. dic-, imperf. dic-, perf. dīx-, plūsquamperf.

dīx-.

Verba: laedere -isse -um; implicāre ---isse -um; plaudere -isse -um; parere -isse -um; retinēre ---isse -um; accendere -isse -um; sedēre ---isse; nūbere -isse; sapere -isse; dēmere ---ijsse ---um; adicere -isse -um; adicere-isse -um.

\booksection{Pēnsum B}

Nōmen est pars --- quae corpus aut rem significat. Aemilia et Iūlia nōmina --- sunt, māter et filia sunt nōmina ---. --- nōminum sunt nōminātīvus, genetīvus, cēt. Gradūs --- sunt trēs: ---, comparātīvus, superlātīvus. Adiectīva quae comparantur --- aut quantitātem significant. Amāre est verbum primae ---. Interiectiō mentis ---, ut laetitiam vel dolōrem, significat. --- est affectus eius quī īrātus est.

Frātrēs geminl tam --- sunt quam ōva.

Synōnyma (vocābula quae idem ferē significant): plānum facere et ---, fortasse et ---, ideō et ---, --- \contr{} citrā et ---, ecce et ---.

Contrāria (vocābula quae rēs contrāriās significant): commūnis et ---, addere et ---.

\booksection{Pēnsum C}

Quae sunt partēs ōrātiōnis? Estne discipulus nōmen proprium? Cāsūs nōminum quī sunt? Quī sunt gradūs comparātiōnis? Prōnōmen quid est? Quot sunt coniugāatiōnes verbōrum? Num:ntus et foris adverbia temporis sunt? Cui cāsuī iungitur inter praepositiō? Quae praepositiōnēs ablātīvō iunguntur? Quibus cāsibus iungitur in praepositiō? Interiectiō quid significat?

\booksection{Vocābula nova}
\noindent\small ōrātiō, coniūnctiō, interiectiō, Mūsa, scamnum, sacerdōs, cāsus, comparātiō, quālitās, quantitās, significātiō, speciēs, īra, affectus, admirātiō, coniugātiō, proprius, appellātīvus, positīvus, optāūvus, cōpulātīvus, disiūnctīvus, explētīvus, causālis, ratiōnālis, inconditus, similis, īnflectere, dēmere, luctārī, explānāre, adnectere, ōrdināre, mentiōnem facere, dumtaxat, tantundem, quidnī, forsitan, sīquidem, quāpropter, proptereā, adversum, cis, ēn, eia, euax, papae, attat, synōnymum.

\bookbackmatter{Notae}

\begin{description}
\item[=] idem atque (pōnitur inter vocābula quae eandem ferē rem significant)
\item[\contr{}] contrārium (pōnitur inter vocābula quae rēs contrāriās
  significant)
\item[$<$] factum ex (pōnitur inter vocābula quōrum alterum ex alterō factum
  est)
\item[$|$] haec nota pōnitur ante litterās quae in dēclīnātiōne adduntur
\end{description}

\begin{multicols}{2}\small\raggedright
\begin{description}[leftmargin=1.6em,itemsep=1pt,parsep=0pt]
\item[abl] ablātīvus
\item[acc] accūsātīvus
\item[āct] āctīvum
\item[a.\,d.] ante diem
\item[adi] adiectīvum
\item[adv] adverbium
\item[cap.] capitulum
\item[cet.] cēterī -ae -a
\item[comp] comparātīvus
\item[con] coniūnctīvus
\item[dat] datīvus
\item[dēcl] dēclīnātiō
\item[dēp] dēpōnēns
\item[f] feminīnum
\item[fut] futūrum
\item[imper] imperātīvus
\item[impf] imperfectum
\item[indēcl] indēclīnābile
\item[īnf] īnfīnītīvus
\item[kal.] kalendae
\item[loc] locātīvus
\item[m, masc.] masculīnum
\item[n, neutr.] neutrum
\item[nōm] nōminātīvus
\item[nōn.] nōnae
\item[pāg.] pāgina
\item[part] participium
\item[pass] passīvum
\item[perf] perfectum
\item[pers] persōna
\item[pl, plūr.] plūrālis
\item[praes] praesēns
\item[prōn] prōnōmen
\item[prp] praepositiō
\item[s.\,d.] salūtem dīcit
\item[sg, sing.] singulāris
\item[sup] superlātīvus
\item[voc] vocātīvus
\end{description}
\end{multicols}


\end{document}
